% Options for packages loaded elsewhere
\PassOptionsToPackage{unicode}{hyperref}
\PassOptionsToPackage{hyphens}{url}
\PassOptionsToPackage{dvipsnames,svgnames,x11names}{xcolor}
\documentclass[
  american,
  11pt,
  letterpaper,
]{article}
\usepackage{xcolor}
\usepackage[margin=0.78in]{geometry}
\usepackage{amsmath,amssymb}
\setcounter{secnumdepth}{5}
\usepackage{iftex}
\ifPDFTeX
  \usepackage[T1]{fontenc}
  \usepackage[utf8]{inputenc}
  \usepackage{textcomp} % provide euro and other symbols
\else % if luatex or xetex
  \usepackage{unicode-math} % this also loads fontspec
  \defaultfontfeatures{Scale=MatchLowercase}
  \defaultfontfeatures[\rmfamily]{Ligatures=TeX,Scale=1}
\fi
\usepackage{lmodern}
\ifPDFTeX\else
  % xetex/luatex font selection
  \setmainfont[]{PingFang SC}
  \setsansfont[]{PingFang SC}
  \setmonofont[]{Menlo}
\fi
% Use upquote if available, for straight quotes in verbatim environments
\IfFileExists{upquote.sty}{\usepackage{upquote}}{}
\IfFileExists{microtype.sty}{% use microtype if available
  \usepackage[]{microtype}
  \UseMicrotypeSet[protrusion]{basicmath} % disable protrusion for tt fonts
}{}
\makeatletter
\@ifundefined{KOMAClassName}{% if non-KOMA class
  \IfFileExists{parskip.sty}{%
    \usepackage{parskip}
  }{% else
    \setlength{\parindent}{0pt}
    \setlength{\parskip}{6pt plus 2pt minus 1pt}}
}{% if KOMA class
  \KOMAoptions{parskip=half}}
\makeatother
\usepackage{color}
\usepackage{fancyvrb}
\newcommand{\VerbBar}{|}
\newcommand{\VERB}{\Verb[commandchars=\\\{\}]}
\DefineVerbatimEnvironment{Highlighting}{Verbatim}{commandchars=\\\{\}}
% Add ',fontsize=\small' for more characters per line
\usepackage{framed}
\definecolor{shadecolor}{RGB}{248,248,248}
\newenvironment{Shaded}{\begin{snugshade}}{\end{snugshade}}
\newcommand{\AlertTok}[1]{\textcolor[rgb]{0.94,0.16,0.16}{#1}}
\newcommand{\AnnotationTok}[1]{\textcolor[rgb]{0.56,0.35,0.01}{\textbf{\textit{#1}}}}
\newcommand{\AttributeTok}[1]{\textcolor[rgb]{0.13,0.29,0.53}{#1}}
\newcommand{\BaseNTok}[1]{\textcolor[rgb]{0.00,0.00,0.81}{#1}}
\newcommand{\BuiltInTok}[1]{#1}
\newcommand{\CharTok}[1]{\textcolor[rgb]{0.31,0.60,0.02}{#1}}
\newcommand{\CommentTok}[1]{\textcolor[rgb]{0.56,0.35,0.01}{\textit{#1}}}
\newcommand{\CommentVarTok}[1]{\textcolor[rgb]{0.56,0.35,0.01}{\textbf{\textit{#1}}}}
\newcommand{\ConstantTok}[1]{\textcolor[rgb]{0.56,0.35,0.01}{#1}}
\newcommand{\ControlFlowTok}[1]{\textcolor[rgb]{0.13,0.29,0.53}{\textbf{#1}}}
\newcommand{\DataTypeTok}[1]{\textcolor[rgb]{0.13,0.29,0.53}{#1}}
\newcommand{\DecValTok}[1]{\textcolor[rgb]{0.00,0.00,0.81}{#1}}
\newcommand{\DocumentationTok}[1]{\textcolor[rgb]{0.56,0.35,0.01}{\textbf{\textit{#1}}}}
\newcommand{\ErrorTok}[1]{\textcolor[rgb]{0.64,0.00,0.00}{\textbf{#1}}}
\newcommand{\ExtensionTok}[1]{#1}
\newcommand{\FloatTok}[1]{\textcolor[rgb]{0.00,0.00,0.81}{#1}}
\newcommand{\FunctionTok}[1]{\textcolor[rgb]{0.13,0.29,0.53}{\textbf{#1}}}
\newcommand{\ImportTok}[1]{#1}
\newcommand{\InformationTok}[1]{\textcolor[rgb]{0.56,0.35,0.01}{\textbf{\textit{#1}}}}
\newcommand{\KeywordTok}[1]{\textcolor[rgb]{0.13,0.29,0.53}{\textbf{#1}}}
\newcommand{\NormalTok}[1]{#1}
\newcommand{\OperatorTok}[1]{\textcolor[rgb]{0.81,0.36,0.00}{\textbf{#1}}}
\newcommand{\OtherTok}[1]{\textcolor[rgb]{0.56,0.35,0.01}{#1}}
\newcommand{\PreprocessorTok}[1]{\textcolor[rgb]{0.56,0.35,0.01}{\textit{#1}}}
\newcommand{\RegionMarkerTok}[1]{#1}
\newcommand{\SpecialCharTok}[1]{\textcolor[rgb]{0.81,0.36,0.00}{\textbf{#1}}}
\newcommand{\SpecialStringTok}[1]{\textcolor[rgb]{0.31,0.60,0.02}{#1}}
\newcommand{\StringTok}[1]{\textcolor[rgb]{0.31,0.60,0.02}{#1}}
\newcommand{\VariableTok}[1]{\textcolor[rgb]{0.00,0.00,0.00}{#1}}
\newcommand{\VerbatimStringTok}[1]{\textcolor[rgb]{0.31,0.60,0.02}{#1}}
\newcommand{\WarningTok}[1]{\textcolor[rgb]{0.56,0.35,0.01}{\textbf{\textit{#1}}}}
\usepackage{longtable,booktabs,array}
\newcounter{none} % for unnumbered tables
\usepackage{calc} % for calculating minipage widths
% Correct order of tables after \paragraph or \subparagraph
\usepackage{etoolbox}
\makeatletter
\patchcmd\longtable{\par}{\if@noskipsec\mbox{}\fi\par}{}{}
\makeatother
% Allow footnotes in longtable head/foot
\IfFileExists{footnotehyper.sty}{\usepackage{footnotehyper}}{\usepackage{footnote}}
\makesavenoteenv{longtable}
\usepackage{graphicx}
\makeatletter
\newsavebox\pandoc@box
\newcommand*\pandocbounded[1]{% scales image to fit in text height/width
  \sbox\pandoc@box{#1}%
  \Gscale@div\@tempa{\textheight}{\dimexpr\ht\pandoc@box+\dp\pandoc@box\relax}%
  \Gscale@div\@tempb{\linewidth}{\wd\pandoc@box}%
  \ifdim\@tempb\p@<\@tempa\p@\let\@tempa\@tempb\fi% select the smaller of both
  \ifdim\@tempa\p@<\p@\scalebox{\@tempa}{\usebox\pandoc@box}%
  \else\usebox{\pandoc@box}%
  \fi%
}
% Set default figure placement to htbp
\def\fps@figure{htbp}
\makeatother
\ifLuaTeX
\usepackage[bidi=basic,shorthands=off]{babel}
\else
\usepackage[bidi=default,shorthands=off]{babel}
\fi
\ifPDFTeX
\else
\babelfont{rm}[]{PingFang SC}
\fi
\ifLuaTeX
  \usepackage{selnolig} % disable illegal ligatures
\fi
\setlength{\emergencystretch}{3em} % prevent overfull lines
\providecommand{\tightlist}{%
  \setlength{\itemsep}{0pt}\setlength{\parskip}{0pt}}
\usepackage{fvextra}
\usepackage{etoolbox}
\usepackage{fancyhdr}
\setkeys{Gin}{keepaspectratio}
\XeTeXlinebreaklocale "zh"
\XeTeXlinebreakskip = 0pt plus 1pt
\fvset{breaklines=true,breakanywhere=true,fontsize=\scriptsize}
\DefineVerbatimEnvironment{Highlighting}{Verbatim}{commandchars=\\\{\},breaklines=true,breakanywhere=true,fontsize=\scriptsize}
\renewenvironment{Shaded}{\par\smallskip}{\par\smallskip}
\AtBeginEnvironment{longtable}{\footnotesize}
\setlength{\emergencystretch}{5em}
\setlength{\headheight}{15pt}
\pagestyle{fancy}
\fancyhf{}
\fancyhead[L]{\small\sffamily Microduck | Basketball + Bridge}
\fancyfoot[L]{\footnotesize\sffamily Simulation evidence edition 1.0}
\fancyfoot[R]{\thepage}
\renewcommand{\headrulewidth}{0.3pt}
\pretocmd{\section}{\clearpage}{}{}
\makeatletter
\renewcommand{\@pnumwidth}{2.2em}
\renewcommand{\@tocrmarg}{3em}
\renewcommand{\maketitle}{%
  \begin{titlepage}
  \vspace*{0.35in}
  {\color{blue!55!black}\rule{\textwidth}{3pt}}\par\bigskip
  {\huge\bfseries\@title\par}
  \vspace{0.35in}
  {\large Author: \@author\par}\medskip
  {\@date\par}
  \vspace{0.5in}
  \includegraphics[width=\textwidth]{assets/evidence/observed-live-scenes.png}\par
  \vspace{0.25in}
  {\small Real simulation evidence. Basketball balance verified; rolling and full bridge crossing remain unaccepted.\par}
  \vfill
  {\small PPO / MuJoCo / Curriculum / ONNX / Reproducible code\par}
  \end{titlepage}
}
\makeatother
\usepackage{bookmark}
\IfFileExists{xurl.sty}{\usepackage{xurl}}{} % add URL line breaks if available
\urlstyle{same}
\hypersetup{
  pdftitle={The Microduck Balance Laboratory},
  pdfauthor={George Hu},
  pdflang={en-US},
  colorlinks=true,
  linkcolor={Maroon},
  filecolor={Maroon},
  citecolor={Blue},
  urlcolor={Blue},
  pdfcreator={LaTeX via pandoc}}

\title{The Microduck Balance Laboratory}
\usepackage{etoolbox}
\makeatletter
\providecommand{\subtitle}[1]{% add subtitle to \maketitle
  \apptocmd{\@title}{\par {\large #1 \par}}{}{}
}
\makeatother
\subtitle{Basketball and Suspended Bridge: Modeling, PPO, Evidence and
Reproduction}
\author{George Hu}
\date{2026-09-11 · v1.0 · Evidence: 2026-09-10}

\begin{document}
\maketitle

\renewcommand*\contentsname{Contents and learning route}
{
\setcounter{tocdepth}{2}
\tableofcontents
}
\section{Two balance laboratories: what this book actually
reproduces}\label{two-balance-laboratories-what-this-book-actually-reproduces}

\subsection{A promise with a boundary}\label{a-promise-with-a-boundary}

Imagine standing on a basketball. Moving your feet changes the ball's
motion; the ball's motion changes the surface under your feet. Now
imagine stepping onto a narrow plank suspended by four cables. Every
footstep changes the support that the next footstep depends on. These
are \textbf{coupled balance problems}, not simply walking animations
placed on unusual scenery.

This book follows the working Microduck repository from physical
assumptions to simulation, observations, actions, PPO updates, exported
policies, evaluation, and the seven-case Duck Viewer. It assumes
algebra, graphs, and elementary probability, not university mechanics or
calculus. Each symbol is introduced before use. Code in the source
appendix is actual implementation, not a fictional miniature presented
as a complete trainer.

The evidence snapshot concerns local experiments recorded on
\textbf{September 10, 2026}. The edition is \textbf{1.0, September 11,
2026}. A version number identifies a document snapshot; it does not mean
that every skill is solved.

{\def\LTcaptype{none} % do not increment counter
\begin{longtable}[]{@{}
  >{\raggedright\arraybackslash}p{(\linewidth - 2\tabcolsep) * \real{0.5000}}
  >{\raggedright\arraybackslash}p{(\linewidth - 2\tabcolsep) * \real{0.5000}}@{}}
\toprule\noalign{}
\begin{minipage}[b]{\linewidth}\raggedright
Question
\end{minipage} & \begin{minipage}[b]{\linewidth}\raggedright
Evidence-backed answer
\end{minipage} \\
\midrule\noalign{}
\endhead
\bottomrule\noalign{}
\endlastfoot
Can a duck balance on a free basketball? & Yes, the fixed-command local
candidate survived all six recorded 60-second nominal trials. \\
Can it drive the ball reliably in the requested direction? & Not
established: zero of three forward-command trials passed the rolling
gate. \\
Does the suspended bridge exist as a working physical simulation? & Yes:
a free plank, four physical cable tendons, platforms, and visible
gantries. \\
Did the final bridge pilot complete an accepted crossing? & No.~It has
32,768 new PPO transitions and zero accepted crossings. \\
Are both examples available in Duck Viewer? & Yes. Seven cases means
five existing cases plus these two, not seven mastered skills. \\
Was either behavior verified on hardware? & No.~All results here are
simulation-only. \\
\end{longtable}
}

The successful basketball behavior began from a \textbf{supplied trained
recurrent actor}. The successful parts of this local experiment are
source import, accurate actor transfer, further PPO training, export
parity, sustained balance, and visible playback. This is not evidence
that a randomly initialized local network learned ball balance in
102,400 steps. Likewise, a successful bridge training command means the
pipeline ran; it does not mean the duck crossed.

``From scratch'' therefore has two meanings. You can reconstruct the
\textbf{software experiment} from source and required assets. You cannot
reconstruct an undocumented upstream training history or replace a
missing mature checkpoint with an assertion. This edition gives the
reproducible local path in full and marks the original policy as an
external prerequisite. The book's \texttt{source/} snapshot and ZIP are
supporting source, not a redistribution of all third-party robot meshes,
software dependencies, or policy weights.

\subsection{Choose a learning route}\label{choose-a-learning-route}

\begin{enumerate}
\def\labelenumi{\arabic{enumi}.}
\tightlist
\item
  \textbf{First afternoon:} read the assumptions, inspect the scene
  diagrams, run preflight, run the pure-Python PPO arithmetic, and
  inspect real contact sheets. No training is necessary to learn the
  data flow.
\item
  \textbf{First experiment:} preserve the supplied artifacts, create a
  fresh run directory, execute a tiny smoke run, and inspect every
  output. A four-step smoke is an integration check only.
\item
  \textbf{Measured reproduction:} run the documented fixed-command
  basketball adaptation or the bridge pilot, export, evaluate
  deterministic policies with zero assistance, then inspect the complete
  rendered episode.
\item
  \textbf{Research extension:} change one physics or curriculum
  parameter, keep the evaluation fixed, run independent seeds, and
  document failures as carefully as successes.
\end{enumerate}

Read the PPO chapter before tuning a learning rate. Read the evaluation
chapter before deciding that a curve proves learning. Read the source
appendix alongside the case chapters rather than starting with thousands
of lines without context.

\subsection{A vocabulary we will keep
consistent}\label{a-vocabulary-we-will-keep-consistent}

{\def\LTcaptype{none} % do not increment counter
\begin{longtable}[]{@{}
  >{\raggedright\arraybackslash}p{(\linewidth - 2\tabcolsep) * \real{0.5000}}
  >{\raggedright\arraybackslash}p{(\linewidth - 2\tabcolsep) * \real{0.5000}}@{}}
\toprule\noalign{}
\begin{minipage}[b]{\linewidth}\raggedright
Term
\end{minipage} & \begin{minipage}[b]{\linewidth}\raggedright
Meaning in this book
\end{minipage} \\
\midrule\noalign{}
\endhead
\bottomrule\noalign{}
\endlastfoot
State & All simulated physical variables, including information the
policy cannot see. \\
Observation & The 61 numbers provided to the actor at one control
step. \\
Action & Fourteen policy outputs interpreted by the robot control
contract. \\
Policy / actor & A function producing an action distribution during
training, and its mean during deterministic evaluation. \\
Critic / value & A learned estimate of future discounted reward, used
for training. \\
Episode & One attempt between reset and termination/time limit. \\
Transition & One observation--action--reward--next-observation
interaction in one environment. \\
Rollout & A block of transitions collected before updating the
network. \\
Checkpoint & Training weights plus whatever metadata that checkpoint
format actually stores. \\
ONNX export & Portable inference graph with the observation normalizer
baked in. \\
Curriculum & A sequence of easier-to-harder physics, spawning, or
judging conditions; not permission to change rewards silently. \\
Parity & Numerical agreement between two implementations on the same
inputs and recurrent-state lifecycle. \\
Skill gate & A measurable acceptance condition, distinct from process
exit or file creation. \\
\end{longtable}
}

\subsection{Keep an experiment
notebook}\label{keep-an-experiment-notebook}

Record the source hashes, command, Python executable, package versions,
actuator mode, seed, spawn distribution, assistance level, randomization
settings, command distribution, step count, and evaluation horizon. Save
the full JSON result, not only a screenshot of its green region.

An example notebook entry is:

\begin{Shaded}
\begin{Highlighting}[]
\FunctionTok{\{}
  \DataTypeTok{"experiment"}\FunctionTok{:} \StringTok{"my{-}independent{-}basketball{-}run"}\FunctionTok{,}
  \DataTypeTok{"initialization"}\FunctionTok{:} \StringTok{"supplied actor; fresh critic and optimizer"}\FunctionTok{,}
  \DataTypeTok{"control\_dt\_s"}\FunctionTok{:} \FloatTok{0.02}\FunctionTok{,}
  \DataTypeTok{"actor\_observations"}\FunctionTok{:} \DecValTok{61}\FunctionTok{,}
  \DataTypeTok{"actions"}\FunctionTok{:} \DecValTok{14}\FunctionTok{,}
  \DataTypeTok{"claim\_to\_test"}\FunctionTok{:} \StringTok{"60{-}second free{-}ball balance"}\FunctionTok{,}
  \DataTypeTok{"claim\_not\_yet\_supported"}\FunctionTok{:} \StringTok{"commanded rolling and steering"}\FunctionTok{,}
  \DataTypeTok{"hardware\_tested"}\FunctionTok{:} \KeywordTok{false}
\FunctionTok{\}}
\end{Highlighting}
\end{Shaded}

This is a template, not an additional measured run. The actual run
settings live beside the checkpoints. Filenames such as \texttt{final},
\texttt{full}, or \texttt{success} are not evidence by themselves.

\subsection{Where everything lives}\label{where-everything-lives}

All shell commands in this book start at the \textbf{workspace root},
unless a block explicitly changes directory. A path beginning
\texttt{rlx/} is relative to that root. Book assets and source snapshots
are relative to
\texttt{docs/\allowbreak{}basketba\allowbreak{}ll-\allowbreak{}bridge-\allowbreak{}book/\allowbreak{}}.

\begin{Shaded}
\begin{Highlighting}[]
\NormalTok{microduck{-}lab/}
\NormalTok{  microduck\_local/       CPU MuJoCo harness, robot contract, lab server}
\NormalTok{  microduck\_rl/          upstream robot/training reference checkout}
\NormalTok{  microduck/             upstream runtime and supplied walking policies}
\NormalTok{  microduck{-}playground/  supplied basketball experiment/assets/checkpoint}
\NormalTok{  rlx/                  local PPO, task environments, ONNX exports}
\NormalTok{  duck{-}viewer/           browser Studio and real scene playback}
\NormalTok{  scripts/              scenario command wrappers and API verification}
\NormalTok{  docs/basketball{-}bridge{-}book/}
\NormalTok{    chapters/en/        English chapter sources}
\NormalTok{    chapters/zh/        Chinese chapter sources}
\NormalTok{    assets/             measured plots, screenshots, schematic diagrams}
\NormalTok{    examples/           executable arithmetic and environment probes}
\NormalTok{    tools/              build, preflight, validation}
\NormalTok{    source/             exact selected implementation snapshot}
\end{Highlighting}
\end{Shaded}

The full workspace is the execution unit. Copying one environment file
to an empty directory will not provide its robot contract, meshes,
actuator code, or training framework.

\section{Mathematical and physical
foundations}\label{mathematical-and-physical-foundations}

\subsection{Units before equations}\label{units-before-equations}

Use metres for length, seconds for time, kilograms for mass, radians for
angles, newtons for force, and newton-metres for torque. A value without
a unit can look reasonable while being wrong by a factor of a thousand.

At a 50 Hz control rate the actor is consulted 50 times each second:

\[\Delta t = 1/50 = 0.02\ \mathrm{s}.\]

A 60-second basketball test contains 3,000 control steps. A 20-second
bridge episode contains 1,000. These are control steps; the physics
engine can perform multiple smaller integration steps inside each one.
More environments increase data per rollout, not the duration of one
simulated second.

Velocity is change of position divided by elapsed time. If the duck
moves 0.03 m in 0.2 s, its average speed is 0.15 m/s. Acceleration is
change of velocity divided by elapsed time. These finite differences are
enough to understand the simulator's job before studying continuous
derivatives.

Forces change momentum. In one dimension, a simple teaching update is

\[v_{t+1}=v_t+(F/m)\Delta t,\qquad x_{t+1}=x_t+v_{t+1}\Delta t.\]

This is an explanatory approximation, \textbf{not} a replacement for
MuJoCo's integration, articulated-body dynamics, and contact solver. The
actual simulator also resolves joints, contacts, friction and
constraints. A torque is a turning effect: force times its perpendicular
lever arm. A foot force displaced from the plank's centre can rotate it
even when the duck remains upright.

\subsection{Why balance and motion
interact}\label{why-balance-and-motion-interact}

On a floor, the support under a planted foot is nearly fixed. On the
ball, a foot tangential force can rotate the ball and move the support.
If an ideal rigid sphere rolls without slip, distance and angle satisfy

\[d=r\theta.\]

With radius 0.12 m, a 90-degree rotation is \(\theta=\pi/2\) radians,
giving \(d\approx0.1885\) m. At 0.15 m/s the ideal angular speed is
\(v/r=1.25\) rad/s. These are scale checks, not guarantees: slipping,
contact deformation and turning break the simple straight-line relation.

For the suspended bridge, gravity pulls the combined masses down, while
cables support and constrain the plank. A cable is not a rigid beam: its
tension acts along its length. The implementation's tendon and limit
settings approximate the desired cable behavior; their exact meaning is
determined by the model, not by how the rendered cable looks. A gantry
can be visible while having no collision role. The visual scene is
evidence of appearance, not a complete description of the dynamics.

A material name such as ``wood'' or ``rubber'' does not automatically
create correct physics. The simulator needs mass, dimensions, inertia,
friction/contact parameters, joint and tendon definitions. The wooden
appearance is a visual material. The 0.45 kg plank mass is a dynamic
parameter. They must be documented separately.

\subsection{Coordinate frames: forward for
whom?}\label{coordinate-frames-forward-for-whom}

The \textbf{world frame} belongs to the scene. The \textbf{body frame}
turns with the duck. If the duck rotates 90 degrees, body-forward no
longer points along world-x. Rewarding world-x velocity while issuing
body-forward commands creates a misleading task.

For a planar heading angle \(\psi\), a body-frame velocity can be
expressed in world coordinates as

\[v_x^w=\cos\psi\,v_x^b-\sin\psi\,v_y^b,\]
\[v_y^w=\sin\psi\,v_x^b+\cos\psi\,v_y^b.\]

Superscripts \(w\) and \(b\) mean world and body, not exponentiation.
With body velocity \((0.15,0)\) and heading 90 degrees, world velocity
is approximately \((0,0.15)\). The basketball evaluator uses an
initial-heading reference for directional progress; the bridge commander
translates navigation information into ordinary body-frame commands.
Neither fact makes world position an actor input.

\subsection{Feedback, memory, and partial
observation}\label{feedback-memory-and-partial-observation}

An \textbf{open-loop} script repeats preset joint targets regardless of
what happens. A \textbf{feedback policy} chooses the next targets using
current sensed information. If a foot slips, the next observation
changes and the policy can respond.

The robot does not receive the entire simulator state. It gets angular
velocity, projected gravity, commands, joint information, previous
actions and reserved command slots. It does not directly receive ball
centre or plank pose. This is partial observation. Two identical current
observations may have different motion histories. The basketball actor
therefore carries LSTM memory; the bridge actor is a transferred
feedforward walking actor. A memory network is not clairvoyance: it can
only summarize evidence that entered its observations.

During learning, random sampled actions help explore nearby behavior.
During deterministic evaluation, the exported actor's mean is used. A
stochastic policy that occasionally catches itself may export to a mean
policy that falls. That is why training reward and exported-policy
testing are different experiments.

\subsection{Arithmetic workshop}\label{arithmetic-workshop}

\begin{Shaded}
\begin{Highlighting}[]
\ImportTok{import}\NormalTok{ math}

\NormalTok{control\_dt }\OperatorTok{=} \FloatTok{0.02}
\NormalTok{radius }\OperatorTok{=} \FloatTok{0.12}
\NormalTok{command\_speed }\OperatorTok{=} \FloatTok{0.15}
\NormalTok{steps }\OperatorTok{=} \BuiltInTok{round}\NormalTok{(}\DecValTok{60} \OperatorTok{/}\NormalTok{ control\_dt)}
\NormalTok{quarter\_turn\_distance }\OperatorTok{=}\NormalTok{ radius }\OperatorTok{*}\NormalTok{ math.pi }\OperatorTok{/} \DecValTok{2}
\NormalTok{ideal\_angular\_speed }\OperatorTok{=}\NormalTok{ command\_speed }\OperatorTok{/}\NormalTok{ radius}
\BuiltInTok{print}\NormalTok{(steps, }\BuiltInTok{round}\NormalTok{(quarter\_turn\_distance, }\DecValTok{4}\NormalTok{), ideal\_angular\_speed)}
\ControlFlowTok{assert}\NormalTok{ steps }\OperatorTok{==} \DecValTok{3000}
\ControlFlowTok{assert} \BuiltInTok{abs}\NormalTok{(ideal\_angular\_speed }\OperatorTok{{-}} \FloatTok{1.25}\NormalTok{) }\OperatorTok{\textless{}} \FloatTok{1e{-}12}
\end{Highlighting}
\end{Shaded}

Expected output: \texttt{3000\ 0.1885\ 1.25}. This is a unit test of
arithmetic, not a balance test.

\textbf{Exercise.} Four environments collect 128 steps before each
update. How many transitions are collected? How much simulated time
elapses in each lane? How many such updates make 32,768 transitions?

\textbf{Answer.} \(4\times128=512\) transitions; each lane advances
\(128\times0.02=2.56\) seconds; \(32768/512=64\) updates. Summing lane
durations gives aggregate simulated experience, not one continuous
655.36-second bridge crossing.

\textbf{Exercise.} Why can a high average reward fail to prove a
60-second hold?

\textbf{Answer.} Short successful segments separated by falls/resets can
have high average reward. A sustained hold requires one unbroken
episode, first-fall accounting, valid foot support, and the entire
requested duration.

\section{Set up a reproducible
workshop}\label{set-up-a-reproducible-workshop}

\subsection{Materials: simulation assets are not a hardware shopping
list}\label{materials-simulation-assets-are-not-a-hardware-shopping-list}

The following is the \textbf{simulation bill of materials}. It
identifies objects and files you need to reproduce the software
experiment. It is not a certified mechanical design, purchasing
recommendation, or construction specification for putting an expensive
robot above a moving plank.

{\def\LTcaptype{none} % do not increment counter
\begin{longtable}[]{@{}
  >{\raggedright\arraybackslash}p{(\linewidth - 4\tabcolsep) * \real{0.3000}}
  >{\raggedleft\arraybackslash}p{(\linewidth - 4\tabcolsep) * \real{0.4000}}
  >{\raggedright\arraybackslash}p{(\linewidth - 4\tabcolsep) * \real{0.3000}}@{}}
\toprule\noalign{}
\begin{minipage}[b]{\linewidth}\raggedright
Item
\end{minipage} & \begin{minipage}[b]{\linewidth}\raggedleft
Quantity
\end{minipage} & \begin{minipage}[b]{\linewidth}\raggedright
Role and limitation
\end{minipage} \\
\midrule\noalign{}
\endhead
\bottomrule\noalign{}
\endlastfoot
Microduck full-collision robot MJCF and referenced meshes & One model
per scene & Fourteen actuated joints; inspect the packaged XML, do not
substitute an unrelated robot. \\
Free basketball & One & Radius 0.12 m, mass 0.62 kg; collision sphere
and separate textured visual mesh. \\
Basketball OBJ and PNG & One each & Supplied playground assets;
appearance is separate from sphere contact physics. \\
Suspended plank & One & Length 1.10 m, width 0.13 m, thickness 0.03 m,
mass 0.45 kg. \\
Cable tendons & Four & Physical constraints/support, not only decorative
lines. \\
Launch and destination platforms & Two & Static support and unambiguous
start/end regions. \\
Visible gantries & Scene components & Show anchor placement;
non-colliding visual structure. \\
Ground plane, lighting, camera & One scene setup & Ground contact is
detected as failure, not hidden by the camera. \\
Supplied basketball actor and reference ONNX & One pair & Required
mature initializer for the measured local experiment. \\
Supplied walking ONNX & One & Exact bridge actor/normalizer
initializer. \\
\end{longtable}
}

No physical timber grade, cable breaking load, fastening specification,
battery limit, servo thermal envelope or safety factor was measured for
this book. Do not infer those from simulated dimensions. Real tests
would need engineered restraints, an exclusion zone, an emergency stop,
torque/current limits, qualified supervision, and separate sim-to-real
validation. The local harness is a prototype, not deployment approval.

\subsection{Computer and software
assumptions}\label{computer-and-software-assumptions}

The recorded path uses an Apple-Silicon Mac. MuJoCo steps the physical
scene on CPU; basketball PPO uses PyTorch CPU; bridge PPO uses
MLX/Metal. ``No CUDA required'' does \textbf{not} mean ``every part runs
without a GPU runtime on every operating system.'' This edition does not
claim a clean Linux reproduction of the bridge learner.

Use Python 3.12, not whichever \texttt{python} happens to be first in
your shell. Node/npm run the viewer. \texttt{uv} manages Python
environments. Pandoc, XeLaTeX, \texttt{pdfinfo}, \texttt{pdftotext}, and
\texttt{pdftoppm} build and inspect the book; they are not needed for
training. Pillow and Matplotlib generate plots, not policy actions.

The measured Python package versions are recorded automatically in
\texttt{environm\allowbreak{}ent.json\allowbreak{}}. This is a
\textbf{recorded environment inventory}, not a tested universal resolver
lock. Source changes outside a Git commit are bound by file SHA-256
values in \texttt{source-\allowbreak{}manifest\allowbreak{}.json}; a
commit hash alone would miss the new uncommitted skill implementations.

\subsection{Acquire the workspace, not just the
PDF}\label{acquire-the-workspace-not-just-the-pdf}

Obtain this workspace checkout from its maintainer, including its
\texttt{rlx} checkout, upstream sibling repositories and the
user-supplied playground references. Do not invent a public download URL
for policy weights. Before installing anything, inspect the local
prerequisites:

\begin{Shaded}
\begin{Highlighting}[]
\BuiltInTok{test} \AttributeTok{{-}f}\NormalTok{ microduck\_local/src/microduck\_local/contract.py}
\BuiltInTok{test} \AttributeTok{{-}f}\NormalTok{ rlx/rlx/mjlab\_microduck/robot/microduck/robot\_allcollisions.xml}
\BuiltInTok{test} \AttributeTok{{-}f}\NormalTok{ microduck/policies/alpha\_walking.onnx}
\BuiltInTok{test} \AttributeTok{{-}f}\NormalTok{ microduck{-}playground/artifacts/basketball/checkpoint.pt}
\BuiltInTok{test} \AttributeTok{{-}f}\NormalTok{ microduck{-}playground/artifacts/basketball/policy.onnx}
\BuiltInTok{test} \AttributeTok{{-}f}\NormalTok{ microduck{-}playground/src/mjlab\_microduck/robot/assets/basketball/basketball.obj}
\BuiltInTok{test} \AttributeTok{{-}f}\NormalTok{ microduck{-}playground/src/mjlab\_microduck/robot/assets/basketball/basketball.png}
\BuiltInTok{command} \AttributeTok{{-}v}\NormalTok{ uv node npm}
\end{Highlighting}
\end{Shaded}

If a test fails, stop at that missing asset. Changing the actor
architecture or replacing its normalizer with zeros does not repair a
missing checkpoint. The bridge scene itself uses the packaged robot
assets; basketball additionally requires the supplied playground
OBJ/PNG. Walking baselines and other harness paths may require the
upstream sibling model checkout.

For a new environment, the repo's documented setup route is:

\begin{Shaded}
\begin{Highlighting}[]
\BuiltInTok{cd}\NormalTok{ rlx}
\ExtensionTok{uv}\NormalTok{ sync }\AttributeTok{{-}{-}python}\NormalTok{ 3.12}
\ExtensionTok{uv}\NormalTok{ pip install }\AttributeTok{{-}{-}python}\NormalTok{ .venv/bin/python }\AttributeTok{{-}e}\NormalTok{ ../microduck\_local}
\ExtensionTok{.venv/bin/python}\NormalTok{ examples/ppo\_microduck\_studio.py }\AttributeTok{{-}{-}help}
\BuiltInTok{cd}\NormalTok{ ..}
\BuiltInTok{cd}\NormalTok{ duck{-}viewer}
\ExtensionTok{npm}\NormalTok{ ci}
\BuiltInTok{cd}\NormalTok{ ..}
\end{Highlighting}
\end{Shaded}

These installation commands require network/package availability. The
book's validation uses the existing recorded environment and does not
claim to have recreated a fresh machine offline. Preserve lockfiles;
compare package versions and test before changing a resolved dependency.
Do not copy a virtual environment between computers.

The recorded executable is
\texttt{rlx/\allowbreak{}.venv-\allowbreak{}microduc\allowbreak{}k/\allowbreak{}bin/\allowbreak{}python}.
The clean setup above creates
\texttt{rlx/\allowbreak{}.venv/\allowbreak{}bin/\allowbreak{}python}.
Substitute the latter consistently in the case chapters if that is your
environment. The following variable makes the distinction explicit:

\begin{Shaded}
\begin{Highlighting}[]
\VariableTok{PYTHON}\OperatorTok{=}\StringTok{"}\VariableTok{$PWD}\StringTok{/rlx/.venv{-}microduck/bin/python"}
\BuiltInTok{test} \AttributeTok{{-}x} \StringTok{"}\VariableTok{$PYTHON}\StringTok{"}
\StringTok{"}\VariableTok{$PYTHON}\StringTok{"}\NormalTok{ docs/basketball{-}bridge{-}book/tools/preflight.py}
\StringTok{"}\VariableTok{$PYTHON}\StringTok{"}\NormalTok{ docs/basketball{-}bridge{-}book/examples/ppo\_arithmetic.py}
\end{Highlighting}
\end{Shaded}

Preflight checks required local assets and reports package versions
without importing Metal. It does not prove that a subsequent Metal or
graphics context can be opened. A restricted sandbox may permit file
reads but reject Metal or CoreGraphics. That is a runtime permission
problem, not proof of a broken actor.

\subsection{Separate preflight, smoke, training, and
evaluation}\label{separate-preflight-smoke-training-and-evaluation}

{\def\LTcaptype{none} % do not increment counter
\begin{longtable}[]{@{}
  >{\raggedright\arraybackslash}p{(\linewidth - 4\tabcolsep) * \real{0.3333}}
  >{\raggedright\arraybackslash}p{(\linewidth - 4\tabcolsep) * \real{0.3333}}
  >{\raggedright\arraybackslash}p{(\linewidth - 4\tabcolsep) * \real{0.3333}}@{}}
\toprule\noalign{}
\begin{minipage}[b]{\linewidth}\raggedright
Stage
\end{minipage} & \begin{minipage}[b]{\linewidth}\raggedright
Smallest useful proof
\end{minipage} & \begin{minipage}[b]{\linewidth}\raggedright
What it cannot establish
\end{minipage} \\
\midrule\noalign{}
\endhead
\bottomrule\noalign{}
\endlastfoot
Preflight & Assets exist, expected files and packages can be located. &
Dynamics, graphics, or policy quality. \\
Contract test & Shapes, resets, curriculum and evaluator behavior are
correct on fixtures/probes. & A trained skill is reliable. \\
Smoke training & A few transitions reach update, checkpoint and ONNX
output. & Skill acquisition or convergence. \\
Pilot training & A bounded real experiment changes weights and records
metrics. & Crossing/steering unless separately evaluated. \\
Deterministic evaluation & Exported policy meets or fails explicit
conditions. & Robustness outside tested conditions. \\
Visual review & Contact, body posture and actual motion agree with
metrics. & An unseen seed or hardware scenario. \\
\end{longtable}
}

Run targeted contract tests before spending a long time training:

\begin{Shaded}
\begin{Highlighting}[]
\ExtensionTok{rlx/.venv{-}microduck/bin/python} \AttributeTok{{-}m}\NormalTok{ pytest }\AttributeTok{{-}q} \DataTypeTok{\textbackslash{}}
\NormalTok{  rlx/tests/test\_basketball.py }\DataTypeTok{\textbackslash{}}
\NormalTok{  rlx/tests/test\_basketball\_policy.py }\DataTypeTok{\textbackslash{}}
\NormalTok{  rlx/tests/test\_basketball\_evaluation.py }\DataTypeTok{\textbackslash{}}
\NormalTok{  rlx/tests/test\_balance\_adapter.py }\DataTypeTok{\textbackslash{}}
\NormalTok{  rlx/tests/test\_bridge.py }\DataTypeTok{\textbackslash{}}
\NormalTok{  rlx/tests/test\_bridge\_evaluation.py }\DataTypeTok{\textbackslash{}}
\NormalTok{  rlx/tests/test\_bridge\_bootstrap.py}
\end{Highlighting}
\end{Shaded}

These tests exercise the named implementation boundary. They do not
assert that every unrelated test in the workspace passes. For this
publication, current test results and book-only checks are saved
separately under \texttt{verifica\allowbreak{}tion/\allowbreak{}}.

\subsection{Before the first step: inspect the
scene}\label{before-the-first-step-inspect-the-scene}

Run the book's environment probe to compile both models, check the
observation/action shape and sample four finite transitions. It applies
zero action and records a construction check; it does not load a learned
policy or prove balance:

\begin{Shaded}
\begin{Highlighting}[]
\ExtensionTok{rlx/.venv{-}microduck/bin/python} \DataTypeTok{\textbackslash{}}
\NormalTok{  docs/basketball{-}bridge{-}book/examples/probe\_scenes.py}
\end{Highlighting}
\end{Shaded}

Each environment owns mutable MuJoCo data. Shared immutable model
resources can save memory, but recurrent hidden states and per-duck
simulation states must never be shared. A viewer can display two ducks
using the same ONNX session while each maintains independent LSTM
memory.

\subsection{Preserve outputs before a new
run}\label{preserve-outputs-before-a-new-run}

Use new output directories, such as
\texttt{rlx/\allowbreak{}runs/\allowbreak{}book/\allowbreak{}basketba\allowbreak{}ll-\allowbreak{}YYYYMMDD\allowbreak{}-\allowbreak{}HHMMSS},
replacing the date token with a fresh run identifier. Check the
destination does not already exist. Never train into the evidence
directories printed in this book.

For each run, preserve at least:

\begin{itemize}
\tightlist
\item
  the executed command, software inventory, source hash and
  source-policy hash;
\item
  checkpoint and associated metadata, final ONNX and parity result;
\item
  full update/episode logs, not only a smoothed reward graph;
\item
  evaluation JSON including failures, seeds, horizon and assistance
  settings;
\item
  video, contact sheet and render metadata bound to the evaluated
  policy.
\end{itemize}

A checkpoint is not automatically a complete optimizer resume. In these
two pathways, source actors are transferred but critics/optimizers are
new. Re-running a command may reproduce the protocol without reproducing
identical floating-point weights. Different numerical libraries, package
versions, scheduling and simulator versions can change trajectories. The
honest goal is a traceable experiment and matching acceptance criteria,
not an unjustified promise of bitwise retraining.

\subsection{Troubleshooting without changing the
question}\label{troubleshooting-without-changing-the-question}

{\def\LTcaptype{none} % do not increment counter
\begin{longtable}[]{@{}
  >{\raggedright\arraybackslash}p{(\linewidth - 4\tabcolsep) * \real{0.3333}}
  >{\raggedright\arraybackslash}p{(\linewidth - 4\tabcolsep) * \real{0.3333}}
  >{\raggedright\arraybackslash}p{(\linewidth - 4\tabcolsep) * \real{0.3333}}@{}}
\toprule\noalign{}
\begin{minipage}[b]{\linewidth}\raggedright
Symptom
\end{minipage} & \begin{minipage}[b]{\linewidth}\raggedright
First check
\end{minipage} & \begin{minipage}[b]{\linewidth}\raggedright
Do not do this
\end{minipage} \\
\midrule\noalign{}
\endhead
\bottomrule\noalign{}
\endlastfoot
Missing ball appearance & OBJ/PNG paths and group-2 visual geoms. &
Replace the physical sphere with decorative animation. \\
Model compiles but viewer is flat ground & Per-policy environment kwargs
and lab scene extraction. & Claim bridge playback from the policy name
alone. \\
Actor output diverges after export & Normalizer divisor, observation
order, h/c reset and ONNX metadata. & Ignore parity because the duck
looks plausible. \\
\texttt{No\ Metal\ device\ available} & Execute MLX in an authorized
local runtime. & Change algorithms silently and call it the same
experiment. \\
CoreGraphics rendering failure & Authorized graphics context and
installed rendering backend. & Substitute a schematic as a measured
rollout screenshot. \\
High reward, no crossing & Contact-valid full-horizon evaluation and
rollout inspection. & Raise a reward coefficient before checking
exploration. \\
JSON/stdout parsing mismatch & Prefer structured
\texttt{training\allowbreak{}-\allowbreak{}metrics.\allowbreak{}jsonl}
and result files. & Treat incomplete log parsing as a flat zero loss
curve. \\
\end{longtable}
}

\textbf{Exercise.} Why preserve both policy hashes and video evidence?
\textbf{Answer:} a good-looking video can come from a different policy.
Hash-bound metadata connects the claim, checkpoint, deterministic
export, evaluation and visible trajectory.

\section{PPO: Learning Carefully from a Short Batch of
Experience}\label{ppo-learning-carefully-from-a-short-batch-of-experience}

Proximal Policy Optimization, usually shortened to \textbf{PPO}, is the
learning algorithm used by the local RLX trainer. PPO does not directly
teach a robot what basketball means. The environment supplies
observations, actions, rewards, and episode boundaries. PPO uses those
records to answer a narrower question:

\begin{quote}
Given the actions the current policy just tried, which actions should
become a little more likely, which should become a little less likely,
and how large a change is safe enough to attempt?
\end{quote}

This chapter develops that answer from high-school algebra, probability
density, and averages. It then maps each idea to the repository's
current implementation:

\begin{itemize}
\tightlist
\item
  \texttt{rlx/\allowbreak{}examples\allowbreak{}/\allowbreak{}ppo\_\allowbreak{}microduc\allowbreak{}k\_\allowbreak{}basketba\allowbreak{}ll.py}:
  the recurrent PyTorch continuation trainer for the released basketball
  actor;
\item
  \texttt{rlx/\allowbreak{}rlx/\allowbreak{}algorith\allowbreak{}ms/\allowbreak{}ppo.py}:
  the generic feedforward MLX PPO implementation;
\item
  \texttt{rlx/\allowbreak{}rlx/\allowbreak{}buffers/\allowbreak{}rollout\_\allowbreak{}buffer.p\allowbreak{}y}:
  the fixed-size rollout storage used by the MLX implementation.
\end{itemize}

The code is the authority for implementation details. PPO papers explain
the general method, but this repository adds choices such as
timeout-aware bootstrapping, a robust critic loss, numerical log-ratio
clipping, and a basketball-specific KL stopping rule.

\subsection{The learning loop in plain
language}\label{the-learning-loop-in-plain-language}

Imagine coaching a player who must act every 20 milliseconds. At each
tick:

\begin{enumerate}
\def\labelenumi{\arabic{enumi}.}
\tightlist
\item
  The policy receives an observation, such as joint state, orientation,
  command, and recent motion information.
\item
  The actor produces a probability distribution over 14 continuous joint
  commands.
\item
  One action is sampled from that distribution.
\item
  The simulator advances and returns a reward plus boundary flags.
\item
  The critic predicts how much discounted future reward was expected
  from the old observation.
\end{enumerate}

PPO saves a short block of these transitions. It then replays the saved
actions through the updated policy and compares:

\[
\text{old log probability}
\quad\text{with}\quad
\text{new log probability}.
\]

If an action turned out better than the critic expected, PPO tries to
increase its probability. If it turned out worse, PPO tries to decrease
its probability. The word \textbf{proximal} means that the update is
deliberately conservative: PPO limits how much benefit it can claim from
moving the probability ratio too far in the helpful direction.

After several optimizer passes, the rollout is discarded. New data must
be collected from the new policy. PPO is therefore \textbf{on-policy}:
it does not keep an unlimited replay database containing actions from
many old policies.

The two implementations in this repository organize that same idea
differently:

{\def\LTcaptype{none} % do not increment counter
\begin{longtable}[]{@{}
  >{\raggedright\arraybackslash}p{(\linewidth - 6\tabcolsep) * \real{0.2500}}
  >{\raggedright\arraybackslash}p{(\linewidth - 6\tabcolsep) * \real{0.2500}}
  >{\raggedright\arraybackslash}p{(\linewidth - 6\tabcolsep) * \real{0.2500}}
  >{\raggedright\arraybackslash}p{(\linewidth - 6\tabcolsep) * \real{0.2500}}@{}}
\toprule\noalign{}
\begin{minipage}[b]{\linewidth}\raggedright
Path
\end{minipage} & \begin{minipage}[b]{\linewidth}\raggedright
Actor memory
\end{minipage} & \begin{minipage}[b]{\linewidth}\raggedright
Update organization
\end{minipage} & \begin{minipage}[b]{\linewidth}\raggedright
Main numerical framework
\end{minipage} \\
\midrule\noalign{}
\endhead
\bottomrule\noalign{}
\endlastfoot
Generic
\texttt{rlx/\allowbreak{}rlx/\allowbreak{}algorith\allowbreak{}ms/\allowbreak{}ppo.py}
& Feedforward & Flatten time and environments, shuffle, use minibatches
& MLX \\
Basketball
\texttt{ppo\_\allowbreak{}microduc\allowbreak{}k\_\allowbreak{}basketba\allowbreak{}ll.py}
& One-layer LSTM & Replay the entire time-by-environment sequence each
epoch & PyTorch \\
\end{longtable}
}

That difference is not cosmetic. A feedforward sample can be trained as
an independent row. An LSTM action depends on the hidden state produced
by earlier observations, so its history must be reconstructed in the
correct order.

\subsection{Actor, critic, and the meaning of
advantage}\label{actor-critic-and-the-meaning-of-advantage}

PPO trains two prediction systems.

The \textbf{actor} chooses actions. For continuous control, it produces
a mean action \(\mu(o)\) and a standard deviation \(\sigma\). The
basketball actor has 14 action dimensions. Its recurrent core first
processes a 61-dimensional observation with an LSTM of hidden size 256,
then an MLP maps the recurrent output through widths 512, 256, and 128
to 14 means.

The \textbf{critic} predicts a scalar value:

\[
V(o_t)\approx
\mathbb E[r_t+\gamma r_{t+1}+\gamma^2r_{t+2}+\cdots].
\]

The discount \(\gamma\) makes distant rewards count less. With
\(\gamma=0.99\), a reward one step away is multiplied by \(0.99\), two
steps away by \(0.99^2\), and so on.

The actor does not train directly from raw reward. It trains from an
\textbf{advantage estimate}:

\[
\widehat A_t \approx
\text{outcome after action }a_t
-
\text{critic's expectation before }a_t.
\]

\begin{itemize}
\tightlist
\item
  \(\widehat A_t>0\): the sampled action did better than the baseline;
  make it more likely.
\item
  \(\widehat A_t<0\): it did worse than the baseline; make it less
  likely.
\item
  \(\widehat A_t\approx0\): the sample provides little evidence for
  changing its probability.
\end{itemize}

``Positive advantage'' does not mean ``positive reward.'' A reward of 2
can have negative advantage if the critic expected 5. A reward of \(-1\)
can have positive advantage if the critic expected \(-4\). Advantage is
a comparison with a baseline.

The basketball critic is feedforward even though the actor is recurrent.
It reads the current normalized 61-vector and predicts one value. The
actor and critic share the actor's frozen observation normalizer, but
the basketball continuation creates a \textbf{fresh critic and fresh
Adam optimizer}. It is an actor warm start, not a full
optimizer-and-critic resume.

\subsection{Continuous actions use probability
density}\label{continuous-actions-use-probability-density}

For one action coordinate, a Gaussian policy has density

\[
p(a\mid\mu,\sigma)
=
\frac{1}{\sigma\sqrt{2\pi}}
\exp\left(
-\frac{(a-\mu)^2}{2\sigma^2}
\right).
\]

The corresponding log density is

\[
\log p(a\mid\mu,\sigma)
=
-\frac12\left[
\frac{(a-\mu)^2}{\sigma^2}
+2\log\sigma
+\log(2\pi)
\right].
\]

For 14 independent Gaussian coordinates, the implementation sums 14 log
densities:

\[
\log\pi(a\mid o)
=
\sum_{j=1}^{14}\log p(a_j\mid\mu_j(o),\sigma_j).
\]

This produces one joint log probability per transition.

\subsubsection{Density can exceed 1}\label{density-can-exceed-1}

Students often learn that a probability cannot exceed 1. That statement
is true for the probability of an event, but a continuous Gaussian
formula returns a \textbf{density}, not the probability of one exact
point. A narrow density can be taller than 1 as long as the total area
under the curve remains 1.

For \(\mu=0\) and \(\sigma=0.1\), the density at the mean is

\[
p(0)=\frac{1}{0.1\sqrt{2\pi}}\approx3.989.
\]

The probability of landing in an interval is the area over that
interval. For a tiny interval of width \(0.01\) near the mean, the
probability is roughly density times width, about
\(3.989\times0.01=0.03989\), still below 1.

The log density can therefore be positive. That is legal. A positive
continuous log density does not mean a probability greater than 1.

The following complete standard-library program reproduces the
calculation:

\begin{Shaded}
\begin{Highlighting}[]
\ImportTok{import}\NormalTok{ math}


\KeywordTok{def}\NormalTok{ gaussian\_log\_density(x: }\BuiltInTok{float}\NormalTok{, mean: }\BuiltInTok{float}\NormalTok{, std: }\BuiltInTok{float}\NormalTok{) }\OperatorTok{{-}\textgreater{}} \BuiltInTok{float}\NormalTok{:}
    \ControlFlowTok{if}\NormalTok{ std }\OperatorTok{\textless{}=} \FloatTok{0.0}\NormalTok{:}
        \ControlFlowTok{raise} \PreprocessorTok{ValueError}\NormalTok{(}\StringTok{"std must be positive"}\NormalTok{)}
\NormalTok{    variance }\OperatorTok{=}\NormalTok{ std }\OperatorTok{*}\NormalTok{ std}
    \ControlFlowTok{return} \OperatorTok{{-}}\FloatTok{0.5} \OperatorTok{*}\NormalTok{ (}
\NormalTok{        (x }\OperatorTok{{-}}\NormalTok{ mean) }\OperatorTok{**} \DecValTok{2} \OperatorTok{/}\NormalTok{ variance}
        \OperatorTok{+} \FloatTok{2.0} \OperatorTok{*}\NormalTok{ math.log(std)}
        \OperatorTok{+}\NormalTok{ math.log(}\FloatTok{2.0} \OperatorTok{*}\NormalTok{ math.pi)}
\NormalTok{    )}


\KeywordTok{def}\NormalTok{ main() }\OperatorTok{{-}\textgreater{}} \VariableTok{None}\NormalTok{:}
\NormalTok{    log\_density }\OperatorTok{=}\NormalTok{ gaussian\_log\_density(}\FloatTok{0.0}\NormalTok{, }\FloatTok{0.0}\NormalTok{, }\FloatTok{0.1}\NormalTok{)}
\NormalTok{    density }\OperatorTok{=}\NormalTok{ math.exp(log\_density)}
\NormalTok{    approximate\_interval\_probability }\OperatorTok{=}\NormalTok{ density }\OperatorTok{*} \FloatTok{0.01}

    \BuiltInTok{print}\NormalTok{(}\SpecialStringTok{f"log density: }\SpecialCharTok{\{}\NormalTok{log\_density}\SpecialCharTok{:.6f\}}\SpecialStringTok{"}\NormalTok{)}
    \BuiltInTok{print}\NormalTok{(}\SpecialStringTok{f"density: }\SpecialCharTok{\{}\NormalTok{density}\SpecialCharTok{:.6f\}}\SpecialStringTok{"}\NormalTok{)}
    \BuiltInTok{print}\NormalTok{(}\SpecialStringTok{f"approximate probability in width 0.01: "}
          \SpecialStringTok{f"}\SpecialCharTok{\{}\NormalTok{approximate\_interval\_probability}\SpecialCharTok{:.6f\}}\SpecialStringTok{"}\NormalTok{)}

    \ControlFlowTok{assert}\NormalTok{ density }\OperatorTok{\textgreater{}} \FloatTok{1.0}
    \ControlFlowTok{assert} \FloatTok{0.0} \OperatorTok{\textless{}}\NormalTok{ approximate\_interval\_probability }\OperatorTok{\textless{}} \FloatTok{1.0}


\ControlFlowTok{if} \VariableTok{\_\_name\_\_} \OperatorTok{==} \StringTok{"\_\_main\_\_"}\NormalTok{:}
\NormalTok{    main()}
\end{Highlighting}
\end{Shaded}

Expected output, up to final-digit rounding:

\begin{Shaded}
\begin{Highlighting}[]
\NormalTok{log density: 1.383647}
\NormalTok{density: 3.989423}
\NormalTok{approximate probability in width 0.01: 0.039894}
\end{Highlighting}
\end{Shaded}

In
\texttt{ppo\_\allowbreak{}microduc\allowbreak{}k\_\allowbreak{}basketba\allowbreak{}ll.py},
\texttt{gaussian\_log\_prob()} implements the same expression with Torch
tensors and sums over the action dimension. \texttt{gaussian\_entropy()}
also sums over dimensions.

\subsection{From log probabilities to the PPO
ratio}\label{from-log-probabilities-to-the-ppo-ratio}

The stored action was sampled under the old policy. PPO asks how its
density changes under the new policy:

\[
r_t(\theta)
=
\frac{\pi_\theta(a_t\mid o_t)}
{\pi_{\theta_{\mathrm{old}}}(a_t\mid o_t)}.
\]

Direct division of tiny densities is numerically awkward, so code uses
log probabilities:

\[
r_t(\theta)
=
\exp\left(
\log\pi_\theta(a_t\mid o_t)
-
\log\pi_{\theta_{\mathrm{old}}}(a_t\mid o_t)
\right).
\]

The ratio is always positive because an exponential is positive.

\begin{itemize}
\tightlist
\item
  \(r=1\): the saved action has the same density.
\item
  \(r=1.2\): its density is 20\% larger.
\item
  \(r=0.7\): its density is 30\% smaller.
\end{itemize}

The ratio is not ``the probability of the action.'' It is a comparison
of new and old densities for the same recorded action.

The generic MLX code clips the \textbf{log ratio} to
\texttt{{[}-20,\ 20{]}} before exponentiating. This is a numerical
overflow guard. It is separate from PPO's much narrower policy clip such
as \([0.8,1.2]\). The basketball trainer currently exponentiates the
unbounded log-ratio directly, then rejects non-finite KL values or
losses.

\subsection{The clipped policy
objective}\label{the-clipped-policy-objective}

PPO's sample-level surrogate is

\[
L_t^{\mathrm{clip}}
=
\min\left(
r_t\widehat A_t,\;
\operatorname{clip}(r_t,1-\epsilon,1+\epsilon)\widehat A_t
\right).
\]

The algorithm wants to maximize this quantity. The implementation
minimizes its negative. The generic MLX source writes the negative form
as a \texttt{maximum}; the basketball source writes the positive form
with \texttt{minimum} and negates the mean. They are algebraically
equivalent.

With \(\epsilon=0.2\), the clipped interval is \([0.8,1.2]\).

\subsubsection{Positive advantage}\label{positive-advantage}

Suppose \(\widehat A=+2\). Increasing the ratio is helpful because it
makes a good action more likely.

\begin{itemize}
\tightlist
\item
  At \(r=1.1\), both terms equal \(2.2\).
\item
  At \(r=1.3\), the ordinary product is \(2.6\), but the clipped product
  is \(1.2\times2=2.4\). PPO uses \(2.4\).
\item
  At \(r=0.7\), the ordinary product is \(1.4\) and the clipped product
  is \(1.6\). PPO uses the smaller value, \(1.4\).
\end{itemize}

The objective stops rewarding an overly large helpful increase, but it
does not erase a harmful decrease.

\subsubsection{Negative advantage}\label{negative-advantage}

Suppose \(\widehat A=-2\). Decreasing the ratio is helpful because it
makes a bad action less likely.

\begin{itemize}
\tightlist
\item
  At \(r=0.7\), the ordinary product is \(-1.4\), while the clipped
  product is \(0.8\times(-2)=-1.6\). PPO uses the smaller value,
  \(-1.6\).
\item
  At \(r=1.3\), the ordinary product is \(-2.6\) and the clipped product
  is \(-2.4\). PPO uses \(-2.6\).
\end{itemize}

For negative advantage, the helpful direction is downward. The plateau
appears below \(0.8\), not above \(1.2\).

This standalone program prints both cases:

\begin{Shaded}
\begin{Highlighting}[]
\KeywordTok{def}\NormalTok{ clipped\_surrogate(ratio: }\BuiltInTok{float}\NormalTok{, advantage: }\BuiltInTok{float}\NormalTok{, epsilon: }\BuiltInTok{float} \OperatorTok{=} \FloatTok{0.2}\NormalTok{) }\OperatorTok{{-}\textgreater{}} \BuiltInTok{float}\NormalTok{:}
\NormalTok{    clipped\_ratio }\OperatorTok{=} \BuiltInTok{min}\NormalTok{(}\BuiltInTok{max}\NormalTok{(ratio, }\FloatTok{1.0} \OperatorTok{{-}}\NormalTok{ epsilon), }\FloatTok{1.0} \OperatorTok{+}\NormalTok{ epsilon)}
    \ControlFlowTok{return} \BuiltInTok{min}\NormalTok{(ratio }\OperatorTok{*}\NormalTok{ advantage, clipped\_ratio }\OperatorTok{*}\NormalTok{ advantage)}


\KeywordTok{def}\NormalTok{ main() }\OperatorTok{{-}\textgreater{}} \VariableTok{None}\NormalTok{:}
    \ControlFlowTok{for}\NormalTok{ advantage }\KeywordTok{in}\NormalTok{ (}\FloatTok{2.0}\NormalTok{, }\OperatorTok{{-}}\FloatTok{2.0}\NormalTok{):}
        \BuiltInTok{print}\NormalTok{(}\SpecialStringTok{f"advantage = }\SpecialCharTok{\{}\NormalTok{advantage}\SpecialCharTok{:+.1f\}}\SpecialStringTok{"}\NormalTok{)}
        \ControlFlowTok{for}\NormalTok{ ratio }\KeywordTok{in}\NormalTok{ (}\FloatTok{0.7}\NormalTok{, }\FloatTok{1.0}\NormalTok{, }\FloatTok{1.3}\NormalTok{):}
\NormalTok{            value }\OperatorTok{=}\NormalTok{ clipped\_surrogate(ratio, advantage)}
            \BuiltInTok{print}\NormalTok{(}\SpecialStringTok{f"  ratio=}\SpecialCharTok{\{}\NormalTok{ratio}\SpecialCharTok{:.1f\}}\SpecialStringTok{, surrogate=}\SpecialCharTok{\{}\NormalTok{value}\SpecialCharTok{:+.1f\}}\SpecialStringTok{"}\NormalTok{)}

    \ControlFlowTok{assert}\NormalTok{ clipped\_surrogate(}\FloatTok{1.3}\NormalTok{, }\FloatTok{2.0}\NormalTok{) }\OperatorTok{==} \FloatTok{2.4}
    \ControlFlowTok{assert}\NormalTok{ clipped\_surrogate(}\FloatTok{0.7}\NormalTok{, }\OperatorTok{{-}}\FloatTok{2.0}\NormalTok{) }\OperatorTok{==} \OperatorTok{{-}}\FloatTok{1.6}


\ControlFlowTok{if} \VariableTok{\_\_name\_\_} \OperatorTok{==} \StringTok{"\_\_main\_\_"}\NormalTok{:}
\NormalTok{    main()}
\end{Highlighting}
\end{Shaded}

Clipping is not a guarantee that the new policy is better. It only
changes the training objective for the sampled batch. The network can
still change on observations absent from that batch, multiple parameters
interact, the critic can be wrong, and the reward can fail to represent
the intended basketball skill.

\begin{figure}
\centering
\includegraphics[width=6.2in,height=5.5in,alt={Schematic PPO clipping curves. This explanatory diagram is not measured training evidence.}]{assets/diagrams/schematic-ppo-clipping.png}
\caption{Schematic PPO clipping curves. This explanatory diagram is not
measured training evidence.}
\end{figure}

\subsection{Generalized Advantage
Estimation}\label{generalized-advantage-estimation}

The one-step temporal-difference residual is

\[
\delta_t
=
r_t+\gamma b_t-V(o_t),
\]

where \(b_t\) is the permitted bootstrap value after the transition.
Generalized Advantage Estimation, or GAE, combines residuals backward:

\[
\widehat A_t
=
\delta_t+\gamma\lambda m_t\widehat A_{t+1}.
\]

Here:

\begin{itemize}
\tightlist
\item
  \(\gamma\) discounts future return;
\item
  \(\lambda\) controls how strongly later residuals affect the current
  estimate;
\item
  \(m_t\) is a continuation mask that prevents a trace from crossing an
  episode reset.
\end{itemize}

With the repository defaults \(\gamma=0.99\) and \(\lambda=0.95\), the
trace multiplier is

\[
\gamma\lambda=0.9405.
\]

A smaller \(\lambda\) relies more on the critic's one-step predictions.
A larger \(\lambda\) carries more sampled future information backward,
usually with more sampling variance. Neither value is universally best.

The critic target is

\[
\widehat R_t=\widehat A_t+V(o_t).
\]

The actor uses \(\widehat A_t\); the critic fits \(\widehat R_t\).

\subsection{Terminal and time-limit masks are
different}\label{terminal-and-time-limit-masks-are-different}

Gymnasium separates two boundaries:

\begin{itemize}
\tightlist
\item
  \textbf{termination}: the modeled task truly ended, such as a failure
  condition;
\item
  \textbf{truncation}: collection stopped for an external reason,
  commonly a time limit, although the underlying task could continue.
\end{itemize}

Both boundaries reset the environment. They do not receive the same
value bootstrap.

For transition \(t\), define:

\[
b_t=
\begin{cases}
0,&\text{if terminated},\\
V(o_{t+1}^{\mathrm{final}}),&\text{if truncated but not terminated},\\
V(o_{t+1}),&\text{if the episode continues}.
\end{cases}
\]

The GAE trace mask is zero for either termination or truncation:

\[
m_t=1-\mathbf 1[\text{terminated}\lor\text{truncated}].
\]

This gives two separate decisions:

\begin{enumerate}
\def\labelenumi{\arabic{enumi}.}
\tightlist
\item
  \textbf{May the current residual bootstrap a value?} A pure timeout
  may; a true terminal may not.
\item
  \textbf{May advantage flow into the next stored transition?} Neither
  boundary may, because the next stored observation can belong to a
  reset episode.
\end{enumerate}

The basketball implementation stores
\texttt{bootstra\allowbreak{}p\_\allowbreak{}values} for every
transition. During collection, it evaluates the critic on the final
pre-reset observation, then explicitly replaces the value with zero
where \texttt{terminated} is true. Its \texttt{compute\_gae()} uses:

\begin{itemize}
\tightlist
\item
  \texttt{\textasciitilde{}terminated} as the bootstrap mask;
\item
  \texttt{\textasciitilde{}(terminated\ \textbar{}\ truncated)} as the
  trace-continuation mask.
\end{itemize}

The generic MLX implementation reaches the same semantics differently.
\texttt{\_truncation\_values()} selects only
\texttt{truncated\ \&\ \textasciitilde{}terminated} environments and
evaluates \texttt{info{[}"terminal\_observation"{]}}. The shared GAE
helper uses the timeout value for a truncation, zero for termination,
and the ordinary next value otherwise. It never bootstraps from an
autoreset observation.

\subsubsection{Worked boundary example}\label{worked-boundary-example}

Suppose:

\[
r_t=1,\quad V(o_t)=2,\quad
V(o_{t+1}^{\mathrm{final}})=5,\quad \gamma=0.99.
\]

If the transition is a true terminal:

\[
\delta_t=1+0-2=-1.
\]

If it is only a time limit:

\[
\delta_t=1+0.99(5)-2=3.95.
\]

If an autoreset observation has value 99, using it would give

\[
1+0.99(99)-2=97.01,
\]

which incorrectly assigns value from the next episode to the previous
one.

The next program computes a three-step trace and proves that a reward in
the reset episode cannot leak backward:

\begin{Shaded}
\begin{Highlighting}[]
\ImportTok{from}\NormalTok{ typing }\ImportTok{import}\NormalTok{ Sequence}


\KeywordTok{def}\NormalTok{ gae(}
\NormalTok{    rewards: Sequence[}\BuiltInTok{float}\NormalTok{],}
\NormalTok{    values: Sequence[}\BuiltInTok{float}\NormalTok{],}
\NormalTok{    bootstraps: Sequence[}\BuiltInTok{float}\NormalTok{],}
\NormalTok{    terminated: Sequence[}\BuiltInTok{bool}\NormalTok{],}
\NormalTok{    truncated: Sequence[}\BuiltInTok{bool}\NormalTok{],}
\NormalTok{    gamma: }\BuiltInTok{float} \OperatorTok{=} \FloatTok{0.99}\NormalTok{,}
\NormalTok{    gae\_lambda: }\BuiltInTok{float} \OperatorTok{=} \FloatTok{0.95}\NormalTok{,}
\NormalTok{) }\OperatorTok{{-}\textgreater{}} \BuiltInTok{tuple}\NormalTok{[}\BuiltInTok{list}\NormalTok{[}\BuiltInTok{float}\NormalTok{], }\BuiltInTok{list}\NormalTok{[}\BuiltInTok{float}\NormalTok{]]:}
\NormalTok{    size }\OperatorTok{=} \BuiltInTok{len}\NormalTok{(rewards)}
    \ControlFlowTok{if} \KeywordTok{not} \BuiltInTok{all}\NormalTok{(}\BuiltInTok{len}\NormalTok{(x) }\OperatorTok{==}\NormalTok{ size }\ControlFlowTok{for}\NormalTok{ x }\KeywordTok{in}\NormalTok{ (}
\NormalTok{        values, bootstraps, terminated, truncated}
\NormalTok{    )):}
        \ControlFlowTok{raise} \PreprocessorTok{ValueError}\NormalTok{(}\StringTok{"all inputs must have equal length"}\NormalTok{)}

\NormalTok{    advantages }\OperatorTok{=}\NormalTok{ [}\FloatTok{0.0}\NormalTok{] }\OperatorTok{*}\NormalTok{ size}
\NormalTok{    trace }\OperatorTok{=} \FloatTok{0.0}
    \ControlFlowTok{for}\NormalTok{ step }\KeywordTok{in} \BuiltInTok{range}\NormalTok{(size }\OperatorTok{{-}} \DecValTok{1}\NormalTok{, }\OperatorTok{{-}}\DecValTok{1}\NormalTok{, }\OperatorTok{{-}}\DecValTok{1}\NormalTok{):}
\NormalTok{        bootstrap\_mask }\OperatorTok{=} \FloatTok{0.0} \ControlFlowTok{if}\NormalTok{ terminated[step] }\ControlFlowTok{else} \FloatTok{1.0}
\NormalTok{        continuation\_mask }\OperatorTok{=}\NormalTok{ (}
            \FloatTok{0.0} \ControlFlowTok{if}\NormalTok{ terminated[step] }\KeywordTok{or}\NormalTok{ truncated[step] }\ControlFlowTok{else} \FloatTok{1.0}
\NormalTok{        )}
\NormalTok{        delta }\OperatorTok{=}\NormalTok{ (}
\NormalTok{            rewards[step]}
            \OperatorTok{+}\NormalTok{ gamma }\OperatorTok{*}\NormalTok{ bootstraps[step] }\OperatorTok{*}\NormalTok{ bootstrap\_mask}
            \OperatorTok{{-}}\NormalTok{ values[step]}
\NormalTok{        )}
\NormalTok{        trace }\OperatorTok{=}\NormalTok{ delta }\OperatorTok{+}\NormalTok{ gamma }\OperatorTok{*}\NormalTok{ gae\_lambda }\OperatorTok{*}\NormalTok{ continuation\_mask }\OperatorTok{*}\NormalTok{ trace}
\NormalTok{        advantages[step] }\OperatorTok{=}\NormalTok{ trace}

\NormalTok{    returns }\OperatorTok{=}\NormalTok{ [a }\OperatorTok{+}\NormalTok{ v }\ControlFlowTok{for}\NormalTok{ a, v }\KeywordTok{in} \BuiltInTok{zip}\NormalTok{(advantages, values)]}
    \ControlFlowTok{return}\NormalTok{ advantages, returns}


\KeywordTok{def}\NormalTok{ main() }\OperatorTok{{-}\textgreater{}} \VariableTok{None}\NormalTok{:}
\NormalTok{    advantages, returns }\OperatorTok{=}\NormalTok{ gae(}
\NormalTok{        rewards}\OperatorTok{=}\NormalTok{[}\FloatTok{0.5}\NormalTok{, }\FloatTok{0.2}\NormalTok{, }\FloatTok{0.0}\NormalTok{],}
\NormalTok{        values}\OperatorTok{=}\NormalTok{[}\FloatTok{1.0}\NormalTok{, }\FloatTok{1.2}\NormalTok{, }\FloatTok{1.1}\NormalTok{],}
\NormalTok{        bootstraps}\OperatorTok{=}\NormalTok{[}\FloatTok{1.2}\NormalTok{, }\FloatTok{1.1}\NormalTok{, }\FloatTok{0.9}\NormalTok{],}
\NormalTok{        terminated}\OperatorTok{=}\NormalTok{[}\VariableTok{False}\NormalTok{, }\VariableTok{False}\NormalTok{, }\VariableTok{False}\NormalTok{],}
\NormalTok{        truncated}\OperatorTok{=}\NormalTok{[}\VariableTok{False}\NormalTok{, }\VariableTok{False}\NormalTok{, }\VariableTok{True}\NormalTok{],}
\NormalTok{    )}
    \BuiltInTok{print}\NormalTok{([}\BuiltInTok{round}\NormalTok{(x, }\DecValTok{6}\NormalTok{) }\ControlFlowTok{for}\NormalTok{ x }\KeywordTok{in}\NormalTok{ advantages])}
    \BuiltInTok{print}\NormalTok{([}\BuiltInTok{round}\NormalTok{(x, }\DecValTok{6}\NormalTok{) }\ControlFlowTok{for}\NormalTok{ x }\KeywordTok{in}\NormalTok{ returns])}

\NormalTok{    changed\_next\_episode\_reward }\OperatorTok{=} \FloatTok{1000.0}
    \ControlFlowTok{assert}\NormalTok{ changed\_next\_episode\_reward }\OperatorTok{\textgreater{}} \FloatTok{0.0}  \CommentTok{\# It is deliberately unused.}
    \ControlFlowTok{assert} \BuiltInTok{round}\NormalTok{(advantages[}\OperatorTok{{-}}\DecValTok{1}\NormalTok{], }\DecValTok{6}\NormalTok{) }\OperatorTok{==} \OperatorTok{{-}}\FloatTok{0.209}
    \ControlFlowTok{assert}\NormalTok{ [}\BuiltInTok{round}\NormalTok{(x, }\DecValTok{6}\NormalTok{) }\ControlFlowTok{for}\NormalTok{ x }\KeywordTok{in}\NormalTok{ advantages] }\OperatorTok{==}\NormalTok{ [}
        \FloatTok{0.586836}\NormalTok{, }\OperatorTok{{-}}\FloatTok{0.107565}\NormalTok{, }\OperatorTok{{-}}\FloatTok{0.209}
\NormalTok{    ]}


\ControlFlowTok{if} \VariableTok{\_\_name\_\_} \OperatorTok{==} \StringTok{"\_\_main\_\_"}\NormalTok{:}
\NormalTok{    main()}
\end{Highlighting}
\end{Shaded}

The expected advantages are:

\begin{Shaded}
\begin{Highlighting}[]
\NormalTok{[0.586836, {-}0.107565, {-}0.209]}
\end{Highlighting}
\end{Shaded}

The timeout transition bootstraps from 0.9, but its continuation mask is
zero, so no later reset-episode advantage enters the trace.

\subsection{Advantage normalization}\label{advantage-normalization}

Raw advantages can have very different scales across tasks or batches.
Both implementations normalize advantages:

\[
\widehat A'_t
=
\frac{\widehat A_t-\operatorname{mean}(\widehat A)}
{\operatorname{std}(\widehat A)+10^{-8}}.
\]

The basketball trainer normalizes once over the complete rollout before
its epoch loop. The generic MLX trainer normalizes inside
\texttt{\_loss\_and\_metrics()} after selecting a minibatch, so each
actual minibatch uses its own mean and standard deviation.

Normalization changes scale and can change a sample's sign relative to
the selected group's mean. It does not change rewards, returns, or old
log probabilities. It also does not mean that an advantage of
\texttt{+2} in a teaching table will literally remain \texttt{+2} in
training.

\subsection{The critic loss}\label{the-critic-loss}

The basketball continuation uses ordinary half squared error:

\[
L_V
=
\frac12\operatorname{mean}\left[
(V_{\mathrm{new}}-\widehat R)^2
\right].
\]

It does not clip the value update.

The generic MLX PPO uses a Huber-style error with threshold 10:

\[
h(e)=
\begin{cases}
\frac12e^2,&|e|\le10,\\
10|e|-50,&|e|>10.
\end{cases}
\]

When \texttt{clip\_value\_loss=True}, it also creates a clipped
prediction:

\[
V_{\mathrm{clip}}
=
V_{\mathrm{old}}+
\operatorname{clip}(V_{\mathrm{new}}-V_{\mathrm{old}},-\epsilon,\epsilon),
\]

then uses the larger of the unclipped and clipped Huber errors. Choosing
the larger prevents the clipped branch from making an excessively
changed value prediction look artificially good.

This is an important implementation difference. ``RLX PPO uses value
clipping'' is true for
\texttt{rlx/\allowbreak{}rlx/\allowbreak{}algorith\allowbreak{}ms/\allowbreak{}ppo.py},
but false for the basketball-specific \texttt{ppo\_update()}.

\subsection{Entropy and why it can be
negative}\label{entropy-and-why-it-can-be-negative}

Entropy encourages exploration. For one Gaussian coordinate,

\[
H
=
\log\sigma+\frac12\log(2\pi e).
\]

For independent coordinates, the implementation sums this value.

This is \textbf{differential entropy}, not discrete Shannon entropy.
Differential entropy may be negative. For example, with \(\sigma=0.03\):

\[
H_{\text{one dimension}}\approx-2.087,
\]

so 14 equal dimensions have total entropy near \(-29.2\).

Nothing is wrong with a negative value. It means the continuous density
is concentrated in a narrow region. It does not mean negative
probability.

Both loss functions subtract the entropy bonus:

\[
L_{\mathrm{total}}
=
L_{\mathrm{policy}}
+c_VL_V
-c_HH.
\]

If \(H\) is negative, the term \(-c_HH\) is positive. The optimizer
still follows the exact differential-entropy objective; the sign should
not be interpreted with discrete-entropy intuition.

The basketball actor stores a directly trained standard-deviation
parameter and clamps it only to a minimum of \texttt{1e-6} when used.
Loading a source actor then fills this parameter with the requested
\texttt{initial\_std}, whose command-line default is \texttt{0.03}. Thus
the continuation deliberately resets exploration scale instead of
retaining the source checkpoint's saved standard deviation.

\subsection{Total loss is not reward, and neither proves
mastery}\label{total-loss-is-not-reward-and-neither-proves-mastery}

The basketball total loss is

\[
L
=
L_{\mathrm{policy}}
+1.0L_V
-0.01H
\quad\text{under its defaults}.
\]

The generic PPO defaults use value coefficient \texttt{0.5} and entropy
coefficient \texttt{0.01}, although example programs can override them.

A falling training loss is not the same as a rising episode reward:

\begin{itemize}
\tightlist
\item
  policy loss measures a clipped probability objective on the latest
  batch;
\item
  value loss measures critic prediction error;
\item
  entropy measures distribution spread;
\item
  reward comes from the environment;
\item
  mastery is an external behavioral claim.
\end{itemize}

The total loss can rise while behavior improves because the critic is
fitting larger return targets. It can fall while the robot learns a
reward loophole. A high reward can still represent holding the ball
incorrectly, exploiting assistance, or failing after the logged horizon.
PPO only optimizes the reward and data it receives.

Therefore a basketball policy requires independent evaluation: complete
episodes, physical task metrics, deterministic exported-policy tests,
and rendered inspection. No scalar PPO metric alone establishes mastery.

\subsection{Recurrent sequence replay in
basketball}\label{recurrent-sequence-replay-in-basketball}

The basketball actor is an LSTM. Its action mean at time \(t\) depends
on:

\[
\mu_t=f(o_t,h_{t-1},c_{t-1}).
\]

If training shuffled individual time steps, the hidden state would no
longer match the one that generated the action. The basketball rollout
therefore stores:

\begin{itemize}
\tightlist
\item
  observations with shape \texttt{{[}time,\ environments,\ 61{]}};
\item
  sampled actions and old log probabilities;
\item
  old critic values, rewards, terminations, truncations, and bootstrap
  values;
\item
  \texttt{episode\_\allowbreak{}starts} for every time step;
\item
  the initial LSTM hidden and cell states for the rollout.
\end{itemize}

During each PPO epoch, \texttt{actor.forward\_sequence()} starts from
the saved initial hidden state and walks forward in time. Before
processing a step whose \texttt{episode\_\allowbreak{}starts} flag is
true, it multiplies the hidden and cell state by zero for that
environment. This reconstructs the recurrent state while preventing
memory from crossing resets.

The update uses the full sequence each epoch. There is no basketball
minibatch loop in \texttt{ppo\_update()}. With defaults of 32 steps and
4 environments, one rollout contains:

\[
32\times4=128\text{ transitions}.
\]

Five epochs present those same 128 transitions to the optimizer up to
five times, unless KL stopping ends earlier. This is 128 new environment
transitions, not 640 new transitions.

The critic is evaluated on a flattened
\texttt{{[}time*environment,\ 61{]}} tensor because it is feedforward.
Its outputs are reshaped back to \texttt{{[}time,\ environment{]}}.

\subsection{Feedforward minibatches in generic MLX
PPO}\label{feedforward-minibatches-in-generic-mlx-ppo}

\texttt{RolloutBuffer.reset()} allocates fixed arrays shaped
\texttt{{[}num\_steps,\ num\_envs,\ ...{]}}. It stores:

{\def\LTcaptype{none} % do not increment counter
\begin{longtable}[]{@{}
  >{\raggedright\arraybackslash}p{(\linewidth - 2\tabcolsep) * \real{0.5000}}
  >{\raggedright\arraybackslash}p{(\linewidth - 2\tabcolsep) * \real{0.5000}}@{}}
\toprule\noalign{}
\begin{minipage}[b]{\linewidth}\raggedright
Field
\end{minipage} & \begin{minipage}[b]{\linewidth}\raggedright
Purpose
\end{minipage} \\
\midrule\noalign{}
\endhead
\bottomrule\noalign{}
\endlastfoot
\texttt{observations} & input used to choose each action \\
\texttt{next\_\allowbreak{}observat\allowbreak{}ions} & returned next
input \\
\texttt{actions} & sampled action \\
\texttt{rewards} & scalar reward \\
\texttt{terminations}, \texttt{truncations} & separate boundary types \\
\texttt{values} & old critic predictions \\
\texttt{log\_probs} & old-policy joint log probabilities \\
\texttt{advantages}, \texttt{returns} & allocated fields, although the
current PPO path passes computed arrays directly to \texttt{update()} \\
\end{longtable}
}

The generic update flattens the first two dimensions. For the library
defaults:

\[
16\text{ steps}\times4096\text{ environments}
=65{,}536\text{ samples}.
\]

With 16 minibatches, each minibatch has 4,096 samples. Two epochs
produce 32 optimizer steps per rollout.

The generic Microduck example does \textbf{not} use those large library
defaults.
\texttt{rlx/\allowbreak{}examples\allowbreak{}/\allowbreak{}ppo\_\allowbreak{}microduc\allowbreak{}k.py}
defines 16 environments, 24 steps, 4 minibatches, and 5 epochs, then
passes those values into \texttt{PPOConfig}. Its rollout has 384
transitions, each minibatch has 96, and one rollout produces 20
optimizer steps.

The MLX update creates a fresh random permutation each epoch. This is
appropriate for its feedforward network because every saved observation
contains the complete input expected by the network.
\texttt{RolloutB\allowbreak{}uffer} stores no recurrent hidden state and
no sequence-start mask; using it unchanged for an LSTM would lose
necessary history.

\subsection{KL divergence: diagnostic in MLX, stopping rule in
basketball}\label{kl-divergence-diagnostic-in-mlx-stopping-rule-in-basketball}

Both implementations calculate

\[
\widehat D_{\mathrm{KL}}
=
\operatorname{mean}\left[(r-1)-\log r\right].
\]

Because samples come from the old policy, this is a common sample
approximation related to
\(D_{\mathrm{KL}}(\pi_{\mathrm{old}}\|\pi_{\mathrm{new}})\). It should
be nonnegative in expectation, although a finite sample and
floating-point arithmetic can produce small irregularities.

The implementations use it differently.

\subsubsection{Generic MLX behavior}\label{generic-mlx-behavior}

\texttt{rlx/\allowbreak{}rlx/\allowbreak{}algorith\allowbreak{}ms/\allowbreak{}ppo.py}:

\begin{itemize}
\tightlist
\item
  clips log ratios to \texttt{{[}-20,20{]}} before exponentiation;
\item
  reports \texttt{approxim\allowbreak{}ate\_\allowbreak{}kl} when
  metrics are requested;
\item
  reports \texttt{clip\_\allowbreak{}fraction\allowbreak{}};
\item
  does \textbf{not} define \texttt{target\_kl};
\item
  does \textbf{not} stop epochs or minibatches based on KL.
\end{itemize}

Its KL is a diagnostic, not a gate.

\subsubsection{Basketball behavior}\label{basketball-behavior}

\texttt{ppo\_\allowbreak{}microduc\allowbreak{}k\_\allowbreak{}basketba\allowbreak{}ll.py}:

\begin{itemize}
\tightlist
\item
  defaults \texttt{target\_kl} to \texttt{0.02};
\item
  computes KL before each epoch's optimizer step;
\item
  can stop with \texttt{epochs\_completed\ ==\ 0} if saved old log
  probabilities are already inconsistent with the current actor;
\item
  recomputes the full recurrent sequence after each optimizer step;
\item
  computes post-update KL and stops further epochs when it exceeds the
  target;
\item
  records \texttt{kl\_\allowbreak{}early\_\allowbreak{}stop}.
\end{itemize}

The pre-update guard matters for stale or mismatched rollout data. The
post-update guard limits repeated reuse of the same short recurrent
rollout. It is still not a mathematical guarantee that every state
changed by less than a fixed amount.

\subsection{Exact defaults and
overrides}\label{exact-defaults-and-overrides}

Do not describe ``the PPO defaults'' without naming the layer.

\subsubsection{\texorpdfstring{Generic \texttt{PPOConfig}
defaults}{Generic PPOConfig defaults}}\label{generic-ppoconfig-defaults}

{\def\LTcaptype{none} % do not increment counter
\begin{longtable}[]{@{}lr@{}}
\toprule\noalign{}
Setting & Value \\
\midrule\noalign{}
\endhead
\bottomrule\noalign{}
\endlastfoot
\texttt{num\_envs} & 4096 \\
\texttt{num\_steps} & 16 \\
\texttt{gamma} & 0.99 \\
\texttt{gae\_lambda} & 0.95 \\
\texttt{num\_\allowbreak{}minibatc\allowbreak{}hes} & 16 \\
\texttt{update\_\allowbreak{}epochs} & 2 \\
\texttt{normaliz\allowbreak{}e\_\allowbreak{}advantag\allowbreak{}es} &
true \\
\texttt{clip\_\allowbreak{}coeffici\allowbreak{}ent} & 0.2 \\
\texttt{clip\_\allowbreak{}value\_\allowbreak{}loss} & true \\
\texttt{entropy\_\allowbreak{}coeffici\allowbreak{}ent} & 0.01 \\
\texttt{value\_\allowbreak{}coeffici\allowbreak{}ent} & 0.5 \\
\texttt{max\_\allowbreak{}grad\_\allowbreak{}norm} & 0.5 \\
\end{longtable}
}

\subsubsection{Generic Microduck example
overrides}\label{generic-microduck-example-overrides}

\texttt{rlx/\allowbreak{}examples\allowbreak{}/\allowbreak{}ppo\_\allowbreak{}microduc\allowbreak{}k.py}
replaces \texttt{num\_envs}, \texttt{num\_steps},
\texttt{num\_\allowbreak{}minibatc\allowbreak{}hes}, and
\texttt{update\_\allowbreak{}epochs} with \texttt{16}, \texttt{24},
\texttt{4}, and \texttt{5}. It keeps gamma \texttt{0.99}, lambda
\texttt{0.95}, clip \texttt{0.2}, and entropy coefficient \texttt{0.01}.
Its optimizer learning-rate default is \texttt{1e-3}, which belongs to
the example arguments rather than \texttt{PPOConfig}.

\subsubsection{Basketball trainer
defaults}\label{basketball-trainer-defaults}

The basketball trainer does not instantiate \texttt{PPOConfig}. Its
command line and function defaults produce:

{\def\LTcaptype{none} % do not increment counter
\begin{longtable}[]{@{}lr@{}}
\toprule\noalign{}
Setting & Value \\
\midrule\noalign{}
\endhead
\bottomrule\noalign{}
\endlastfoot
environments & 4 \\
rollout steps & 32 \\
updates & 5 \\
PPO epochs per rollout & 5 \\
gamma & 0.99 \\
GAE lambda & 0.95 \\
clip coefficient & 0.2 \\
value coefficient & 1.0 \\
entropy coefficient & 0.01 \\
gradient norm limit & 1.0 \\
Adam learning rate & \texttt{2e-5} \\
target KL & 0.02 \\
reset exploration standard deviation & 0.03 \\
\end{longtable}
}

Its environment construction also overrides training conditions
explicitly:

\begin{itemize}
\tightlist
\item
  episode limit: 10 seconds;
\item
  actuator: \texttt{"bam"} by default;
\item
  observation noise: off;
\item
  domain randomization: off;
\item
  pushes: off;
\item
  curriculum: off;
\item
  hold assistance: \texttt{0.0};
\item
  command: fixed \texttt{(0.08,\ 0.0,\ 0.0)} unless
  \texttt{-\allowbreak{}-\allowbreak{}randomiz\allowbreak{}ed-\allowbreak{}commands\allowbreak{}}
  is selected.
\end{itemize}

These choices describe the default local continuation experiment. They
are not universal PPO recommendations and do not reproduce the complete
upstream training recipe.

\subsection{One complete update, end to
end}\label{one-complete-update-end-to-end}

For the default basketball run:

\begin{enumerate}
\def\labelenumi{\arabic{enumi}.}
\tightlist
\item
  Reset four environments with deterministic seed offsets.
\item
  Initialize LSTM hidden and cell tensors with shape
  \texttt{{[}1,\ 4,\ 256{]}}.
\item
  Collect 32 control steps, giving 128 transitions.
\item
  At every step, zero recurrent state for environments whose current
  observation starts a new episode.
\item
  Sample 14-dimensional Gaussian actions and save their joint old log
  probabilities.
\item
  Step each environment independently.
\item
  Evaluate final pre-reset observations for bootstrap values; force
  terminal bootstraps to zero.
\item
  Reset ended environments and mark the next observations as episode
  starts.
\item
  Run backward GAE with termination and truncation masks.
\item
  Normalize the 128 advantages together.
\item
  Replay the entire actor sequence from the saved initial LSTM state.
\item
  Check pre-update approximate KL.
\item
  Compute clipped policy loss, squared critic loss, and Gaussian
  entropy.
\item
  Backpropagate, clip the combined actor-and-critic gradient norm to
  1.0, and step Adam.
\item
  Replay the sequence again, measure post-update KL, and decide whether
  another epoch is allowed.
\end{enumerate}

The command-line defaults request five such rollout/update cycles:

\[
4\times32\times5=640
\text{ local environment transitions}.
\]

That small default is suitable as a short continuation or smoke-scale
run, not evidence of basketball mastery.

\subsection{Common interpretation
mistakes}\label{common-interpretation-mistakes}

\textbf{Mistake: ``The Gaussian returned probability 3.9, which is
impossible.''}\\
Correction: it returned density. Probability is area over an interval.

\textbf{Mistake: ``A positive log probability is invalid.''}\\
Correction: continuous log density may be positive when the density
exceeds 1.

\textbf{Mistake: ``Clipping keeps every ratio between 0.8 and 1.2.''}\\
Correction: the objective uses a clipped comparison. The actual ratio
can lie outside the interval.

\textbf{Mistake: ``A timeout gets zero future value because the
environment reset.''}\\
Correction: a pure time limit bootstraps from the final pre-reset
observation, then stops the trace.

\textbf{Mistake: ``Negative entropy means negative probability.''}\\
Correction: continuous differential entropy can be negative.

\textbf{Mistake: ``Lower total loss means higher reward.''}\\
Correction: total loss mixes policy, critic, and entropy terms on one
batch.

\textbf{Mistake: ``High reward means the skill is mastered.''}\\
Correction: reward is a designed training signal. Mastery requires
independent physical criteria and inspection.

\textbf{Mistake: ``The basketball updater uses MLX minibatches.''}\\
Correction: it is a separate PyTorch recurrent updater that replays full
sequences.

\textbf{Mistake: ``Both PPO implementations stop at KL 0.02.''}\\
Correction: only the basketball path has that target and early stopping.
Generic MLX PPO reports KL without stopping.

\subsection{Exercises}\label{exercises}

\begin{enumerate}
\def\labelenumi{\arabic{enumi}.}
\tightlist
\item
  A Gaussian has mean 0 and standard deviation 0.2. Compute its density
  at the mean. Can the answer exceed 1?
\item
  An action has old log density \texttt{-3.0} and new log density
  \texttt{-2.8}. Compute the PPO ratio.
\item
  With advantage \texttt{+3}, ratio \texttt{1.4}, and clip coefficient
  \texttt{0.2}, compute the unclipped product, clipped product, and
  selected surrogate.
\item
  Repeat Exercise 3 with advantage \texttt{-3}.
\item
  A transition has reward 2, current value 4, final-observation value 5,
  and gamma 0.9. Compute its TD residual for a true terminal and for a
  pure timeout.
\item
  Why must an LSTM rollout save its initial hidden state and
  episode-start masks?
\item
  Under generic \texttt{PPOConfig} defaults, how many optimizer steps
  occur per rollout?
\item
  Under basketball defaults, how many new transitions are collected in
  one rollout, and how many times can they be presented to the
  optimizer?
\item
  A training log shows total loss falling and reward rising. Give two
  reasons this still does not prove basketball mastery.
\item
  State the KL behavior difference between the generic MLX and
  basketball implementations in one sentence.
\end{enumerate}

\subsection{Answers}\label{answers}

\begin{enumerate}
\def\labelenumi{\arabic{enumi}.}
\tightlist
\item
  At the mean, density is \(1/(0.2\sqrt{2\pi})\approx1.995\). Yes. It is
  density, and its total area is still 1.
\item
  \(r=\exp(-2.8-(-3.0))=\exp(0.2)\approx1.2214\).
\item
  Unclipped: \(1.4\times3=4.2\). Clipped ratio: 1.2. Clipped product:
  3.6. PPO selects the smaller surrogate, 3.6.
\item
  Unclipped: \(1.4\times(-3)=-4.2\). Clipped product:
  \(1.2\times(-3)=-3.6\). PPO selects \(-4.2\) because it is smaller.
  Clipping does not protect movement in the harmful direction.
\item
  Terminal: \(2-4=-2\). Timeout: \(2+0.9(5)-4=2.5\).
\item
  The same observation can produce a different action mean under a
  different hidden state. Sequence replay must reconstruct the state
  that corresponds to the stored action, and reset masks must prevent
  memory from crossing episode boundaries.
\item
  \texttt{update\_epochs\ ×\ num\_minibatches\ =\ 2\ ×\ 16\ =\ 32}
  optimizer steps.
\item
  \texttt{4\ ×\ 32\ =\ 128} new transitions. Up to five full-sequence
  optimizer presentations occur, unless KL early stopping ends the loop
  sooner.
\item
  The reward may omit important physical requirements or contain an
  exploitable shortcut; evaluation may cover only short or assisted
  episodes. In addition, loss and reward are not the same quantity.
\item
  Generic MLX PPO reports approximate KL as a metric but never stops on
  it, while basketball checks KL before and after each full-sequence
  epoch and stops further epochs when it exceeds the default target
  \texttt{0.02}.
\end{enumerate}

\subsection{References}\label{references}

\begin{itemize}
\tightlist
\item
  John Schulman, Filip Wolski, Prafulla Dhariwal, Alec Radford, and Oleg
  Klimov. \textbf{Proximal Policy Optimization Algorithms}.
  arXiv:1707.06347, 2017. \url{https://arxiv.org/abs/1707.06347}
\item
  John Schulman, Philipp Moritz, Sergey Levine, Michael Jordan, and
  Pieter Abbeel. \textbf{High-Dimensional Continuous Control Using
  Generalized Advantage Estimation}. arXiv:1506.02438, 2015.
  \url{https://arxiv.org/abs/1506.02438}
\end{itemize}

The paper metadata and abstracts were checked against locally retrieved
arXiv pages. The repository implementation statements in this chapter
were checked separately against the local source files named at the
beginning.

\section{Basketball: Blind Recurrent Balance on a Free
Ball}\label{basketball-blind-recurrent-balance-on-a-free-ball}

\subsection{Learning objectives}\label{learning-objectives}

After this chapter, you should be able to:

\begin{itemize}
\tightlist
\item
  reconstruct the basketball scene from the numbers in the repository;
\item
  explain every slice of the shared 61-value observation and all 14
  actions;
\item
  describe why the actor needs recurrent memory even though the ball is
  absent from its input;
\item
  follow hidden and cell state through reset, rollout collection, PPO
  replay, ONNX export, evaluation, and Studio playback;
\item
  calculate every local reward term from the implemented formulas and
  coefficients;
\item
  explain the six-level physical-assistance curriculum without calling
  it a reward curriculum;
\item
  distinguish a released-actor warm start from a full optimizer resume
  and from training truly from scratch;
\item
  run the parser-verified training, evaluation, rendering, and
  Studio-adapter commands;
\item
  apply the physical evaluation gate and report the retained result
  honestly: \textbf{6/6 sustained balance, 0/3 commanded rolling};
\item
  state the limits: steering is not solved, lateral and yaw steering are
  not assessed, bridge behavior is not established, and no hardware
  claim is made.
\end{itemize}

\subsection{Evidence and claim
boundary}\label{evidence-and-claim-boundary}

This chapter uses repository code and retained repository artifacts
only. The main sources are:

{\def\LTcaptype{none} % do not increment counter
\begin{longtable}[]{@{}
  >{\raggedright\arraybackslash}p{(\linewidth - 2\tabcolsep) * \real{0.5000}}
  >{\raggedright\arraybackslash}p{(\linewidth - 2\tabcolsep) * \real{0.5000}}@{}}
\toprule\noalign{}
\begin{minipage}[b]{\linewidth}\raggedright
Question
\end{minipage} & \begin{minipage}[b]{\linewidth}\raggedright
Repository source
\end{minipage} \\
\midrule\noalign{}
\endhead
\bottomrule\noalign{}
\endlastfoot
Shared observation, joint order, timing, and default pose &
\texttt{microduc\allowbreak{}k\_\allowbreak{}local/\allowbreak{}src/\allowbreak{}microduc\allowbreak{}k\_\allowbreak{}local/\allowbreak{}contract\allowbreak{}.py} \\
Action application, one-step joint-velocity lag, and observation
assembly &
\texttt{microduc\allowbreak{}k\_\allowbreak{}local/\allowbreak{}src/\allowbreak{}microduc\allowbreak{}k\_\allowbreak{}local/\allowbreak{}walk\_\allowbreak{}env.py} \\
Basketball scene, assistance, termination, rewards, and metrics &
\texttt{rlx/\allowbreak{}rlx/\allowbreak{}environm\allowbreak{}ents/\allowbreak{}basketba\allowbreak{}ll.py} \\
LSTM actor, local critic, source verification, ONNX export, and parity &
\texttt{rlx/\allowbreak{}rlx/\allowbreak{}models/\allowbreak{}basketba\allowbreak{}ll.py} \\
Local recurrent PPO and direct CLI &
\texttt{rlx/\allowbreak{}examples\allowbreak{}/\allowbreak{}ppo\_\allowbreak{}microduc\allowbreak{}k\_\allowbreak{}basketba\allowbreak{}ll.py} \\
First-fall physical evaluator &
\texttt{rlx/\allowbreak{}scripts/\allowbreak{}eval\_\allowbreak{}basketba\allowbreak{}ll\_\allowbreak{}local.py\allowbreak{}} \\
Studio train/eval/render/import adapter &
\texttt{rlx/\allowbreak{}examples\allowbreak{}/\allowbreak{}ppo\_\allowbreak{}microduc\allowbreak{}k\_\allowbreak{}balance.\allowbreak{}py} \\
Studio recurrent-policy loader &
\texttt{microduc\allowbreak{}k\_\allowbreak{}local/\allowbreak{}src/\allowbreak{}microduc\allowbreak{}k\_\allowbreak{}local/\allowbreak{}studio\_\allowbreak{}policies\allowbreak{}.py} \\
Human-readable result summary &
\texttt{docs/\allowbreak{}basketba\allowbreak{}ll-\allowbreak{}showcase\allowbreak{}/\allowbreak{}RESULTS.\allowbreak{}md} \\
Retained direct evidence &
\texttt{docs/\allowbreak{}basketba\allowbreak{}ll-\allowbreak{}showcase\allowbreak{}/\allowbreak{}evidence\allowbreak{}/\allowbreak{}evaluati\allowbreak{}on.json} \\
Latest retained Studio evidence &
\texttt{rlx/\allowbreak{}runs/\allowbreak{}studio/\allowbreak{}basketba\allowbreak{}ll/\allowbreak{}basketba\allowbreak{}ll-\allowbreak{}balance-\allowbreak{}01/\allowbreak{}evaluati\allowbreak{}on.json} \\
\end{longtable}
}

The retained artifacts were produced on September 10, 2026. The latest
Studio report evaluates seeds 101, 102, and 103 for 60 seconds under
zero and \texttt{0.15\ m/s} forward commands. All six trials sustained
balance, but none of the three commanded trials passed controlled
rolling. The adapter also states that lateral and yaw command trials are
not implemented. Therefore:

\begin{quote}
The demonstrated result is unassisted sustained balance on a free
basketball. It is not controlled rolling, general steering, bridge
traversal, climbing onto the ball, recovery from the floor, or hardware
validation.
\end{quote}

The source basketball release in
\texttt{microduc\allowbreak{}k-\allowbreak{}playgrou\allowbreak{}nd}
reports a much larger upstream evaluation, but that published evaluation
is not a result of this local trainer. This chapter keeps source
evidence, local evidence, and Studio evidence separate.

\textbf{Block labels used in this chapter}

\begin{itemize}
\tightlist
\item
  \textbf{Source excerpt:} an actual repository excerpt shown for study.
  Surrounding imports or class context may be omitted, so it is not a
  standalone executable program.
\item
  \textbf{Executable command:} a shell command runnable from the
  workspace root when the documented environment and assets are
  available. Training commands create new artifacts.
\item
  \textbf{Recorded text/data:} a literal path, sequence, or metadata
  value from repository evidence. It is not a command.
\end{itemize}

\subsection{Build the physical scene from
numbers}\label{build-the-physical-scene-from-numbers}

\subsubsection{Units}\label{units}

The implementation uses the usual MuJoCo metric quantities, made
explicit by variable and metric names:

{\def\LTcaptype{none} % do not increment counter
\begin{longtable}[]{@{}
  >{\raggedright\arraybackslash}p{(\linewidth - 4\tabcolsep) * \real{0.3333}}
  >{\raggedright\arraybackslash}p{(\linewidth - 4\tabcolsep) * \real{0.3333}}
  >{\raggedright\arraybackslash}p{(\linewidth - 4\tabcolsep) * \real{0.3333}}@{}}
\toprule\noalign{}
\begin{minipage}[b]{\linewidth}\raggedright
Quantity
\end{minipage} & \begin{minipage}[b]{\linewidth}\raggedright
Unit
\end{minipage} & \begin{minipage}[b]{\linewidth}\raggedright
Example
\end{minipage} \\
\midrule\noalign{}
\endhead
\bottomrule\noalign{}
\endlastfoot
length and position & metre, \texttt{m} & ball radius
\texttt{0.12\ m} \\
mass & kilogram, \texttt{kg} & ball mass \texttt{0.62\ kg} \\
time & second, \texttt{s} & control interval \texttt{0.02\ s} \\
linear velocity & metres per second, \texttt{m/s} & command
\texttt{0.15\ m/s} \\
angle & radian internally & joint offsets and yaw command \\
angular velocity & radians per second, \texttt{rad/s} & yaw and ball
rotation rates \\
reported tilt & degree & evaluator limits tilt to less than
\texttt{50\ deg} \\
force & newton, \texttt{N} & assistance spring and damping forces \\
torque & newton-metre, \texttt{N\ m} & assistance torques \\
\end{longtable}
}

A radian is an angle measured by arc length divided by radius. For
orientation, \texttt{pi\ rad\ =\ 180\ deg}. The policy's 14 action
values are joint-angle offsets in radians, not motor torques.

\subsubsection{Solver, floor, ball, and visual
mesh}\label{solver-floor-ball-and-visual-mesh}

The scene starts from the full-collision Microduck robot and adds a
floor and a free body:

\begin{figure}
\centering
\includegraphics[width=6.2in,height=5.5in,alt={Schematic basketball geometry using the implemented radius and spawn heights}]{assets/diagrams/schematic-basketball-geometry.png}
\caption{Schematic basketball geometry using the implemented radius and
spawn heights}
\end{figure}

\textbf{Source excerpt (not standalone):}
\texttt{rlx/\allowbreak{}rlx/\allowbreak{}environm\allowbreak{}ents/\allowbreak{}basketba\allowbreak{}ll.py}

\begin{Shaded}
\begin{Highlighting}[]
\CommentTok{\# rlx/rlx/environments/basketball.py}
\NormalTok{BALL\_RADIUS }\OperatorTok{=} \FloatTok{0.12}
\NormalTok{BALL\_MASS }\OperatorTok{=} \FloatTok{0.62}

\NormalTok{spec.option.timestep }\OperatorTok{=}\NormalTok{ C.PHYSICS\_DT}
\NormalTok{spec.option.integrator }\OperatorTok{=}\NormalTok{ mujoco.mjtIntegrator.mjINT\_IMPLICITFAST}
\NormalTok{spec.option.cone }\OperatorTok{=}\NormalTok{ mujoco.mjtCone.mjCONE\_PYRAMIDAL}
\NormalTok{spec.option.iterations }\OperatorTok{=} \DecValTok{10}
\NormalTok{spec.option.ls\_iterations }\OperatorTok{=} \DecValTok{20}

\NormalTok{spec.worldbody.add\_geom(}
\NormalTok{    name}\OperatorTok{=}\StringTok{"floor"}\NormalTok{, }\BuiltInTok{type}\OperatorTok{=}\NormalTok{mujoco.mjtGeom.mjGEOM\_PLANE,}
\NormalTok{    size}\OperatorTok{=}\NormalTok{[}\DecValTok{0}\NormalTok{, }\DecValTok{0}\NormalTok{, }\FloatTok{0.05}\NormalTok{], material}\OperatorTok{=}\StringTok{"basketball\_ground\_mat"}\NormalTok{,}
\NormalTok{    friction}\OperatorTok{=}\NormalTok{[}\DecValTok{1}\NormalTok{, }\FloatTok{0.005}\NormalTok{, }\FloatTok{0.0001}\NormalTok{],}
\NormalTok{)}

\NormalTok{ball }\OperatorTok{=}\NormalTok{ spec.worldbody.add\_body(}
\NormalTok{    name}\OperatorTok{=}\StringTok{"basketball"}\NormalTok{, pos}\OperatorTok{=}\NormalTok{[}\DecValTok{0}\NormalTok{, }\DecValTok{0}\NormalTok{, BALL\_RADIUS }\OperatorTok{+} \FloatTok{0.001}\NormalTok{]}
\NormalTok{)}
\NormalTok{ball.add\_freejoint(name}\OperatorTok{=}\StringTok{"basketball\_freejoint"}\NormalTok{)}
\NormalTok{ball.add\_geom(}
\NormalTok{    name}\OperatorTok{=}\StringTok{"ball\_sphere"}\NormalTok{, }\BuiltInTok{type}\OperatorTok{=}\NormalTok{mujoco.mjtGeom.mjGEOM\_SPHERE,}
\NormalTok{    size}\OperatorTok{=}\NormalTok{[BALL\_RADIUS, }\DecValTok{0}\NormalTok{, }\DecValTok{0}\NormalTok{], mass}\OperatorTok{=}\NormalTok{BALL\_MASS,}
\NormalTok{    friction}\OperatorTok{=}\NormalTok{[}\FloatTok{1.2}\NormalTok{, }\FloatTok{0.01}\NormalTok{, }\FloatTok{0.001}\NormalTok{], priority}\OperatorTok{=}\DecValTok{1}\NormalTok{, condim}\OperatorTok{=}\DecValTok{4}\NormalTok{,}
\NormalTok{)}
\end{Highlighting}
\end{Shaded}

The sphere's centre begins at \texttt{z\ =\ 0.121\ m}, one millimetre
above a radius of \texttt{0.12\ m}. Its physical geom has mass and
contact. A second mesh geom supplies the basketball appearance but has
\texttt{mass=0}, \texttt{contype=0}, and \texttt{conaffinity=0}; it
cannot secretly support the robot. The texture and mesh come from
\texttt{microduc\allowbreak{}k-\allowbreak{}playgrou\allowbreak{}nd/\allowbreak{}src/\allowbreak{}mjlab\_\allowbreak{}microduc\allowbreak{}k/\allowbreak{}robot/\allowbreak{}assets/\allowbreak{}basketba\allowbreak{}ll/\allowbreak{}}.

The floor uses checkerboard texture colours and friction
\texttt{{[}1,\ 0.005,\ 0.0001{]}}. The ball uses
\texttt{{[}1.2,\ 0.01,\ 0.001{]}}, contact dimension 4, and priority 1.
These are simulator contact parameters, not measured certificates for a
particular real floor or ball.

The free joint gives the basketball seven configuration values: three
position coordinates plus a four-value quaternion. It also gives six
velocity values: three linear and three angular. Tests lock the complete
model at \texttt{nq=28}, \texttt{nv=26}, and \texttt{nu=14}.

\subsubsection{Timing}\label{timing}

The physics step is

\[
\Delta t_{\text{physics}}=0.005\text{ s}.
\]

Each policy action is held for four physics steps:

\[
\Delta t_{\text{control}}
=4(0.005)
=0.02\text{ s},
\]

so the policy runs at

\[
f_{\text{control}}=\frac{1}{0.02}=50\text{ Hz}.
\]

A 60-second evaluation therefore contains exactly

\[
\frac{60}{0.02}=3000
\]

control steps per trial.

\subsubsection{Spawn geometry}\label{spawn-geometry}

The ball centre is at \texttt{0.121\ m}. The robot root is placed at

\[
z_{\text{root}}
=2R+0.128
=2(0.12)+0.128
=0.368\text{ m}.
\]

The root is therefore

\[
0.368-0.121=0.247\text{ m}
\]

above the ball centre and about

\[
0.368-(0.121+0.12)=0.127\text{ m}
\]

above the ball's top. The reset test locks
\texttt{root\_\allowbreak{}above\_\allowbreak{}ball\_\allowbreak{}m} to
\texttt{0.247}.

The robot starts upright near the apex. Its horizontal root position is
sampled within \texttt{+-0.01\ m}, roll and pitch within
\texttt{+-2\ deg}, yaw across \texttt{{[}-pi,\ pi{]}} when random yaw is
enabled, and every joint within \texttt{+-0.05\ rad} of the default
pose. The ball starts at rest. This task does not teach mounting the
ball or recovering from the ground.

\subsection{The exact 61-observation and 14-action
contract}\label{the-exact-61-observation-and-14-action-contract}

\subsubsection{Observation slices}\label{observation-slices}

The actor always receives a \texttt{float32{[}61{]}} vector:

{\def\LTcaptype{none} % do not increment counter
\begin{longtable}[]{@{}lrrl@{}}
\toprule\noalign{}
Python slice & Indices & Width & Meaning \\
\midrule\noalign{}
\endhead
\bottomrule\noalign{}
\endlastfoot
\texttt{0:3} & 0-2 & 3 & trunk angular velocity \\
\texttt{3:6} & 3-5 & 3 & gravity projected into the trunk frame \\
\texttt{6:20} & 6-19 & 14 & joint position minus the default pose \\
\texttt{20:34} & 20-33 & 14 & joint velocity, delayed one control
step \\
\texttt{34:48} & 34-47 & 14 & previous raw actor action \\
\texttt{48:51} & 48-50 & 3 & forward, lateral, and yaw velocity
command \\
\texttt{51:55} & 51-54 & 4 & head-command padding, forced to zero \\
\texttt{55:61} & 55-60 & 6 & body-command padding, forced to zero \\
\end{longtable}
}

The assembly is direct:

\textbf{Source excerpt (not standalone):}
\texttt{microduc\allowbreak{}k\_\allowbreak{}local/\allowbreak{}src/\allowbreak{}microduc\allowbreak{}k\_\allowbreak{}local/\allowbreak{}walk\_\allowbreak{}env.py}

\begin{Shaded}
\begin{Highlighting}[]
\CommentTok{\# microduck\_local/src/microduck\_local/walk\_env.py}
\NormalTok{obs }\OperatorTok{=}\NormalTok{ np.empty(C.OBS\_DIM, np.float32)}
\NormalTok{obs[}\DecValTok{0}\NormalTok{:}\DecValTok{3}\NormalTok{] }\OperatorTok{=}\NormalTok{ gyro}
\NormalTok{obs[}\DecValTok{3}\NormalTok{:}\DecValTok{6}\NormalTok{] }\OperatorTok{=}\NormalTok{ gravity}
\NormalTok{obs[}\DecValTok{6}\NormalTok{:}\DecValTok{20}\NormalTok{] }\OperatorTok{=}\NormalTok{ joint\_pos}
\NormalTok{obs[}\DecValTok{20}\NormalTok{:}\DecValTok{34}\NormalTok{] }\OperatorTok{=}\NormalTok{ joint\_vel}
\NormalTok{obs[}\DecValTok{34}\NormalTok{:}\DecValTok{48}\NormalTok{] }\OperatorTok{=} \VariableTok{self}\NormalTok{.last\_action}
\NormalTok{obs[}\DecValTok{48}\NormalTok{:}\DecValTok{51}\NormalTok{] }\OperatorTok{=} \VariableTok{self}\NormalTok{.twist\_cmd}
\NormalTok{obs[}\DecValTok{51}\NormalTok{:}\DecValTok{55}\NormalTok{] }\OperatorTok{=} \VariableTok{self}\NormalTok{.head\_cmd}
\NormalTok{obs[}\DecValTok{55}\NormalTok{:}\DecValTok{61}\NormalTok{] }\OperatorTok{=} \VariableTok{self}\NormalTok{.body\_cmd}
\end{Highlighting}
\end{Shaded}

Basketball's command sampler sets the last ten values to zero. Both the
trainer and tests reject nonzero \texttt{51:61} values. Moving the ball
without changing robot proprioception leaves the observation unchanged
in the focused test. The actor is therefore \textbf{blind to explicit
ball state}: no ball position, velocity, rotation, radius, contact
label, or world position appears in the 61 values.

The word ``blind'' does not mean information-free. Contact with a moving
ball changes body angular velocity, projected gravity, joint state, and
the action-response history. The LSTM can use that time sequence to
infer useful hidden dynamics.

\subsubsection{Joint and action order}\label{joint-and-action-order}

The 14 action entries use this fixed servo order:

{\def\LTcaptype{none} % do not increment counter
\begin{longtable}[]{@{}rlr@{}}
\toprule\noalign{}
Action index & Joint & Default pose, rad \\
\midrule\noalign{}
\endhead
\bottomrule\noalign{}
\endlastfoot
0 & left hip yaw & 0.0000 \\
1 & left hip roll & -0.0873 \\
2 & left hip pitch & -0.4579 \\
3 & left knee & -0.0049 \\
4 & left ankle & 0.4530 \\
5 & neck pitch & 0.3491 \\
6 & head pitch & 0.3491 \\
7 & head yaw & 0.0000 \\
8 & head roll & 0.0000 \\
9 & right hip yaw & 0.0000 \\
10 & right hip roll & 0.0873 \\
11 & right hip pitch & 0.4579 \\
12 & right knee & 0.0049 \\
13 & right ankle & -0.4530 \\
\end{longtable}
}

For raw actor output \(\mathbf a_t\in\mathbb R^{14}\), the applied
position target is

\[
\mathbf q^{\text{target}}_t
=\mathbf q^{\text{default}}
+\operatorname{clip}(\mathbf a_t,-4,4).
\]

The observation stores the \textbf{raw} action, not the clipped action.
This keeps the next observation and action-change penalty aware of what
the network requested. Under BAM actuation, the position target drives
the voltage-based servo model, including its inherited 3-6 physics-step
bus delay. Four physics steps make one control step, so that internal
delay is shorter than two control intervals.

\subsection{The recurrent actor and its state
lifecycle}\label{the-recurrent-actor-and-its-state-lifecycle}

\subsubsection{Network}\label{network}

The local actor is:

\[
61
\rightarrow \operatorname{LSTM}(256)
\rightarrow 512
\rightarrow 256
\rightarrow 128
\rightarrow 14,
\]

with ELU activations in the multilayer perceptron. The diagonal Gaussian
exploration scale has 14 trainable values and starts at the requested
\texttt{initial\_std}, usually \texttt{0.03}. It is a direct
standard-deviation parameter, clamped to at least \texttt{1e-6}, not a
log-standard-deviation parameter.

The observation normalizer computes

\[
\widehat{\mathbf o}
=\frac{\mathbf o-\boldsymbol\mu}
{\boldsymbol\sigma+0.01}.
\]

The imported source normalizer is frozen. Its stored standard deviation
for the ten zero-padded head/body slots is zero, so the \texttt{+0.01}
denominator keeps those divisions finite. The exported ONNX contains the
normalizer.

The local critic is different from the released upstream training critic
described in the source model card. The local
\texttt{Basketba\allowbreak{}llCritic\allowbreak{}} takes only the same
normalized 61 values and is a fresh feedforward
\texttt{61\ -\textgreater{}\ 512\ -\textgreater{}\ 256\ -\textgreater{}\ 128\ -\textgreater{}\ 1}
network. It has no ball-state input.

\subsubsection{One-step inference}\label{one-step-inference}

\textbf{Source excerpt (not standalone):}
\texttt{rlx/\allowbreak{}rlx/\allowbreak{}models/\allowbreak{}basketba\allowbreak{}ll.py}

\begin{Shaded}
\begin{Highlighting}[]
\CommentTok{\# rlx/rlx/models/basketball.py}
\KeywordTok{def}\NormalTok{ forward(}\VariableTok{self}\NormalTok{, observations, h\_in, c\_in):}
\NormalTok{    normalized }\OperatorTok{=} \VariableTok{self}\NormalTok{.obs\_normalizer(observations)}
\NormalTok{    recurrent, (h\_out, c\_out) }\OperatorTok{=} \VariableTok{self}\NormalTok{.rnn.rnn(}
\NormalTok{        normalized.unsqueeze(}\DecValTok{0}\NormalTok{),}
\NormalTok{        (h\_in, c\_in),}
\NormalTok{    )}
\NormalTok{    actions }\OperatorTok{=} \VariableTok{self}\NormalTok{.mlp(recurrent.squeeze(}\DecValTok{0}\NormalTok{))}
    \ControlFlowTok{return}\NormalTok{ actions, h\_out, c\_out}
\end{Highlighting}
\end{Shaded}

The ONNX interface is exact:

\begin{itemize}
\tightlist
\item
  inputs: \texttt{obs\ {[}1,61{]}}, \texttt{h\_in\ {[}1,1,256{]}},
  \texttt{c\_in\ {[}1,1,256{]}};
\item
  outputs: \texttt{actions\ {[}1,14{]}}, \texttt{h\_out\ {[}1,1,256{]}},
  \texttt{c\_out\ {[}1,1,256{]}};
\item
  data type: \texttt{float32}.
\end{itemize}

\subsubsection{When memory is carried and
reset}\label{when-memory-is-carried-and-reset}

At the beginning of a trial or episode, hidden state \(\mathbf h\) and
cell state \(\mathbf c\) are zero. At every ordinary 50 Hz step:

\begin{enumerate}
\def\labelenumi{\arabic{enumi}.}
\tightlist
\item
  send \texttt{obs}, \texttt{h\_in}, and \texttt{c\_in} to the actor;
\item
  apply \texttt{actions};
\item
  carry \texttt{h\_out} and \texttt{c\_out} into the next call.
\end{enumerate}

An ordinary command change does not reset memory. A trial reset, episode
boundary, policy activation/switch, or recovery boundary must reset both
tensors.

Training stores an \texttt{episode\_starts{[}time,\ env{]}} mask. During
collection and during every PPO replay, the code multiplies hidden and
cell state by zero where an episode starts:

\textbf{Source excerpt (not standalone):}
\texttt{rlx/\allowbreak{}rlx/\allowbreak{}models/\allowbreak{}basketba\allowbreak{}ll.py}

\begin{Shaded}
\begin{Highlighting}[]
\CommentTok{\# rlx/rlx/models/basketball.py}
\ControlFlowTok{for}\NormalTok{ step }\KeywordTok{in} \BuiltInTok{range}\NormalTok{(observations.shape[}\DecValTok{0}\NormalTok{]):}
\NormalTok{    carry }\OperatorTok{=}\NormalTok{ (}\OperatorTok{\textasciitilde{}}\NormalTok{episode\_starts[step].}\BuiltInTok{bool}\NormalTok{()).to(observations.dtype)}
\NormalTok{    carry }\OperatorTok{=}\NormalTok{ carry.view(}\DecValTok{1}\NormalTok{, }\OperatorTok{{-}}\DecValTok{1}\NormalTok{, }\DecValTok{1}\NormalTok{)}
\NormalTok{    hidden }\OperatorTok{=}\NormalTok{ hidden }\OperatorTok{*}\NormalTok{ carry}
\NormalTok{    cell }\OperatorTok{=}\NormalTok{ cell }\OperatorTok{*}\NormalTok{ carry}
\NormalTok{    actions, hidden, cell }\OperatorTok{=} \VariableTok{self}\NormalTok{(observations[step], hidden, cell)}
\NormalTok{    outputs.append(actions)}
\end{Highlighting}
\end{Shaded}

This matters because recurrent PPO cannot treat stored observations as
unrelated rows. It recomputes the full time sequence from the rollout's
initial hidden and cell tensors, applying the same episode masks. The
evaluator similarly calls \texttt{policy.reset()} once at each trial
start and carries state until first fall or the exact horizon.

The exporter runs a 40-step PyTorch/ONNX parity test with a state reset
at step 20. The retained fixed-command actor's maximum absolute action
error was \texttt{1.43e-6}; the mixed-command actor's was
\texttt{2.15e-6}.

\subsection{The exact local reward}\label{the-exact-local-reward}

\subsubsection{Symbols}\label{symbols}

Let:

\begin{itemize}
\tightlist
\item
  \(\mathbf v_{xy}\) be body-frame forward and lateral trunk velocity;
\item
  \(\mathbf c_{xy}\) be the commanded forward and lateral velocity;
\item
  \(\boldsymbol\omega\) be trunk gyroscope values;
\item
  \(c_\omega\) be commanded yaw rate;
\item
  \(\mathbf g\) be gravity projected into the trunk frame;
\item
  \(\mathbf d_{xy}\) be trunk position minus ball position in the
  horizontal plane;
\item
  \(g_i\) be foot \(i\)'s radial gap from
  \texttt{BALL\_RADIUS\ +\ 0.012};
\item
  \(h\) be trunk height above the ball centre;
\item
  \(\mathbf v_b\) be horizontal ball velocity;
\item
  \(\mathbf a_t-\mathbf a_{t-1}\) be the raw action change.
\end{itemize}

Every weighted term is multiplied by \texttt{CTRL\_DT\ =\ 0.02}. The
reward is:

\[
r_t=0.02\sum_k w_k f_k.
\]

\subsubsection{Implemented formulas and
weights}\label{implemented-formulas-and-weights}

{\def\LTcaptype{none} % do not increment counter
\begin{longtable}[]{@{}
  >{\raggedright\arraybackslash}p{(\linewidth - 4\tabcolsep) * \real{0.3000}}
  >{\raggedright\arraybackslash}p{(\linewidth - 4\tabcolsep) * \real{0.3000}}
  >{\raggedleft\arraybackslash}p{(\linewidth - 4\tabcolsep) * \real{0.4000}}@{}}
\toprule\noalign{}
\begin{minipage}[b]{\linewidth}\raggedright
Term
\end{minipage} & \begin{minipage}[b]{\linewidth}\raggedright
Unweighted formula \(f_k\)
\end{minipage} & \begin{minipage}[b]{\linewidth}\raggedleft
Weight \(w_k\)
\end{minipage} \\
\midrule\noalign{}
\endhead
\bottomrule\noalign{}
\endlastfoot
linear tracking & \(\exp(-\|\mathbf v_{xy}-\mathbf c_{xy}\|^2/0.1)\) &
1.0 \\
yaw tracking & \(\exp(-(\omega_z-c_\omega)^2/0.5)\) & 0.5 \\
upright & \(\exp(-\|(g_x,g_y)\|^2/0.09)\) & 1.0 \\
centred & \(\exp(-\|\mathbf d_{xy}\|^2/0.0016)\) & 2.0 \\
feet on ball & \(\frac12\sum_i\exp(-g_i^2/0.0004)\) & 1.0 \\
height & \(\exp(-(h-0.245)^2/0.0009)\) & 1.0 \\
ball-speed penalty & \(-\min(\|\mathbf v_b\|^2,100)\) & 0.05 \\
angular-velocity penalty & \(-\min(\omega_x^2+\omega_y^2,100)\) &
0.05 \\
action-rate penalty & \(-\min(\|\mathbf a_t-\mathbf a_{t-1}\|^2,100)\) &
0.2 \\
\end{longtable}
}

The source is compact enough to compare directly:

\textbf{Source excerpt (not standalone):}
\texttt{rlx/\allowbreak{}rlx/\allowbreak{}environm\allowbreak{}ents/\allowbreak{}basketba\allowbreak{}ll.py}

\begin{Shaded}
\begin{Highlighting}[]
\CommentTok{\# rlx/rlx/environments/basketball.py}
\NormalTok{terms }\OperatorTok{=}\NormalTok{ \{}
    \StringTok{"linear\_tracking"}\NormalTok{: math.exp(}
        \OperatorTok{{-}}\BuiltInTok{float}\NormalTok{(np.}\BuiltInTok{sum}\NormalTok{((}\VariableTok{self}\NormalTok{.body\_lin\_vel()[:}\DecValTok{2}\NormalTok{] }\OperatorTok{{-}} \VariableTok{self}\NormalTok{.twist\_cmd[:}\DecValTok{2}\NormalTok{]) }\OperatorTok{**} \DecValTok{2}\NormalTok{)) }\OperatorTok{/} \FloatTok{.1}
\NormalTok{    ),}
    \StringTok{"yaw\_tracking"}\NormalTok{: math.exp(}
        \OperatorTok{{-}}\BuiltInTok{float}\NormalTok{((}\VariableTok{self}\NormalTok{.\_gyro[}\DecValTok{2}\NormalTok{] }\OperatorTok{{-}} \VariableTok{self}\NormalTok{.twist\_cmd[}\DecValTok{2}\NormalTok{]) }\OperatorTok{**} \DecValTok{2}\NormalTok{) }\OperatorTok{/} \FloatTok{.5}
\NormalTok{    ),}
    \StringTok{"upright"}\NormalTok{: math.exp(}\OperatorTok{{-}}\BuiltInTok{float}\NormalTok{(np.}\BuiltInTok{sum}\NormalTok{(gravity[:}\DecValTok{2}\NormalTok{] }\OperatorTok{**} \DecValTok{2}\NormalTok{)) }\OperatorTok{/} \FloatTok{.09}\NormalTok{),}
    \StringTok{"centered"}\NormalTok{: math.exp(}\OperatorTok{{-}}\BuiltInTok{float}\NormalTok{(np.}\BuiltInTok{sum}\NormalTok{(relative[:}\DecValTok{2}\NormalTok{] }\OperatorTok{**} \DecValTok{2}\NormalTok{)) }\OperatorTok{/} \FloatTok{.0016}\NormalTok{),}
    \StringTok{"feet\_on\_ball"}\NormalTok{: }\BuiltInTok{float}\NormalTok{(np.mean(np.exp(}\OperatorTok{{-}}\NormalTok{gaps }\OperatorTok{**} \DecValTok{2} \OperatorTok{/} \FloatTok{.0004}\NormalTok{))),}
    \StringTok{"height"}\NormalTok{: math.exp(}\OperatorTok{{-}}\BuiltInTok{float}\NormalTok{((relative[}\DecValTok{2}\NormalTok{] }\OperatorTok{{-}}\NormalTok{ BALL\_RADIUS }\OperatorTok{{-}} \FloatTok{.125}\NormalTok{) }\OperatorTok{**} \DecValTok{2}\NormalTok{) }\OperatorTok{/} \FloatTok{.0009}\NormalTok{),}
    \StringTok{"ball\_speed\_penalty"}\NormalTok{: }\OperatorTok{{-}}\BuiltInTok{float}\NormalTok{(}\BuiltInTok{min}\NormalTok{(np.}\BuiltInTok{sum}\NormalTok{(ball\_velocity }\OperatorTok{**} \DecValTok{2}\NormalTok{), }\DecValTok{100}\NormalTok{)),}
    \StringTok{"angular\_velocity\_penalty"}\NormalTok{: }\OperatorTok{{-}}\BuiltInTok{float}\NormalTok{(}\BuiltInTok{min}\NormalTok{(np.}\BuiltInTok{sum}\NormalTok{(}\VariableTok{self}\NormalTok{.\_gyro[:}\DecValTok{2}\NormalTok{] }\OperatorTok{**} \DecValTok{2}\NormalTok{), }\DecValTok{100}\NormalTok{)),}
    \StringTok{"action\_rate\_penalty"}\NormalTok{: }\OperatorTok{{-}}\BuiltInTok{float}\NormalTok{(}
        \BuiltInTok{min}\NormalTok{(np.}\BuiltInTok{sum}\NormalTok{((}\VariableTok{self}\NormalTok{.last\_action }\OperatorTok{{-}} \VariableTok{self}\NormalTok{.prev\_action) }\OperatorTok{**} \DecValTok{2}\NormalTok{), }\DecValTok{100}\NormalTok{)}
\NormalTok{    ),}
\NormalTok{\}}
\NormalTok{weighted }\OperatorTok{=}\NormalTok{ \{}
\NormalTok{    key: value }\OperatorTok{*}\NormalTok{ REWARD\_WEIGHTS[key] }\OperatorTok{*}\NormalTok{ CTRL\_DT}
    \ControlFlowTok{for}\NormalTok{ key, value }\KeywordTok{in}\NormalTok{ terms.items()}
\NormalTok{\}}
\end{Highlighting}
\end{Shaded}

The reward uses privileged simulator state for the ball, contacts, and
geometry. That is allowed for training rewards, but none of it enters
the actor observation.

\subsubsection{Worked reward example}\label{worked-reward-example}

Suppose one step has:

\begin{itemize}
\tightlist
\item
  command \texttt{0.15\ m/s} forward and measured forward speed
  \texttt{0.01\ m/s};
\item
  no lateral error;
\item
  horizontal root-ball offset \texttt{0.02\ m};
\item
  reset-like height \texttt{h=0.247\ m};
\item
  raw action-change squared sum \texttt{0.04};
\item
  all other positive terms equal to 1 and other penalties zero.
\end{itemize}

Linear tracking:

\[
f_{\text{lin}}
=\exp\left(-\frac{(0.01-0.15)^2}{0.1}\right)
=\exp(-0.196)
\approx0.8220.
\]

Its weighted contribution is

\[
0.8220(1.0)(0.02)\approx0.01644.
\]

Centred:

\[
f_{\text{center}}
=\exp\left(-\frac{0.02^2}{0.0016}\right)
=e^{-0.25}
\approx0.7788,
\]

so the contribution is

\[
0.7788(2.0)(0.02)\approx0.03115.
\]

Height:

\[
f_{\text{height}}
=\exp\left(-\frac{(0.247-0.245)^2}{0.0009}\right)
\approx0.9956,
\]

giving approximately \texttt{0.01991}.

Action-rate penalty:

\[
(-0.04)(0.2)(0.02)=-0.00016.
\]

If every positive term were exactly 1 and every penalty were zero, the
maximum per-step total would be

\[
0.02(1+0.5+1+2+1+1)=0.13.
\]

This is a dense shaping score, not the success test. The evaluator
deliberately ignores reward totals.

\subsection{Physical-assistance
curriculum}\label{physical-assistance-curriculum}

\subsubsection{Ladder and transition
rule}\label{ladder-and-transition-rule}

The only curriculum levels are:

\textbf{Recorded text/data (not executable):} the implemented assistance
ladder.

\begin{Shaded}
\begin{Highlighting}[]
\NormalTok{1.0 {-}\textgreater{} 0.5 {-}\textgreater{} 0.25 {-}\textgreater{} 0.1 {-}\textgreater{} 0.03 {-}\textgreater{} 0.0}
\end{Highlighting}
\end{Shaded}

After a completed episode:

\begin{itemize}
\tightlist
\item
  duration at least \texttt{6\ s}: move one level toward zero
  assistance;
\item
  duration below \texttt{1.5\ s}: move one level toward stronger
  assistance;
\item
  otherwise: remain at the same level.
\end{itemize}

The transition is applied on the next reset and only if the previous
episode actually finished. Reward weights never change.

\subsubsection{What assistance does}\label{what-assistance-does}

The assistance method first clears every external force. At
\texttt{hold=0}, it returns immediately. For nonzero hold, it applies:

\textbf{Source excerpt (not standalone):}
\texttt{rlx/\allowbreak{}rlx/\allowbreak{}environm\allowbreak{}ents/\allowbreak{}basketba\allowbreak{}ll.py}

\begin{Shaded}
\begin{Highlighting}[]
\CommentTok{\# rlx/rlx/environments/basketball.py}
\NormalTok{ball\_force }\OperatorTok{=} \OperatorTok{{-}}\DecValTok{400} \OperatorTok{*}\NormalTok{ (ball\_pos }\OperatorTok{{-}} \VariableTok{self}\NormalTok{.\_ball\_anchor) }\OperatorTok{{-}} \DecValTok{20} \OperatorTok{*}\NormalTok{ ball\_vel[:}\DecValTok{3}\NormalTok{]}
\NormalTok{ball\_force[}\DecValTok{2}\NormalTok{] }\OperatorTok{=} \OperatorTok{{-}}\DecValTok{20} \OperatorTok{*}\NormalTok{ ball\_vel[}\DecValTok{2}\NormalTok{]}
\VariableTok{self}\NormalTok{.data.xfrc\_applied[}\VariableTok{self}\NormalTok{.ball\_body\_id, :}\DecValTok{3}\NormalTok{] }\OperatorTok{=} \VariableTok{self}\NormalTok{.hold }\OperatorTok{*}\NormalTok{ ball\_force}

\VariableTok{self}\NormalTok{.data.xfrc\_applied[}\VariableTok{self}\NormalTok{.ball\_body\_id, }\DecValTok{3}\NormalTok{:] }\OperatorTok{=}\NormalTok{ (}
    \OperatorTok{{-}}\VariableTok{self}\NormalTok{.hold }\OperatorTok{*} \FloatTok{.2} \OperatorTok{*}\NormalTok{ (ball\_rotation }\OperatorTok{@}\NormalTok{ ball\_vel[}\DecValTok{3}\NormalTok{:])}
\NormalTok{)}

\NormalTok{duck\_force }\OperatorTok{=} \OperatorTok{{-}}\DecValTok{40} \OperatorTok{*}\NormalTok{ (}
    \VariableTok{self}\NormalTok{.data.xpos[}\VariableTok{self}\NormalTok{.trunk\_body\_id] }\OperatorTok{{-}}\NormalTok{ ball\_pos}
\NormalTok{) }\OperatorTok{{-}} \DecValTok{4} \OperatorTok{*} \VariableTok{self}\NormalTok{.data.qvel[:}\DecValTok{3}\NormalTok{]}
\NormalTok{duck\_force[}\DecValTok{2}\NormalTok{] }\OperatorTok{=} \DecValTok{0}
\VariableTok{self}\NormalTok{.data.xfrc\_applied[}\VariableTok{self}\NormalTok{.trunk\_body\_id, :}\DecValTok{3}\NormalTok{] }\OperatorTok{=} \VariableTok{self}\NormalTok{.hold }\OperatorTok{*}\NormalTok{ duck\_force}

\VariableTok{self}\NormalTok{.data.xfrc\_applied[}\VariableTok{self}\NormalTok{.trunk\_body\_id, }\DecValTok{3}\NormalTok{:] }\OperatorTok{=} \VariableTok{self}\NormalTok{.hold }\OperatorTok{*}\NormalTok{ (}
    \DecValTok{3} \OperatorTok{*}\NormalTok{ np.cross(up, [}\DecValTok{0}\NormalTok{, }\DecValTok{0}\NormalTok{, }\DecValTok{1}\NormalTok{]) }\OperatorTok{{-}} \FloatTok{.08} \OperatorTok{*}\NormalTok{ angular\_world}
\NormalTok{)}
\end{Highlighting}
\end{Shaded}

The ball receives horizontal spring-damper support toward its anchor,
vertical damping without a vertical position spring, and rotational
damping. The trunk receives horizontal centring force plus an
upright-restoring and angular-damping torque.

\textbf{Worked assistance example.} If the ball is displaced
\texttt{0.02\ m} in positive x with zero velocity, its unscaled x force
is

\[
-400(0.02)=-8\text{ N}.
\]

At \texttt{hold=0.25}, the applied x force is \texttt{-2\ N}. At
\texttt{hold=0}, it is exactly zero.

The retained curriculum smoke ran 7,168 environment steps with two
environments, completed 39 episodes, and exercised both promotion and
demotion. It ended with environments between \texttt{hold=0.25} and
\texttt{hold=0.5}. It did \textbf{not} graduate to zero and does not
demonstrate learning free-ball balance from scratch.

\subsection{Training: warm start, PPO, and the scratch
limit}\label{training-warm-start-ppo-and-the-scratch-limit}

\subsubsection{What is loaded}\label{what-is-loaded}

The direct trainer defaults to:

\textbf{Recorded text/data (not executable):} the default
source-checkpoint path.

\begin{Shaded}
\begin{Highlighting}[]
\NormalTok{microduck{-}playground/artifacts/basketball/checkpoint.pt}
\end{Highlighting}
\end{Shaded}

It verifies the released source checkpoint and adjacent ONNX with fixed
SHA256 values. It then loads
\texttt{actor\_\allowbreak{}state\_\allowbreak{}dict}, freezes the
source observation normalizer, and overwrites the actor's exploration
standard deviation with
\texttt{-\allowbreak{}-\allowbreak{}initial-\allowbreak{}std}.

It does \textbf{not} restore the source critic, Adam optimizer, GPU
curriculum counter, or exact mjlab simulator state. Instead:

\textbf{Source excerpt (not standalone):}
\texttt{rlx/\allowbreak{}examples\allowbreak{}/\allowbreak{}ppo\_\allowbreak{}microduc\allowbreak{}k\_\allowbreak{}basketba\allowbreak{}ll.py}

\begin{Shaded}
\begin{Highlighting}[]
\CommentTok{\# rlx/examples/ppo\_microduck\_basketball.py}
\NormalTok{actor, source\_metadata }\OperatorTok{=}\NormalTok{ load\_source\_actor(...)}
\NormalTok{critic }\OperatorTok{=}\NormalTok{ BasketballCritic(actor.obs\_normalizer)}
\NormalTok{optimizer }\OperatorTok{=}\NormalTok{ torch.optim.Adam(parameters, lr}\OperatorTok{=}\NormalTok{args.learning\_rate)}
\end{Highlighting}
\end{Shaded}

The retained summaries correctly label this:

\textbf{Recorded text/data (not executable):} retained run metadata.

\begin{Shaded}
\begin{Highlighting}[]
\NormalTok{actor warm{-}start with fresh local critic and optimizer}
\NormalTok{full\_upstream\_resume: false}
\end{Highlighting}
\end{Shaded}

A local checkpoint may be used as another actor warm start, but it must
have the local format, frozen-normalization contract, an adjacent ONNX,
and recurrent parity. This is still not a full upstream resume.

\subsubsection{What ``scratch'' would
require}\label{what-scratch-would-require}

The local CLI has no \texttt{-\/-scratch} or random-actor option.
\texttt{train()} always calls \texttt{load\_source\_actor()}. The Studio
adapter also defaults to the released checkpoint or an
\texttt{-\/-init-from} checkpoint. Therefore the repository contains no
retained evidence that this local reward, curriculum, and CPU PPO
implementation can discover basketball balance from a random actor.

The mature b11 actor is the source of the demonstrated balance. Local
runs adapt it. The assistance smoke also begins from that actor. A truly
scratch experiment would need a deliberate code path, a declared random
initialization, a smoke test, and an unassisted first-fall evaluation.
None is claimed here.

\subsubsection{Rollout and PPO details}\label{rollout-and-ppo-details}

For \texttt{N} environments and \texttt{T} rollout steps, one update
collects

\[
N\times T
\]

transitions. The retained fixed-command run used \texttt{N=8},
\texttt{T=128}, and 100 updates:

\[
8\times128\times100=102{,}400
\]

new local transitions.

The actor samples

\[
\mathbf a_t
=\boldsymbol\mu_t+\boldsymbol\epsilon_t\odot\boldsymbol\sigma,
\qquad
\boldsymbol\epsilon_t\sim\mathcal N(\mathbf0,\mathbf I).
\]

GAE uses \texttt{gamma=0.99} and \texttt{lambda=0.95}. A true
termination does not bootstrap. A time-limit truncation uses the final
state's critic value in the one-step delta but stops the multi-step
trace across the reset. Advantages are normalized.

The PPO update uses one full recurrent sequence minibatch, five epochs,
clip coefficient \texttt{0.2}, value coefficient \texttt{1.0}, entropy
coefficient \texttt{0.01}, maximum gradient norm \texttt{1.0}, and
target KL \texttt{0.02}. It recomputes the complete LSTM sequence each
epoch. KL checking can stop later epochs, but a completed optimizer step
can still leave post-update KL above the target.

The fixed-command run used learning rate \texttt{2e-6}, initial standard
deviation \texttt{0.03}, and command \texttt{{[}0.15,0,0{]}}. The
mixed-command run separately used 153,600 steps, learning rate
\texttt{1e-5}, standard deviation \texttt{0.08}, and the direct
trainer's randomized sampler: exact zero 25\% of the time, otherwise
uniform forward \texttt{{[}-0.15,0.15{]}\ m/s}, lateral
\texttt{{[}-0.1,0.1{]}\ m/s}, and yaw \texttt{{[}-0.5,0.5{]}\ rad/s}.
They are independent warm starts, not a single 256,000-step lineage.

\begin{figure}
\centering
\includegraphics[width=6.2in,height=5.5in,alt={Observed PPO diagnostics for the two independent basketball continuation runs; the empty reward panel records that per-update reward history was not saved}]{assets/evidence/basketball-training-curves.png}
\caption{Observed PPO diagnostics for the two independent basketball
continuation runs; the empty reward panel records that per-update reward
history was not saved}
\end{figure}

\subsection{Parser-verified command line
workflow}\label{parser-verified-command-line-workflow}

All commands below run from the workspace root. Their option spellings
were checked against the current parsers. Training commands create new
artifacts and can take substantial time.

\subsubsection{Focused contract tests}\label{focused-contract-tests}

\textbf{Executable command (workspace root):}

\begin{Shaded}
\begin{Highlighting}[]
\ExtensionTok{rlx/.venv{-}microduck/bin/python} \AttributeTok{{-}m}\NormalTok{ pytest }\DataTypeTok{\textbackslash{}}
\NormalTok{  rlx/tests/test\_basketball.py }\DataTypeTok{\textbackslash{}}
\NormalTok{  rlx/tests/test\_basketball\_policy.py }\DataTypeTok{\textbackslash{}}
\NormalTok{  rlx/tests/test\_basketball\_evaluation.py }\AttributeTok{{-}q}
\end{Highlighting}
\end{Shaded}

\subsubsection{Direct five-update smoke}\label{direct-five-update-smoke}

Use a new output directory; the trainer refuses to overwrite a nonempty
run.

\textbf{Executable command (workspace root):}

\begin{Shaded}
\begin{Highlighting}[]
\FunctionTok{bash}\NormalTok{ scripts/train{-}basketball.sh }\DataTypeTok{\textbackslash{}}
  \AttributeTok{{-}{-}output}\NormalTok{ rlx/artifacts/basketball{-}book{-}smoke }\DataTypeTok{\textbackslash{}}
  \AttributeTok{{-}{-}num{-}envs}\NormalTok{ 4 }\DataTypeTok{\textbackslash{}}
  \AttributeTok{{-}{-}updates}\NormalTok{ 5 }\DataTypeTok{\textbackslash{}}
  \AttributeTok{{-}{-}steps}\NormalTok{ 32 }\DataTypeTok{\textbackslash{}}
  \AttributeTok{{-}{-}learning{-}rate}\NormalTok{ 2e{-}6 }\DataTypeTok{\textbackslash{}}
  \AttributeTok{{-}{-}initial{-}std}\NormalTok{ 0.03 }\DataTypeTok{\textbackslash{}}
  \AttributeTok{{-}{-}hold}\NormalTok{ 0 }\DataTypeTok{\textbackslash{}}
  \AttributeTok{{-}{-}no{-}curriculum} \DataTypeTok{\textbackslash{}}
  \AttributeTok{{-}{-}command}\NormalTok{ 0.15 0 0}
\end{Highlighting}
\end{Shaded}

This requests \texttt{4\ x\ 32\ x\ 5\ =\ 640} new transitions. It is a
wiring check, not learning evidence.

\subsubsection{Reproduce the fixed-command local adaptation
shape}\label{reproduce-the-fixed-command-local-adaptation-shape}

\textbf{Executable command (workspace root):}

\begin{Shaded}
\begin{Highlighting}[]
\FunctionTok{bash}\NormalTok{ scripts/train{-}basketball.sh }\DataTypeTok{\textbackslash{}}
  \AttributeTok{{-}{-}output}\NormalTok{ rlx/artifacts/basketball{-}book{-}fixed }\DataTypeTok{\textbackslash{}}
  \AttributeTok{{-}{-}num{-}envs}\NormalTok{ 8 }\DataTypeTok{\textbackslash{}}
  \AttributeTok{{-}{-}updates}\NormalTok{ 100 }\DataTypeTok{\textbackslash{}}
  \AttributeTok{{-}{-}steps}\NormalTok{ 128 }\DataTypeTok{\textbackslash{}}
  \AttributeTok{{-}{-}learning{-}rate}\NormalTok{ 2e{-}6 }\DataTypeTok{\textbackslash{}}
  \AttributeTok{{-}{-}target{-}kl}\NormalTok{ 0.02 }\DataTypeTok{\textbackslash{}}
  \AttributeTok{{-}{-}save{-}interval}\NormalTok{ 25 }\DataTypeTok{\textbackslash{}}
  \AttributeTok{{-}{-}initial{-}std}\NormalTok{ 0.03 }\DataTypeTok{\textbackslash{}}
  \AttributeTok{{-}{-}hold}\NormalTok{ 0 }\DataTypeTok{\textbackslash{}}
  \AttributeTok{{-}{-}no{-}curriculum} \DataTypeTok{\textbackslash{}}
  \AttributeTok{{-}{-}command}\NormalTok{ 0.15 0 0}
\end{Highlighting}
\end{Shaded}

\subsubsection{Honest direct evaluation}\label{honest-direct-evaluation}

The repeated \texttt{-\/-command} option is intentional.

\textbf{Executable command (workspace root):}

\begin{Shaded}
\begin{Highlighting}[]
\ExtensionTok{rlx/.venv{-}microduck/bin/python}\NormalTok{ rlx/scripts/eval\_basketball\_local.py }\DataTypeTok{\textbackslash{}}
  \AttributeTok{{-}{-}policy}\NormalTok{ rlx/artifacts/basketball{-}book{-}fixed/policy.onnx }\DataTypeTok{\textbackslash{}}
  \AttributeTok{{-}{-}output}\NormalTok{ rlx/artifacts/basketball{-}book{-}fixed{-}eval }\DataTypeTok{\textbackslash{}}
  \AttributeTok{{-}{-}seconds}\NormalTok{ 60 }\DataTypeTok{\textbackslash{}}
  \AttributeTok{{-}{-}seeds}\NormalTok{ 101 202 303 }\DataTypeTok{\textbackslash{}}
  \AttributeTok{{-}{-}command}\NormalTok{ 0 0 0 }\DataTypeTok{\textbackslash{}}
  \AttributeTok{{-}{-}command}\NormalTok{ 0.15 0 0 }\DataTypeTok{\textbackslash{}}
  \AttributeTok{{-}{-}actuator}\NormalTok{ bam }\DataTypeTok{\textbackslash{}}
  \AttributeTok{{-}{-}render}
\end{Highlighting}
\end{Shaded}

The evaluator returns exit status 1 when the aggregate physical gate
fails, while still writing \texttt{evaluati\allowbreak{}on.json} and
failure media. A nonzero status is expected when commanded rolling
fails.

\subsubsection{Studio-compatible
training}\label{studio-compatible-training}

The adapter rounds requested timesteps up to a whole
\texttt{num\_envs\ x\ num\_steps} recurrent batch. This command divides
exactly:

\textbf{Executable command (workspace root):}

\begin{Shaded}
\begin{Highlighting}[]
\ExtensionTok{rlx/.venv{-}microduck/bin/python}\NormalTok{ rlx/examples/ppo\_microduck\_balance.py train }\DataTypeTok{\textbackslash{}}
  \AttributeTok{{-}{-}recipe}\NormalTok{ basketball }\DataTypeTok{\textbackslash{}}
  \AttributeTok{{-}{-}output{-}dir}\NormalTok{ rlx/runs/studio/basketball/basketball{-}book{-}fixed }\DataTypeTok{\textbackslash{}}
  \AttributeTok{{-}{-}total{-}timesteps}\NormalTok{ 102400 }\DataTypeTok{\textbackslash{}}
  \AttributeTok{{-}{-}num{-}envs}\NormalTok{ 8 }\DataTypeTok{\textbackslash{}}
  \AttributeTok{{-}{-}num{-}steps}\NormalTok{ 128 }\DataTypeTok{\textbackslash{}}
  \AttributeTok{{-}{-}num{-}minibatches}\NormalTok{ 1 }\DataTypeTok{\textbackslash{}}
  \AttributeTok{{-}{-}learning{-}rate}\NormalTok{ 2e{-}6 }\DataTypeTok{\textbackslash{}}
  \AttributeTok{{-}{-}initial{-}std}\NormalTok{ 0.03 }\DataTypeTok{\textbackslash{}}
  \AttributeTok{{-}{-}target{-}kl}\NormalTok{ 0.02 }\DataTypeTok{\textbackslash{}}
  \AttributeTok{{-}{-}hold}\NormalTok{ 0 }\DataTypeTok{\textbackslash{}}
  \AttributeTok{{-}{-}no{-}curriculum}
\end{Highlighting}
\end{Shaded}

The adapter fixes the command to \texttt{{[}0.15,0,0{]}}, BAM actuation,
a 10-second training horizon, frozen observation normalization, one
full-sequence minibatch, and the standalone learner's PPO constants.

\subsubsection{Studio evaluation and
rendering}\label{studio-evaluation-and-rendering}

\textbf{Executable command (workspace root):}

\begin{Shaded}
\begin{Highlighting}[]
\ExtensionTok{rlx/.venv{-}microduck/bin/python}\NormalTok{ rlx/examples/ppo\_microduck\_balance.py eval }\DataTypeTok{\textbackslash{}}
  \AttributeTok{{-}{-}recipe}\NormalTok{ basketball }\DataTypeTok{\textbackslash{}}
  \AttributeTok{{-}{-}policy}\NormalTok{ rlx/runs/studio/basketball/basketball{-}book{-}fixed/policy.onnx }\DataTypeTok{\textbackslash{}}
  \AttributeTok{{-}{-}output{-}dir}\NormalTok{ rlx/runs/studio/basketball/basketball{-}book{-}fixed }\DataTypeTok{\textbackslash{}}
  \AttributeTok{{-}{-}num{-}envs}\NormalTok{ 3 }\DataTypeTok{\textbackslash{}}
  \AttributeTok{{-}{-}seed}\NormalTok{ 101 }\DataTypeTok{\textbackslash{}}
  \AttributeTok{{-}{-}max{-}episode{-}s}\NormalTok{ 60 }\DataTypeTok{\textbackslash{}}
  \AttributeTok{{-}{-}evaluation{-}mode}\NormalTok{ skill}
\end{Highlighting}
\end{Shaded}

The adapter returns exit status 2 while full steering is unassessed or
rolling fails. It evaluates both zero command and \texttt{0.15\ m/s}
forward command for seeds 101, 102, and 103.

\textbf{Executable command (workspace root):}

\begin{Shaded}
\begin{Highlighting}[]
\ExtensionTok{rlx/.venv{-}microduck/bin/python}\NormalTok{ rlx/examples/ppo\_microduck\_balance.py render }\DataTypeTok{\textbackslash{}}
  \AttributeTok{{-}{-}recipe}\NormalTok{ basketball }\DataTypeTok{\textbackslash{}}
  \AttributeTok{{-}{-}policy}\NormalTok{ rlx/runs/studio/basketball/basketball{-}book{-}fixed/policy.onnx }\DataTypeTok{\textbackslash{}}
  \AttributeTok{{-}{-}output{-}dir}\NormalTok{ rlx/runs/studio/basketball/basketball{-}book{-}fixed }\DataTypeTok{\textbackslash{}}
  \AttributeTok{{-}{-}output}\NormalTok{ rlx/runs/studio/basketball/basketball{-}book{-}fixed/render }\DataTypeTok{\textbackslash{}}
  \AttributeTok{{-}{-}episodes}\NormalTok{ 1 }\DataTypeTok{\textbackslash{}}
  \AttributeTok{{-}{-}seed}\NormalTok{ 101 }\DataTypeTok{\textbackslash{}}
  \AttributeTok{{-}{-}render{-}seconds}\NormalTok{ 60 }\DataTypeTok{\textbackslash{}}
  \AttributeTok{{-}{-}width}\NormalTok{ 1280 }\DataTypeTok{\textbackslash{}}
  \AttributeTok{{-}{-}height}\NormalTok{ 720}
\end{Highlighting}
\end{Shaded}

Studio's renderer currently records the commanded-forward case only. A
failed rolling render can still be valuable evidence of sustained
balance and drift.

\subsection{Evaluation: balance is not
rolling}\label{evaluation-balance-is-not-rolling}

\subsubsection{First-fall accounting}\label{first-fall-accounting}

Evaluation does not auto-reset. It stops a trial at the first terminal
state or at the exact horizon. A complete trial requires:

\begin{itemize}
\tightlist
\item
  exactly the required number of samples;
\item
  no invalid physical metrics;
\item
  no termination or truncation;
\item
  zero automatic resets.
\end{itemize}

The environment terminates for low root height, root-ball offset above
\texttt{0.16\ m}, tilt above \texttt{55\ deg}, robot-floor contact,
non-foot robot-ball contact, or nonfinite physics. The evaluator is
stricter on accepted geometry: maximum tilt must remain below
\texttt{50\ deg}, and maximum root-ball offset below \texttt{0.16\ m}.

\subsubsection{Sustained balance}\label{sustained-balance}

A trial counts as sustained balance only when all are true:

\[
\begin{aligned}
&\text{duration complete}\ \land\ \text{foot-contact fraction}\ge0.5\\
&\land\ \text{no body-ball contact}\ \land\ \text{no robot-floor contact}\\
&\land\ \text{hold always zero}\ \land\ \text{stable geometry}.
\end{aligned}
\]

Balance does not require the ball to remain at one world position. A
policy may balance while the ball drifts or turns.

\subsubsection{Controlled rolling}\label{controlled-rolling}

For a nonzero planar command, controlled rolling additionally requires:

\begin{enumerate}
\def\labelenumi{\arabic{enumi}.}
\tightlist
\item
  enough signed ball progress along the initial command direction;
\item
  low forward and lateral body-velocity error;
\item
  low yaw-rate error;
\item
  enough integrated ball rotation.
\end{enumerate}

For command speed \(s\), duration \(T\), and ball radius \(R=0.12\):

\[
d_{\min}=\max(0.25,\;0.35sT),
\]

\[
e_{\text{linear,max}}=\max(0.04,\;0.6s),
\]

\[
e_{\text{yaw,max}}=0.5,
\]

\[
\theta_{\min}
=\frac{\max(0.25,\;0.35sT)}{R}(0.5).
\]

\textbf{Worked 60-second gate.} With \texttt{s=0.15\ m/s}:

\[
sT=0.15(60)=9.0\text{ m},
\]

\[
d_{\min}=0.35(9.0)=3.15\text{ m},
\]

\[
e_{\text{linear,max}}=0.6(0.15)=0.09\text{ m/s},
\]

\[
\theta_{\min}=\frac{3.15}{0.12}(0.5)=13.125\text{ rad}.
\]

Ball travel alone is insufficient. The evaluator labels substantial
travel as random movement when signed command progress is too small
relative to the required distance or to total travel.

\subsubsection{Latest retained Studio
result}\label{latest-retained-studio-result}

The latest retained Studio evaluation uses BAM, action delay, random
initial yaw, no pushes, no observation noise, no mass/friction domain
randomization, no assistance, and no automatic reset.

{\def\LTcaptype{none} % do not increment counter
\begin{longtable}[]{@{}
  >{\raggedright\arraybackslash}p{(\linewidth - 8\tabcolsep) * \real{0.1667}}
  >{\raggedleft\arraybackslash}p{(\linewidth - 8\tabcolsep) * \real{0.2222}}
  >{\raggedleft\arraybackslash}p{(\linewidth - 8\tabcolsep) * \real{0.2222}}
  >{\raggedleft\arraybackslash}p{(\linewidth - 8\tabcolsep) * \real{0.2222}}
  >{\raggedright\arraybackslash}p{(\linewidth - 8\tabcolsep) * \real{0.1667}}@{}}
\toprule\noalign{}
\begin{minipage}[b]{\linewidth}\raggedright
Case
\end{minipage} & \begin{minipage}[b]{\linewidth}\raggedleft
Trials
\end{minipage} & \begin{minipage}[b]{\linewidth}\raggedleft
Sustained balance
\end{minipage} & \begin{minipage}[b]{\linewidth}\raggedleft
Controlled rolling
\end{minipage} & \begin{minipage}[b]{\linewidth}\raggedright
Case verdict
\end{minipage} \\
\midrule\noalign{}
\endhead
\bottomrule\noalign{}
\endlastfoot
command \texttt{{[}0,0,0{]}} & 3 & 3/3 & not the case criterion &
pass \\
command \texttt{{[}0.15,0,0{]}} & 3 & 3/3 & 0/3 & fail \\
all trials & 6 & \textbf{6/6} & \textbf{0/3 commanded} & full task
fail \\
\end{longtable}
}

Every commanded trial failed both command-directed progress and velocity
tracking. The adapter reports:

\begin{figure}
\centering
\includegraphics[width=6.2in,height=5.5in,alt={Observed zero-command 60-second rollout, seed 101}]{assets/evidence/observed-basketball-zero-command-contact-sheet.png}
\caption{Observed zero-command 60-second rollout, seed 101}
\end{figure}

\begin{figure}
\centering
\includegraphics[width=6.2in,height=5.5in,alt={Observed 0.15 m/s forward-command attempt, seed 101; balance continues but the rolling gate fails}]{assets/evidence/observed-basketball-forward-command-contact-sheet.png}
\caption{Observed 0.15 m/s forward-command attempt, seed 101; balance
continues but the rolling gate fails}
\end{figure}

\textbf{Source excerpt (not standalone):}
\texttt{rlx/\allowbreak{}examples\allowbreak{}/\allowbreak{}ppo\_\allowbreak{}microduc\allowbreak{}k\_\allowbreak{}balance.\allowbreak{}py}

\begin{Shaded}
\begin{Highlighting}[]
\CommentTok{\# rlx/examples/ppo\_microduck\_balance.py}
\NormalTok{forward\_rolling\_passed }\OperatorTok{=}\NormalTok{ (}
    \BuiltInTok{bool}\NormalTok{(rolling\_cases) }\KeywordTok{and} \BuiltInTok{all}\NormalTok{(case[}\StringTok{"passed"}\NormalTok{] }\ControlFlowTok{for}\NormalTok{ case }\KeywordTok{in}\NormalTok{ rolling\_cases)}
\NormalTok{)}
\NormalTok{rolling\_passed }\OperatorTok{=} \VariableTok{False}
\NormalTok{success }\OperatorTok{=} \VariableTok{False}
\NormalTok{failures.append(}
    \StringTok{"full steering is not assessed: lateral and yaw command trials are not implemented"}
\NormalTok{)}
\end{Highlighting}
\end{Shaded}

The hard-coded fail-closed result is deliberate. Even a future
forward-only pass would not prove full steering. The current evidence
does not pass forward rolling either.

\subsection{Reading artifacts without
overclaiming}\label{reading-artifacts-without-overclaiming}

Use the files for different questions:

{\def\LTcaptype{none} % do not increment counter
\begin{longtable}[]{@{}
  >{\raggedright\arraybackslash}p{(\linewidth - 2\tabcolsep) * \real{0.5000}}
  >{\raggedright\arraybackslash}p{(\linewidth - 2\tabcolsep) * \real{0.5000}}@{}}
\toprule\noalign{}
\begin{minipage}[b]{\linewidth}\raggedright
Artifact
\end{minipage} & \begin{minipage}[b]{\linewidth}\raggedright
What it can establish
\end{minipage} \\
\midrule\noalign{}
\endhead
\bottomrule\noalign{}
\endlastfoot
\texttt{summary.json} & source hashes, settings, local steps, actor
change, parity, PPO diagnostics \\
\texttt{checkpoi\allowbreak{}nt.pt} & local actor, fresh local critic,
optimizer, and run metadata \\
\texttt{policy.onnx} & deterministic recurrent deployment graph with
frozen normalizer \\
\texttt{evaluati\allowbreak{}on.json} & every physical trial, criteria,
failures, and aggregate verdict \\
\texttt{rollout.mp4} & visible motion for one rendered seed and
command \\
\texttt{contact-\allowbreak{}sheet.pn\allowbreak{}g} & spaced visual
inspection across the rollout \\
Studio \texttt{result.json} & adapter-level pass/fail and provenance \\
\end{longtable}
}

Do not infer:

\begin{itemize}
\tightlist
\item
  steering from rotating ball texture;
\item
  straight commanded travel from total path length;
\item
  scratch learning from a curriculum smoke;
\item
  exact upstream continuation from actor warm-start training;
\item
  hardware safety from CPU MuJoCo;
\item
  bridge traversal from basketball balance.
\end{itemize}

The correct stopping statement for the retained policy is:

\begin{quote}
It sustains unassisted balance for all six retained 60-second Studio
trials. It fails all three commanded-forward rolling trials. Lateral and
yaw steering are unassessed. Bridge behavior and hardware behavior are
outside this evidence.
\end{quote}

\subsection{Exercises}\label{exercises-1}

\begin{enumerate}
\def\labelenumi{\arabic{enumi}.}
\tightlist
\item
  The ball radius is \texttt{0.12\ m}. What is the ball diameter?
\item
  At 50 Hz, how many control steps occur in 12 seconds?
\item
  Which observation indices hold the yaw command?
\item
  Does observation index 55 contain ball position?
\item
  A raw action at joint 4 is \texttt{0.20\ rad}. Ignoring clipping, what
  ankle target is applied when the default is \texttt{0.4530\ rad}?
\item
  Compute the centred reward's unweighted value for horizontal offset
  \texttt{0.04\ m}.
\item
  At \texttt{hold=0.5}, a stationary ball is displaced \texttt{0.01\ m}
  in x. What x assistance force is applied to the ball?
\item
  A curriculum environment at \texttt{hold=0.25} completes a
  \texttt{6.4\ s} episode. What is the next hold?
\item
  How many new transitions are collected by 8 environments, 128 steps,
  and 25 updates?
\item
  For \texttt{0.15\ m/s} over 60 seconds, why is \texttt{1.30\ m} signed
  progress a failure even if the robot remains upright?
\item
  A trial has 60-second duration, 90\% foot contact, zero hold, no
  forbidden contacts, \texttt{4\ deg} maximum tilt, and \texttt{0.02\ m}
  maximum offset. The ball travels \texttt{4\ m} sideways under a
  forward command. Does it pass sustained balance? Does it pass
  controlled rolling?
\item
  Why is the retained 7,168-step curriculum smoke not evidence of
  scratch learning?
\end{enumerate}

\subsection{Answers}\label{answers-1}

\begin{enumerate}
\def\labelenumi{\arabic{enumi}.}
\tightlist
\item
  \(2R=2(0.12)=0.24\text{ m}\).
\item
  \(12/0.02=600\) steps.
\item
  Index 50, the third value in slice \texttt{48:51}.
\item
  No.~Slice \texttt{55:61} is zero body-command padding. The actor
  receives no explicit ball state.
\item
  \(0.4530+0.20=0.6530\text{ rad}\).
\item
  \(\exp(-0.04^2/0.0016)=\exp(-1)\approx0.3679\).
\item
  Unscaled force is \(-400(0.01)=-4\text{ N}\); half hold applies
  \texttt{-2\ N}.
\item
  The ladder moves one step toward zero:
  \texttt{0.25\ -\textgreater{}\ 0.1}.
\item
  \(8\times128\times25=25{,}600\) transitions.
\item
  The gate requires at least \texttt{3.15\ m} signed progress, so
  \texttt{1.30\ m} is below target. Upright survival alone is balance,
  not controlled rolling.
\item
  It passes sustained balance because those conditions satisfy the
  balance gate. It fails controlled rolling because sideways travel does
  not provide enough signed forward progress and likely fails velocity
  tracking.
\item
  The smoke started from the mature released actor, used physical
  assistance, and ended above zero assistance. The local CLI did not
  initialize a random actor, and no unassisted scratch policy was
  evaluated.
\end{enumerate}

\section{Suspended Bridge: Learning to Walk on a Moving
Support}\label{suspended-bridge-learning-to-walk-on-a-moving-support}

\subsection{Learning objectives}\label{learning-objectives-1}

After this chapter, you should be able to:

\begin{enumerate}
\def\labelenumi{\arabic{enumi}.}
\tightlist
\item
  describe the simulated bridge geometry, four tendon cables, free
  plank, platforms, and non-colliding gantry with exact dimensions;
\item
  distinguish a simulated bill of materials from a certified physical
  construction plan;
\item
  trace the unchanged 61-observation and 14-action deployment contract;
\item
  explain why the navigation commander may use world state while the
  learned actor receives only the standard command slots;
\item
  derive every inherited locomotion reward term and its actual weight;
\item
  explain why the curriculum changes physics and spawn position, but
  never reward weights;
\item
  describe the walking-policy transfer: transferred actor and
  normalizer, frozen normalization statistics, fresh critic, and fresh
  optimizer;
\item
  reproduce the current build, train, export, evaluation, and render
  commands; and
\item
  apply the strict full-20-second crossing gate without turning partial
  progress into a successful result.
\end{enumerate}

The most important result comes first:

\begin{quote}
The recorded \texttt{bridge-\allowbreak{}studio-\allowbreak{}02} pilot
trained for 32,768 PPO transitions and did \textbf{not} cross the
suspended bridge. Its best measured progress was 1.175 m from the launch
position, below the 1.60 m progress threshold and without a
contact-validated two-foot landing. This chapter documents a working
simulation and training protocol, not a mastered skill.
\end{quote}

\subsection{Evidence and claim
boundary}\label{evidence-and-claim-boundary-1}

This chapter is grounded in the current repository implementation:

{\def\LTcaptype{none} % do not increment counter
\begin{longtable}[]{@{}
  >{\raggedright\arraybackslash}p{(\linewidth - 2\tabcolsep) * \real{0.5000}}
  >{\raggedright\arraybackslash}p{(\linewidth - 2\tabcolsep) * \real{0.5000}}@{}}
\toprule\noalign{}
\begin{minipage}[b]{\linewidth}\raggedright
Question
\end{minipage} & \begin{minipage}[b]{\linewidth}\raggedright
Source
\end{minipage} \\
\midrule\noalign{}
\endhead
\bottomrule\noalign{}
\endlastfoot
Bridge geometry, commander, curriculum, contacts, and metrics &
\texttt{rlx/\allowbreak{}rlx/\allowbreak{}environm\allowbreak{}ents/\allowbreak{}bridge.p\allowbreak{}y} \\
Strict crossing gate &
\texttt{rlx/\allowbreak{}rlx/\allowbreak{}environm\allowbreak{}ents/\allowbreak{}bridge\_\allowbreak{}evaluati\allowbreak{}on.py} \\
Walking actor transfer &
\texttt{rlx/\allowbreak{}rlx/\allowbreak{}models/\allowbreak{}bridge\_\allowbreak{}bootstra\allowbreak{}p.py} \\
Studio train/eval/render/export parser &
\texttt{rlx/\allowbreak{}examples\allowbreak{}/\allowbreak{}ppo\_\allowbreak{}microduc\allowbreak{}k\_\allowbreak{}studio.p\allowbreak{}y} \\
Shared observation and action order &
\texttt{microduc\allowbreak{}k\_\allowbreak{}local/\allowbreak{}src/\allowbreak{}microduc\allowbreak{}k\_\allowbreak{}local/\allowbreak{}contract\allowbreak{}.py} \\
Inherited reward equations &
\texttt{microduc\allowbreak{}k\_\allowbreak{}local/\allowbreak{}src/\allowbreak{}microduc\allowbreak{}k\_\allowbreak{}local/\allowbreak{}walk\_\allowbreak{}env.py} \\
Recorded pilot and deterministic evaluation &
\texttt{docs/\allowbreak{}bridge-\allowbreak{}showcase\allowbreak{}/\allowbreak{}}
and
\texttt{rlx/\allowbreak{}runs/\allowbreak{}studio/\allowbreak{}bridge/\allowbreak{}bridge-\allowbreak{}studio-\allowbreak{}02/\allowbreak{}} \\
\end{longtable}
}

The local training playbook adds three interpretation rules:

\begin{itemize}
\tightlist
\item
  never change the 61/14 interface for one task;
\item
  do not treat reward or PPO completion as proof of behavior; and
\item
  inspect a deterministic exported rollout before claiming a skill.
\end{itemize}

The bridge model is a \textbf{local MuJoCo prototype}. Its dimensions,
masses, friction coefficients, tendon stiffness, damping, and limits are
simulator parameters. They are not measurements from a constructed
bridge, cable ratings, load calculations, a safety factor, or a hardware
certification. The render material names such as
\texttt{bridge\_\allowbreak{}plank\_\allowbreak{}mat} describe color and
appearance, not a specified species of wood or grade of metal.

\subsection{What is physically
simulated?}\label{what-is-physically-simulated}

The task places Microduck on a raised launch platform. A narrow plank
spans the gap to a second platform. The plank is not welded, hinged, or
kinematically animated. It has a MuJoCo free joint, so contact forces
can translate it and rotate it in roll, pitch, and yaw. Four spatial
tendons connect moving attachment sites on the plank to fixed world
anchors.

\subsubsection{Simulated bill of
materials}\label{simulated-bill-of-materials}

This table is a simulator inventory, not a shopping or construction
list.

{\def\LTcaptype{none} % do not increment counter
\begin{longtable}[]{@{}
  >{\raggedleft\arraybackslash}p{(\linewidth - 4\tabcolsep) * \real{0.4000}}
  >{\raggedright\arraybackslash}p{(\linewidth - 4\tabcolsep) * \real{0.3000}}
  >{\raggedright\arraybackslash}p{(\linewidth - 4\tabcolsep) * \real{0.3000}}@{}}
\toprule\noalign{}
\begin{minipage}[b]{\linewidth}\raggedleft
Count
\end{minipage} & \begin{minipage}[b]{\linewidth}\raggedright
Simulated item
\end{minipage} & \begin{minipage}[b]{\linewidth}\raggedright
Exact model values
\end{minipage} \\
\midrule\noalign{}
\endhead
\bottomrule\noalign{}
\endlastfoot
1 & Free plank & \texttt{1.10\ x\ 0.13\ x\ 0.03\ m}, mass
\texttt{0.45\ kg}, center initially at \texttt{(0,\ 0,\ 0.225)\ m} \\
2 & Fixed platforms & each \texttt{0.60\ x\ 0.48\ x\ 0.24\ m}, centers
at \texttt{x=-0.85\ m} and \texttt{x=+0.85\ m} \\
4 & Spatial tendon cables & stiffness \texttt{320}, damping
\texttt{2.0}, render width \texttt{0.0015\ m}, length-limited \\
4 & Fixed cable anchors & \texttt{(x,y,z)=(+/-0.45,+/-0.17,0.77)\ m} \\
4 & Plank attachment sites & local
\texttt{(x,y,z)=(+/-0.45,+/-0.057,0)\ m} \\
2 & Gantry crossbars & \texttt{0.028\ x\ 0.48\ x\ 0.028\ m}, centered at
\texttt{x=+/-0.45\ m}, \texttt{z=0.784\ m} \\
4 & Gantry posts & \texttt{0.028\ x\ 0.028\ x\ 0.77\ m}, at
\texttt{x=+/-0.45\ m}, \texttt{y=+/-0.225\ m} \\
1 & Ground plane & top at world \texttt{z=0}, below the platform tops \\
\end{longtable}
}

MuJoCo box \texttt{size} values are half-extents. For example, the plank
source uses half-extents \texttt{(0.55,\ 0.065,\ 0.015)}, which gives
the full dimensions in the table.

The start platform spans \texttt{x=-1.15} to \texttt{-0.55\ m}. The
plank spans \texttt{x=-0.55} to \texttt{+0.55\ m}. The end platform
spans \texttt{x=+0.55} to \texttt{+1.15\ m}. Their boundaries touch
exactly in the nominal model. All three top surfaces are at
\texttt{z=0.24\ m}; the plank center is therefore
\texttt{0.24-0.015=0.225\ m}.

\subsubsection{The four cables}\label{the-four-cables}

At each of the two longitudinal positions, \texttt{x=-0.45\ m} and
\texttt{x=+0.45\ m}, one left and one right cable connect the plank to
the gantry. The lateral anchor-to-attachment offset is

\[
\!0.17-0.057=0.113\ \text{m}
\]

on the left and the same magnitude on the right. The vertical offset is

\[
0.77-0.225=0.545\ \text{m}.
\]

The nominal straight-line length is therefore

\[
L=\sqrt{0.113^2+0.545^2}
  =0.556591412079\ \text{m}.
\]

The source sets the upper spring length to \texttt{L-0.0015}, or about
\texttt{0.555091\ m}, and the maximum limited length to
\texttt{L+0.008}, or about \texttt{0.564591\ m}. The tendons are
massless MuJoCo constraints with elastic and damping parameters; the
code does not model cable strands, knots, connectors, fatigue, breaking
strength, or anchor pull-out.

A simplified linear-spring estimate gives an initial extension of
\texttt{0.0015\ m} and approximately

\[
F=k\Delta L=320(0.0015)=0.48
\]

in the simulator's force units before damping and constraint effects.
This is an instructional estimate, not a physical load rating.

\subsubsection{The gantry is visible but cannot catch the
robot}\label{the-gantry-is-visible-but-cannot-catch-the-robot}

The crossbars and posts use \texttt{contype=0} and
\texttt{conaffinity=0}. MuJoCo renders them, but collision detection
does not let the robot stand on them, lean on them, or receive a rescue
contact. The four tendon constraints are physical simulation elements;
the rectangular gantry geometry is decorative, non-colliding context.

The essential model construction is shown in this abridged source
excerpt. It omits the surrounding loop and is not a standalone program:

\begin{Shaded}
\begin{Highlighting}[]
\NormalTok{plank }\OperatorTok{=}\NormalTok{ spec.worldbody.add\_body(}
\NormalTok{    name}\OperatorTok{=}\StringTok{"bridge\_plank"}\NormalTok{,}
\NormalTok{    pos}\OperatorTok{=}\NormalTok{(}\DecValTok{0}\NormalTok{, }\DecValTok{0}\NormalTok{, }\FloatTok{0.225}\NormalTok{),}
\NormalTok{)}
\NormalTok{plank.add\_freejoint(name}\OperatorTok{=}\StringTok{"bridge\_freejoint"}\NormalTok{)}
\NormalTok{plank.add\_geom(}
\NormalTok{    name}\OperatorTok{=}\StringTok{"bridge\_surface"}\NormalTok{,}
    \BuiltInTok{type}\OperatorTok{=}\NormalTok{mujoco.mjtGeom.mjGEOM\_BOX,}
\NormalTok{    size}\OperatorTok{=}\NormalTok{(}\FloatTok{0.55}\NormalTok{, }\FloatTok{0.065}\NormalTok{, }\FloatTok{0.015}\NormalTok{),}
\NormalTok{    mass}\OperatorTok{=}\FloatTok{0.45}\NormalTok{,}
\NormalTok{    friction}\OperatorTok{=}\NormalTok{(}\FloatTok{1.15}\NormalTok{, }\FloatTok{0.006}\NormalTok{, }\FloatTok{0.0002}\NormalTok{),}
\NormalTok{    condim}\OperatorTok{=}\DecValTok{4}\NormalTok{,}
\NormalTok{)}

\NormalTok{tendon }\OperatorTok{=}\NormalTok{ spec.add\_tendon(}
\NormalTok{    stiffness}\OperatorTok{=}\FloatTok{320.0}\NormalTok{,}
\NormalTok{    damping}\OperatorTok{=}\FloatTok{2.0}\NormalTok{,}
\NormalTok{    springlength}\OperatorTok{=}\NormalTok{(}\FloatTok{0.0}\NormalTok{, cable\_length }\OperatorTok{{-}} \FloatTok{0.0015}\NormalTok{),}
\NormalTok{    limited}\OperatorTok{=}\VariableTok{True}\NormalTok{,}
    \BuiltInTok{range}\OperatorTok{=}\NormalTok{(}\FloatTok{0.0}\NormalTok{, cable\_length }\OperatorTok{+} \FloatTok{0.008}\NormalTok{),}
\NormalTok{    width}\OperatorTok{=}\FloatTok{0.0015}\NormalTok{,}
\NormalTok{)}
\NormalTok{tendon.wrap\_site(}\SpecialStringTok{f"}\SpecialCharTok{\{}\NormalTok{name}\SpecialCharTok{\}}\SpecialStringTok{\_anchor"}\NormalTok{)}
\NormalTok{tendon.wrap\_site(}\SpecialStringTok{f"}\SpecialCharTok{\{}\NormalTok{name}\SpecialCharTok{\}}\SpecialStringTok{\_attach"}\NormalTok{)}
\end{Highlighting}
\end{Shaded}

The floor friction is \texttt{(1.0,\ 0.005,\ 0.0001)}, platform friction
is \texttt{(1.1,\ 0.005,\ 0.0001)}, and plank friction is
\texttt{(1.15,\ 0.006,\ 0.0002)}. These coefficients are chosen
local-training values.

\subsection{Contact, motion, and
failure}\label{contact-motion-and-failure}

The simulation runs with a \texttt{0.005\ s} physics step and four
physics substeps per policy action:

\[
\Delta t_{\text{control}}=4(0.005)=0.02\ \text{s},
\]

so the actor runs at 50 Hz. A default 20-second episode contains 1,000
control steps.

The action has 14 values in the shared joint order. It is a joint-target
offset:

\[
\mathbf q^{\text{target}}
=\mathbf q^{\text{default}}+\operatorname{clip}(\mathbf a,-4,4).
\]

The plank moves only through simulated dynamics: gravity, tendon forces,
contacts, and optional curriculum forces. No control step teleports the
plank or robot forward. The reset function does place the robot at a
declared curriculum spawn and gives the plank a small random initial
roll, pitch, and angular velocity.

The bridge episode terminates when any of these conditions occurs:

\begin{itemize}
\tightlist
\item
  any robot body contacts the ground plane;
\item
  trunk tilt exceeds 70 degrees;
\item
  trunk height falls below \texttt{0.10\ m};
\item
  the trunk moves beyond
  \texttt{\textbar{}y\textbar{}\ \textgreater{}\ 0.36\ m} while neither
  foot has valid support;
\item
  physics becomes non-finite.
\end{itemize}

Reaching the ordinary time limit is a truncation, not a fall. Valid foot
support may come from the start platform, plank, or end platform. A
hand, head, trunk, or other non-foot body touching one of those support
surfaces is tracked as \texttt{nonfoot\_\allowbreak{}support}.

\subsection{The observation contract: privileged commander, ordinary
actor}\label{the-observation-contract-privileged-commander-ordinary-actor}

\subsubsection{The actor still receives exactly 61
values}\label{the-actor-still-receives-exactly-61-values}

The bridge actor has no special bridge-state input:

{\def\LTcaptype{none} % do not increment counter
\begin{longtable}[]{@{}lrl@{}}
\toprule\noalign{}
Slice & Width & Meaning \\
\midrule\noalign{}
\endhead
\bottomrule\noalign{}
\endlastfoot
\texttt{0:3} & 3 & trunk angular velocity \\
\texttt{3:6} & 3 & gravity projected into the trunk frame \\
\texttt{6:20} & 14 & joint position relative to
\texttt{DEFAULT\_POSE} \\
\texttt{20:34} & 14 & joint velocity delayed by one control step \\
\texttt{34:48} & 14 & previous raw actor action \\
\texttt{48:51} & 3 & body-frame forward, lateral, and yaw-rate
command \\
\texttt{51:55} & 4 & head command, zero for Bridge \\
\texttt{55:61} & 6 & body command, zero for Bridge \\
\end{longtable}
}

Bridge position, bridge velocity, cable length, plank roll and pitch,
world position, contact labels, crossing state, and assistance level are
absent. A test moves the plank laterally by \texttt{0.03\ m} without
changing the robot and confirms that the 61-value observation remains
byte-for-byte unchanged.

The actor and critic are separate
\texttt{61\ -\textgreater{}\ 512\ -\textgreater{}\ 256\ -\textgreater{}\ 128}
ELU networks, but both receive the same normalized 61-vector. This
implementation does \textbf{not} provide an asymmetric privileged
critic. Privileged simulator quantities are used by reward calculation,
curriculum decisions, logs, and evaluation.

\subsubsection{Why the commander is
privileged}\label{why-the-commander-is-privileged}

The navigation commander runs outside the learned actor. It reads the
trunk's world \texttt{x}, world lateral displacement \texttt{y}, and
world heading. It converts those quantities into the normal three
command slots. The actor sees the resulting command, not the world
measurements that produced it.

Let the horizontal forward heading be

\[
\mathbf h=(h_x,h_y,0),
\]

and the body-left direction be

\[
\mathbf s=(-h_y,h_x,0).
\]

Before the destination stop, the desired world velocity is

\[
\mathbf d=
\left(
0.25,\
\operatorname{clip}(-1.2y,-0.12,0.12),\
0
\right).
\]

The heading error is

\[
\psi=\operatorname{atan2}(h_y,h_x).
\]

The published body-frame command is

\[
\mathbf c=
\left(
\mathbf d\cdot\mathbf h,\
\mathbf d\cdot\mathbf s,\
\operatorname{clip}(-1.5\psi,-0.6,0.6)
\right).
\]

The implementation is:

\begin{Shaded}
\begin{Highlighting}[]
\NormalTok{reached\_destination }\OperatorTok{=}\NormalTok{ (}
\NormalTok{    trunk\_world\_x }\OperatorTok{\textgreater{}=}\NormalTok{ END\_PLATFORM\_X }\OperatorTok{{-}} \FloatTok{0.05}
\NormalTok{)  }\CommentTok{\# 0.85 {-} 0.05 = 0.80 m}

\NormalTok{desired\_world }\OperatorTok{=}\NormalTok{ np.array((}
    \FloatTok{0.0} \ControlFlowTok{if}\NormalTok{ reached\_destination }\ControlFlowTok{else}\NormalTok{ fixed\_command[}\DecValTok{0}\NormalTok{],}
\NormalTok{    np.clip(}\OperatorTok{{-}}\FloatTok{1.2} \OperatorTok{*}\NormalTok{ trunk\_world\_y, }\OperatorTok{{-}}\FloatTok{0.12}\NormalTok{, }\FloatTok{0.12}\NormalTok{),}
    \FloatTok{0.0}\NormalTok{,}
\NormalTok{))}

\NormalTok{twist\_cmd[:] }\OperatorTok{=}\NormalTok{ (}
\NormalTok{    desired\_world }\OperatorTok{@}\NormalTok{ heading,}
\NormalTok{    desired\_world }\OperatorTok{@}\NormalTok{ side,}
\NormalTok{    np.clip(}\OperatorTok{{-}}\FloatTok{1.5} \OperatorTok{*}\NormalTok{ yaw\_error, }\OperatorTok{{-}}\FloatTok{0.6}\NormalTok{, }\FloatTok{0.6}\NormalTok{),}
\NormalTok{)}
\end{Highlighting}
\end{Shaded}

This commander is not a physical assist: it applies no force and writes
no joint action. However, it is privileged protocol logic because it
uses world-frame localization unavailable inside the 61-vector. A
complete system claim must therefore include the commander, not
attribute all navigation to the actor.

When trunk \texttt{x\ \textgreater{}=\ 0.80\ m}, the current commander
sets desired world-forward speed to zero while retaining lateral and yaw
corrections. This destination stop distinguishes the final
\texttt{bridge-\allowbreak{}studio-\allowbreak{}02} protocol from
\texttt{bridge-\allowbreak{}studio-\allowbreak{}01}, which preceded that
stop behavior. The two runs are independent 32,768-transition warm
starts; they are not one 65,536-transition training chain.

\subsection{Reward: inherited walking, not a crossing
jackpot}\label{reward-inherited-walking-not-a-crossing-jackpot}

\texttt{BridgeEnv} does not override
\texttt{\_\allowbreak{}compute\_\allowbreak{}reward}. The bridge recipe
also declares an empty set of bridge reward override keys. Therefore all
curriculum stages use the exact base walking reward:

\[
r_t=
r_{\text{lin}}+
r_{\text{ang}}+
r_{\text{upright}}+
r_{\text{pose}}+
r_{\text{head}}+
r_{\text{air}}+
r_{\Delta a}+
r_{\omega xy}.
\]

There is no reward for world \texttt{x}, bridge progress, plank angle,
cable tension, end-platform contact, or
\texttt{bridge\_\allowbreak{}crossed}.

Let:

\begin{itemize}
\tightlist
\item
  \(\mathbf c=(c_x,c_y,c_\omega)\) be \texttt{twist\_cmd};
\item
  \(\mathbf v_b=(v_x,v_y,v_z)\) be true trunk linear velocity in the
  body frame;
\item
  \(\boldsymbol\omega=(\omega_x,\omega_y,\omega_z)\) be true trunk
  angular velocity;
\item
  \(\mathbf g_b=(g_x,g_y,g_z)\) be projected gravity;
\item
  \(\delta q_j=q_j-q_j^{\text{default}}\);
\item
  \(h_j\) be the head command, zero in Bridge.
\end{itemize}

The exact implemented terms and weights are:

\subsubsection{Linear velocity tracking}\label{linear-velocity-tracking}

\[
e_v=(c_x-v_x)^2+(c_y-v_y)^2+v_z^2,
\]

\[
r_{\text{lin}}=2\exp(-e_v/0.1).
\]

\subsubsection{Angular velocity
tracking}\label{angular-velocity-tracking}

\[
e_\omega=(c_\omega-\omega_z)^2+\omega_x^2+\omega_y^2,
\]

\[
r_{\text{ang}}=2\exp(-e_\omega/0.5).
\]

\subsubsection{Uprightness}\label{uprightness}

\[
r_{\text{upright}}
=2\exp[-(g_x^2+g_y^2)/0.05].
\]

\subsubsection{Default leg pose}\label{default-leg-pose}

For the ten leg joints,

\[
r_{\text{pose}}
=\exp\left[
-\frac{\sum_{j\in\text{legs}}\delta q_j^2}{0.5}
\right].
\]

\subsubsection{Head pose}\label{head-pose}

For the four head joints,

\[
r_{\text{head}}
=2\left[
\frac14\sum_{j\in\text{head}}
\exp\left(-\left(\frac{\delta q_j-h_j}{0.5}\right)^2\right)
\right].
\]

\subsubsection{Dense foot air time}\label{dense-foot-air-time}

Each unsupported foot accumulates air time in \texttt{0.02\ s}
increments. If the motion command magnitude exceeds \texttt{0.01}, a
foot contributes one while

\[
0.125<\tau_f<0.300\ \text{s}.
\]

Thus

\[
r_{\text{air}}
=3\sum_{f\in\{L,R\}}
\mathbb 1(0.125<\tau_f<0.300).
\]

Bridge overrides foot-contact detection so the start platform, plank,
and end platform all count as valid support for this inherited term.

\subsubsection{Action-rate penalty}\label{action-rate-penalty}

\[
r_{\Delta a}
=-w_{\Delta a}
\sum_{j=1}^{14}(a_{t,j}-a_{t-1,j})^2.
\]

The per-environment lifetime schedule is:

\begin{Shaded}
\begin{Highlighting}[]
\NormalTok{environment steps:  0     12,000  18,000  24,000  30,000  36,000}
\NormalTok{weight:             0.1   0.2     0.4     0.6     0.8     1.0}
\end{Highlighting}
\end{Shaded}

The recorded pilot used four environments and 32,768 total transitions,
or 8,192 control steps per environment. It therefore remained on weight
\texttt{0.1} throughout that pilot.

\subsubsection{Roll/pitch angular-velocity
penalty}\label{rollpitch-angular-velocity-penalty}

\[
r_{\omega xy}
=-0.05(\omega_x^2+\omega_y^2).
\]

Both penalty terms are non-positive by construction.

\subsubsection{Worked numerical example}\label{worked-numerical-example}

Suppose one step has:

\begin{Shaded}
\begin{Highlighting}[]
\NormalTok{command             = (0.25, 0, 0)}
\NormalTok{body velocity       = (0.15, 0, 0.05) m/s}
\NormalTok{angular velocity    = (0.10, 0.05, 0.10) rad/s}
\NormalTok{projected gravity   = (0.10, 0, approximately {-}0.995)}
\NormalTok{leg squared error   = 0.20}
\NormalTok{head errors         = (0.10, {-}0.10, 0, 0.20) rad}
\NormalTok{one foot air time   = 0.20 s; the other foot is supported}
\NormalTok{action delta sumsq  = 0.08}
\NormalTok{action{-}rate weight  = 0.1}
\end{Highlighting}
\end{Shaded}

Then:

{\def\LTcaptype{none} % do not increment counter
\begin{longtable}[]{@{}lr@{}}
\toprule\noalign{}
Term & Value \\
\midrule\noalign{}
\endhead
\bottomrule\noalign{}
\endlastfoot
\texttt{track\_\allowbreak{}lin\_\allowbreak{}vel} &
\texttt{1.764994} \\
\texttt{track\_\allowbreak{}ang\_\allowbreak{}vel} &
\texttt{1.911995} \\
\texttt{upright} & \texttt{1.637462} \\
\texttt{pose} & \texttt{0.670320} \\
\texttt{head\_pose} & \texttt{1.886861} \\
\texttt{feet\_\allowbreak{}air\_\allowbreak{}time} &
\texttt{3.000000} \\
\texttt{action\_\allowbreak{}rate\_\allowbreak{}penalty} &
\texttt{-0.008000} \\
\texttt{ang\_\allowbreak{}vel\_\allowbreak{}xy\_\allowbreak{}penalty} &
\texttt{-0.000625} \\
\textbf{total} & \textbf{\texttt{10.863007}} \\
\end{longtable}
}

This high reward still says nothing about crossing. The robot could earn
it while walking on an early assisted spawn or before reaching the end
platform.

\subsection{Physics-only curriculum}\label{physics-only-curriculum}

The curriculum has five stages:

{\def\LTcaptype{none} % do not increment counter
\begin{longtable}[]{@{}
  >{\raggedleft\arraybackslash}p{(\linewidth - 6\tabcolsep) * \real{0.2667}}
  >{\raggedleft\arraybackslash}p{(\linewidth - 6\tabcolsep) * \real{0.2667}}
  >{\raggedleft\arraybackslash}p{(\linewidth - 6\tabcolsep) * \real{0.2667}}
  >{\raggedright\arraybackslash}p{(\linewidth - 6\tabcolsep) * \real{0.2000}}@{}}
\toprule\noalign{}
\begin{minipage}[b]{\linewidth}\raggedleft
Stage
\end{minipage} & \begin{minipage}[b]{\linewidth}\raggedleft
Assistance \(\alpha\)
\end{minipage} & \begin{minipage}[b]{\linewidth}\raggedleft
Robot spawn \texttt{x}
\end{minipage} & \begin{minipage}[b]{\linewidth}\raggedright
Interpretation
\end{minipage} \\
\midrule\noalign{}
\endhead
\bottomrule\noalign{}
\endlastfoot
0 & \texttt{1.0} & \texttt{0.00\ m} & starts near plank center with full
plank stabilization \\
1 & \texttt{0.6} & \texttt{-0.30\ m} & longer plank approach, reduced
stabilization \\
2 & \texttt{0.3} & \texttt{-0.55\ m} & begins at the
plank/start-platform boundary \\
3 & \texttt{0.1} & \texttt{-0.75\ m} & begins on the launch platform \\
4 & \texttt{0.0} & \texttt{-0.90\ m} & exact unassisted evaluation
spawn \\
\end{longtable}
}

Assistance acts only on the plank. For plank displacement
\(\Delta\mathbf p=\mathbf p-\mathbf p_0\) and linear velocity
\(\mathbf v\):

\[
\mathbf F_{xy}
=\alpha(-45\Delta\mathbf p_{xy}-4\mathbf v_{xy}),
\]

\[
F_z=\alpha(-25\Delta p_z-2.5v_z).
\]

For plank up-axis \(\mathbf z_p\), world up \(\mathbf z_w\), and world
angular velocity \(\boldsymbol\omega_w\):

\[
\boldsymbol\tau
=\alpha\left[
3.5(\mathbf z_p\times\mathbf z_w)
-0.18\boldsymbol\omega_w
\right].
\]

At \texttt{alpha=0}, the applied force array is exactly zero. The
curriculum never pushes the robot forward, changes an action, modifies a
reward coefficient, or grants crossing credit.

After an episode:

\begin{Shaded}
\begin{Highlighting}[]
\ControlFlowTok{if}\NormalTok{ crossed:}
\NormalTok{    stage }\OperatorTok{=} \BuiltInTok{min}\NormalTok{(stage }\OperatorTok{+} \DecValTok{1}\NormalTok{, }\DecValTok{4}\NormalTok{)}
\ControlFlowTok{elif}\NormalTok{ elapsed\_s }\OperatorTok{\textless{}} \FloatTok{2.0}\NormalTok{:}
\NormalTok{    stage }\OperatorTok{=} \BuiltInTok{max}\NormalTok{(stage }\OperatorTok{{-}} \DecValTok{1}\NormalTok{, }\DecValTok{0}\NormalTok{)}
\CommentTok{\# Otherwise keep the same stage.}
\end{Highlighting}
\end{Shaded}

This follows the playbook's exploration rule: if the target state never
occurs, change physics or starts so it can be sampled; do not invent a
sparse jackpot. The easiest rung did produce one contact-valid crossing
in a diagnostic \texttt{alpha\_\allowbreak{}walking.\allowbreak{}onnx}
probe, but that attempt later fell. It is exploration evidence, not a
strict skill pass.

\subsection{Transfer from the shipped walking
actor}\label{transfer-from-the-shipped-walking-actor}

A full Bridge run without \texttt{-\/-init-from} automatically creates
an initializer from
\texttt{microduc\allowbreak{}k/\allowbreak{}policies\allowbreak{}/\allowbreak{}alpha\_\allowbreak{}walking.\allowbreak{}onnx}
when \texttt{total\_timesteps\ \textgreater{}\ 4}.

The bootstrap validates:

\begin{itemize}
\tightlist
\item
  input name \texttt{obs}, shape \texttt{(1,61)};
\item
  output name \texttt{actions}, shape \texttt{(1,14)};
\item
  exact ONNX node sequence;
\item
  expected tensor names, shapes, and \texttt{float32} types;
\item
  finite normalizer mean and positive standard deviation; and
\item
  unchanged source hash during transfer.
\end{itemize}

It creates a new MLX actor-critic. The four actor linear layers are
copied from the ONNX actor. The observation mean and variance are copied
from the ONNX normalizer. The critic is freshly seeded.
\texttt{actor\_\allowbreak{}log\_\allowbreak{}std} is newly initialized
from \texttt{initial\_std=0.03}. The training caller creates a new Adam
optimizer.

The actor is \textbf{not frozen}. It is a warm start and remains
trainable. The observation normalization statistics are frozen so the
transferred actor keeps the coordinate system in which the walking
weights were trained.

Before saving, the code compares the transferred MLX actor with the
source ONNX on 32 representative observations. The allowed maximum
absolute action error is \texttt{1e-4}. The recorded
\texttt{bridge-\allowbreak{}studio-\allowbreak{}02} initializer measured
\texttt{1.0728836e-6}.

The resulting continuation is accurately summarized as:

\begin{Shaded}
\begin{Highlighting}[]
\NormalTok{actor mean weights:       transferred, then trainable}
\NormalTok{observation normalizer:   transferred and frozen}
\NormalTok{critic weights:           fresh seeded MLX initialization}
\NormalTok{actor exploration std:    fresh, initial value 0.03}
\NormalTok{optimizer state:          fresh Adam}
\end{Highlighting}
\end{Shaded}

It is not an exact PPO resume from the upstream walking run.

\subsection{Reproduce the workflow}\label{reproduce-the-workflow}

Run these commands from the workspace root. They use the existing
\texttt{rlx/\allowbreak{}.venv-\allowbreak{}microduc\allowbreak{}k}
environment through
\texttt{scripts/\allowbreak{}train-\allowbreak{}bridge.s\allowbreak{}h}.

\subsubsection{Build and inspect the simulated
scene}\label{build-and-inspect-the-simulated-scene}

The Studio parser has no \texttt{build} subcommand.
\texttt{bridge\_model()} builds the MuJoCo model at runtime. This
checked command constructs one environment and verifies the public
dimensions:

\begin{Shaded}
\begin{Highlighting}[]
\ExtensionTok{rlx/.venv{-}microduck/bin/python} \AttributeTok{{-}} \OperatorTok{\textless{}\textless{}\textquotesingle{}PY\textquotesingle{}}
\StringTok{from rlx.environments.bridge import BridgeEnv}

\StringTok{env = BridgeEnv(}
\StringTok{    actuator="xml",}
\StringTok{    curriculum=False,}
\StringTok{    obs\_noise=False,}
\StringTok{    domain\_rand=False,}
\StringTok{    action\_delay=False,}
\StringTok{    random\_yaw=False,}
\StringTok{    seed=7,}
\StringTok{)}
\StringTok{obs, info = env.reset(seed=7)}
\StringTok{print(obs.shape, env.action\_space.shape, env.model.ntendon)}
\StringTok{print(round(info["bridge\_position"][2], 3), info["assistance"])}
\StringTok{env.close()}
\OperatorTok{PY}
\end{Highlighting}
\end{Shaded}

Expected output:

\begin{Shaded}
\begin{Highlighting}[]
\NormalTok{(61,) (14,) 4}
\NormalTok{0.225 0.0}
\end{Highlighting}
\end{Shaded}

\subsubsection{Train the bounded pilot
recipe}\label{train-the-bounded-pilot-recipe}

\begin{Shaded}
\begin{Highlighting}[]
\ExtensionTok{./scripts/train{-}bridge.sh}\NormalTok{ train }\DataTypeTok{\textbackslash{}}
  \AttributeTok{{-}{-}bridge{-}curriculum} \DataTypeTok{\textbackslash{}}
  \AttributeTok{{-}{-}total{-}timesteps}\NormalTok{ 32768 }\DataTypeTok{\textbackslash{}}
  \AttributeTok{{-}{-}num{-}envs}\NormalTok{ 4 }\DataTypeTok{\textbackslash{}}
  \AttributeTok{{-}{-}num{-}steps}\NormalTok{ 128 }\DataTypeTok{\textbackslash{}}
  \AttributeTok{{-}{-}num{-}minibatches}\NormalTok{ 4 }\DataTypeTok{\textbackslash{}}
  \AttributeTok{{-}{-}update{-}epochs}\NormalTok{ 2 }\DataTypeTok{\textbackslash{}}
  \AttributeTok{{-}{-}learning{-}rate}\NormalTok{ 1e{-}5 }\DataTypeTok{\textbackslash{}}
  \AttributeTok{{-}{-}gamma}\NormalTok{ 0.99 }\DataTypeTok{\textbackslash{}}
  \AttributeTok{{-}{-}gae{-}lambda}\NormalTok{ 0.95 }\DataTypeTok{\textbackslash{}}
  \AttributeTok{{-}{-}clip{-}coefficient}\NormalTok{ 0.1 }\DataTypeTok{\textbackslash{}}
  \AttributeTok{{-}{-}entropy{-}coefficient}\NormalTok{ 0 }\DataTypeTok{\textbackslash{}}
  \AttributeTok{{-}{-}max{-}grad{-}norm}\NormalTok{ 0.5 }\DataTypeTok{\textbackslash{}}
  \AttributeTok{{-}{-}initial{-}std}\NormalTok{ 0.03 }\DataTypeTok{\textbackslash{}}
  \AttributeTok{{-}{-}backend}\NormalTok{ dummy }\DataTypeTok{\textbackslash{}}
  \AttributeTok{{-}{-}actuator}\NormalTok{ xml }\DataTypeTok{\textbackslash{}}
  \AttributeTok{{-}{-}max{-}episode{-}s}\NormalTok{ 20 }\DataTypeTok{\textbackslash{}}
  \AttributeTok{{-}{-}seed}\NormalTok{ 7 }\DataTypeTok{\textbackslash{}}
  \AttributeTok{{-}{-}no{-}normalize{-}rewards} \DataTypeTok{\textbackslash{}}
  \AttributeTok{{-}{-}no{-}domain{-}rand} \DataTypeTok{\textbackslash{}}
  \AttributeTok{{-}{-}no{-}obs{-}noise} \DataTypeTok{\textbackslash{}}
  \AttributeTok{{-}{-}no{-}action{-}delay} \DataTypeTok{\textbackslash{}}
  \AttributeTok{{-}{-}no{-}random{-}yaw} \DataTypeTok{\textbackslash{}}
  \AttributeTok{{-}{-}output{-}dir}\NormalTok{ rlx/runs/studio/bridge/my{-}bridge}
\end{Highlighting}
\end{Shaded}

Training automatically writes
\texttt{bridge.s\allowbreak{}afetenso\allowbreak{}rs}, its JSON sidecar,
\texttt{bridge.onnx}, an initial snapshot, and
\texttt{training\allowbreak{}-\allowbreak{}metrics.\allowbreak{}jsonl}.
Periodic snapshots additionally depend on
\texttt{-\allowbreak{}-\allowbreak{}checkpoi\allowbreak{}nt-\allowbreak{}interval\allowbreak{}}.

The batch accounting is exact:

\[
4\text{ envs}\times128\text{ steps}=512
\text{ transitions per rollout}.
\]

\[
32768/512=64\text{ rollouts}.
\]

Each rollout uses four minibatches and two epochs:

\[
64\times4\times2=512
\text{ optimizer minibatch steps}.
\]

\subsubsection{Export explicitly}\label{export-explicitly}

Training exports ONNX by default. To rebuild only the deterministic ONNX
actor:

\begin{Shaded}
\begin{Highlighting}[]
\ExtensionTok{./scripts/train{-}bridge.sh}\NormalTok{ export }\DataTypeTok{\textbackslash{}}
  \AttributeTok{{-}{-}checkpoint}\NormalTok{ rlx/runs/studio/bridge/my{-}bridge/bridge.safetensors }\DataTypeTok{\textbackslash{}}
  \AttributeTok{{-}{-}output}\NormalTok{ rlx/runs/studio/bridge/my{-}bridge/bridge.onnx}
\end{Highlighting}
\end{Shaded}

The export contains the observation normalizer and actor mean. It does
not export the critic or stochastic action sampling.

\subsubsection{Evaluate the strict skill
gate}\label{evaluate-the-strict-skill-gate}

\begin{Shaded}
\begin{Highlighting}[]
\ExtensionTok{./scripts/train{-}bridge.sh}\NormalTok{ eval }\DataTypeTok{\textbackslash{}}
  \AttributeTok{{-}{-}backend}\NormalTok{ dummy }\DataTypeTok{\textbackslash{}}
  \AttributeTok{{-}{-}num{-}envs}\NormalTok{ 3 }\DataTypeTok{\textbackslash{}}
  \AttributeTok{{-}{-}seed}\NormalTok{ 101 }\DataTypeTok{\textbackslash{}}
  \AttributeTok{{-}{-}policy}\NormalTok{ rlx/runs/studio/bridge/my{-}bridge/bridge.onnx }\DataTypeTok{\textbackslash{}}
  \AttributeTok{{-}{-}evaluation{-}mode}\NormalTok{ skill }\DataTypeTok{\textbackslash{}}
  \AttributeTok{{-}{-}eval{-}steps}\NormalTok{ 1000 }\DataTypeTok{\textbackslash{}}
  \AttributeTok{{-}{-}max{-}episode{-}s}\NormalTok{ 20 }\DataTypeTok{\textbackslash{}}
  \AttributeTok{{-}{-}actuator}\NormalTok{ xml }\DataTypeTok{\textbackslash{}}
  \AttributeTok{{-}{-}no{-}domain{-}rand} \DataTypeTok{\textbackslash{}}
  \AttributeTok{{-}{-}no{-}obs{-}noise} \DataTypeTok{\textbackslash{}}
  \AttributeTok{{-}{-}no{-}action{-}delay} \DataTypeTok{\textbackslash{}}
  \AttributeTok{{-}{-}no{-}random{-}yaw}
\end{Highlighting}
\end{Shaded}

Bridge skill evaluation rejects the \texttt{fork} backend because it
needs control-step metrics. Use \texttt{dummy} or \texttt{subproc}.
Skill mode also requires complete episode-horizon multiples and at least
1,000 steps per 20-second episode. The command exits with status 2 when
the skill result fails, even if the inference pipeline itself stayed
finite.

\subsubsection{Render and look}\label{render-and-look}

\begin{Shaded}
\begin{Highlighting}[]
\ExtensionTok{./scripts/train{-}bridge.sh}\NormalTok{ render }\DataTypeTok{\textbackslash{}}
  \AttributeTok{{-}{-}policy}\NormalTok{ rlx/runs/studio/bridge/my{-}bridge/bridge.onnx }\DataTypeTok{\textbackslash{}}
  \AttributeTok{{-}{-}output}\NormalTok{ rlx/runs/studio/bridge/my{-}bridge/render }\DataTypeTok{\textbackslash{}}
  \AttributeTok{{-}{-}render{-}seconds}\NormalTok{ 20 }\DataTypeTok{\textbackslash{}}
  \AttributeTok{{-}{-}episodes}\NormalTok{ 1 }\DataTypeTok{\textbackslash{}}
  \AttributeTok{{-}{-}seed}\NormalTok{ 101 }\DataTypeTok{\textbackslash{}}
  \AttributeTok{{-}{-}actuator}\NormalTok{ xml }\DataTypeTok{\textbackslash{}}
  \AttributeTok{{-}{-}camera}\NormalTok{ side }\DataTypeTok{\textbackslash{}}
  \AttributeTok{{-}{-}width}\NormalTok{ 640 }\DataTypeTok{\textbackslash{}}
  \AttributeTok{{-}{-}height}\NormalTok{ 360 }\DataTypeTok{\textbackslash{}}
  \AttributeTok{{-}{-}fps}\NormalTok{ 25 }\DataTypeTok{\textbackslash{}}
  \AttributeTok{{-}{-}sheet{-}frames}\NormalTok{ 12}
\end{Highlighting}
\end{Shaded}

Rendering forces domain randomization, observation noise, action delay,
random yaw, and bridge assistance off. If the robot falls before 20
seconds, the renderer resets and continues recording; the overlay
reports attempt and reset counts. A 20-second video may therefore
contain several failed attempts rather than one complete episode.

\subsection{The strict full-20-second crossing
rule}\label{the-strict-full-20-second-crossing-rule}

Environment-level
\texttt{crossing\allowbreak{}\_\allowbreak{}accepted\allowbreak{}}
becomes sticky only when all of these are true at one control step:

\begin{enumerate}
\def\labelenumi{\arabic{enumi}.}
\tightlist
\item
  trunk \texttt{x\ \textgreater{}=\ END\_PLATFORM\_X-0.08\ =\ 0.77\ m};
\item
  both feet contact the end platform;
\item
  neither foot contacts the ground;
\item
  no robot body contacts the ground;
\item
  no non-foot body supports the robot; and
\item
  at least one earlier foot contact with the suspended plank was
  observed.
\end{enumerate}

The separate \texttt{BridgeEv\allowbreak{}aluation\allowbreak{}} then
requires more:

\begin{itemize}
\tightlist
\item
  the episode starts with launch support near progress zero;
\item
  maximum progress is at least \texttt{1.60\ m};
\item
  assistance is zero on every measured step;
\item
  no floor or non-foot support occurs;
\item
  the contact-validated crossing flag is observed;
\item
  the episode continues to a 1,000-step truncation without termination;
  and
\item
  every evaluated lane and episode passes.
\end{itemize}

Progress is measured from the final launch spawn:

\[
p=x_{\text{trunk}}-(-0.90)=x_{\text{trunk}}+0.90.
\]

Thus the \texttt{1.60\ m} progress threshold alone corresponds to
\texttt{x\_trunk\ \textgreater{}=\ 0.70\ m}. It is deliberately
insufficient: the stricter crossing flag still requires
\texttt{x\ \textgreater{}=\ 0.77\ m}, two end-platform feet, prior plank
contact, and clean support history.

Crossing early and falling later fails because the required 20-second
horizon did not complete. Teleporting directly to the end platform fails
because plank contact history is missing. Walking on the floor fails
even if the final \texttt{x} coordinate is large.

\subsection{What the current pilot actually
achieved}\label{what-the-current-pilot-actually-achieved}

The latest recorded final pilot is
\texttt{rlx/\allowbreak{}runs/\allowbreak{}studio/\allowbreak{}bridge/\allowbreak{}bridge-\allowbreak{}studio-\allowbreak{}02/\allowbreak{}},
produced on September 10, 2026. Its checkpoint records:

{\def\LTcaptype{none} % do not increment counter
\begin{longtable}[]{@{}lr@{}}
\toprule\noalign{}
Item & Value \\
\midrule\noalign{}
\endhead
\bottomrule\noalign{}
\endlastfoot
New PPO transitions & \texttt{32,768} \\
Environments & \texttt{4} \\
Rollout length & \texttt{128} \\
Minibatches & \texttt{4} \\
Update epochs & \texttt{2} \\
Seed & \texttt{7} \\
Actuator & XML \\
Observation normalizer & transferred and frozen \\
Reward normalization & off \\
Curriculum & physics/spawn only \\
\end{longtable}
}

The deterministic skill evaluation used three lanes seeded 101, 102, and
103, 1,000 control steps per lane, no randomization/noise/delay/yaw, XML
actuators, zero assistance, and the destination-stop commander.

The report contains 11 attempts because terminated lanes reset while the
3,000-transition evaluation continued. None completed the full 20-second
horizon. None crossed. Eight attempts reached the plank. The best
attempt made \texttt{1.174915875\ m} of progress from the launch spawn,
about

\[
1.60-1.174915875=0.425084125\ \text{m}
\]

short of the evaluator's progress threshold. No attempt recorded feet on
the end platform. One attempt contacted the floor. The pipeline was
finite, but the skill status was \texttt{failed}.

The correct summary is:

\begin{quote}
The pilot learned or retained partial approach and plank-walking
behavior, but no successful unassisted crossing was observed.
\end{quote}

It is incorrect to say that 32,768 steps ``completed the curriculum,''
that a high return proves crossing, or that the render shows one
continuous 20-second crossing. The retained video explicitly includes
resets and partial attempts.

\subsection{Exercises}\label{exercises-2}

\subsubsection{Exercise 1: Cable
geometry}\label{exercise-1-cable-geometry}

Using anchor \texttt{(0.45,0.17,0.77)} and nominal attachment
\texttt{(0.45,0.057,0.225)}, calculate the cable length.

\subsubsection{Exercise 2: Commander
correction}\label{exercise-2-commander-correction}

Assume the trunk is aligned with world \texttt{+x}, at
\texttt{y=0.08\ m}, and has not reached the destination. What command is
published?

\subsubsection{Exercise 3: Observation
privilege}\label{exercise-3-observation-privilege}

The plank rolls five degrees while the robot state and command remain
exactly unchanged. Which observation indices change because of the plank
roll?

\subsubsection{Exercise 4: Curriculum
transition}\label{exercise-4-curriculum-transition}

An environment is at stage 3. Determine the next stage after:

\begin{enumerate}
\def\labelenumi{\arabic{enumi}.}
\tightlist
\item
  a clean crossing at 8 seconds;
\item
  a failure at 1.5 seconds;
\item
  a failure at 6 seconds.
\end{enumerate}

\subsubsection{Exercise 5: PPO
accounting}\label{exercise-5-ppo-accounting}

For 32,768 transitions, four environments, 128 rollout steps, four
minibatches, and two epochs, calculate rollouts, transitions per
minibatch, and optimizer minibatch steps.

\subsubsection{Exercise 6: Honest
crossing}\label{exercise-6-honest-crossing}

For each case, state pass or fail:

\begin{enumerate}
\def\labelenumi{\arabic{enumi}.}
\tightlist
\item
  progress 1.65 m, crossed flag true at step 700, no invalid contacts,
  then truncation at step 1,000;
\item
  the same crossing, followed by a fall at step 850;
\item
  two feet on the end platform with no earlier plank contact;
\item
  progress 1.59 m with a true crossing flag and full horizon;
\item
  a perfect 20-second episode with assistance \texttt{0.1} for one step.
\end{enumerate}

\subsubsection{Exercise 7: Result
interpretation}\label{exercise-7-result-interpretation}

The policy earns a larger mean return than its initializer but has
\texttt{bridge\_crossed.max\ =\ 0}. What may be claimed?

\subsection{Answers}\label{answers-2}

\subsubsection{Answer 1}\label{answer-1}

The longitudinal difference is zero, lateral difference is
\texttt{0.113\ m}, and vertical difference is \texttt{0.545\ m}:

\[
L=\sqrt{0^2+0.113^2+0.545^2}=0.556591412079\ \text{m}.
\]

\subsubsection{Answer 2}\label{answer-2}

With heading \texttt{(1,0,0)}, side \texttt{(0,1,0)}, and
\texttt{y=0.08},

\[
d_y=\operatorname{clip}(-1.2(0.08),-0.12,0.12)=-0.096.
\]

The yaw error is zero, so the command is

\begin{Shaded}
\begin{Highlighting}[]
\NormalTok{(0.25, {-}0.096, 0.0)}
\end{Highlighting}
\end{Shaded}

before floating-point conversion.

\subsubsection{Answer 3}\label{answer-3}

None change directly. There is no plank-state block in the 61-vector.
Future robot motion caused by contact may change IMU, joints, or
previous action, but the isolated plank pose is privileged.

\subsubsection{Answer 4}\label{answer-4}

\begin{enumerate}
\def\labelenumi{\arabic{enumi}.}
\tightlist
\item
  crossing: stage 4;
\item
  failure before 2 seconds: stage 2;
\item
  later failure: remain at stage 3.
\end{enumerate}

\subsubsection{Answer 5}\label{answer-5}

There are \texttt{4*128=512} transitions per rollout and
\texttt{32768/512=64} rollouts. Each minibatch contains
\texttt{512/4=128} transitions. Two epochs over four minibatches give
\texttt{8} optimizer steps per rollout and \texttt{64*8=512} total.

\subsubsection{Answer 6}\label{answer-6}

\begin{enumerate}
\def\labelenumi{\arabic{enumi}.}
\tightlist
\item
  pass, assuming launch support and metric coverage are valid;
\item
  fail: the episode terminated before the required horizon;
\item
  fail: suspended-plank foot-contact history is required;
\item
  fail: progress is below \texttt{1.60\ m};
\item
  fail: any assisted step invalidates unassisted evaluation.
\end{enumerate}

\subsubsection{Answer 7}\label{answer-7}

You may claim improved return under the inherited locomotion objective,
if the comparison is otherwise controlled. You may not claim a bridge
crossing or a mastered bridge skill. Deterministic physical evaluation
remains the deciding evidence.

\section{Read the evidence, not the
celebration}\label{read-the-evidence-not-the-celebration}

\subsection{The evidence chain}\label{the-evidence-chain}

\begin{figure}
\centering
\includegraphics[width=6.2in,height=5.5in,alt={Evidence-preparation schematic: retained run records pass through schema checks and CSV extraction, then become observed plots or copied frames, and finally hash-bound book assets. This is not a measured rollout.}]{assets/diagrams/schematic-evidence-pipeline.png}
\caption{Evidence-preparation schematic: retained run records pass
through schema checks and CSV extraction, then become observed plots or
copied frames, and finally hash-bound book assets. This is not a
measured rollout.}
\end{figure}

A training program can execute perfectly while learning the wrong
behavior. A reward can improve while balance deteriorates. An exported
model can differ from its training actor. A viewer can show the wrong
scene. Therefore the final claim must pass several independent checks:
correct model and inputs, real parameter updates, export parity, honest
unassisted evaluation, and actual visual inspection.

The word \textbf{independent} matters. If you define ``success'' as high
reward and then show the same reward as proof, you have checked one
definition twice. Contact validity, no-fall duration, directional
displacement and complete bridge crossing are separate observable tests
of the requested physical task.

\subsection{Read real reward and loss
curves}\label{read-real-reward-and-loss-curves}

\begin{figure}
\centering
\includegraphics[width=6.2in,height=5.5in,alt={Real recorded basketball training losses, with an explicit missing-history panel for reward. The two policies are independent warm starts, not a single continuous training run.}]{assets/evidence/basketball-training-curves.png}
\caption{Real recorded basketball training losses, with an explicit
missing-history panel for reward. The two policies are independent warm
starts, not a single continuous training run.}
\end{figure}

The fixed-command basketball \texttt{summary.json} has 100 update
records; the mixed-command run has 150. They include total loss, policy
loss, value loss, entropy, approximate KL and gradient norm.
\textbf{They do not include historical per-update reward or episode
return.} This is an evidence gap, not permission to draw an attractive
rising curve. The missing-reward panel is intentional.

Both runs begin from the supplied actor with fresh critics and
optimizers. Plotting the fixed run followed by the mixed run on one
cumulative axis would falsely imply continuation. The correct horizontal
axes end at 102,400 and 153,600 transitions separately. Each
fixed-command update corresponds to \(8\times128=1024\) new transitions.

The separate basketball reward diagnostic supplied with this edition
records deterministic per-step rewards from the existing policy. It is a
\textbf{new diagnostic trajectory}, not a recovered historical training
curve and not evidence of further optimization. Its CSV and receipt
identify the policy, conditions and actual horizon. Do not splice it
into the old training history.

\begin{figure}
\centering
\includegraphics[width=6.2in,height=5.5in,alt={New deterministic basketball reward diagnostic, not historical PPO training reward. The source policy is unchanged.}]{assets/evidence/basketball-reward-trace.png}
\caption{New deterministic basketball reward diagnostic, not historical
PPO training reward. The source policy is unchanged.}
\end{figure}

\begin{figure}
\centering
\includegraphics[width=6.2in,height=5.5in,alt={Real bridge pilot telemetry: collection mean reward and completed-episode returns are distinct from PPO loss components.}]{assets/evidence/bridge-training-curves.png}
\caption{Real bridge pilot telemetry: collection mean reward and
completed-episode returns are distinct from PPO loss components.}
\end{figure}

The bridge log contains collection events, update events and
completed-episode events. Collection \texttt{mean\_reward} is a
per-transition average.
\texttt{mean\_\allowbreak{}raw\_\allowbreak{}return} is the reward
summed over an episode, then averaged over completed episodes in that
event. An episode that lasts longer can have a larger return even
without better per-step behavior. An update \texttt{mean\_loss} averages
optimizer minibatch objectives. These three numbers do not have
interchangeable meanings.

The bridge's total/value-loss panel uses a logarithmic vertical axis.
Equal vertical distances represent equal multiplicative changes, not
equal additive changes. A reduction from 1,000 to 100 has the same
log-axis distance as 100 to 10. Negative policy losses cannot be placed
on the same ordinary log axis and belong in their own panel.

Small policy loss is not proof of a solved policy. Advantage
normalization deliberately centres advantages near zero; ratios close to
one can make their average surrogate close to zero. A large value loss
early in these runs is also unsurprising because the actor is mature but
the critic is newly initialized. Entropy can be negative for continuous
densities, as explained in the PPO chapter. A target KL is not a
guarantee that a single optimizer step never overshoots.

\textbf{A good curve caption names five things:} the run, exact recorded
field, horizontal-axis unit, any transformation/smoothing, and the claim
it does not prove. The book's CSV extracts preserve the unsmoothed
numbers. You can regenerate figures without running training.

\subsection{Deterministic evaluation: what
passed?}\label{deterministic-evaluation-what-passed}

\begin{figure}
\centering
\includegraphics[width=6.2in,height=5.5in,alt={Recorded evaluation outcomes, not training scores. Basketball balance succeeds in the tested trials while controlled rolling and accepted bridge crossing do not.}]{assets/evidence/evaluation-outcomes.png}
\caption{Recorded evaluation outcomes, not training scores. Basketball
balance succeeds in the tested trials while controlled rolling and
accepted bridge crossing do not.}
\end{figure}

{\def\LTcaptype{none} % do not increment counter
\begin{longtable}[]{@{}
  >{\raggedright\arraybackslash}p{(\linewidth - 4\tabcolsep) * \real{0.3333}}
  >{\raggedright\arraybackslash}p{(\linewidth - 4\tabcolsep) * \real{0.3333}}
  >{\raggedright\arraybackslash}p{(\linewidth - 4\tabcolsep) * \real{0.3333}}@{}}
\toprule\noalign{}
\begin{minipage}[b]{\linewidth}\raggedright
Recorded experiment
\end{minipage} & \begin{minipage}[b]{\linewidth}\raggedright
Result
\end{minipage} & \begin{minipage}[b]{\linewidth}\raggedright
Correct interpretation
\end{minipage} \\
\midrule\noalign{}
\endhead
\bottomrule\noalign{}
\endlastfoot
Basketball fixed adaptation, six 60-second trials & 6/6 sustained
survival & Verified local nominal-condition free-ball balance. \\
Same candidate, three forward-command trials & 0/3 rolling passes &
Movement is not reliable commanded rolling. \\
Basketball mixed adaptation & 6/6 survival, 0/3 rolling passes & A
separate candidate, not further steps of the fixed candidate. \\
Bridge final pilot & 0 accepted crossings & Training/export work; the
requested physical task remains unsolved. \\
Seven-case API/Studio checks & Seven scenarios reachable and rendered &
Product integration, not seven skill certificates. \\
\end{longtable}
}

The direct basketball report uses seeds 101, 202 and 303. The later
Studio evaluation uses 101, 102 and 103. Both contain zero/forward
command pairs, but they are separate evidence sets. Do not relabel one
report's seeds using another report's screenshot. The minimum
direct-trial foot-ball support coverage was about 99.93\%, and maximum
tilt about 4.74 degrees; these describe that report's candidate and
conditions, not an untested distribution.

\textbf{Historical source boundary.} The retained direct basketball
report predates the current source snapshot: its evaluator/environment
hashes begin \texttt{fcfa65}/\texttt{2eda2a}, while this edition's
supplied evaluator/environment begin \texttt{bbe735}/\texttt{6bff54}.
The old report is preserved, not presented as a fresh execution of these
exact source bytes. The new 10-second reward diagnostic binds the
current environment hash, but is not a replacement for the six-trial
60-second acceptance battery. To test the current snapshot, run the
documented evaluator into a new directory and retain its new hashes.
Protocol reproduction and byte-identical historical reconstruction are
different claims.

The direct fixed run's forward tracking mean absolute error is about
0.138 m/s against a 0.15 m/s command. The mixed run's is about 0.160
m/s. These are poor command-following results even though balance is
good. A small fixed-versus-source change near 0.140 m/s does not
establish a robust improvement.

For bridge, 1.174916 m maximum progress is measured relative to the
launch position. It is \textbf{not} ``1.174916 m of a 1.10 m bridge,
therefore done.'' The duck must leave its platform, make valid plank
contact, reach the destination with both feet, avoid forbidden support
and survive the required full episode. The final evaluation has three
vector lanes and 3,000 total transitions, with resets producing multiple
attempts and trailing fragments. Those fragments are not extra
successful 20-second trials.

Three or six successes in nominal tests do not establish broad
reliability. The current tests do not exhaust friction, mass, cable
stiffness, command directions, pushes, noisy sensors, actuator
temperature, real manufacturing variation or hardware conditions.
General steering coverage is explicitly absent for lateral/yaw commands.
Report the tested set, not an invented population success percentage.

\subsection{Visual review: inspect consecutive
frames}\label{visual-review-inspect-consecutive-frames}

\begin{figure}
\centering
\includegraphics[width=6.2in,height=5.5in,alt={Measured zero-command basketball rollout, seed 101. This is evidence of sustained balance in the recorded trial, not commanded locomotion.}]{assets/evidence/observed-basketball-zero-command-contact-sheet.png}
\caption{Measured zero-command basketball rollout, seed 101. This is
evidence of sustained balance in the recorded trial, not commanded
locomotion.}
\end{figure}

Check more than whether the duck is above the orange ball. Is its trunk
upright? Are feet actually supporting it? Is the ground or another body
part supporting the robot? Are the textures rotating consistently with
physical motion? Did the clock continue without a reset? Was the camera
following so tightly that drift became invisible?

\begin{figure}
\centering
\includegraphics[width=6.2in,height=5.5in,alt={Measured forward-command attempt. The duck balances and the ball moves, but the directional rolling gate fails.}]{assets/evidence/observed-basketball-forward-command-contact-sheet.png}
\caption{Measured forward-command attempt. The duck balances and the
ball moves, but the directional rolling gate fails.}
\end{figure}

A contact sheet samples a video. It can miss fast foot chatter or an
event between frames. When uncertain, inspect the MP4 at the original
frame rate or render at the 50 Hz control rate. Use the numerical
contact/termination record alongside images. ``Looks successful'' and
``metrics look successful'' should agree before accepting a skill.

\begin{figure}
\centering
\includegraphics[width=6.2in,height=5.5in,alt={Bridge diagnostic sequence from the final pilot. The 20-second video contains resets; it is not one uninterrupted successful crossing.}]{assets/evidence/observed-bridge-rollout-contact-sheet.png}
\caption{Bridge diagnostic sequence from the final pilot. The 20-second
video contains resets; it is not one uninterrupted successful crossing.}
\end{figure}

The retained bridge video spans 20 seconds and contains two resets. A
later attempt reaches partway onto the plank. A reset returns the robot
to a fresh spawn; it cannot be hidden and counted as continuous
progress. The book includes this failed sequence because seeing where a
policy loses contact is useful for curriculum design.

Videos remain in the workspace rather than being embedded in the PDF:

\begin{itemize}
\tightlist
\item
  \texttt{docs/\allowbreak{}basketba\allowbreak{}ll-\allowbreak{}showcase\allowbreak{}/\allowbreak{}evidence\allowbreak{}/\allowbreak{}render/\allowbreak{}cmd-\allowbreak{}p0\_\allowbreak{}000-\allowbreak{}p0\_\allowbreak{}000-\allowbreak{}p0\_\allowbreak{}000/\allowbreak{}seed-\allowbreak{}101/\allowbreak{}rollout.\allowbreak{}mp4}
\item
  \texttt{docs/\allowbreak{}basketba\allowbreak{}ll-\allowbreak{}showcase\allowbreak{}/\allowbreak{}evidence\allowbreak{}/\allowbreak{}render/\allowbreak{}cmd-\allowbreak{}p0\_\allowbreak{}150-\allowbreak{}p0\_\allowbreak{}000-\allowbreak{}p0\_\allowbreak{}000/\allowbreak{}seed-\allowbreak{}101/\allowbreak{}rollout.\allowbreak{}mp4}
\item
  \texttt{rlx/\allowbreak{}runs/\allowbreak{}studio/\allowbreak{}bridge/\allowbreak{}bridge-\allowbreak{}studio-\allowbreak{}02/\allowbreak{}render/\allowbreak{}ep0.mp4}
\end{itemize}

\subsection{Run both scenes in Duck
Viewer}\label{run-both-scenes-in-duck-viewer}

\begin{figure}
\centering
\includegraphics[width=6.2in,height=5.5in,alt={Observed seven-case Studio screen. Card availability is separate from learned-skill acceptance.}]{assets/evidence/observed-seven-cases.png}
\caption{Observed seven-case Studio screen. Card availability is
separate from learned-skill acceptance.}
\end{figure}

The seven cases are \textbf{Dance, Swing, Running, Stilts, Backflip,
Basketball, Bridge}. The basketball card must say balance-only when
rolling is unaccepted. The bridge card must not turn a pilot video into
an accepted-crossing preview.

Start a new instance only if there is no appropriate existing service.
Keep each server in its own terminal; these are long-running processes.
Use available ports instead of stopping someone else's server.

\begin{Shaded}
\begin{Highlighting}[]
\BuiltInTok{export} \VariableTok{MICRODUCK\_STUDIO\_PYTHON\_DIRECT}\OperatorTok{=}\StringTok{"}\VariableTok{$PWD}\StringTok{/rlx/.venv{-}microduck/bin/python"}
\BuiltInTok{cd}\NormalTok{ duck{-}viewer}
\ExtensionTok{node}\NormalTok{ node\_modules/next/dist/bin/next dev }\AttributeTok{{-}{-}hostname}\NormalTok{ 127.0.0.1 }\AttributeTok{{-}p}\NormalTok{ 63317}
\end{Highlighting}
\end{Shaded}

In a second terminal, from the workspace root:

\begin{Shaded}
\begin{Highlighting}[]
\BuiltInTok{cd}\NormalTok{ microduck\_local}
\VariableTok{MICRODUCK\_ACTUATOR}\OperatorTok{=}\NormalTok{bam }\ExtensionTok{.venv/bin/duck{-}lab} \AttributeTok{{-}{-}port}\NormalTok{ 8788 }\DataTypeTok{\textbackslash{}}
\NormalTok{  runs/first{-}gait ../microduck/policies/alpha\_walking.onnx}
\end{Highlighting}
\end{Shaded}

This second command requires the local harness's own virtual environment
and the named baseline paths. If
\texttt{runs/\allowbreak{}first-\allowbreak{}gait} is not available on a
clean checkout, omit that positional baseline and use the supplied
walking ONNX only. Create the harness environment with
\texttt{uv\ sync\ -\/-python\ 3.12} inside
\texttt{microduc\allowbreak{}k\_\allowbreak{}local} if necessary.
Per-policy custom-scene settings select basketball or bridge; the
default actuator environment variable is not a substitute for those
settings.

Open \texttt{http://127.0.0.1:63317}, select the experiment and its
exported policy, and check the live scene. Browser availability alone is
not enough: a saved preview and a live physics scene are different
surfaces.

\begin{figure}
\centering
\includegraphics[width=6.2in,height=5.5in,alt={Observed live custom scenes: basketball at left and suspended bridge with platforms and gantries at right.}]{assets/evidence/observed-live-scenes.png}
\caption{Observed live custom scenes: basketball at left and suspended
bridge with platforms and gantries at right.}
\end{figure}

Two earlier integration defects are instructive. First, the lab
extracted group-2 visual geometry, so a correct physics object placed
only in group 0 was invisible. Second, the CLI policy loader needed each
policy's custom environment arguments; otherwise a named bridge policy
ran on the wrong scene. The fix was correct scene routing and visual
extraction, not adding decorative CSS objects. LSTM cache sharing also
had to preserve private h/c state for every duck.

Use the repository's readiness and API checks with the services running:

\begin{Shaded}
\begin{Highlighting}[]
\FunctionTok{bash}\NormalTok{ restart{-}lab.sh }\AttributeTok{{-}{-}readiness{-}only}
\ExtensionTok{node}\NormalTok{ scripts/verify{-}seven{-}cases.mjs}
\ExtensionTok{node}\NormalTok{ scripts/verify{-}balance{-}studio{-}api.mjs}
\end{Highlighting}
\end{Shaded}

The API verifier invokes real evaluation/render operations and may
create evidence. The optional
\texttt{-\allowbreak{}-\allowbreak{}smoke-\allowbreak{}train} variant
performs real tiny training jobs, not a read-only check. Those
integration checks are separate from the static book build. They may
need browser/graphics access and sufficient time.

\subsection{A research plan for the unresolved
skills}\label{a-research-plan-for-the-unresolved-skills}

First ask whether any rollout contains the missing behavior. If not,
changing a reward coefficient cannot teach a state the learner never
visits. For bridge, investigate early-spawn practice, plank
stabilization, step placement and transition off the launch platform.
For basketball, separate balance retention, forward command tracking and
lateral/yaw steering coverage.

Keep the nominal accepted gate frozen while changing one training
condition. Log assistance on every attempt and evaluate with it off.
Compare a transferred walking or source actor, a zero-action baseline,
and the newly trained actor. A platform-standing baseline proves
platform standing, not plank balance. A partial-spawn crossing is
exploration evidence, not a full launch-to-destination solution.

Do not promise a number of additional steps sufficient to solve the
task. Training time depends on exploration, dynamics, initialization and
optimization, not only the requested total. A scientifically useful next
run may fail while revealing the exact contact transition the curriculum
must make learnable.

\subsection{Rebuild and inspect this
book}\label{rebuild-and-inspect-this-book}

\begin{Shaded}
\begin{Highlighting}[]
\ExtensionTok{rlx/.venv{-}microduck/bin/python}\NormalTok{ docs/basketball{-}bridge{-}book/tools/prepare\_assets.py}
\ExtensionTok{rlx/.venv{-}microduck/bin/python}\NormalTok{ docs/basketball{-}bridge{-}book/tools/build\_book.py}
\ExtensionTok{rlx/.venv{-}microduck/bin/python}\NormalTok{ docs/basketball{-}bridge{-}book/tools/validate\_book.py}
\end{Highlighting}
\end{Shaded}

The build uses existing dependencies; it does not install software or
retrain either policy. Markdown is the copyable source edition. PDF adds
wrapped code, page numbers, equations and a table of contents.
\texttt{supporti\allowbreak{}ng-\allowbreak{}source.z\allowbreak{}ip}
and \texttt{source-\allowbreak{}manifest\allowbreak{}.json} bind the
implementation snapshot. The validation report lists syntax checks,
source/asset hashes, PDF page/text checks and any remaining warnings.

\subsection{Final practical
examination}\label{final-practical-examination}

\begin{enumerate}
\def\labelenumi{\arabic{enumi}.}
\tightlist
\item
  Explain why 61 observations do not mean 61 independent sensors.
\item
  Show where the exported normalizer is applied and why changing its
  denominator breaks parity.
\item
  Compute the transitions in one rollout and distinguish simulated time
  from wall time.
\item
  Explain why a timeout bootstraps value but stops the GAE trace across
  reset.
\item
  Identify one observed loss curve, one observed reward curve, and one
  schematic.
\item
  Report basketball balance and rolling separately, with seeds and
  horizon.
\item
  Explain why the bridge's best progress number does not prove a
  crossing.
\item
  Start both viewer scenes without changing another user's running
  policies.
\item
  State which assets prevent an empty-folder, offline reproduction.
\item
  Propose one bounded next experiment without altering the acceptance
  gate.
\end{enumerate}

\textbf{Answer guide.} Observations include derived gravity, commands
and past actions; normalization is part of the actor contract;
transitions equal lanes times rollout length; timeout is not physical
failure but reset must break trace propagation; plots require
provenance; the measured balance passes while rolling fails;
launch-relative progress includes approach; route per-policy scenes and
preserve recurrent memory; supplied weights/meshes and dependencies
remain prerequisites; change one curriculum or physical condition and
keep deterministic full-horizon evaluation fixed.

\section{References and provenance}\label{references-and-provenance}

The implementation snapshot is the authority for what these examples
actually execute. General algorithm and simulator references explain the
ideas; they are not evidence of either trained skill.

\begin{enumerate}
\def\labelenumi{\arabic{enumi}.}
\tightlist
\item
  John Schulman, Filip Wolski, Prafulla Dhariwal, Alec Radford and Oleg
  Klimov. \textbf{Proximal Policy Optimization Algorithms}. 2017.
  arXiv:1707.06347. Original PPO paper; clipped surrogate and repeated
  sample reuse. Official record checked for this edition:
  \texttt{https://arxiv.org/abs/1707.06347}.
\item
  John Schulman, Philipp Moritz, Sergey Levine, Michael Jordan and
  Pieter Abbeel. \textbf{High-Dimensional Continuous Control Using
  Generalized Advantage Estimation}. arXiv:1506.02438. The GAE
  reference; this book separately explains the repository's
  termination/truncation masks.
  \texttt{https://arxiv.org/abs/1506.02438}.
\item
  \textbf{MuJoCo XML Reference}, official documentation. Body, free
  joint, geom, material, spatial tendon, springlength and limit
  semantics. Consult documentation matching the installed MuJoCo
  version; this snapshot records MuJoCo 3.10.0.
  \texttt{https://mujoco.readthedocs.io/en/stable/XMLreference.html}.
\item
  \texttt{microduc\allowbreak{}k\_\allowbreak{}local/\allowbreak{}AGENTS.m\allowbreak{}d}:
  the local training playbook, fixed observation contract, physics-only
  curricula, and export--evaluate--look discipline.
\item
  \texttt{microduc\allowbreak{}k-\allowbreak{}playgrou\allowbreak{}nd/\allowbreak{}experime\allowbreak{}nts/\allowbreak{}basketba\allowbreak{}ll/\allowbreak{}}
  and
  \texttt{microduc\allowbreak{}k-\allowbreak{}playgrou\allowbreak{}nd/\allowbreak{}artifact\allowbreak{}s/\allowbreak{}basketba\allowbreak{}ll/\allowbreak{}}:
  user-supplied reference design and mature actor, including
  provenance/license notices.
\item
  \texttt{microduc\allowbreak{}k-\allowbreak{}playgrou\allowbreak{}nd/\allowbreak{}experime\allowbreak{}nts/\allowbreak{}swing/\allowbreak{}README.m\allowbreak{}d}:
  suspended-support reference; a bridge is not a renamed swing task.
\item
  \texttt{docs/\allowbreak{}basketba\allowbreak{}ll-\allowbreak{}showcase\allowbreak{}/\allowbreak{}RESULTS.\allowbreak{}md}
  and
  \texttt{docs/\allowbreak{}bridge-\allowbreak{}showcase\allowbreak{}/\allowbreak{}RESULTS.\allowbreak{}md}:
  retained local results and limitations.
\item
  \texttt{assets/\allowbreak{}provenan\allowbreak{}ce.json},
  \texttt{assets/\allowbreak{}SHA256SU\allowbreak{}MS},
  \texttt{environm\allowbreak{}ent.json\allowbreak{}},
  \texttt{source-\allowbreak{}manifest\allowbreak{}.json}:
  machine-readable image/data, runtime and implementation fingerprints.
\end{enumerate}

No external paper's text is reproduced wholesale. Equations,
explanations and code discussions are tied to the local implementation.
Existing copied source retains its original content and accompanying
applicable notices; this educational snapshot does not change
third-party licensing.

\section{Appendix: complete core source and version
fingerprints}\label{appendix-complete-core-source-and-version-fingerprints}

Seven core modules follow in full, preserving their original imports,
comments and implementation. Exact files and SHA-256 manifest accompany
the book. Other adapters, the complete Studio CLI, basketball evaluator,
shared models, viewer interfaces and tests are in \texttt{source/} and
\texttt{supporti\allowbreak{}ng-\allowbreak{}source.z\allowbreak{}ip}.
This snapshot is not a standalone package: execution requires the
complete workspace, robot assets and recorded dependencies. Copy code
from Markdown or \texttt{.py} files; visual line wrapping in the PDF is
not extra source syntax.

\subsection{\texorpdfstring{A1.
\texttt{rlx/\allowbreak{}rlx/\allowbreak{}environm\allowbreak{}ents/\allowbreak{}basketba\allowbreak{}ll.py}}{A1. }}\label{a1.-rlxrlxenvironmentsbasketball.py}

SHA-256:
\texttt{6bff54a8\allowbreak{}b1f7c856\allowbreak{}ec00cce5\allowbreak{}9bb2b9b3\allowbreak{}c528d3c1\allowbreak{}6054d1f7\allowbreak{}30446fe6\allowbreak{}8c65c757\allowbreak{}}

\begin{Shaded}
\begin{Highlighting}[]
\CommentTok{"""CPU MuJoCo basketball prototype grounded in the playground b11 release."""}

\ImportTok{from}\NormalTok{ \_\_future\_\_ }\ImportTok{import}\NormalTok{ annotations}

\ImportTok{import}\NormalTok{ math}
\ImportTok{from}\NormalTok{ pathlib }\ImportTok{import}\NormalTok{ Path}
\ImportTok{from}\NormalTok{ types }\ImportTok{import}\NormalTok{ MappingProxyType}

\ImportTok{import}\NormalTok{ mujoco}
\ImportTok{import}\NormalTok{ numpy }\ImportTok{as}\NormalTok{ np}

\ImportTok{from}\NormalTok{ microduck\_local }\ImportTok{import}\NormalTok{ contract }\ImportTok{as}\NormalTok{ C}
\ImportTok{from}\NormalTok{ microduck\_local.walk\_env }\ImportTok{import}\NormalTok{ MicroduckWalkEnv}

\NormalTok{ROOT }\OperatorTok{=}\NormalTok{ Path(}\VariableTok{\_\_file\_\_}\NormalTok{).resolve().parents[}\DecValTok{3}\NormalTok{]}
\NormalTok{REFERENCE }\OperatorTok{=}\NormalTok{ ROOT }\OperatorTok{/} \StringTok{"microduck{-}playground"}
\NormalTok{BALL\_RADIUS }\OperatorTok{=} \FloatTok{0.12}
\NormalTok{BALL\_MASS }\OperatorTok{=} \FloatTok{0.62}
\NormalTok{CTRL\_DT }\OperatorTok{=}\NormalTok{ C.CTRL\_DT}
\NormalTok{HOLD\_LEVELS }\OperatorTok{=}\NormalTok{ (}\FloatTok{1.0}\NormalTok{, }\FloatTok{0.5}\NormalTok{, }\FloatTok{0.25}\NormalTok{, }\FloatTok{0.1}\NormalTok{, }\FloatTok{0.03}\NormalTok{, }\FloatTok{0.0}\NormalTok{)}
\NormalTok{REWARD\_WEIGHTS }\OperatorTok{=}\NormalTok{ MappingProxyType(\{}
    \StringTok{"linear\_tracking"}\NormalTok{: }\FloatTok{1.0}\NormalTok{,}
    \StringTok{"yaw\_tracking"}\NormalTok{: }\FloatTok{0.5}\NormalTok{,}
    \StringTok{"upright"}\NormalTok{: }\FloatTok{1.0}\NormalTok{,}
    \StringTok{"centered"}\NormalTok{: }\FloatTok{2.0}\NormalTok{,}
    \StringTok{"feet\_on\_ball"}\NormalTok{: }\FloatTok{1.0}\NormalTok{,}
    \StringTok{"height"}\NormalTok{: }\FloatTok{1.0}\NormalTok{,}
    \StringTok{"ball\_speed\_penalty"}\NormalTok{: }\FloatTok{0.05}\NormalTok{,}
    \StringTok{"angular\_velocity\_penalty"}\NormalTok{: }\FloatTok{0.05}\NormalTok{,}
    \StringTok{"action\_rate\_penalty"}\NormalTok{: }\FloatTok{0.2}\NormalTok{,}
\NormalTok{\})}


\KeywordTok{def}\NormalTok{ basketball\_model(actuator: }\BuiltInTok{str} \OperatorTok{=} \StringTok{"bam"}\NormalTok{) }\OperatorTok{{-}\textgreater{}}\NormalTok{ mujoco.MjModel:}
    \ControlFlowTok{if}\NormalTok{ actuator }\KeywordTok{not} \KeywordTok{in}\NormalTok{ (}\StringTok{"bam"}\NormalTok{, }\StringTok{"xml"}\NormalTok{):}
        \ControlFlowTok{raise} \PreprocessorTok{ValueError}\NormalTok{(}\StringTok{"actuator must be bam or xml"}\NormalTok{)}
\NormalTok{    robot\_dir }\OperatorTok{=}\NormalTok{ Path(}\VariableTok{\_\_file\_\_}\NormalTok{).resolve().parents[}\DecValTok{1}\NormalTok{] }\OperatorTok{/} \StringTok{"mjlab\_microduck/robot/microduck"}
\NormalTok{    asset\_dir }\OperatorTok{=}\NormalTok{ REFERENCE }\OperatorTok{/} \StringTok{"src/mjlab\_microduck/robot/assets/basketball"}
\NormalTok{    spec }\OperatorTok{=}\NormalTok{ mujoco.MjSpec.from\_file(}\BuiltInTok{str}\NormalTok{(robot\_dir }\OperatorTok{/} \StringTok{"robot\_allcollisions.xml"}\NormalTok{))}
\NormalTok{    spec.option.timestep }\OperatorTok{=}\NormalTok{ C.PHYSICS\_DT}
\NormalTok{    spec.option.integrator }\OperatorTok{=}\NormalTok{ mujoco.mjtIntegrator.mjINT\_IMPLICITFAST}
\NormalTok{    spec.option.cone }\OperatorTok{=}\NormalTok{ mujoco.mjtCone.mjCONE\_PYRAMIDAL}
\NormalTok{    spec.option.iterations }\OperatorTok{=} \DecValTok{10}
\NormalTok{    spec.option.ls\_iterations }\OperatorTok{=} \DecValTok{20}
\NormalTok{    ground\_texture }\OperatorTok{=}\NormalTok{ spec.add\_texture(name}\OperatorTok{=}\StringTok{"basketball\_ground"}\NormalTok{)}
\NormalTok{    ground\_texture.}\BuiltInTok{type} \OperatorTok{=}\NormalTok{ mujoco.mjtTexture.mjTEXTURE\_2D}
\NormalTok{    ground\_texture.builtin }\OperatorTok{=}\NormalTok{ mujoco.mjtBuiltin.mjBUILTIN\_CHECKER}
\NormalTok{    ground\_texture.rgb1 }\OperatorTok{=}\NormalTok{ [}\FloatTok{.22}\NormalTok{, }\FloatTok{.27}\NormalTok{, }\FloatTok{.31}\NormalTok{]}
\NormalTok{    ground\_texture.rgb2 }\OperatorTok{=}\NormalTok{ [}\FloatTok{.16}\NormalTok{, }\FloatTok{.20}\NormalTok{, }\FloatTok{.24}\NormalTok{]}
\NormalTok{    ground\_texture.width }\OperatorTok{=} \DecValTok{256}
\NormalTok{    ground\_texture.height }\OperatorTok{=} \DecValTok{256}
\NormalTok{    ground\_material }\OperatorTok{=}\NormalTok{ spec.add\_material(name}\OperatorTok{=}\StringTok{"basketball\_ground\_mat"}\NormalTok{)}
\NormalTok{    ground\_textures }\OperatorTok{=} \BuiltInTok{list}\NormalTok{(ground\_material.textures)}
\NormalTok{    ground\_textures[}\BuiltInTok{int}\NormalTok{(mujoco.mjtTextureRole.mjTEXROLE\_RGB)] }\OperatorTok{=} \StringTok{"basketball\_ground"}
\NormalTok{    ground\_material.textures }\OperatorTok{=}\NormalTok{ ground\_textures}
\NormalTok{    ground\_material.texuniform }\OperatorTok{=} \VariableTok{True}
\NormalTok{    ground\_material.texrepeat }\OperatorTok{=}\NormalTok{ [}\DecValTok{4}\NormalTok{, }\DecValTok{4}\NormalTok{]}
\NormalTok{    spec.worldbody.add\_geom(}
\NormalTok{        name}\OperatorTok{=}\StringTok{"floor"}\NormalTok{, }\BuiltInTok{type}\OperatorTok{=}\NormalTok{mujoco.mjtGeom.mjGEOM\_PLANE,}
\NormalTok{        size}\OperatorTok{=}\NormalTok{[}\DecValTok{0}\NormalTok{, }\DecValTok{0}\NormalTok{, }\FloatTok{0.05}\NormalTok{], material}\OperatorTok{=}\StringTok{"basketball\_ground\_mat"}\NormalTok{,}
\NormalTok{        friction}\OperatorTok{=}\NormalTok{[}\DecValTok{1}\NormalTok{, }\FloatTok{0.005}\NormalTok{, }\FloatTok{0.0001}\NormalTok{],}
\NormalTok{    )}
\NormalTok{    spec.worldbody.add\_light(pos}\OperatorTok{=}\NormalTok{[}\DecValTok{0}\NormalTok{, }\OperatorTok{{-}}\DecValTok{1}\NormalTok{, }\DecValTok{3}\NormalTok{], }\BuiltInTok{dir}\OperatorTok{=}\NormalTok{[}\DecValTok{0}\NormalTok{, }\DecValTok{0}\NormalTok{, }\OperatorTok{{-}}\DecValTok{1}\NormalTok{])}
\NormalTok{    texture }\OperatorTok{=}\NormalTok{ spec.add\_texture(name}\OperatorTok{=}\StringTok{"basketball\_tex"}\NormalTok{)}
\NormalTok{    texture.}\BuiltInTok{type} \OperatorTok{=}\NormalTok{ mujoco.mjtTexture.mjTEXTURE\_2D}
\NormalTok{    texture.}\BuiltInTok{file} \OperatorTok{=} \BuiltInTok{str}\NormalTok{(asset\_dir }\OperatorTok{/} \StringTok{"basketball.png"}\NormalTok{)}
\NormalTok{    material }\OperatorTok{=}\NormalTok{ spec.add\_material(name}\OperatorTok{=}\StringTok{"basketball\_mat"}\NormalTok{)}
\NormalTok{    textures }\OperatorTok{=} \BuiltInTok{list}\NormalTok{(material.textures)}
\NormalTok{    textures[}\BuiltInTok{int}\NormalTok{(mujoco.mjtTextureRole.mjTEXROLE\_RGB)] }\OperatorTok{=} \StringTok{"basketball\_tex"}
\NormalTok{    material.textures }\OperatorTok{=}\NormalTok{ textures}
\NormalTok{    material.specular }\OperatorTok{=} \FloatTok{0.15}
\NormalTok{    spec.add\_mesh(name}\OperatorTok{=}\StringTok{"basketball\_mesh"}\NormalTok{, }\BuiltInTok{file}\OperatorTok{=}\BuiltInTok{str}\NormalTok{(asset\_dir }\OperatorTok{/} \StringTok{"basketball.obj"}\NormalTok{))}
\NormalTok{    ball }\OperatorTok{=}\NormalTok{ spec.worldbody.add\_body(name}\OperatorTok{=}\StringTok{"basketball"}\NormalTok{, pos}\OperatorTok{=}\NormalTok{[}\DecValTok{0}\NormalTok{, }\DecValTok{0}\NormalTok{, BALL\_RADIUS }\OperatorTok{+} \FloatTok{0.001}\NormalTok{])}
\NormalTok{    ball.add\_freejoint(name}\OperatorTok{=}\StringTok{"basketball\_freejoint"}\NormalTok{)}
\NormalTok{    ball.add\_geom(}
\NormalTok{        name}\OperatorTok{=}\StringTok{"ball\_sphere"}\NormalTok{, }\BuiltInTok{type}\OperatorTok{=}\NormalTok{mujoco.mjtGeom.mjGEOM\_SPHERE,}
\NormalTok{        size}\OperatorTok{=}\NormalTok{[BALL\_RADIUS, }\DecValTok{0}\NormalTok{, }\DecValTok{0}\NormalTok{], mass}\OperatorTok{=}\NormalTok{BALL\_MASS, friction}\OperatorTok{=}\NormalTok{[}\FloatTok{1.2}\NormalTok{, }\FloatTok{0.01}\NormalTok{, }\FloatTok{0.001}\NormalTok{],}
\NormalTok{        priority}\OperatorTok{=}\DecValTok{1}\NormalTok{, condim}\OperatorTok{=}\DecValTok{4}\NormalTok{, rgba}\OperatorTok{=}\NormalTok{[}\DecValTok{1}\NormalTok{, }\FloatTok{0.45}\NormalTok{, }\FloatTok{0.02}\NormalTok{, }\DecValTok{0}\NormalTok{], group}\OperatorTok{=}\DecValTok{3}\NormalTok{,}
\NormalTok{    )}
\NormalTok{    ball.add\_geom(}
\NormalTok{        name}\OperatorTok{=}\StringTok{"ball\_visual"}\NormalTok{, }\BuiltInTok{type}\OperatorTok{=}\NormalTok{mujoco.mjtGeom.mjGEOM\_MESH,}
\NormalTok{        meshname}\OperatorTok{=}\StringTok{"basketball\_mesh"}\NormalTok{, material}\OperatorTok{=}\StringTok{"basketball\_mat"}\NormalTok{,}
\NormalTok{        mass}\OperatorTok{=}\DecValTok{0}\NormalTok{, contype}\OperatorTok{=}\DecValTok{0}\NormalTok{, conaffinity}\OperatorTok{=}\DecValTok{0}\NormalTok{, group}\OperatorTok{=}\DecValTok{2}\NormalTok{,}
\NormalTok{    )}
\NormalTok{    spec.add\_key(}
\NormalTok{        name}\OperatorTok{=}\StringTok{"STAND"}\NormalTok{, qpos}\OperatorTok{=}\NormalTok{np.concatenate((}
\NormalTok{            [}\DecValTok{0}\NormalTok{, }\DecValTok{0}\NormalTok{, }\DecValTok{2} \OperatorTok{*}\NormalTok{ BALL\_RADIUS }\OperatorTok{+} \FloatTok{0.128}\NormalTok{, }\DecValTok{1}\NormalTok{, }\DecValTok{0}\NormalTok{, }\DecValTok{0}\NormalTok{, }\DecValTok{0}\NormalTok{],}
\NormalTok{            C.DEFAULT\_POSE, [}\DecValTok{0}\NormalTok{, }\DecValTok{0}\NormalTok{, BALL\_RADIUS }\OperatorTok{+} \FloatTok{0.001}\NormalTok{, }\DecValTok{1}\NormalTok{, }\DecValTok{0}\NormalTok{, }\DecValTok{0}\NormalTok{, }\DecValTok{0}\NormalTok{],}
\NormalTok{        )), ctrl}\OperatorTok{=}\NormalTok{C.DEFAULT\_POSE,}
\NormalTok{    )}
\NormalTok{    model }\OperatorTok{=}\NormalTok{ spec.}\BuiltInTok{compile}\NormalTok{()}
\NormalTok{    model.vis.global\_.offwidth }\OperatorTok{=} \DecValTok{1280}
\NormalTok{    model.vis.global\_.offheight }\OperatorTok{=} \DecValTok{960}
    \ControlFlowTok{return}\NormalTok{ model}


\KeywordTok{def}\NormalTok{ curriculum\_level(stage: }\BuiltInTok{int}\NormalTok{, seconds: }\BuiltInTok{float}\NormalTok{) }\OperatorTok{{-}\textgreater{}} \BuiltInTok{int}\NormalTok{:}
    \ControlFlowTok{if} \KeywordTok{not} \DecValTok{0} \OperatorTok{\textless{}=}\NormalTok{ stage }\OperatorTok{\textless{}} \BuiltInTok{len}\NormalTok{(HOLD\_LEVELS) }\KeywordTok{or} \KeywordTok{not}\NormalTok{ math.isfinite(seconds) }\KeywordTok{or}\NormalTok{ seconds }\OperatorTok{\textless{}} \DecValTok{0}\NormalTok{:}
        \ControlFlowTok{raise} \PreprocessorTok{ValueError}\NormalTok{(}\StringTok{"invalid curriculum stage or episode duration"}\NormalTok{)}
    \ControlFlowTok{if}\NormalTok{ seconds }\OperatorTok{\textgreater{}=} \DecValTok{6}\NormalTok{:}
        \ControlFlowTok{return} \BuiltInTok{min}\NormalTok{(stage }\OperatorTok{+} \DecValTok{1}\NormalTok{, }\BuiltInTok{len}\NormalTok{(HOLD\_LEVELS) }\OperatorTok{{-}} \DecValTok{1}\NormalTok{)}
    \ControlFlowTok{if}\NormalTok{ seconds }\OperatorTok{\textless{}} \FloatTok{1.5}\NormalTok{:}
        \ControlFlowTok{return} \BuiltInTok{max}\NormalTok{(stage }\OperatorTok{{-}} \DecValTok{1}\NormalTok{, }\DecValTok{0}\NormalTok{)}
    \ControlFlowTok{return}\NormalTok{ stage}


\KeywordTok{class}\NormalTok{ BasketballEnv(MicroduckWalkEnv):}
    \CommentTok{"""A blind 61D actor steps on a sphere; all assistance is explicit in info.}

\CommentTok{    This is a local actuator/physics adaptation, not a bitwise mjlab port.}
\CommentTok{    The sealed bounded reward subset is unchanged across curriculum stages.}
\CommentTok{    """}

    \KeywordTok{def} \FunctionTok{\_\_init\_\_}\NormalTok{(}
        \VariableTok{self}\NormalTok{, }\OperatorTok{*}\NormalTok{, max\_episode\_s: }\BuiltInTok{float} \OperatorTok{=} \DecValTok{10}\NormalTok{, hold: }\BuiltInTok{float} \OperatorTok{=} \DecValTok{0}\NormalTok{,}
\NormalTok{        curriculum: }\BuiltInTok{bool} \OperatorTok{=} \VariableTok{False}\NormalTok{, command: }\BuiltInTok{tuple}\NormalTok{[}\BuiltInTok{float}\NormalTok{, }\BuiltInTok{float}\NormalTok{, }\BuiltInTok{float}\NormalTok{] }\OperatorTok{|} \VariableTok{None} \OperatorTok{=} \VariableTok{None}\NormalTok{,}
\NormalTok{        seed: }\BuiltInTok{int} \OperatorTok{|} \VariableTok{None} \OperatorTok{=} \VariableTok{None}\NormalTok{, actuator: }\BuiltInTok{str} \OperatorTok{=} \StringTok{"bam"}\NormalTok{, obs\_noise: }\BuiltInTok{bool} \OperatorTok{=} \VariableTok{False}\NormalTok{,}
\NormalTok{        domain\_rand: }\BuiltInTok{bool} \OperatorTok{=} \VariableTok{False}\NormalTok{, pushes: }\BuiltInTok{bool} \OperatorTok{=} \VariableTok{False}\NormalTok{, random\_yaw: }\BuiltInTok{bool} \OperatorTok{=} \VariableTok{True}\NormalTok{,}
\NormalTok{        model: mujoco.MjModel }\OperatorTok{|} \VariableTok{None} \OperatorTok{=} \VariableTok{None}\NormalTok{,}
\NormalTok{    ):}
        \ControlFlowTok{if} \KeywordTok{not}\NormalTok{ math.isfinite(hold) }\KeywordTok{or} \KeywordTok{not} \DecValTok{0} \OperatorTok{\textless{}=}\NormalTok{ hold }\OperatorTok{\textless{}=} \DecValTok{1}\NormalTok{:}
            \ControlFlowTok{raise} \PreprocessorTok{ValueError}\NormalTok{(}\StringTok{"hold must be finite and between zero and one"}\NormalTok{)}
        \ControlFlowTok{if} \KeywordTok{not}\NormalTok{ math.isfinite(max\_episode\_s) }\KeywordTok{or}\NormalTok{ max\_episode\_s }\OperatorTok{\textless{}=} \DecValTok{0}\NormalTok{:}
            \ControlFlowTok{raise} \PreprocessorTok{ValueError}\NormalTok{(}\StringTok{"episode length must be positive and finite"}\NormalTok{)}
        \ControlFlowTok{if}\NormalTok{ curriculum }\KeywordTok{and}\NormalTok{ hold }\KeywordTok{not} \KeywordTok{in}\NormalTok{ HOLD\_LEVELS:}
            \ControlFlowTok{raise} \PreprocessorTok{ValueError}\NormalTok{(}\StringTok{"curriculum hold must be one of HOLD\_LEVELS"}\NormalTok{)}
        \VariableTok{self}\NormalTok{.hold }\OperatorTok{=} \BuiltInTok{float}\NormalTok{(hold)}
        \VariableTok{self}\NormalTok{.curriculum }\OperatorTok{=}\NormalTok{ curriculum}
        \VariableTok{self}\NormalTok{.stage }\OperatorTok{=}\NormalTok{ HOLD\_LEVELS.index(hold) }\ControlFlowTok{if}\NormalTok{ hold }\KeywordTok{in}\NormalTok{ HOLD\_LEVELS }\ControlFlowTok{else} \DecValTok{0}
        \VariableTok{self}\NormalTok{.fixed\_command }\OperatorTok{=} \VariableTok{self}\NormalTok{.\_validate\_command(command)}
        \VariableTok{self}\NormalTok{.pushes }\OperatorTok{=}\NormalTok{ pushes}
        \VariableTok{self}\NormalTok{.\_episode\_finished }\OperatorTok{=} \VariableTok{False}
        \BuiltInTok{super}\NormalTok{().}\FunctionTok{\_\_init\_\_}\NormalTok{(}
\NormalTok{            max\_episode\_s}\OperatorTok{=}\NormalTok{max\_episode\_s, command\_resample\_s}\OperatorTok{=}\DecValTok{5}\NormalTok{,}
\NormalTok{            obs\_noise}\OperatorTok{=}\NormalTok{obs\_noise, domain\_rand}\OperatorTok{=}\NormalTok{domain\_rand, action\_delay}\OperatorTok{=}\VariableTok{True}\NormalTok{,}
\NormalTok{            random\_yaw}\OperatorTok{=}\NormalTok{random\_yaw, seed}\OperatorTok{=}\NormalTok{seed, actuator\_force}\OperatorTok{=}\NormalTok{actuator,}
\NormalTok{            model}\OperatorTok{=}\NormalTok{model }\ControlFlowTok{if}\NormalTok{ model }\KeywordTok{is} \KeywordTok{not} \VariableTok{None} \ControlFlowTok{else}\NormalTok{ basketball\_model(actuator),}
\NormalTok{            terminate\_on\_fall}\OperatorTok{=}\VariableTok{False}\NormalTok{, height\_termination}\OperatorTok{=}\VariableTok{False}\NormalTok{,}
\NormalTok{        )}
        \VariableTok{self}\NormalTok{.ball\_body\_id }\OperatorTok{=} \VariableTok{self}\NormalTok{.model.body(}\StringTok{"basketball"}\NormalTok{).}\BuiltInTok{id}
        \VariableTok{self}\NormalTok{.ball\_geom\_id }\OperatorTok{=} \VariableTok{self}\NormalTok{.model.geom(}\StringTok{"ball\_sphere"}\NormalTok{).}\BuiltInTok{id}
        \VariableTok{self}\NormalTok{.ball\_qpos\_adr }\OperatorTok{=} \BuiltInTok{int}\NormalTok{(}\VariableTok{self}\NormalTok{.model.joint(}\StringTok{"basketball\_freejoint"}\NormalTok{).qposadr[}\DecValTok{0}\NormalTok{])}
        \VariableTok{self}\NormalTok{.ball\_qvel\_adr }\OperatorTok{=} \BuiltInTok{int}\NormalTok{(}\VariableTok{self}\NormalTok{.model.joint(}\StringTok{"basketball\_freejoint"}\NormalTok{).dofadr[}\DecValTok{0}\NormalTok{])}
        \VariableTok{self}\NormalTok{.foot\_site\_ids }\OperatorTok{=}\NormalTok{ [}\VariableTok{self}\NormalTok{.model.site(name).}\BuiltInTok{id} \ControlFlowTok{for}\NormalTok{ name }\KeywordTok{in}\NormalTok{ (}\StringTok{"left\_foot"}\NormalTok{, }\StringTok{"right\_foot"}\NormalTok{)]}
        \VariableTok{self}\NormalTok{.\_ball\_anchor }\OperatorTok{=}\NormalTok{ np.array([}\DecValTok{0}\NormalTok{, }\DecValTok{0}\NormalTok{, BALL\_RADIUS }\OperatorTok{+} \FloatTok{0.001}\NormalTok{])}
        \VariableTok{self}\NormalTok{.\_next\_push }\OperatorTok{=}\NormalTok{ math.inf}

    \AttributeTok{@staticmethod}
    \KeywordTok{def}\NormalTok{ \_validate\_command(command):}
        \ControlFlowTok{if}\NormalTok{ command }\KeywordTok{is} \VariableTok{None}\NormalTok{:}
            \ControlFlowTok{return} \VariableTok{None}
\NormalTok{        array }\OperatorTok{=}\NormalTok{ np.asarray(command, dtype}\OperatorTok{=}\NormalTok{np.float32)}
        \ControlFlowTok{if}\NormalTok{ array.shape }\OperatorTok{!=}\NormalTok{ (}\DecValTok{3}\NormalTok{,) }\KeywordTok{or} \KeywordTok{not}\NormalTok{ np.isfinite(array).}\BuiltInTok{all}\NormalTok{():}
            \ControlFlowTok{raise} \PreprocessorTok{ValueError}\NormalTok{(}\StringTok{"command must contain three finite numbers"}\NormalTok{)}
        \ControlFlowTok{return}\NormalTok{ array.copy()}

    \KeywordTok{def}\NormalTok{ set\_command(}\VariableTok{self}\NormalTok{, command: }\BuiltInTok{tuple}\NormalTok{[}\BuiltInTok{float}\NormalTok{, }\BuiltInTok{float}\NormalTok{, }\BuiltInTok{float}\NormalTok{]) }\OperatorTok{{-}\textgreater{}} \VariableTok{None}\NormalTok{:}
\NormalTok{        validated }\OperatorTok{=} \VariableTok{self}\NormalTok{.\_validate\_command(command)}
        \ControlFlowTok{if}\NormalTok{ validated }\KeywordTok{is} \VariableTok{None}\NormalTok{:}
            \ControlFlowTok{raise} \PreprocessorTok{ValueError}\NormalTok{(}\StringTok{"set\_command requires three numbers"}\NormalTok{)}
        \VariableTok{self}\NormalTok{.fixed\_command }\OperatorTok{=}\NormalTok{ validated}
        \VariableTok{self}\NormalTok{.twist\_cmd[:] }\OperatorTok{=}\NormalTok{ validated}

    \KeywordTok{def}\NormalTok{ \_sample\_commands(}\VariableTok{self}\NormalTok{) }\OperatorTok{{-}\textgreater{}} \VariableTok{None}\NormalTok{:}
        \ControlFlowTok{if} \VariableTok{self}\NormalTok{.fixed\_command }\KeywordTok{is} \KeywordTok{not} \VariableTok{None}\NormalTok{:}
            \VariableTok{self}\NormalTok{.twist\_cmd[:] }\OperatorTok{=} \VariableTok{self}\NormalTok{.fixed\_command}
        \ControlFlowTok{elif} \VariableTok{self}\NormalTok{.\_rng.random() }\OperatorTok{\textless{}} \FloatTok{0.25}\NormalTok{:}
            \VariableTok{self}\NormalTok{.twist\_cmd[:] }\OperatorTok{=} \DecValTok{0}
        \ControlFlowTok{else}\NormalTok{:}
            \VariableTok{self}\NormalTok{.twist\_cmd[:] }\OperatorTok{=} \VariableTok{self}\NormalTok{.\_rng.uniform([}\OperatorTok{{-}}\FloatTok{.15}\NormalTok{, }\OperatorTok{{-}}\FloatTok{.1}\NormalTok{, }\OperatorTok{{-}}\FloatTok{.5}\NormalTok{], [}\FloatTok{.15}\NormalTok{, }\FloatTok{.1}\NormalTok{, }\FloatTok{.5}\NormalTok{])}
        \VariableTok{self}\NormalTok{.head\_cmd[:] }\OperatorTok{=} \DecValTok{0}
        \VariableTok{self}\NormalTok{.body\_cmd[:] }\OperatorTok{=} \DecValTok{0}

    \KeywordTok{def}\NormalTok{ reset(}\VariableTok{self}\NormalTok{, }\OperatorTok{*}\NormalTok{, seed: }\BuiltInTok{int} \OperatorTok{|} \VariableTok{None} \OperatorTok{=} \VariableTok{None}\NormalTok{, options: }\BuiltInTok{dict} \OperatorTok{|} \VariableTok{None} \OperatorTok{=} \VariableTok{None}\NormalTok{):}
        \ControlFlowTok{if} \VariableTok{self}\NormalTok{.curriculum }\KeywordTok{and} \VariableTok{self}\NormalTok{.\_episode\_finished:}
            \VariableTok{self}\NormalTok{.stage }\OperatorTok{=}\NormalTok{ curriculum\_level(}\VariableTok{self}\NormalTok{.stage, }\VariableTok{self}\NormalTok{.step\_count }\OperatorTok{*}\NormalTok{ CTRL\_DT)}
            \VariableTok{self}\NormalTok{.hold }\OperatorTok{=}\NormalTok{ HOLD\_LEVELS[}\VariableTok{self}\NormalTok{.stage]}
        \VariableTok{self}\NormalTok{.\_episode\_finished }\OperatorTok{=} \VariableTok{False}
        \BuiltInTok{super}\NormalTok{().reset(seed}\OperatorTok{=}\NormalTok{seed, options}\OperatorTok{=}\NormalTok{options)}
        \VariableTok{self}\NormalTok{.data.qpos[:}\DecValTok{2}\NormalTok{] }\OperatorTok{=} \VariableTok{self}\NormalTok{.\_rng.uniform(}\OperatorTok{{-}}\FloatTok{.01}\NormalTok{, }\FloatTok{.01}\NormalTok{, }\DecValTok{2}\NormalTok{)}
        \VariableTok{self}\NormalTok{.data.qpos[}\DecValTok{2}\NormalTok{] }\OperatorTok{=} \DecValTok{2} \OperatorTok{*}\NormalTok{ BALL\_RADIUS }\OperatorTok{+} \FloatTok{.128}
\NormalTok{        yaw }\OperatorTok{=} \VariableTok{self}\NormalTok{.\_rng.uniform(}\OperatorTok{{-}}\NormalTok{math.pi, math.pi) }\ControlFlowTok{if} \VariableTok{self}\NormalTok{.random\_yaw }\ControlFlowTok{else} \DecValTok{0}
\NormalTok{        roll, pitch }\OperatorTok{=} \VariableTok{self}\NormalTok{.\_rng.uniform(}\OperatorTok{{-}}\NormalTok{math.radians(}\DecValTok{2}\NormalTok{), math.radians(}\DecValTok{2}\NormalTok{), }\DecValTok{2}\NormalTok{)}
\NormalTok{        yaw\_quat }\OperatorTok{=}\NormalTok{ np.array([math.cos(yaw }\OperatorTok{/} \DecValTok{2}\NormalTok{), }\DecValTok{0}\NormalTok{, }\DecValTok{0}\NormalTok{, math.sin(yaw }\OperatorTok{/} \DecValTok{2}\NormalTok{)])}
\NormalTok{        pitch\_quat }\OperatorTok{=}\NormalTok{ np.array([math.cos(pitch }\OperatorTok{/} \DecValTok{2}\NormalTok{), }\DecValTok{0}\NormalTok{, math.sin(pitch }\OperatorTok{/} \DecValTok{2}\NormalTok{), }\DecValTok{0}\NormalTok{])}
\NormalTok{        roll\_quat }\OperatorTok{=}\NormalTok{ np.array([math.cos(roll }\OperatorTok{/} \DecValTok{2}\NormalTok{), math.sin(roll }\OperatorTok{/} \DecValTok{2}\NormalTok{), }\DecValTok{0}\NormalTok{, }\DecValTok{0}\NormalTok{])}
\NormalTok{        rotation }\OperatorTok{=}\NormalTok{ np.empty(}\DecValTok{4}\NormalTok{)}
\NormalTok{        mujoco.mju\_mulQuat(rotation, yaw\_quat, pitch\_quat)}
\NormalTok{        mujoco.mju\_mulQuat(}\VariableTok{self}\NormalTok{.data.qpos[}\DecValTok{3}\NormalTok{:}\DecValTok{7}\NormalTok{], rotation, roll\_quat)}
        \VariableTok{self}\NormalTok{.data.qpos[}\VariableTok{self}\NormalTok{.joint\_qpos\_adr] }\OperatorTok{=}\NormalTok{ C.DEFAULT\_POSE }\OperatorTok{+} \VariableTok{self}\NormalTok{.\_rng.uniform(}\OperatorTok{{-}}\FloatTok{.05}\NormalTok{, }\FloatTok{.05}\NormalTok{, }\DecValTok{14}\NormalTok{)}
        \VariableTok{self}\NormalTok{.data.qpos[}\VariableTok{self}\NormalTok{.ball\_qpos\_adr:}\VariableTok{self}\NormalTok{.ball\_qpos\_adr }\OperatorTok{+} \DecValTok{7}\NormalTok{] }\OperatorTok{=}\NormalTok{ [}\DecValTok{0}\NormalTok{, }\DecValTok{0}\NormalTok{, BALL\_RADIUS }\OperatorTok{+} \FloatTok{.001}\NormalTok{, }\DecValTok{1}\NormalTok{, }\DecValTok{0}\NormalTok{, }\DecValTok{0}\NormalTok{, }\DecValTok{0}\NormalTok{]}
        \VariableTok{self}\NormalTok{.data.qvel[:] }\OperatorTok{=} \DecValTok{0}
        \VariableTok{self}\NormalTok{.data.xfrc\_applied[:] }\OperatorTok{=} \DecValTok{0}
        \VariableTok{self}\NormalTok{.data.ctrl[:] }\OperatorTok{=} \VariableTok{self}\NormalTok{.data.qpos[}\VariableTok{self}\NormalTok{.joint\_qpos\_adr]}
\NormalTok{        mujoco.mj\_forward(}\VariableTok{self}\NormalTok{.model, }\VariableTok{self}\NormalTok{.data)}
        \ControlFlowTok{if} \VariableTok{self}\NormalTok{.bam }\KeywordTok{is} \KeywordTok{not} \VariableTok{None}\NormalTok{:}
            \VariableTok{self}\NormalTok{.bam.reset()}
        \VariableTok{self}\NormalTok{.prev\_joint\_vel[:] }\OperatorTok{=} \DecValTok{0}
        \VariableTok{self}\NormalTok{.\_next\_push }\OperatorTok{=} \VariableTok{self}\NormalTok{.\_rng.uniform(}\FloatTok{1.5}\NormalTok{, }\DecValTok{3}\NormalTok{) }\ControlFlowTok{if} \VariableTok{self}\NormalTok{.pushes }\ControlFlowTok{else}\NormalTok{ math.inf}
        \ControlFlowTok{return} \VariableTok{self}\NormalTok{.\_get\_obs(), }\VariableTok{self}\NormalTok{.metrics()}

    \KeywordTok{def}\NormalTok{ \_apply\_assistance(}\VariableTok{self}\NormalTok{) }\OperatorTok{{-}\textgreater{}} \VariableTok{None}\NormalTok{:}
        \VariableTok{self}\NormalTok{.data.xfrc\_applied[:] }\OperatorTok{=} \DecValTok{0}
        \ControlFlowTok{if} \VariableTok{self}\NormalTok{.hold }\OperatorTok{==} \DecValTok{0}\NormalTok{:}
            \ControlFlowTok{return}
\NormalTok{        ball\_pos }\OperatorTok{=} \VariableTok{self}\NormalTok{.data.xpos[}\VariableTok{self}\NormalTok{.ball\_body\_id]}
\NormalTok{        ball\_vel }\OperatorTok{=} \VariableTok{self}\NormalTok{.data.qvel[}\VariableTok{self}\NormalTok{.ball\_qvel\_adr:}\VariableTok{self}\NormalTok{.ball\_qvel\_adr }\OperatorTok{+} \DecValTok{6}\NormalTok{]}
\NormalTok{        ball\_force }\OperatorTok{=} \OperatorTok{{-}}\DecValTok{400} \OperatorTok{*}\NormalTok{ (ball\_pos }\OperatorTok{{-}} \VariableTok{self}\NormalTok{.\_ball\_anchor) }\OperatorTok{{-}} \DecValTok{20} \OperatorTok{*}\NormalTok{ ball\_vel[:}\DecValTok{3}\NormalTok{]}
\NormalTok{        ball\_force[}\DecValTok{2}\NormalTok{] }\OperatorTok{=} \OperatorTok{{-}}\DecValTok{20} \OperatorTok{*}\NormalTok{ ball\_vel[}\DecValTok{2}\NormalTok{]}
        \VariableTok{self}\NormalTok{.data.xfrc\_applied[}\VariableTok{self}\NormalTok{.ball\_body\_id, :}\DecValTok{3}\NormalTok{] }\OperatorTok{=} \VariableTok{self}\NormalTok{.hold }\OperatorTok{*}\NormalTok{ ball\_force}
\NormalTok{        ball\_rotation }\OperatorTok{=} \VariableTok{self}\NormalTok{.data.xmat[}\VariableTok{self}\NormalTok{.ball\_body\_id].reshape(}\DecValTok{3}\NormalTok{, }\DecValTok{3}\NormalTok{)}
        \VariableTok{self}\NormalTok{.data.xfrc\_applied[}\VariableTok{self}\NormalTok{.ball\_body\_id, }\DecValTok{3}\NormalTok{:] }\OperatorTok{=} \OperatorTok{{-}}\VariableTok{self}\NormalTok{.hold }\OperatorTok{*} \FloatTok{.2} \OperatorTok{*}\NormalTok{ (ball\_rotation }\OperatorTok{@}\NormalTok{ ball\_vel[}\DecValTok{3}\NormalTok{:])}
\NormalTok{        rotation }\OperatorTok{=} \VariableTok{self}\NormalTok{.data.xmat[}\VariableTok{self}\NormalTok{.trunk\_body\_id].reshape(}\DecValTok{3}\NormalTok{, }\DecValTok{3}\NormalTok{)}
\NormalTok{        up }\OperatorTok{=}\NormalTok{ rotation[:, }\DecValTok{2}\NormalTok{]}
\NormalTok{        angular\_world }\OperatorTok{=}\NormalTok{ rotation }\OperatorTok{@} \VariableTok{self}\NormalTok{.data.qvel[}\DecValTok{3}\NormalTok{:}\DecValTok{6}\NormalTok{]}
\NormalTok{        duck\_force }\OperatorTok{=} \OperatorTok{{-}}\DecValTok{40} \OperatorTok{*}\NormalTok{ (}\VariableTok{self}\NormalTok{.data.xpos[}\VariableTok{self}\NormalTok{.trunk\_body\_id] }\OperatorTok{{-}}\NormalTok{ ball\_pos) }\OperatorTok{{-}} \DecValTok{4} \OperatorTok{*} \VariableTok{self}\NormalTok{.data.qvel[:}\DecValTok{3}\NormalTok{]}
\NormalTok{        duck\_force[}\DecValTok{2}\NormalTok{] }\OperatorTok{=} \DecValTok{0}
        \VariableTok{self}\NormalTok{.data.xfrc\_applied[}\VariableTok{self}\NormalTok{.trunk\_body\_id, :}\DecValTok{3}\NormalTok{] }\OperatorTok{=} \VariableTok{self}\NormalTok{.hold }\OperatorTok{*}\NormalTok{ duck\_force}
        \VariableTok{self}\NormalTok{.data.xfrc\_applied[}\VariableTok{self}\NormalTok{.trunk\_body\_id, }\DecValTok{3}\NormalTok{:] }\OperatorTok{=} \VariableTok{self}\NormalTok{.hold }\OperatorTok{*}\NormalTok{ (}
            \DecValTok{3} \OperatorTok{*}\NormalTok{ np.cross(up, [}\DecValTok{0}\NormalTok{, }\DecValTok{0}\NormalTok{, }\DecValTok{1}\NormalTok{]) }\OperatorTok{{-}} \FloatTok{.08} \OperatorTok{*}\NormalTok{ angular\_world}
\NormalTok{        )}

    \KeywordTok{def}\NormalTok{ step(}\VariableTok{self}\NormalTok{, action):}
\NormalTok{        action }\OperatorTok{=}\NormalTok{ np.asarray(action, dtype}\OperatorTok{=}\NormalTok{np.float32)}
        \ControlFlowTok{if}\NormalTok{ action.shape }\OperatorTok{!=}\NormalTok{ (}\DecValTok{14}\NormalTok{,) }\KeywordTok{or} \KeywordTok{not}\NormalTok{ np.isfinite(action).}\BuiltInTok{all}\NormalTok{():}
            \ControlFlowTok{raise} \PreprocessorTok{ValueError}\NormalTok{(}\StringTok{"actions must be finite with shape (14,)"}\NormalTok{)}
        \ControlFlowTok{if} \VariableTok{self}\NormalTok{.\_episode\_finished:}
            \ControlFlowTok{raise} \PreprocessorTok{RuntimeError}\NormalTok{(}\StringTok{"reset is required after a terminal episode"}\NormalTok{)}
        \VariableTok{self}\NormalTok{.\_apply\_assistance()}
        \ControlFlowTok{if} \VariableTok{self}\NormalTok{.step\_count }\OperatorTok{*}\NormalTok{ CTRL\_DT }\OperatorTok{\textgreater{}=} \VariableTok{self}\NormalTok{.\_next\_push:}
            \VariableTok{self}\NormalTok{.data.qvel[:}\DecValTok{2}\NormalTok{] }\OperatorTok{=} \VariableTok{self}\NormalTok{.\_rng.uniform(}\OperatorTok{{-}}\FloatTok{.09}\NormalTok{, }\FloatTok{.09}\NormalTok{, }\DecValTok{2}\NormalTok{)}
            \VariableTok{self}\NormalTok{.\_next\_push }\OperatorTok{+=} \VariableTok{self}\NormalTok{.\_rng.uniform(}\FloatTok{1.5}\NormalTok{, }\DecValTok{3}\NormalTok{)}
\NormalTok{        observation, reward, terminated, truncated, info }\OperatorTok{=} \BuiltInTok{super}\NormalTok{().step(action)}
\NormalTok{        metrics }\OperatorTok{=} \VariableTok{self}\NormalTok{.metrics()}
\NormalTok{        reasons }\OperatorTok{=}\NormalTok{ []}
        \ControlFlowTok{if}\NormalTok{ metrics[}\StringTok{"root\_above\_ball\_m"}\NormalTok{] }\OperatorTok{\textless{}}\NormalTok{ BALL\_RADIUS }\OperatorTok{+} \FloatTok{.05}\NormalTok{:}
\NormalTok{            reasons.append(}\StringTok{"height"}\NormalTok{)}
        \ControlFlowTok{if}\NormalTok{ metrics[}\StringTok{"root\_ball\_offset\_m"}\NormalTok{] }\OperatorTok{\textgreater{}} \FloatTok{.16}\NormalTok{:}
\NormalTok{            reasons.append(}\StringTok{"offset"}\NormalTok{)}
        \ControlFlowTok{if}\NormalTok{ metrics[}\StringTok{"tilt\_deg"}\NormalTok{] }\OperatorTok{\textgreater{}} \DecValTok{55}\NormalTok{:}
\NormalTok{            reasons.append(}\StringTok{"tilt"}\NormalTok{)}
        \ControlFlowTok{if}\NormalTok{ metrics[}\StringTok{"robot\_floor\_contact"}\NormalTok{]:}
\NormalTok{            reasons.append(}\StringTok{"robot\_floor\_contact"}\NormalTok{)}
        \ControlFlowTok{if}\NormalTok{ metrics[}\StringTok{"body\_ball\_contact"}\NormalTok{]:}
\NormalTok{            reasons.append(}\StringTok{"body\_ball\_contact"}\NormalTok{)}
        \ControlFlowTok{if} \KeywordTok{not}\NormalTok{ np.isfinite(}\VariableTok{self}\NormalTok{.data.qpos).}\BuiltInTok{all}\NormalTok{() }\KeywordTok{or} \KeywordTok{not}\NormalTok{ np.isfinite(}\VariableTok{self}\NormalTok{.data.qvel).}\BuiltInTok{all}\NormalTok{():}
\NormalTok{            reasons.append(}\StringTok{"nonfinite\_physics"}\NormalTok{)}
\NormalTok{        terminated }\OperatorTok{=}\NormalTok{ terminated }\KeywordTok{or} \BuiltInTok{bool}\NormalTok{(reasons)}
        \VariableTok{self}\NormalTok{.\_episode\_finished }\OperatorTok{=}\NormalTok{ terminated }\KeywordTok{or}\NormalTok{ truncated}
\NormalTok{        info.update(metrics)}
\NormalTok{        info[}\StringTok{"termination\_reasons"}\NormalTok{] }\OperatorTok{=}\NormalTok{ reasons}
        \ControlFlowTok{if} \VariableTok{self}\NormalTok{.\_episode\_finished:}
\NormalTok{            info[}\StringTok{"episode\_rewards"}\NormalTok{] }\OperatorTok{=} \BuiltInTok{dict}\NormalTok{(}\VariableTok{self}\NormalTok{.reward\_sums)}
        \ControlFlowTok{return}\NormalTok{ observation, reward, terminated, truncated, info}

    \KeywordTok{def}\NormalTok{ \_compute\_reward(}\VariableTok{self}\NormalTok{):}
\NormalTok{        relative }\OperatorTok{=} \VariableTok{self}\NormalTok{.data.xpos[}\VariableTok{self}\NormalTok{.trunk\_body\_id] }\OperatorTok{{-}} \VariableTok{self}\NormalTok{.data.xpos[}\VariableTok{self}\NormalTok{.ball\_body\_id]}
\NormalTok{        gravity }\OperatorTok{=} \VariableTok{self}\NormalTok{.\_projected\_gravity()}
\NormalTok{        feet }\OperatorTok{=} \VariableTok{self}\NormalTok{.data.site\_xpos[}\VariableTok{self}\NormalTok{.foot\_site\_ids]}
\NormalTok{        gaps }\OperatorTok{=}\NormalTok{ np.linalg.norm(feet }\OperatorTok{{-}} \VariableTok{self}\NormalTok{.data.xpos[}\VariableTok{self}\NormalTok{.ball\_body\_id], axis}\OperatorTok{=}\DecValTok{1}\NormalTok{) }\OperatorTok{{-}}\NormalTok{ (BALL\_RADIUS }\OperatorTok{+} \FloatTok{.012}\NormalTok{)}
\NormalTok{        ball\_velocity }\OperatorTok{=} \VariableTok{self}\NormalTok{.data.qvel[}\VariableTok{self}\NormalTok{.ball\_qvel\_adr:}\VariableTok{self}\NormalTok{.ball\_qvel\_adr }\OperatorTok{+} \DecValTok{2}\NormalTok{]}
\NormalTok{        terms }\OperatorTok{=}\NormalTok{ \{}
            \StringTok{"linear\_tracking"}\NormalTok{: math.exp(}\OperatorTok{{-}}\BuiltInTok{float}\NormalTok{(np.}\BuiltInTok{sum}\NormalTok{((}\VariableTok{self}\NormalTok{.body\_lin\_vel()[:}\DecValTok{2}\NormalTok{] }\OperatorTok{{-}} \VariableTok{self}\NormalTok{.twist\_cmd[:}\DecValTok{2}\NormalTok{]) }\OperatorTok{**} \DecValTok{2}\NormalTok{)) }\OperatorTok{/} \FloatTok{.1}\NormalTok{),}
            \StringTok{"yaw\_tracking"}\NormalTok{: math.exp(}\OperatorTok{{-}}\BuiltInTok{float}\NormalTok{((}\VariableTok{self}\NormalTok{.\_gyro[}\DecValTok{2}\NormalTok{] }\OperatorTok{{-}} \VariableTok{self}\NormalTok{.twist\_cmd[}\DecValTok{2}\NormalTok{]) }\OperatorTok{**} \DecValTok{2}\NormalTok{) }\OperatorTok{/} \FloatTok{.5}\NormalTok{),}
            \StringTok{"upright"}\NormalTok{: math.exp(}\OperatorTok{{-}}\BuiltInTok{float}\NormalTok{(np.}\BuiltInTok{sum}\NormalTok{(gravity[:}\DecValTok{2}\NormalTok{] }\OperatorTok{**} \DecValTok{2}\NormalTok{)) }\OperatorTok{/} \FloatTok{.09}\NormalTok{),}
            \StringTok{"centered"}\NormalTok{: math.exp(}\OperatorTok{{-}}\BuiltInTok{float}\NormalTok{(np.}\BuiltInTok{sum}\NormalTok{(relative[:}\DecValTok{2}\NormalTok{] }\OperatorTok{**} \DecValTok{2}\NormalTok{)) }\OperatorTok{/} \FloatTok{.0016}\NormalTok{),}
            \StringTok{"feet\_on\_ball"}\NormalTok{: }\BuiltInTok{float}\NormalTok{(np.mean(np.exp(}\OperatorTok{{-}}\NormalTok{gaps }\OperatorTok{**} \DecValTok{2} \OperatorTok{/} \FloatTok{.0004}\NormalTok{))),}
            \StringTok{"height"}\NormalTok{: math.exp(}\OperatorTok{{-}}\BuiltInTok{float}\NormalTok{((relative[}\DecValTok{2}\NormalTok{] }\OperatorTok{{-}}\NormalTok{ BALL\_RADIUS }\OperatorTok{{-}} \FloatTok{.125}\NormalTok{) }\OperatorTok{**} \DecValTok{2}\NormalTok{) }\OperatorTok{/} \FloatTok{.0009}\NormalTok{),}
            \StringTok{"ball\_speed\_penalty"}\NormalTok{: }\OperatorTok{{-}}\BuiltInTok{float}\NormalTok{(}\BuiltInTok{min}\NormalTok{(np.}\BuiltInTok{sum}\NormalTok{(ball\_velocity }\OperatorTok{**} \DecValTok{2}\NormalTok{), }\DecValTok{100}\NormalTok{)),}
            \StringTok{"angular\_velocity\_penalty"}\NormalTok{: }\OperatorTok{{-}}\BuiltInTok{float}\NormalTok{(}\BuiltInTok{min}\NormalTok{(np.}\BuiltInTok{sum}\NormalTok{(}\VariableTok{self}\NormalTok{.\_gyro[:}\DecValTok{2}\NormalTok{] }\OperatorTok{**} \DecValTok{2}\NormalTok{), }\DecValTok{100}\NormalTok{)),}
            \StringTok{"action\_rate\_penalty"}\NormalTok{: }\OperatorTok{{-}}\BuiltInTok{float}\NormalTok{(}\BuiltInTok{min}\NormalTok{(np.}\BuiltInTok{sum}\NormalTok{((}\VariableTok{self}\NormalTok{.last\_action }\OperatorTok{{-}} \VariableTok{self}\NormalTok{.prev\_action) }\OperatorTok{**} \DecValTok{2}\NormalTok{), }\DecValTok{100}\NormalTok{)),}
\NormalTok{        \}}
\NormalTok{        weighted }\OperatorTok{=}\NormalTok{ \{key: value }\OperatorTok{*}\NormalTok{ REWARD\_WEIGHTS[key] }\OperatorTok{*}\NormalTok{ CTRL\_DT }\ControlFlowTok{for}\NormalTok{ key, value }\KeywordTok{in}\NormalTok{ terms.items()\}}
\NormalTok{        reward }\OperatorTok{=} \BuiltInTok{sum}\NormalTok{(weighted.values())}
        \ControlFlowTok{if} \KeywordTok{not}\NormalTok{ math.isfinite(reward):}
            \ControlFlowTok{raise} \PreprocessorTok{FloatingPointError}\NormalTok{(}\StringTok{"nonfinite basketball reward"}\NormalTok{)}
        \ControlFlowTok{return}\NormalTok{ reward, weighted}

    \KeywordTok{def}\NormalTok{ metrics(}\VariableTok{self}\NormalTok{) }\OperatorTok{{-}\textgreater{}} \BuiltInTok{dict}\NormalTok{:}
\NormalTok{        ball\_pos }\OperatorTok{=} \VariableTok{self}\NormalTok{.data.xpos[}\VariableTok{self}\NormalTok{.ball\_body\_id]}
\NormalTok{        relative }\OperatorTok{=} \VariableTok{self}\NormalTok{.data.xpos[}\VariableTok{self}\NormalTok{.trunk\_body\_id] }\OperatorTok{{-}}\NormalTok{ ball\_pos}
\NormalTok{        foot\_geoms }\OperatorTok{=} \BuiltInTok{set}\NormalTok{(}\VariableTok{self}\NormalTok{.foot\_geoms.values())}
\NormalTok{        touching }\OperatorTok{=} \BuiltInTok{set}\NormalTok{()}
\NormalTok{        floor\_contact }\OperatorTok{=} \VariableTok{False}
\NormalTok{        body\_ball\_contact }\OperatorTok{=} \VariableTok{False}
        \ControlFlowTok{for}\NormalTok{ contact }\KeywordTok{in} \VariableTok{self}\NormalTok{.data.contact:}
\NormalTok{            pair }\OperatorTok{=}\NormalTok{ \{}\BuiltInTok{int}\NormalTok{(contact.geom1), }\BuiltInTok{int}\NormalTok{(contact.geom2)\}}
            \ControlFlowTok{if} \VariableTok{self}\NormalTok{.ball\_geom\_id }\KeywordTok{in}\NormalTok{ pair:}
\NormalTok{                other }\OperatorTok{=} \BuiltInTok{next}\NormalTok{(}\BuiltInTok{iter}\NormalTok{(pair }\OperatorTok{{-}}\NormalTok{ \{}\VariableTok{self}\NormalTok{.ball\_geom\_id\}), }\VariableTok{self}\NormalTok{.ball\_geom\_id)}
                \ControlFlowTok{if}\NormalTok{ other }\KeywordTok{in}\NormalTok{ foot\_geoms:}
\NormalTok{                    touching.add(other)}
                \ControlFlowTok{elif} \VariableTok{self}\NormalTok{.model.geom\_bodyid[other] }\KeywordTok{not} \KeywordTok{in}\NormalTok{ (}\DecValTok{0}\NormalTok{, }\VariableTok{self}\NormalTok{.ball\_body\_id):}
\NormalTok{                    body\_ball\_contact }\OperatorTok{=} \VariableTok{True}
            \ControlFlowTok{if} \VariableTok{self}\NormalTok{.floor\_geom }\KeywordTok{in}\NormalTok{ pair:}
\NormalTok{                other }\OperatorTok{=} \BuiltInTok{next}\NormalTok{(}\BuiltInTok{iter}\NormalTok{(pair }\OperatorTok{{-}}\NormalTok{ \{}\VariableTok{self}\NormalTok{.floor\_geom\}), }\VariableTok{self}\NormalTok{.floor\_geom)}
                \ControlFlowTok{if} \VariableTok{self}\NormalTok{.model.geom\_bodyid[other] }\KeywordTok{not} \KeywordTok{in}\NormalTok{ (}\DecValTok{0}\NormalTok{, }\VariableTok{self}\NormalTok{.ball\_body\_id):}
\NormalTok{                    floor\_contact }\OperatorTok{=} \VariableTok{True}
\NormalTok{        ball\_velocity }\OperatorTok{=} \VariableTok{self}\NormalTok{.data.qvel[}\VariableTok{self}\NormalTok{.ball\_qvel\_adr:}\VariableTok{self}\NormalTok{.ball\_qvel\_adr }\OperatorTok{+} \DecValTok{6}\NormalTok{]}
        \ControlFlowTok{return}\NormalTok{ \{}
            \StringTok{"tilt\_deg"}\NormalTok{: math.degrees(math.acos(}\BuiltInTok{float}\NormalTok{(np.clip(}\OperatorTok{{-}}\VariableTok{self}\NormalTok{.\_projected\_gravity()[}\DecValTok{2}\NormalTok{], }\OperatorTok{{-}}\DecValTok{1}\NormalTok{, }\DecValTok{1}\NormalTok{)))),}
            \StringTok{"root\_ball\_offset\_m"}\NormalTok{: }\BuiltInTok{float}\NormalTok{(np.linalg.norm(relative[:}\DecValTok{2}\NormalTok{])),}
            \StringTok{"root\_above\_ball\_m"}\NormalTok{: }\BuiltInTok{float}\NormalTok{(relative[}\DecValTok{2}\NormalTok{]),}
            \StringTok{"ball\_speed\_mps"}\NormalTok{: }\BuiltInTok{float}\NormalTok{(np.linalg.norm(ball\_velocity[:}\DecValTok{2}\NormalTok{])),}
            \StringTok{"ball\_position"}\NormalTok{: ball\_pos.tolist(),}
            \StringTok{"ball\_angular\_speed\_rad\_s"}\NormalTok{: }\BuiltInTok{float}\NormalTok{(np.linalg.norm(ball\_velocity[}\DecValTok{3}\NormalTok{:])),}
            \StringTok{"body\_forward\_mps"}\NormalTok{: }\BuiltInTok{float}\NormalTok{(}\VariableTok{self}\NormalTok{.body\_lin\_vel()[}\DecValTok{0}\NormalTok{]),}
            \StringTok{"body\_lateral\_mps"}\NormalTok{: }\BuiltInTok{float}\NormalTok{(}\VariableTok{self}\NormalTok{.body\_lin\_vel()[}\DecValTok{1}\NormalTok{]),}
            \StringTok{"body\_yaw\_rate\_rad\_s"}\NormalTok{: }\BuiltInTok{float}\NormalTok{(}\VariableTok{self}\NormalTok{.\_gyro[}\DecValTok{2}\NormalTok{]),}
            \StringTok{"foot\_ball\_contacts"}\NormalTok{: }\BuiltInTok{len}\NormalTok{(touching),}
            \StringTok{"body\_ball\_contact"}\NormalTok{: body\_ball\_contact,}
            \StringTok{"robot\_floor\_contact"}\NormalTok{: floor\_contact,}
            \StringTok{"hold"}\NormalTok{: }\VariableTok{self}\NormalTok{.hold,}
            \StringTok{"elapsed\_s"}\NormalTok{: }\VariableTok{self}\NormalTok{.step\_count }\OperatorTok{*}\NormalTok{ CTRL\_DT,}
\NormalTok{        \}}
\end{Highlighting}
\end{Shaded}

\subsection{\texorpdfstring{A2.
\texttt{rlx/\allowbreak{}rlx/\allowbreak{}models/\allowbreak{}basketba\allowbreak{}ll.py}}{A2. }}\label{a2.-rlxrlxmodelsbasketball.py}

SHA-256:
\texttt{df79b783\allowbreak{}e643ed40\allowbreak{}ea19fe58\allowbreak{}3d05dec7\allowbreak{}a5a116cf\allowbreak{}74cced0e\allowbreak{}48776149\allowbreak{}f8ea4e06\allowbreak{}}

\begin{Shaded}
\begin{Highlighting}[]
\CommentTok{"""Torch recurrent policy support for the released blind basketball actor."""}

\ImportTok{from}\NormalTok{ \_\_future\_\_ }\ImportTok{import}\NormalTok{ annotations}

\ImportTok{import}\NormalTok{ hashlib}
\ImportTok{from}\NormalTok{ pathlib }\ImportTok{import}\NormalTok{ Path}
\ImportTok{from}\NormalTok{ typing }\ImportTok{import}\NormalTok{ Any}

\ImportTok{import}\NormalTok{ numpy }\ImportTok{as}\NormalTok{ np}
\ImportTok{import}\NormalTok{ torch}
\ImportTok{from}\NormalTok{ torch }\ImportTok{import}\NormalTok{ nn}

\NormalTok{OBSERVATION\_DIM }\OperatorTok{=} \DecValTok{61}
\NormalTok{ACTION\_DIM }\OperatorTok{=} \DecValTok{14}
\NormalTok{HIDDEN\_SIZE }\OperatorTok{=} \DecValTok{256}
\NormalTok{HIDDEN\_DIMS }\OperatorTok{=}\NormalTok{ (}\DecValTok{512}\NormalTok{, }\DecValTok{256}\NormalTok{, }\DecValTok{128}\NormalTok{)}
\NormalTok{NORMALIZER\_EPSILON }\OperatorTok{=} \FloatTok{0.01}
\NormalTok{SOURCE\_CHECKPOINT\_SHA256 }\OperatorTok{=}\NormalTok{ (}
    \StringTok{"57a322eff092cb71e7cba831791232f0cc4fd45ede0441cad6d87813b0033e41"}
\NormalTok{)}
\NormalTok{SOURCE\_ONNX\_SHA256 }\OperatorTok{=}\NormalTok{ (}
    \StringTok{"e105148b160b3a86170621648215be38dd5018e955930ce182394f4513c7569b"}
\NormalTok{)}
\NormalTok{LOCAL\_CHECKPOINT\_FORMAT }\OperatorTok{=} \StringTok{"rlx.basketball.local\_recurrent\_ppo.v1"}


\KeywordTok{class}\NormalTok{ ObservationNormalizer(nn.Module):}
    \KeywordTok{def} \FunctionTok{\_\_init\_\_}\NormalTok{(}\VariableTok{self}\NormalTok{) }\OperatorTok{{-}\textgreater{}} \VariableTok{None}\NormalTok{:}
        \BuiltInTok{super}\NormalTok{().}\FunctionTok{\_\_init\_\_}\NormalTok{()}
        \VariableTok{self}\NormalTok{.register\_buffer(}\StringTok{"\_mean"}\NormalTok{, torch.zeros(}\DecValTok{1}\NormalTok{, OBSERVATION\_DIM))}
        \VariableTok{self}\NormalTok{.register\_buffer(}\StringTok{"\_var"}\NormalTok{, torch.ones(}\DecValTok{1}\NormalTok{, OBSERVATION\_DIM))}
        \VariableTok{self}\NormalTok{.register\_buffer(}\StringTok{"\_std"}\NormalTok{, torch.ones(}\DecValTok{1}\NormalTok{, OBSERVATION\_DIM))}
        \VariableTok{self}\NormalTok{.register\_buffer(}\StringTok{"count"}\NormalTok{, torch.tensor(}\FloatTok{0.0}\NormalTok{))}

    \KeywordTok{def}\NormalTok{ forward(}\VariableTok{self}\NormalTok{, observations: torch.Tensor) }\OperatorTok{{-}\textgreater{}}\NormalTok{ torch.Tensor:}
        \ControlFlowTok{return}\NormalTok{ (observations }\OperatorTok{{-}} \VariableTok{self}\NormalTok{.\_mean) }\OperatorTok{/}\NormalTok{ (}\VariableTok{self}\NormalTok{.\_std }\OperatorTok{+}\NormalTok{ NORMALIZER\_EPSILON)}


\KeywordTok{class}\NormalTok{ RecurrentCore(nn.Module):}
    \KeywordTok{def} \FunctionTok{\_\_init\_\_}\NormalTok{(}\VariableTok{self}\NormalTok{) }\OperatorTok{{-}\textgreater{}} \VariableTok{None}\NormalTok{:}
        \BuiltInTok{super}\NormalTok{().}\FunctionTok{\_\_init\_\_}\NormalTok{()}
        \VariableTok{self}\NormalTok{.rnn }\OperatorTok{=}\NormalTok{ nn.LSTM(OBSERVATION\_DIM, HIDDEN\_SIZE, num\_layers}\OperatorTok{=}\DecValTok{1}\NormalTok{)}


\KeywordTok{class}\NormalTok{ GaussianDistribution(nn.Module):}
    \KeywordTok{def} \FunctionTok{\_\_init\_\_}\NormalTok{(}\VariableTok{self}\NormalTok{, initial\_std: }\BuiltInTok{float}\NormalTok{) }\OperatorTok{{-}\textgreater{}} \VariableTok{None}\NormalTok{:}
        \BuiltInTok{super}\NormalTok{().}\FunctionTok{\_\_init\_\_}\NormalTok{()}
        \ControlFlowTok{if} \KeywordTok{not}\NormalTok{ np.isfinite(initial\_std) }\KeywordTok{or}\NormalTok{ initial\_std }\OperatorTok{\textless{}=} \FloatTok{0.0}\NormalTok{:}
            \ControlFlowTok{raise} \PreprocessorTok{ValueError}\NormalTok{(}\StringTok{"initial\_std must be finite and positive"}\NormalTok{)}
        \VariableTok{self}\NormalTok{.std\_param }\OperatorTok{=}\NormalTok{ nn.Parameter(torch.full((ACTION\_DIM,), }\BuiltInTok{float}\NormalTok{(initial\_std)))}


\KeywordTok{class}\NormalTok{ BasketballActor(nn.Module):}
    \KeywordTok{def} \FunctionTok{\_\_init\_\_}\NormalTok{(}\VariableTok{self}\NormalTok{, }\OperatorTok{*}\NormalTok{, initial\_std: }\BuiltInTok{float} \OperatorTok{=} \FloatTok{0.03}\NormalTok{) }\OperatorTok{{-}\textgreater{}} \VariableTok{None}\NormalTok{:}
        \BuiltInTok{super}\NormalTok{().}\FunctionTok{\_\_init\_\_}\NormalTok{()}
        \VariableTok{self}\NormalTok{.obs\_normalizer }\OperatorTok{=}\NormalTok{ ObservationNormalizer()}
        \VariableTok{self}\NormalTok{.distribution }\OperatorTok{=}\NormalTok{ GaussianDistribution(initial\_std)}
        \VariableTok{self}\NormalTok{.rnn }\OperatorTok{=}\NormalTok{ RecurrentCore()}
\NormalTok{        layers: }\BuiltInTok{list}\NormalTok{[nn.Module] }\OperatorTok{=}\NormalTok{ []}
\NormalTok{        previous }\OperatorTok{=}\NormalTok{ HIDDEN\_SIZE}
        \ControlFlowTok{for}\NormalTok{ width }\KeywordTok{in}\NormalTok{ HIDDEN\_DIMS:}
\NormalTok{            layers.extend((nn.Linear(previous, width), nn.ELU()))}
\NormalTok{            previous }\OperatorTok{=}\NormalTok{ width}
\NormalTok{        layers.append(nn.Linear(previous, ACTION\_DIM))}
        \VariableTok{self}\NormalTok{.mlp }\OperatorTok{=}\NormalTok{ nn.Sequential(}\OperatorTok{*}\NormalTok{layers)}

    \KeywordTok{def}\NormalTok{ initial\_state(}
        \VariableTok{self}\NormalTok{,}
\NormalTok{        batch\_size: }\BuiltInTok{int}\NormalTok{,}
        \OperatorTok{*}\NormalTok{,}
\NormalTok{        device: torch.device }\OperatorTok{|} \BuiltInTok{str} \OperatorTok{|} \VariableTok{None} \OperatorTok{=} \VariableTok{None}\NormalTok{,}
\NormalTok{        dtype: torch.dtype }\OperatorTok{|} \VariableTok{None} \OperatorTok{=} \VariableTok{None}\NormalTok{,}
\NormalTok{    ) }\OperatorTok{{-}\textgreater{}} \BuiltInTok{tuple}\NormalTok{[torch.Tensor, torch.Tensor]:}
        \ControlFlowTok{if}\NormalTok{ batch\_size }\OperatorTok{\textless{}} \DecValTok{1}\NormalTok{:}
            \ControlFlowTok{raise} \PreprocessorTok{ValueError}\NormalTok{(}\StringTok{"batch\_size must be positive"}\NormalTok{)}
\NormalTok{        parameter }\OperatorTok{=} \BuiltInTok{next}\NormalTok{(}\VariableTok{self}\NormalTok{.parameters())}
\NormalTok{        target\_device }\OperatorTok{=}\NormalTok{ parameter.device }\ControlFlowTok{if}\NormalTok{ device }\KeywordTok{is} \VariableTok{None} \ControlFlowTok{else}\NormalTok{ device}
\NormalTok{        target\_dtype }\OperatorTok{=}\NormalTok{ parameter.dtype }\ControlFlowTok{if}\NormalTok{ dtype }\KeywordTok{is} \VariableTok{None} \ControlFlowTok{else}\NormalTok{ dtype}
\NormalTok{        shape }\OperatorTok{=}\NormalTok{ (}\DecValTok{1}\NormalTok{, batch\_size, HIDDEN\_SIZE)}
        \ControlFlowTok{return}\NormalTok{ (}
\NormalTok{            torch.zeros(shape, device}\OperatorTok{=}\NormalTok{target\_device, dtype}\OperatorTok{=}\NormalTok{target\_dtype),}
\NormalTok{            torch.zeros(shape, device}\OperatorTok{=}\NormalTok{target\_device, dtype}\OperatorTok{=}\NormalTok{target\_dtype),}
\NormalTok{        )}

    \KeywordTok{def}\NormalTok{ forward(}
        \VariableTok{self}\NormalTok{,}
\NormalTok{        observations: torch.Tensor,}
\NormalTok{        h\_in: torch.Tensor,}
\NormalTok{        c\_in: torch.Tensor,}
\NormalTok{    ) }\OperatorTok{{-}\textgreater{}} \BuiltInTok{tuple}\NormalTok{[torch.Tensor, torch.Tensor, torch.Tensor]:}
\NormalTok{        normalized }\OperatorTok{=} \VariableTok{self}\NormalTok{.obs\_normalizer(observations)}
\NormalTok{        recurrent, (h\_out, c\_out) }\OperatorTok{=} \VariableTok{self}\NormalTok{.rnn.rnn(}
\NormalTok{            normalized.unsqueeze(}\DecValTok{0}\NormalTok{),}
\NormalTok{            (h\_in, c\_in),}
\NormalTok{        )}
\NormalTok{        actions }\OperatorTok{=} \VariableTok{self}\NormalTok{.mlp(recurrent.squeeze(}\DecValTok{0}\NormalTok{))}
        \ControlFlowTok{return}\NormalTok{ actions, h\_out, c\_out}

    \KeywordTok{def}\NormalTok{ forward\_sequence(}
        \VariableTok{self}\NormalTok{,}
\NormalTok{        observations: torch.Tensor,}
\NormalTok{        h\_in: torch.Tensor,}
\NormalTok{        c\_in: torch.Tensor,}
\NormalTok{        episode\_starts: torch.Tensor,}
\NormalTok{    ) }\OperatorTok{{-}\textgreater{}} \BuiltInTok{tuple}\NormalTok{[torch.Tensor, torch.Tensor, torch.Tensor]:}
        \ControlFlowTok{if}\NormalTok{ observations.ndim }\OperatorTok{!=} \DecValTok{3} \KeywordTok{or}\NormalTok{ observations.shape[}\OperatorTok{{-}}\DecValTok{1}\NormalTok{] }\OperatorTok{!=}\NormalTok{ OBSERVATION\_DIM:}
            \ControlFlowTok{raise} \PreprocessorTok{ValueError}\NormalTok{(}\StringTok{"observations must have shape [time, batch, 61]"}\NormalTok{)}
        \ControlFlowTok{if}\NormalTok{ episode\_starts.shape }\OperatorTok{!=}\NormalTok{ observations.shape[:}\DecValTok{2}\NormalTok{]:}
            \ControlFlowTok{raise} \PreprocessorTok{ValueError}\NormalTok{(}\StringTok{"episode\_starts must have shape [time, batch]"}\NormalTok{)}
\NormalTok{        hidden }\OperatorTok{=}\NormalTok{ h\_in}
\NormalTok{        cell }\OperatorTok{=}\NormalTok{ c\_in}
\NormalTok{        outputs: }\BuiltInTok{list}\NormalTok{[torch.Tensor] }\OperatorTok{=}\NormalTok{ []}
        \ControlFlowTok{for}\NormalTok{ step }\KeywordTok{in} \BuiltInTok{range}\NormalTok{(observations.shape[}\DecValTok{0}\NormalTok{]):}
\NormalTok{            carry }\OperatorTok{=}\NormalTok{ (}\OperatorTok{\textasciitilde{}}\NormalTok{episode\_starts[step].}\BuiltInTok{bool}\NormalTok{()).to(observations.dtype)}
\NormalTok{            carry }\OperatorTok{=}\NormalTok{ carry.view(}\DecValTok{1}\NormalTok{, }\OperatorTok{{-}}\DecValTok{1}\NormalTok{, }\DecValTok{1}\NormalTok{)}
\NormalTok{            hidden }\OperatorTok{=}\NormalTok{ hidden }\OperatorTok{*}\NormalTok{ carry}
\NormalTok{            cell }\OperatorTok{=}\NormalTok{ cell }\OperatorTok{*}\NormalTok{ carry}
\NormalTok{            actions, hidden, cell }\OperatorTok{=} \VariableTok{self}\NormalTok{(observations[step], hidden, cell)}
\NormalTok{            outputs.append(actions)}
        \ControlFlowTok{return}\NormalTok{ torch.stack(outputs), hidden, cell}

    \KeywordTok{def}\NormalTok{ action\_std(}\VariableTok{self}\NormalTok{) }\OperatorTok{{-}\textgreater{}}\NormalTok{ torch.Tensor:}
        \ControlFlowTok{return} \VariableTok{self}\NormalTok{.distribution.std\_param.clamp\_min(}\FloatTok{1e{-}6}\NormalTok{)}


\KeywordTok{class}\NormalTok{ BasketballCritic(nn.Module):}
    \KeywordTok{def} \FunctionTok{\_\_init\_\_}\NormalTok{(}\VariableTok{self}\NormalTok{, normalizer: ObservationNormalizer) }\OperatorTok{{-}\textgreater{}} \VariableTok{None}\NormalTok{:}
        \BuiltInTok{super}\NormalTok{().}\FunctionTok{\_\_init\_\_}\NormalTok{()}
        \VariableTok{self}\NormalTok{.normalizer }\OperatorTok{=}\NormalTok{ normalizer}
\NormalTok{        layers: }\BuiltInTok{list}\NormalTok{[nn.Module] }\OperatorTok{=}\NormalTok{ []}
\NormalTok{        previous }\OperatorTok{=}\NormalTok{ OBSERVATION\_DIM}
        \ControlFlowTok{for}\NormalTok{ width }\KeywordTok{in}\NormalTok{ HIDDEN\_DIMS:}
\NormalTok{            layers.extend((nn.Linear(previous, width), nn.ELU()))}
\NormalTok{            previous }\OperatorTok{=}\NormalTok{ width}
\NormalTok{        layers.append(nn.Linear(previous, }\DecValTok{1}\NormalTok{))}
        \VariableTok{self}\NormalTok{.mlp }\OperatorTok{=}\NormalTok{ nn.Sequential(}\OperatorTok{*}\NormalTok{layers)}

    \KeywordTok{def}\NormalTok{ forward(}\VariableTok{self}\NormalTok{, observations: torch.Tensor) }\OperatorTok{{-}\textgreater{}}\NormalTok{ torch.Tensor:}
        \ControlFlowTok{with}\NormalTok{ torch.no\_grad():}
\NormalTok{            normalized }\OperatorTok{=} \VariableTok{self}\NormalTok{.normalizer(observations)}
        \ControlFlowTok{return} \VariableTok{self}\NormalTok{.mlp(normalized).squeeze(}\OperatorTok{{-}}\DecValTok{1}\NormalTok{)}


\KeywordTok{def}\NormalTok{ sha256\_file(path: }\BuiltInTok{str} \OperatorTok{|}\NormalTok{ Path) }\OperatorTok{{-}\textgreater{}} \BuiltInTok{str}\NormalTok{:}
\NormalTok{    digest }\OperatorTok{=}\NormalTok{ hashlib.sha256()}
    \ControlFlowTok{with}\NormalTok{ Path(path).}\BuiltInTok{open}\NormalTok{(}\StringTok{"rb"}\NormalTok{) }\ImportTok{as}\NormalTok{ stream:}
        \ControlFlowTok{for}\NormalTok{ block }\KeywordTok{in} \BuiltInTok{iter}\NormalTok{(}\KeywordTok{lambda}\NormalTok{: stream.read(}\DecValTok{1024} \OperatorTok{*} \DecValTok{1024}\NormalTok{), }\StringTok{b""}\NormalTok{):}
\NormalTok{            digest.update(block)}
    \ControlFlowTok{return}\NormalTok{ digest.hexdigest()}


\KeywordTok{def}\NormalTok{ verify\_source\_artifacts(}
\NormalTok{    checkpoint: }\BuiltInTok{str} \OperatorTok{|}\NormalTok{ Path,}
\NormalTok{    onnx\_path: }\BuiltInTok{str} \OperatorTok{|}\NormalTok{ Path }\OperatorTok{|} \VariableTok{None} \OperatorTok{=} \VariableTok{None}\NormalTok{,}
\NormalTok{) }\OperatorTok{{-}\textgreater{}} \BuiltInTok{dict}\NormalTok{[}\BuiltInTok{str}\NormalTok{, }\BuiltInTok{str}\NormalTok{]:}
\NormalTok{    checkpoint\_path }\OperatorTok{=}\NormalTok{ Path(checkpoint)}
\NormalTok{    checkpoint\_hash }\OperatorTok{=}\NormalTok{ sha256\_file(checkpoint\_path)}
    \ControlFlowTok{if}\NormalTok{ checkpoint\_hash }\OperatorTok{!=}\NormalTok{ SOURCE\_CHECKPOINT\_SHA256:}
        \ControlFlowTok{raise} \PreprocessorTok{ValueError}\NormalTok{(}
            \SpecialStringTok{f"unexpected basketball checkpoint SHA256: }\SpecialCharTok{\{}\NormalTok{checkpoint\_hash}\SpecialCharTok{\}}\SpecialStringTok{"}
\NormalTok{        )}
\NormalTok{    result }\OperatorTok{=}\NormalTok{ \{}\StringTok{"checkpoint\_sha256"}\NormalTok{: checkpoint\_hash\}}
\NormalTok{    candidate }\OperatorTok{=}\NormalTok{ Path(onnx\_path) }\ControlFlowTok{if}\NormalTok{ onnx\_path }\KeywordTok{is} \KeywordTok{not} \VariableTok{None} \ControlFlowTok{else}\NormalTok{ checkpoint\_path.with\_name(}
        \StringTok{"policy.onnx"}
\NormalTok{    )}
    \ControlFlowTok{if} \KeywordTok{not}\NormalTok{ candidate.is\_file():}
        \ControlFlowTok{raise} \PreprocessorTok{FileNotFoundError}\NormalTok{(}\SpecialStringTok{f"required source ONNX is missing: }\SpecialCharTok{\{}\NormalTok{candidate}\SpecialCharTok{\}}\SpecialStringTok{"}\NormalTok{)}
\NormalTok{    onnx\_hash }\OperatorTok{=}\NormalTok{ sha256\_file(candidate)}
    \ControlFlowTok{if}\NormalTok{ onnx\_hash }\OperatorTok{!=}\NormalTok{ SOURCE\_ONNX\_SHA256:}
        \ControlFlowTok{raise} \PreprocessorTok{ValueError}\NormalTok{(}\SpecialStringTok{f"unexpected basketball ONNX SHA256: }\SpecialCharTok{\{}\NormalTok{onnx\_hash}\SpecialCharTok{\}}\SpecialStringTok{"}\NormalTok{)}
\NormalTok{    result[}\StringTok{"onnx\_sha256"}\NormalTok{] }\OperatorTok{=}\NormalTok{ onnx\_hash}
    \ControlFlowTok{return}\NormalTok{ result}


\KeywordTok{def}\NormalTok{ load\_source\_actor(}
\NormalTok{    checkpoint: }\BuiltInTok{str} \OperatorTok{|}\NormalTok{ Path,}
    \OperatorTok{*}\NormalTok{,}
\NormalTok{    initial\_std: }\BuiltInTok{float} \OperatorTok{=} \FloatTok{0.03}\NormalTok{,}
\NormalTok{    verify\_hash: }\BuiltInTok{bool} \OperatorTok{=} \VariableTok{True}\NormalTok{,}
\NormalTok{    allow\_local: }\BuiltInTok{bool} \OperatorTok{=} \VariableTok{False}\NormalTok{,}
\NormalTok{) }\OperatorTok{{-}\textgreater{}} \BuiltInTok{tuple}\NormalTok{[BasketballActor, }\BuiltInTok{dict}\NormalTok{[}\BuiltInTok{str}\NormalTok{, Any]]:}
\NormalTok{    checkpoint\_path }\OperatorTok{=}\NormalTok{ Path(checkpoint)}
\NormalTok{    local\_candidate }\OperatorTok{=}\NormalTok{ (}
\NormalTok{        allow\_local }\KeywordTok{and}\NormalTok{ verify\_hash}
        \KeywordTok{and}\NormalTok{ sha256\_file(checkpoint\_path) }\OperatorTok{!=}\NormalTok{ SOURCE\_CHECKPOINT\_SHA256}
\NormalTok{    )}
    \ControlFlowTok{if}\NormalTok{ local\_candidate:}
\NormalTok{        policy\_path }\OperatorTok{=}\NormalTok{ checkpoint\_path.with\_name(}\StringTok{"policy.onnx"}\NormalTok{)}
\NormalTok{        hashes }\OperatorTok{=}\NormalTok{ \{}
            \StringTok{"checkpoint\_sha256"}\NormalTok{: sha256\_file(checkpoint\_path),}
            \StringTok{"onnx\_sha256"}\NormalTok{: sha256\_file(policy\_path),}
\NormalTok{        \}}
    \ControlFlowTok{else}\NormalTok{:}
\NormalTok{        hashes }\OperatorTok{=}\NormalTok{ verify\_source\_artifacts(checkpoint\_path) }\ControlFlowTok{if}\NormalTok{ verify\_hash }\ControlFlowTok{else}\NormalTok{ \{\}}
\NormalTok{    payload }\OperatorTok{=}\NormalTok{ torch.load(checkpoint\_path, map\_location}\OperatorTok{=}\StringTok{"cpu"}\NormalTok{, weights\_only}\OperatorTok{=}\VariableTok{True}\NormalTok{)}
    \ControlFlowTok{if} \KeywordTok{not} \BuiltInTok{isinstance}\NormalTok{(payload, }\BuiltInTok{dict}\NormalTok{) }\KeywordTok{or} \StringTok{"actor\_state\_dict"} \KeywordTok{not} \KeywordTok{in}\NormalTok{ payload:}
        \ControlFlowTok{raise} \PreprocessorTok{ValueError}\NormalTok{(}\StringTok{"basketball checkpoint is missing actor\_state\_dict"}\NormalTok{)}
\NormalTok{    source\_steps }\OperatorTok{=} \DecValTok{0}
    \ControlFlowTok{if}\NormalTok{ local\_candidate:}
        \ControlFlowTok{if}\NormalTok{ payload.get(}\StringTok{"format"}\NormalTok{) }\OperatorTok{!=}\NormalTok{ LOCAL\_CHECKPOINT\_FORMAT:}
            \ControlFlowTok{raise} \PreprocessorTok{ValueError}\NormalTok{(}\StringTok{"unsupported local basketball checkpoint format"}\NormalTok{)}
        \ControlFlowTok{if}\NormalTok{ (}
\NormalTok{            payload.get(}\StringTok{"normalization\_frozen"}\NormalTok{) }\KeywordTok{is} \KeywordTok{not} \VariableTok{True}
            \KeywordTok{or}\NormalTok{ payload.get(}\StringTok{"actor\_observation\_dim"}\NormalTok{) }\OperatorTok{!=}\NormalTok{ OBSERVATION\_DIM}
            \KeywordTok{or}\NormalTok{ payload.get(}\StringTok{"critic\_observation\_dim"}\NormalTok{) }\OperatorTok{!=}\NormalTok{ OBSERVATION\_DIM}
\NormalTok{        ):}
            \ControlFlowTok{raise} \PreprocessorTok{ValueError}\NormalTok{(}\StringTok{"invalid local basketball normalization/observation contract"}\NormalTok{)}
\NormalTok{        source\_steps }\OperatorTok{=}\NormalTok{ payload.get(}\StringTok{"total\_steps"}\NormalTok{, payload.get(}\StringTok{"local\_steps"}\NormalTok{))}
        \ControlFlowTok{if} \BuiltInTok{type}\NormalTok{(source\_steps) }\KeywordTok{is} \KeywordTok{not} \BuiltInTok{int} \KeywordTok{or}\NormalTok{ source\_steps }\OperatorTok{\textless{}} \DecValTok{0}\NormalTok{:}
            \ControlFlowTok{raise} \PreprocessorTok{ValueError}\NormalTok{(}\StringTok{"local basketball checkpoint requires nonnegative training steps"}\NormalTok{)}
\NormalTok{    actor }\OperatorTok{=}\NormalTok{ BasketballActor(initial\_std}\OperatorTok{=}\NormalTok{initial\_std)}
\NormalTok{    actor.load\_state\_dict(payload[}\StringTok{"actor\_state\_dict"}\NormalTok{], strict}\OperatorTok{=}\VariableTok{True}\NormalTok{)}
    \ControlFlowTok{if} \BuiltInTok{any}\NormalTok{(}\KeywordTok{not}\NormalTok{ torch.isfinite(value).}\BuiltInTok{all}\NormalTok{() }\ControlFlowTok{for}\NormalTok{ value }\KeywordTok{in}\NormalTok{ actor.state\_dict().values()):}
        \ControlFlowTok{raise} \PreprocessorTok{ValueError}\NormalTok{(}\StringTok{"basketball actor contains non{-}finite parameters"}\NormalTok{)}
\NormalTok{    parity }\OperatorTok{=} \VariableTok{None}
    \ControlFlowTok{if}\NormalTok{ local\_candidate:}
\NormalTok{        parity }\OperatorTok{=}\NormalTok{ compare\_actor\_to\_onnx(actor, policy\_path)}
\NormalTok{        error }\OperatorTok{=} \BuiltInTok{max}\NormalTok{(}\BuiltInTok{float}\NormalTok{(parity[}\StringTok{"max\_abs\_error"}\NormalTok{]), }\BuiltInTok{float}\NormalTok{(parity[}\StringTok{"max\_state\_abs\_error"}\NormalTok{]))}
        \ControlFlowTok{if} \KeywordTok{not}\NormalTok{ np.isfinite(error) }\KeywordTok{or}\NormalTok{ error }\OperatorTok{\textgreater{}} \FloatTok{3e{-}5}\NormalTok{:}
            \ControlFlowTok{raise} \PreprocessorTok{ValueError}\NormalTok{(}\SpecialStringTok{f"local basketball checkpoint/ONNX parity failed: }\SpecialCharTok{\{}\NormalTok{error}\SpecialCharTok{\}}\SpecialStringTok{"}\NormalTok{)}
        \ControlFlowTok{if}\NormalTok{ hashes }\OperatorTok{!=}\NormalTok{ \{}
            \StringTok{"checkpoint\_sha256"}\NormalTok{: sha256\_file(checkpoint\_path),}
            \StringTok{"onnx\_sha256"}\NormalTok{: sha256\_file(policy\_path),}
\NormalTok{        \}:}
            \ControlFlowTok{raise} \PreprocessorTok{RuntimeError}\NormalTok{(}\StringTok{"local basketball source changed during validation"}\NormalTok{)}
    \ControlFlowTok{with}\NormalTok{ torch.no\_grad():}
\NormalTok{        actor.distribution.std\_param.fill\_(}\BuiltInTok{float}\NormalTok{(initial\_std))}
\NormalTok{    actor.obs\_normalizer.requires\_grad\_(}\VariableTok{False}\NormalTok{)}
\NormalTok{    metadata }\OperatorTok{=}\NormalTok{ \{}
        \StringTok{"source\_iteration"}\NormalTok{: }\BuiltInTok{int}\NormalTok{(payload.get(}\StringTok{"iter"}\NormalTok{, }\OperatorTok{{-}}\DecValTok{1}\NormalTok{)),}
        \StringTok{"source\_hashes"}\NormalTok{: hashes,}
        \StringTok{"source\_actor\_keys"}\NormalTok{: }\BuiltInTok{sorted}\NormalTok{(payload[}\StringTok{"actor\_state\_dict"}\NormalTok{]),}
        \StringTok{"source\_kind"}\NormalTok{: }\StringTok{"local\_actor\_warm\_start"} \ControlFlowTok{if}\NormalTok{ local\_candidate }\ControlFlowTok{else} \StringTok{"released\_actor"}\NormalTok{,}
        \StringTok{"source\_training\_steps"}\NormalTok{: source\_steps,}
        \StringTok{"source\_onnx\_parity"}\NormalTok{: parity,}
\NormalTok{    \}}
    \ControlFlowTok{return}\NormalTok{ actor, metadata}


\KeywordTok{def}\NormalTok{ export\_recurrent\_onnx(}
\NormalTok{    actor: BasketballActor,}
\NormalTok{    output: }\BuiltInTok{str} \OperatorTok{|}\NormalTok{ Path,}
\NormalTok{) }\OperatorTok{{-}\textgreater{}}\NormalTok{ Path:}
\NormalTok{    target }\OperatorTok{=}\NormalTok{ Path(output)}
\NormalTok{    target.parent.mkdir(parents}\OperatorTok{=}\VariableTok{True}\NormalTok{, exist\_ok}\OperatorTok{=}\VariableTok{True}\NormalTok{)}
\NormalTok{    actor }\OperatorTok{=}\NormalTok{ actor.cpu().}\BuiltInTok{eval}\NormalTok{()}
\NormalTok{    h\_in, c\_in }\OperatorTok{=}\NormalTok{ actor.initial\_state(}\DecValTok{1}\NormalTok{, device}\OperatorTok{=}\StringTok{"cpu"}\NormalTok{, dtype}\OperatorTok{=}\NormalTok{torch.float32)}
\NormalTok{    torch.onnx.export(}
\NormalTok{        actor,}
\NormalTok{        (torch.zeros(}\DecValTok{1}\NormalTok{, OBSERVATION\_DIM), h\_in, c\_in),}
\NormalTok{        target,}
\NormalTok{        input\_names}\OperatorTok{=}\NormalTok{[}\StringTok{"obs"}\NormalTok{, }\StringTok{"h\_in"}\NormalTok{, }\StringTok{"c\_in"}\NormalTok{],}
\NormalTok{        output\_names}\OperatorTok{=}\NormalTok{[}\StringTok{"actions"}\NormalTok{, }\StringTok{"h\_out"}\NormalTok{, }\StringTok{"c\_out"}\NormalTok{],}
\NormalTok{        opset\_version}\OperatorTok{=}\DecValTok{17}\NormalTok{,}
\NormalTok{        do\_constant\_folding}\OperatorTok{=}\VariableTok{True}\NormalTok{,}
\NormalTok{        dynamo}\OperatorTok{=}\VariableTok{False}\NormalTok{,}
\NormalTok{    )}
    \ControlFlowTok{return}\NormalTok{ target}


\KeywordTok{def}\NormalTok{ compare\_actor\_to\_onnx(}
\NormalTok{    actor: BasketballActor,}
\NormalTok{    onnx\_path: }\BuiltInTok{str} \OperatorTok{|}\NormalTok{ Path,}
    \OperatorTok{*}\NormalTok{,}
\NormalTok{    steps: }\BuiltInTok{int} \OperatorTok{=} \DecValTok{40}\NormalTok{,}
\NormalTok{    reset\_at: }\BuiltInTok{int} \OperatorTok{=} \DecValTok{20}\NormalTok{,}
\NormalTok{    seed: }\BuiltInTok{int} \OperatorTok{=} \DecValTok{42}\NormalTok{,}
\NormalTok{) }\OperatorTok{{-}\textgreater{}} \BuiltInTok{dict}\NormalTok{[}\BuiltInTok{str}\NormalTok{, }\BuiltInTok{float} \OperatorTok{|} \BuiltInTok{int}\NormalTok{]:}
    \ImportTok{import}\NormalTok{ onnxruntime }\ImportTok{as}\NormalTok{ ort}

    \ControlFlowTok{if}\NormalTok{ steps }\OperatorTok{\textless{}} \DecValTok{1} \KeywordTok{or} \KeywordTok{not} \DecValTok{0} \OperatorTok{\textless{}=}\NormalTok{ reset\_at }\OperatorTok{\textless{}}\NormalTok{ steps:}
        \ControlFlowTok{raise} \PreprocessorTok{ValueError}\NormalTok{(}\StringTok{"reset\_at must identify a step in the comparison"}\NormalTok{)}
\NormalTok{    generator }\OperatorTok{=}\NormalTok{ torch.Generator().manual\_seed(seed)}
\NormalTok{    options }\OperatorTok{=}\NormalTok{ ort.SessionOptions()}
\NormalTok{    options.intra\_op\_num\_threads }\OperatorTok{=} \DecValTok{1}
\NormalTok{    options.inter\_op\_num\_threads }\OperatorTok{=} \DecValTok{1}
\NormalTok{    session }\OperatorTok{=}\NormalTok{ ort.InferenceSession(}
        \BuiltInTok{str}\NormalTok{(onnx\_path),}
\NormalTok{        sess\_options}\OperatorTok{=}\NormalTok{options,}
\NormalTok{        providers}\OperatorTok{=}\NormalTok{[}\StringTok{"CPUExecutionProvider"}\NormalTok{],}
\NormalTok{    )}
    \ControlFlowTok{if}\NormalTok{ [item.name }\ControlFlowTok{for}\NormalTok{ item }\KeywordTok{in}\NormalTok{ session.get\_inputs()] }\OperatorTok{!=}\NormalTok{ [}\StringTok{"obs"}\NormalTok{, }\StringTok{"h\_in"}\NormalTok{, }\StringTok{"c\_in"}\NormalTok{]:}
        \ControlFlowTok{raise} \PreprocessorTok{ValueError}\NormalTok{(}\StringTok{"unexpected recurrent ONNX input contract"}\NormalTok{)}
    \ControlFlowTok{if}\NormalTok{ [item.name }\ControlFlowTok{for}\NormalTok{ item }\KeywordTok{in}\NormalTok{ session.get\_outputs()] }\OperatorTok{!=}\NormalTok{ [}
        \StringTok{"actions"}\NormalTok{,}
        \StringTok{"h\_out"}\NormalTok{,}
        \StringTok{"c\_out"}\NormalTok{,}
\NormalTok{    ]:}
        \ControlFlowTok{raise} \PreprocessorTok{ValueError}\NormalTok{(}\StringTok{"unexpected recurrent ONNX output contract"}\NormalTok{)}
\NormalTok{    actor }\OperatorTok{=}\NormalTok{ actor.cpu().}\BuiltInTok{eval}\NormalTok{()}
\NormalTok{    hidden, cell }\OperatorTok{=}\NormalTok{ actor.initial\_state(}\DecValTok{1}\NormalTok{, device}\OperatorTok{=}\StringTok{"cpu"}\NormalTok{, dtype}\OperatorTok{=}\NormalTok{torch.float32)}
\NormalTok{    onnx\_hidden }\OperatorTok{=}\NormalTok{ hidden.numpy()}
\NormalTok{    onnx\_cell }\OperatorTok{=}\NormalTok{ cell.numpy()}
\NormalTok{    max\_error }\OperatorTok{=} \FloatTok{0.0}
\NormalTok{    max\_state\_error }\OperatorTok{=} \FloatTok{0.0}
    \ControlFlowTok{with}\NormalTok{ torch.inference\_mode():}
        \ControlFlowTok{for}\NormalTok{ step }\KeywordTok{in} \BuiltInTok{range}\NormalTok{(steps):}
            \ControlFlowTok{if}\NormalTok{ step }\OperatorTok{==}\NormalTok{ reset\_at:}
\NormalTok{                hidden.zero\_()}
\NormalTok{                cell.zero\_()}
\NormalTok{                onnx\_hidden.fill(}\FloatTok{0.0}\NormalTok{)}
\NormalTok{                onnx\_cell.fill(}\FloatTok{0.0}\NormalTok{)}
\NormalTok{            observation }\OperatorTok{=}\NormalTok{ torch.randn(}
                \DecValTok{1}\NormalTok{,}
\NormalTok{                OBSERVATION\_DIM,}
\NormalTok{                generator}\OperatorTok{=}\NormalTok{generator,}
\NormalTok{            ) }\OperatorTok{*} \FloatTok{0.2}
\NormalTok{            observation[:, }\OperatorTok{{-}}\DecValTok{10}\NormalTok{:] }\OperatorTok{=} \FloatTok{0.0}
\NormalTok{            expected, hidden, cell }\OperatorTok{=}\NormalTok{ actor(observation, hidden, cell)}
\NormalTok{            actual, onnx\_hidden, onnx\_cell }\OperatorTok{=}\NormalTok{ session.run(}
                \VariableTok{None}\NormalTok{,}
\NormalTok{                \{}
                    \StringTok{"obs"}\NormalTok{: observation.numpy(),}
                    \StringTok{"h\_in"}\NormalTok{: onnx\_hidden,}
                    \StringTok{"c\_in"}\NormalTok{: onnx\_cell,}
\NormalTok{                \},}
\NormalTok{            )}
            \ControlFlowTok{for}\NormalTok{ index, (torch\_value, onnx\_value) }\KeywordTok{in} \BuiltInTok{enumerate}\NormalTok{((}
\NormalTok{                (expected, actual), (hidden, onnx\_hidden), (cell, onnx\_cell),}
\NormalTok{            )):}
\NormalTok{                expected\_array }\OperatorTok{=}\NormalTok{ torch\_value.numpy()}
                \ControlFlowTok{if} \KeywordTok{not}\NormalTok{ np.isfinite(expected\_array).}\BuiltInTok{all}\NormalTok{() }\KeywordTok{or} \KeywordTok{not}\NormalTok{ np.isfinite(onnx\_value).}\BuiltInTok{all}\NormalTok{():}
                    \ControlFlowTok{raise} \PreprocessorTok{ValueError}\NormalTok{(}\StringTok{"non{-}finite recurrent ONNX parity output"}\NormalTok{)}
\NormalTok{                error }\OperatorTok{=} \BuiltInTok{float}\NormalTok{(np.}\BuiltInTok{max}\NormalTok{(np.}\BuiltInTok{abs}\NormalTok{(onnx\_value }\OperatorTok{{-}}\NormalTok{ expected\_array)))}
                \ControlFlowTok{if}\NormalTok{ index }\OperatorTok{==} \DecValTok{0}\NormalTok{:}
\NormalTok{                    max\_error }\OperatorTok{=} \BuiltInTok{max}\NormalTok{(max\_error, error)}
                \ControlFlowTok{else}\NormalTok{:}
\NormalTok{                    max\_state\_error }\OperatorTok{=} \BuiltInTok{max}\NormalTok{(max\_state\_error, error)}
    \ControlFlowTok{return}\NormalTok{ \{}
        \StringTok{"steps"}\NormalTok{: steps, }\StringTok{"reset\_at"}\NormalTok{: reset\_at, }\StringTok{"max\_abs\_error"}\NormalTok{: max\_error,}
        \StringTok{"max\_state\_abs\_error"}\NormalTok{: max\_state\_error,}
\NormalTok{    \}}
\end{Highlighting}
\end{Shaded}

\subsection{\texorpdfstring{A3.
\texttt{rlx/\allowbreak{}examples\allowbreak{}/\allowbreak{}ppo\_\allowbreak{}microduc\allowbreak{}k\_\allowbreak{}basketba\allowbreak{}ll.py}}{A3. }}\label{a3.-rlxexamplesppo_microduck_basketball.py}

SHA-256:
\texttt{d465af0f\allowbreak{}34751d61\allowbreak{}ed6b8ea2\allowbreak{}33178f9b\allowbreak{}70472ab8\allowbreak{}39c5724f\allowbreak{}aef811e4\allowbreak{}206b9b71\allowbreak{}}

\begin{Shaded}
\begin{Highlighting}[]
\CommentTok{"""Continue the released recurrent basketball actor with local CPU PPO."""}

\ImportTok{from}\NormalTok{ \_\_future\_\_ }\ImportTok{import}\NormalTok{ annotations}

\ImportTok{import}\NormalTok{ argparse}
\ImportTok{from}\NormalTok{ dataclasses }\ImportTok{import}\NormalTok{ dataclass}
\ImportTok{from}\NormalTok{ datetime }\ImportTok{import}\NormalTok{ datetime, timezone}
\ImportTok{import}\NormalTok{ json}
\ImportTok{import}\NormalTok{ math}
\ImportTok{from}\NormalTok{ pathlib }\ImportTok{import}\NormalTok{ Path}
\ImportTok{from}\NormalTok{ typing }\ImportTok{import}\NormalTok{ Any, Sequence}

\ImportTok{import}\NormalTok{ numpy }\ImportTok{as}\NormalTok{ np}
\ImportTok{import}\NormalTok{ torch}
\ImportTok{from}\NormalTok{ torch }\ImportTok{import}\NormalTok{ nn}

\ImportTok{from}\NormalTok{ rlx.models.basketball }\ImportTok{import}\NormalTok{ (}
\NormalTok{    LOCAL\_CHECKPOINT\_FORMAT,}
\NormalTok{    OBSERVATION\_DIM,}
\NormalTok{    BasketballActor,}
\NormalTok{    BasketballCritic,}
\NormalTok{    compare\_actor\_to\_onnx,}
\NormalTok{    export\_recurrent\_onnx,}
\NormalTok{    load\_source\_actor,}
\NormalTok{    sha256\_file,}
\NormalTok{)}

\NormalTok{CTRL\_DT }\OperatorTok{=} \FloatTok{0.02}
\NormalTok{DEFAULT\_CHECKPOINT }\OperatorTok{=}\NormalTok{ (}
\NormalTok{    Path(}\VariableTok{\_\_file\_\_}\NormalTok{).resolve().parents[}\DecValTok{2}\NormalTok{]}
    \OperatorTok{/} \StringTok{"microduck{-}playground"}
    \OperatorTok{/} \StringTok{"artifacts"}
    \OperatorTok{/} \StringTok{"basketball"}
    \OperatorTok{/} \StringTok{"checkpoint.pt"}
\NormalTok{)}


\AttributeTok{@dataclass}
\KeywordTok{class}\NormalTok{ Rollout:}
\NormalTok{    observations: torch.Tensor}
\NormalTok{    actions: torch.Tensor}
\NormalTok{    old\_log\_probs: torch.Tensor}
\NormalTok{    old\_values: torch.Tensor}
\NormalTok{    rewards: torch.Tensor}
\NormalTok{    terminated: torch.Tensor}
\NormalTok{    truncated: torch.Tensor}
\NormalTok{    bootstrap\_values: torch.Tensor}
\NormalTok{    episode\_starts: torch.Tensor}
\NormalTok{    initial\_hidden: torch.Tensor}
\NormalTok{    initial\_cell: torch.Tensor}
\NormalTok{    next\_hidden: torch.Tensor}
\NormalTok{    next\_cell: torch.Tensor}
\NormalTok{    next\_episode\_starts: torch.Tensor}
\NormalTok{    infos: }\BuiltInTok{list}\NormalTok{[}\BuiltInTok{list}\NormalTok{[}\BuiltInTok{dict}\NormalTok{[}\BuiltInTok{str}\NormalTok{, Any]]]}


\KeywordTok{def}\NormalTok{ gaussian\_log\_prob(}
\NormalTok{    actions: torch.Tensor,}
\NormalTok{    means: torch.Tensor,}
\NormalTok{    std: torch.Tensor,}
\NormalTok{) }\OperatorTok{{-}\textgreater{}}\NormalTok{ torch.Tensor:}
\NormalTok{    variance }\OperatorTok{=}\NormalTok{ std.square()}
    \ControlFlowTok{return}\NormalTok{ (}
        \OperatorTok{{-}}\FloatTok{0.5}
        \OperatorTok{*}\NormalTok{ (}
\NormalTok{            (actions }\OperatorTok{{-}}\NormalTok{ means).square() }\OperatorTok{/}\NormalTok{ variance}
            \OperatorTok{+} \FloatTok{2.0} \OperatorTok{*}\NormalTok{ torch.log(std)}
            \OperatorTok{+}\NormalTok{ math.log(}\FloatTok{2.0} \OperatorTok{*}\NormalTok{ math.pi)}
\NormalTok{        )}
\NormalTok{    ).}\BuiltInTok{sum}\NormalTok{(dim}\OperatorTok{={-}}\DecValTok{1}\NormalTok{)}


\KeywordTok{def}\NormalTok{ gaussian\_entropy(std: torch.Tensor) }\OperatorTok{{-}\textgreater{}}\NormalTok{ torch.Tensor:}
    \ControlFlowTok{return}\NormalTok{ (torch.log(std) }\OperatorTok{+} \FloatTok{0.5} \OperatorTok{*}\NormalTok{ math.log(}\FloatTok{2.0} \OperatorTok{*}\NormalTok{ math.pi }\OperatorTok{*}\NormalTok{ math.e)).}\BuiltInTok{sum}\NormalTok{()}


\KeywordTok{def}\NormalTok{ compute\_gae(}
\NormalTok{    rewards: torch.Tensor,}
\NormalTok{    values: torch.Tensor,}
\NormalTok{    bootstrap\_values: torch.Tensor,}
\NormalTok{    terminated: torch.Tensor,}
\NormalTok{    truncated: torch.Tensor,}
    \OperatorTok{*}\NormalTok{,}
\NormalTok{    gamma: }\BuiltInTok{float} \OperatorTok{=} \FloatTok{0.99}\NormalTok{,}
\NormalTok{    gae\_lambda: }\BuiltInTok{float} \OperatorTok{=} \FloatTok{0.95}\NormalTok{,}
\NormalTok{) }\OperatorTok{{-}\textgreater{}} \BuiltInTok{tuple}\NormalTok{[torch.Tensor, torch.Tensor]:}
    \ControlFlowTok{if} \KeywordTok{not}\NormalTok{ (}
\NormalTok{        rewards.shape}
        \OperatorTok{==}\NormalTok{ values.shape}
        \OperatorTok{==}\NormalTok{ bootstrap\_values.shape}
        \OperatorTok{==}\NormalTok{ terminated.shape}
        \OperatorTok{==}\NormalTok{ truncated.shape}
\NormalTok{    ):}
        \ControlFlowTok{raise} \PreprocessorTok{ValueError}\NormalTok{(}\StringTok{"GAE tensors must share [time, env] shape"}\NormalTok{)}
\NormalTok{    advantages }\OperatorTok{=}\NormalTok{ torch.zeros\_like(rewards)}
\NormalTok{    trace }\OperatorTok{=}\NormalTok{ torch.zeros\_like(rewards[}\DecValTok{0}\NormalTok{])}
    \ControlFlowTok{for}\NormalTok{ step }\KeywordTok{in} \BuiltInTok{range}\NormalTok{(rewards.shape[}\DecValTok{0}\NormalTok{] }\OperatorTok{{-}} \DecValTok{1}\NormalTok{, }\OperatorTok{{-}}\DecValTok{1}\NormalTok{, }\OperatorTok{{-}}\DecValTok{1}\NormalTok{):}
\NormalTok{        bootstrap\_mask }\OperatorTok{=}\NormalTok{ (}\OperatorTok{\textasciitilde{}}\NormalTok{terminated[step]).to(rewards.dtype)}
\NormalTok{        continuation\_mask }\OperatorTok{=}\NormalTok{ (}\OperatorTok{\textasciitilde{}}\NormalTok{(terminated[step] }\OperatorTok{|}\NormalTok{ truncated[step])).to(}
\NormalTok{            rewards.dtype}
\NormalTok{        )}
\NormalTok{        delta }\OperatorTok{=}\NormalTok{ (}
\NormalTok{            rewards[step]}
            \OperatorTok{+}\NormalTok{ gamma }\OperatorTok{*}\NormalTok{ bootstrap\_values[step] }\OperatorTok{*}\NormalTok{ bootstrap\_mask}
            \OperatorTok{{-}}\NormalTok{ values[step]}
\NormalTok{        )}
\NormalTok{        trace }\OperatorTok{=}\NormalTok{ delta }\OperatorTok{+}\NormalTok{ gamma }\OperatorTok{*}\NormalTok{ gae\_lambda }\OperatorTok{*}\NormalTok{ continuation\_mask }\OperatorTok{*}\NormalTok{ trace}
\NormalTok{        advantages[step] }\OperatorTok{=}\NormalTok{ trace}
    \ControlFlowTok{return}\NormalTok{ advantages, advantages }\OperatorTok{+}\NormalTok{ values}


\KeywordTok{def}\NormalTok{ \_validate\_observations(observations: Any, num\_envs: }\BuiltInTok{int}\NormalTok{) }\OperatorTok{{-}\textgreater{}}\NormalTok{ np.ndarray:}
\NormalTok{    array }\OperatorTok{=}\NormalTok{ np.asarray(observations, dtype}\OperatorTok{=}\NormalTok{np.float32)}
\NormalTok{    expected }\OperatorTok{=}\NormalTok{ (num\_envs, OBSERVATION\_DIM)}
    \ControlFlowTok{if}\NormalTok{ array.shape }\OperatorTok{!=}\NormalTok{ expected:}
        \ControlFlowTok{raise} \PreprocessorTok{ValueError}\NormalTok{(}\SpecialStringTok{f"observations must have shape }\SpecialCharTok{\{}\NormalTok{expected}\SpecialCharTok{\}}\SpecialStringTok{, got }\SpecialCharTok{\{}\NormalTok{array}\SpecialCharTok{.}\NormalTok{shape}\SpecialCharTok{\}}\SpecialStringTok{"}\NormalTok{)}
    \ControlFlowTok{if} \KeywordTok{not}\NormalTok{ np.isfinite(array).}\BuiltInTok{all}\NormalTok{():}
        \ControlFlowTok{raise} \PreprocessorTok{RuntimeError}\NormalTok{(}\StringTok{"basketball observations contain non{-}finite values"}\NormalTok{)}
    \ControlFlowTok{if}\NormalTok{ np.count\_nonzero(array[:, }\DecValTok{51}\NormalTok{:}\DecValTok{61}\NormalTok{]):}
        \ControlFlowTok{raise} \PreprocessorTok{RuntimeError}\NormalTok{(}\StringTok{"basketball actor head/body padding must remain zero"}\NormalTok{)}
    \ControlFlowTok{return}\NormalTok{ array}


\KeywordTok{def}\NormalTok{ \_validate\_step(}
\NormalTok{    observation: Any,}
\NormalTok{    reward: Any,}
\NormalTok{    terminated: Any,}
\NormalTok{    truncated: Any,}
\NormalTok{    info: Any,}
\NormalTok{) }\OperatorTok{{-}\textgreater{}} \BuiltInTok{tuple}\NormalTok{[np.ndarray, }\BuiltInTok{float}\NormalTok{, }\BuiltInTok{bool}\NormalTok{, }\BuiltInTok{bool}\NormalTok{, }\BuiltInTok{dict}\NormalTok{[}\BuiltInTok{str}\NormalTok{, Any]]:}
\NormalTok{    obs }\OperatorTok{=}\NormalTok{ np.asarray(observation, dtype}\OperatorTok{=}\NormalTok{np.float32)}
    \ControlFlowTok{if}\NormalTok{ obs.shape }\OperatorTok{!=}\NormalTok{ (OBSERVATION\_DIM,) }\KeywordTok{or} \KeywordTok{not}\NormalTok{ np.isfinite(obs).}\BuiltInTok{all}\NormalTok{():}
        \ControlFlowTok{raise} \PreprocessorTok{RuntimeError}\NormalTok{(}\StringTok{"basketball step returned an invalid observation"}\NormalTok{)}
    \ControlFlowTok{if}\NormalTok{ np.count\_nonzero(obs[}\DecValTok{51}\NormalTok{:}\DecValTok{61}\NormalTok{]):}
        \ControlFlowTok{raise} \PreprocessorTok{RuntimeError}\NormalTok{(}\StringTok{"basketball actor received nonzero head/body padding"}\NormalTok{)}
\NormalTok{    value }\OperatorTok{=} \BuiltInTok{float}\NormalTok{(reward)}
    \ControlFlowTok{if} \KeywordTok{not}\NormalTok{ math.isfinite(value):}
        \ControlFlowTok{raise} \PreprocessorTok{RuntimeError}\NormalTok{(}\StringTok{"basketball step returned a non{-}finite reward"}\NormalTok{)}
    \ControlFlowTok{if} \KeywordTok{not} \BuiltInTok{isinstance}\NormalTok{(info, }\BuiltInTok{dict}\NormalTok{):}
        \ControlFlowTok{raise} \PreprocessorTok{TypeError}\NormalTok{(}\StringTok{"basketball info must be a dictionary"}\NormalTok{)}
    \ControlFlowTok{return}\NormalTok{ obs, value, }\BuiltInTok{bool}\NormalTok{(terminated), }\BuiltInTok{bool}\NormalTok{(truncated), info}


\KeywordTok{def}\NormalTok{ make\_environments(args: argparse.Namespace) }\OperatorTok{{-}\textgreater{}} \BuiltInTok{list}\NormalTok{[Any]:}
    \ImportTok{from}\NormalTok{ rlx.environments.basketball }\ImportTok{import}\NormalTok{ BasketballEnv, basketball\_model}

\NormalTok{    shared\_model }\OperatorTok{=}\NormalTok{ (}
\NormalTok{        basketball\_model(actuator}\OperatorTok{=}\NormalTok{args.actuator) }\ControlFlowTok{if}\NormalTok{ args.actuator }\OperatorTok{==} \StringTok{"xml"} \ControlFlowTok{else} \VariableTok{None}
\NormalTok{    )}
\NormalTok{    command }\OperatorTok{=} \VariableTok{None} \ControlFlowTok{if}\NormalTok{ args.command }\KeywordTok{is} \VariableTok{None} \ControlFlowTok{else} \BuiltInTok{tuple}\NormalTok{(}\BuiltInTok{float}\NormalTok{(x) }\ControlFlowTok{for}\NormalTok{ x }\KeywordTok{in}\NormalTok{ args.command)}
\NormalTok{    environments }\OperatorTok{=}\NormalTok{ []}
    \ControlFlowTok{for}\NormalTok{ index }\KeywordTok{in} \BuiltInTok{range}\NormalTok{(args.num\_envs):}
\NormalTok{        kwargs }\OperatorTok{=}\NormalTok{ \{}
            \StringTok{"max\_episode\_s"}\NormalTok{: }\FloatTok{10.0}\NormalTok{,}
            \StringTok{"hold"}\NormalTok{: args.hold,}
            \StringTok{"curriculum"}\NormalTok{: args.curriculum,}
            \StringTok{"command"}\NormalTok{: command,}
            \StringTok{"seed"}\NormalTok{: args.seed }\OperatorTok{+}\NormalTok{ index,}
            \StringTok{"actuator"}\NormalTok{: args.actuator,}
            \StringTok{"obs\_noise"}\NormalTok{: }\VariableTok{False}\NormalTok{,}
            \StringTok{"domain\_rand"}\NormalTok{: }\VariableTok{False}\NormalTok{,}
            \StringTok{"pushes"}\NormalTok{: args.pushes,}
\NormalTok{        \}}
        \ControlFlowTok{if}\NormalTok{ shared\_model }\KeywordTok{is} \KeywordTok{not} \VariableTok{None}\NormalTok{:}
\NormalTok{            kwargs[}\StringTok{"model"}\NormalTok{] }\OperatorTok{=}\NormalTok{ shared\_model}
\NormalTok{        environments.append(}
\NormalTok{            BasketballEnv(}\OperatorTok{**}\NormalTok{kwargs)}
\NormalTok{        )}
    \ControlFlowTok{return}\NormalTok{ environments}


\KeywordTok{def}\NormalTok{ reset\_environments(}
\NormalTok{    environments: Sequence[Any],}
    \OperatorTok{*}\NormalTok{,}
\NormalTok{    seed: }\BuiltInTok{int}\NormalTok{,}
\NormalTok{) }\OperatorTok{{-}\textgreater{}} \BuiltInTok{tuple}\NormalTok{[np.ndarray, }\BuiltInTok{list}\NormalTok{[}\BuiltInTok{dict}\NormalTok{[}\BuiltInTok{str}\NormalTok{, Any]]]:}
\NormalTok{    observations: }\BuiltInTok{list}\NormalTok{[np.ndarray] }\OperatorTok{=}\NormalTok{ []}
\NormalTok{    infos: }\BuiltInTok{list}\NormalTok{[}\BuiltInTok{dict}\NormalTok{[}\BuiltInTok{str}\NormalTok{, Any]] }\OperatorTok{=}\NormalTok{ []}
    \ControlFlowTok{for}\NormalTok{ index, environment }\KeywordTok{in} \BuiltInTok{enumerate}\NormalTok{(environments):}
\NormalTok{        observation, info }\OperatorTok{=}\NormalTok{ environment.reset(seed}\OperatorTok{=}\NormalTok{seed }\OperatorTok{+}\NormalTok{ index)}
\NormalTok{        observations.append(np.asarray(observation, dtype}\OperatorTok{=}\NormalTok{np.float32))}
        \ControlFlowTok{if} \KeywordTok{not} \BuiltInTok{isinstance}\NormalTok{(info, }\BuiltInTok{dict}\NormalTok{):}
            \ControlFlowTok{raise} \PreprocessorTok{TypeError}\NormalTok{(}\StringTok{"basketball reset info must be a dictionary"}\NormalTok{)}
\NormalTok{        infos.append(info)}
    \ControlFlowTok{return}\NormalTok{ \_validate\_observations(observations, }\BuiltInTok{len}\NormalTok{(environments)), infos}


\KeywordTok{def}\NormalTok{ collect\_rollout(}
\NormalTok{    actor: BasketballActor,}
\NormalTok{    critic: BasketballCritic,}
\NormalTok{    environments: Sequence[Any],}
\NormalTok{    observations: np.ndarray,}
\NormalTok{    hidden: torch.Tensor,}
\NormalTok{    cell: torch.Tensor,}
\NormalTok{    episode\_starts: torch.Tensor,}
    \OperatorTok{*}\NormalTok{,}
\NormalTok{    steps: }\BuiltInTok{int}\NormalTok{,}
\NormalTok{    generator: torch.Generator,}
\NormalTok{    seed: }\BuiltInTok{int}\NormalTok{,}
\NormalTok{    reset\_counts: np.ndarray,}
\NormalTok{) }\OperatorTok{{-}\textgreater{}} \BuiltInTok{tuple}\NormalTok{[Rollout, np.ndarray, torch.Tensor, torch.Tensor, torch.Tensor]:}
    \ControlFlowTok{if}\NormalTok{ steps }\OperatorTok{\textless{}} \DecValTok{1}\NormalTok{:}
        \ControlFlowTok{raise} \PreprocessorTok{ValueError}\NormalTok{(}\StringTok{"steps must be positive"}\NormalTok{)}
\NormalTok{    num\_envs }\OperatorTok{=} \BuiltInTok{len}\NormalTok{(environments)}
\NormalTok{    initial\_hidden }\OperatorTok{=}\NormalTok{ hidden.detach().clone()}
\NormalTok{    initial\_cell }\OperatorTok{=}\NormalTok{ cell.detach().clone()}
\NormalTok{    stored: }\BuiltInTok{dict}\NormalTok{[}\BuiltInTok{str}\NormalTok{, }\BuiltInTok{list}\NormalTok{[torch.Tensor]] }\OperatorTok{=}\NormalTok{ \{}
\NormalTok{        name: []}
        \ControlFlowTok{for}\NormalTok{ name }\KeywordTok{in}\NormalTok{ (}
            \StringTok{"observations"}\NormalTok{,}
            \StringTok{"actions"}\NormalTok{,}
            \StringTok{"old\_log\_probs"}\NormalTok{,}
            \StringTok{"old\_values"}\NormalTok{,}
            \StringTok{"rewards"}\NormalTok{,}
            \StringTok{"terminated"}\NormalTok{,}
            \StringTok{"truncated"}\NormalTok{,}
            \StringTok{"bootstrap\_values"}\NormalTok{,}
            \StringTok{"episode\_starts"}\NormalTok{,}
\NormalTok{        )}
\NormalTok{    \}}
\NormalTok{    all\_infos: }\BuiltInTok{list}\NormalTok{[}\BuiltInTok{list}\NormalTok{[}\BuiltInTok{dict}\NormalTok{[}\BuiltInTok{str}\NormalTok{, Any]]] }\OperatorTok{=}\NormalTok{ []}
\NormalTok{    actor.}\BuiltInTok{eval}\NormalTok{()}
\NormalTok{    critic.}\BuiltInTok{eval}\NormalTok{()}
    \ControlFlowTok{for}\NormalTok{ \_ }\KeywordTok{in} \BuiltInTok{range}\NormalTok{(steps):}
\NormalTok{        observation\_tensor }\OperatorTok{=}\NormalTok{ torch.from\_numpy(observations)}
        \ControlFlowTok{with}\NormalTok{ torch.inference\_mode():}
\NormalTok{            carry }\OperatorTok{=}\NormalTok{ (}\OperatorTok{\textasciitilde{}}\NormalTok{episode\_starts.}\BuiltInTok{bool}\NormalTok{()).to(observation\_tensor.dtype).view(}
                \DecValTok{1}\NormalTok{, num\_envs, }\DecValTok{1}
\NormalTok{            )}
\NormalTok{            hidden }\OperatorTok{=}\NormalTok{ hidden }\OperatorTok{*}\NormalTok{ carry}
\NormalTok{            cell }\OperatorTok{=}\NormalTok{ cell }\OperatorTok{*}\NormalTok{ carry}
\NormalTok{            means, next\_hidden, next\_cell }\OperatorTok{=}\NormalTok{ actor(}
\NormalTok{                observation\_tensor,}
\NormalTok{                hidden,}
\NormalTok{                cell,}
\NormalTok{            )}
\NormalTok{            std }\OperatorTok{=}\NormalTok{ actor.action\_std()}
\NormalTok{            actions }\OperatorTok{=}\NormalTok{ means }\OperatorTok{+}\NormalTok{ torch.randn(}
\NormalTok{                means.shape,}
\NormalTok{                generator}\OperatorTok{=}\NormalTok{generator,}
\NormalTok{                dtype}\OperatorTok{=}\NormalTok{means.dtype,}
\NormalTok{            ) }\OperatorTok{*}\NormalTok{ std}
\NormalTok{            log\_probs }\OperatorTok{=}\NormalTok{ gaussian\_log\_prob(actions, means, std)}
\NormalTok{            values }\OperatorTok{=}\NormalTok{ critic(observation\_tensor)}
\NormalTok{        next\_observations: }\BuiltInTok{list}\NormalTok{[np.ndarray] }\OperatorTok{=}\NormalTok{ []}
\NormalTok{        rewards: }\BuiltInTok{list}\NormalTok{[}\BuiltInTok{float}\NormalTok{] }\OperatorTok{=}\NormalTok{ []}
\NormalTok{        terminations: }\BuiltInTok{list}\NormalTok{[}\BuiltInTok{bool}\NormalTok{] }\OperatorTok{=}\NormalTok{ []}
\NormalTok{        truncations: }\BuiltInTok{list}\NormalTok{[}\BuiltInTok{bool}\NormalTok{] }\OperatorTok{=}\NormalTok{ []}
\NormalTok{        step\_infos: }\BuiltInTok{list}\NormalTok{[}\BuiltInTok{dict}\NormalTok{[}\BuiltInTok{str}\NormalTok{, Any]] }\OperatorTok{=}\NormalTok{ []}
\NormalTok{        final\_observations: }\BuiltInTok{list}\NormalTok{[np.ndarray] }\OperatorTok{=}\NormalTok{ []}
        \ControlFlowTok{for}\NormalTok{ index, environment }\KeywordTok{in} \BuiltInTok{enumerate}\NormalTok{(environments):}
\NormalTok{            result }\OperatorTok{=}\NormalTok{ \_validate\_step(}\OperatorTok{*}\NormalTok{environment.step(actions[index].numpy()))}
\NormalTok{            next\_observation, reward, terminated, truncated, info }\OperatorTok{=}\NormalTok{ result}
\NormalTok{            final\_observations.append(next\_observation)}
\NormalTok{            rewards.append(reward)}
\NormalTok{            terminations.append(terminated)}
\NormalTok{            truncations.append(truncated)}
\NormalTok{            step\_infos.append(info)}
            \ControlFlowTok{if}\NormalTok{ terminated }\KeywordTok{or}\NormalTok{ truncated:}
\NormalTok{                reset\_counts[index] }\OperatorTok{+=} \DecValTok{1}
\NormalTok{                reset\_observation, reset\_info }\OperatorTok{=}\NormalTok{ environment.reset(}
\NormalTok{                    seed}\OperatorTok{=}\NormalTok{seed }\OperatorTok{+}\NormalTok{ index }\OperatorTok{+} \BuiltInTok{int}\NormalTok{(reset\_counts[index]) }\OperatorTok{*}\NormalTok{ num\_envs}
\NormalTok{                )}
\NormalTok{                reset\_observation }\OperatorTok{=}\NormalTok{ np.asarray(reset\_observation, dtype}\OperatorTok{=}\NormalTok{np.float32)}
                \ControlFlowTok{if} \KeywordTok{not} \BuiltInTok{isinstance}\NormalTok{(reset\_info, }\BuiltInTok{dict}\NormalTok{):}
                    \ControlFlowTok{raise} \PreprocessorTok{TypeError}\NormalTok{(}\StringTok{"basketball reset info must be a dictionary"}\NormalTok{)}
\NormalTok{                info }\OperatorTok{=}\NormalTok{ \{}\OperatorTok{**}\NormalTok{info, }\StringTok{"reset\_info"}\NormalTok{: reset\_info\}}
\NormalTok{                step\_infos[}\OperatorTok{{-}}\DecValTok{1}\NormalTok{] }\OperatorTok{=}\NormalTok{ info}
\NormalTok{                next\_observations.append(reset\_observation)}
            \ControlFlowTok{else}\NormalTok{:}
\NormalTok{                next\_observations.append(next\_observation)}
\NormalTok{        next\_batch }\OperatorTok{=}\NormalTok{ \_validate\_observations(next\_observations, num\_envs)}
\NormalTok{        final\_batch }\OperatorTok{=}\NormalTok{ \_validate\_observations(final\_observations, num\_envs)}
\NormalTok{        terminated\_tensor }\OperatorTok{=}\NormalTok{ torch.tensor(terminations, dtype}\OperatorTok{=}\NormalTok{torch.}\BuiltInTok{bool}\NormalTok{)}
\NormalTok{        truncated\_tensor }\OperatorTok{=}\NormalTok{ torch.tensor(truncations, dtype}\OperatorTok{=}\NormalTok{torch.}\BuiltInTok{bool}\NormalTok{)}
        \ControlFlowTok{with}\NormalTok{ torch.inference\_mode():}
\NormalTok{            bootstrap\_values }\OperatorTok{=}\NormalTok{ critic(torch.from\_numpy(final\_batch))}
\NormalTok{            bootstrap\_values }\OperatorTok{=}\NormalTok{ torch.where(}
\NormalTok{                terminated\_tensor,}
\NormalTok{                torch.zeros\_like(bootstrap\_values),}
\NormalTok{                bootstrap\_values,}
\NormalTok{            )}
\NormalTok{        stored[}\StringTok{"observations"}\NormalTok{].append(observation\_tensor)}
\NormalTok{        stored[}\StringTok{"actions"}\NormalTok{].append(actions)}
\NormalTok{        stored[}\StringTok{"old\_log\_probs"}\NormalTok{].append(log\_probs)}
\NormalTok{        stored[}\StringTok{"old\_values"}\NormalTok{].append(values)}
\NormalTok{        stored[}\StringTok{"rewards"}\NormalTok{].append(torch.tensor(rewards, dtype}\OperatorTok{=}\NormalTok{torch.float32))}
\NormalTok{        stored[}\StringTok{"terminated"}\NormalTok{].append(terminated\_tensor)}
\NormalTok{        stored[}\StringTok{"truncated"}\NormalTok{].append(truncated\_tensor)}
\NormalTok{        stored[}\StringTok{"bootstrap\_values"}\NormalTok{].append(bootstrap\_values)}
\NormalTok{        stored[}\StringTok{"episode\_starts"}\NormalTok{].append(episode\_starts.clone())}
\NormalTok{        all\_infos.append(step\_infos)}
\NormalTok{        observations }\OperatorTok{=}\NormalTok{ next\_batch}
\NormalTok{        hidden }\OperatorTok{=}\NormalTok{ next\_hidden}
\NormalTok{        cell }\OperatorTok{=}\NormalTok{ next\_cell}
\NormalTok{        episode\_starts }\OperatorTok{=}\NormalTok{ terminated\_tensor }\OperatorTok{|}\NormalTok{ truncated\_tensor}
\NormalTok{    rollout }\OperatorTok{=}\NormalTok{ Rollout(}
        \OperatorTok{**}\NormalTok{\{name: torch.stack(values) }\ControlFlowTok{for}\NormalTok{ name, values }\KeywordTok{in}\NormalTok{ stored.items()\},}
\NormalTok{        initial\_hidden}\OperatorTok{=}\NormalTok{initial\_hidden,}
\NormalTok{        initial\_cell}\OperatorTok{=}\NormalTok{initial\_cell,}
\NormalTok{        next\_hidden}\OperatorTok{=}\NormalTok{hidden.detach(),}
\NormalTok{        next\_cell}\OperatorTok{=}\NormalTok{cell.detach(),}
\NormalTok{        next\_episode\_starts}\OperatorTok{=}\NormalTok{episode\_starts,}
\NormalTok{        infos}\OperatorTok{=}\NormalTok{all\_infos,}
\NormalTok{    )}
    \ControlFlowTok{return}\NormalTok{ rollout, observations, hidden.detach(), cell.detach(), episode\_starts}


\KeywordTok{def}\NormalTok{ ppo\_update(}
\NormalTok{    actor: BasketballActor,}
\NormalTok{    critic: BasketballCritic,}
\NormalTok{    optimizer: torch.optim.Optimizer,}
\NormalTok{    rollout: Rollout,}
    \OperatorTok{*}\NormalTok{,}
\NormalTok{    epochs: }\BuiltInTok{int} \OperatorTok{=} \DecValTok{5}\NormalTok{,}
\NormalTok{    clip\_coefficient: }\BuiltInTok{float} \OperatorTok{=} \FloatTok{0.2}\NormalTok{,}
\NormalTok{    value\_coefficient: }\BuiltInTok{float} \OperatorTok{=} \FloatTok{1.0}\NormalTok{,}
\NormalTok{    entropy\_coefficient: }\BuiltInTok{float} \OperatorTok{=} \FloatTok{0.01}\NormalTok{,}
\NormalTok{    max\_grad\_norm: }\BuiltInTok{float} \OperatorTok{=} \FloatTok{1.0}\NormalTok{,}
\NormalTok{    target\_kl: }\BuiltInTok{float} \OperatorTok{=} \FloatTok{0.02}\NormalTok{,}
\NormalTok{) }\OperatorTok{{-}\textgreater{}} \BuiltInTok{dict}\NormalTok{[}\BuiltInTok{str}\NormalTok{, }\BuiltInTok{float}\NormalTok{]:}
\NormalTok{    advantages, returns }\OperatorTok{=}\NormalTok{ compute\_gae(}
\NormalTok{        rollout.rewards,}
\NormalTok{        rollout.old\_values,}
\NormalTok{        rollout.bootstrap\_values,}
\NormalTok{        rollout.terminated,}
\NormalTok{        rollout.truncated,}
\NormalTok{    )}
\NormalTok{    advantages }\OperatorTok{=}\NormalTok{ (advantages }\OperatorTok{{-}}\NormalTok{ advantages.mean()) }\OperatorTok{/}\NormalTok{ (}
\NormalTok{        advantages.std(unbiased}\OperatorTok{=}\VariableTok{False}\NormalTok{) }\OperatorTok{+} \FloatTok{1e{-}8}
\NormalTok{    )}
\NormalTok{    actor.train()}
\NormalTok{    critic.train()}
\NormalTok{    metrics: }\BuiltInTok{dict}\NormalTok{[}\BuiltInTok{str}\NormalTok{, }\BuiltInTok{float}\NormalTok{] }\OperatorTok{=}\NormalTok{ \{\}}
    \ControlFlowTok{for}\NormalTok{ epoch }\KeywordTok{in} \BuiltInTok{range}\NormalTok{(epochs):}
\NormalTok{        means, \_, \_ }\OperatorTok{=}\NormalTok{ actor.forward\_sequence(}
\NormalTok{            rollout.observations,}
\NormalTok{            rollout.initial\_hidden,}
\NormalTok{            rollout.initial\_cell,}
\NormalTok{            rollout.episode\_starts,}
\NormalTok{        )}
\NormalTok{        std }\OperatorTok{=}\NormalTok{ actor.action\_std()}
\NormalTok{        log\_probs }\OperatorTok{=}\NormalTok{ gaussian\_log\_prob(rollout.actions, means, std)}
\NormalTok{        log\_ratio }\OperatorTok{=}\NormalTok{ log\_probs }\OperatorTok{{-}}\NormalTok{ rollout.old\_log\_probs}
\NormalTok{        ratio }\OperatorTok{=}\NormalTok{ torch.exp(log\_ratio)}
\NormalTok{        approximate\_kl }\OperatorTok{=}\NormalTok{ ((ratio }\OperatorTok{{-}} \DecValTok{1}\NormalTok{) }\OperatorTok{{-}}\NormalTok{ log\_ratio).mean()}
        \ControlFlowTok{if} \KeywordTok{not}\NormalTok{ torch.isfinite(approximate\_kl):}
            \ControlFlowTok{raise} \PreprocessorTok{RuntimeError}\NormalTok{(}\StringTok{"PPO KL became non{-}finite"}\NormalTok{)}
        \ControlFlowTok{if} \BuiltInTok{float}\NormalTok{(approximate\_kl.detach()) }\OperatorTok{\textgreater{}}\NormalTok{ target\_kl:}
\NormalTok{            metrics[}\StringTok{"approx\_kl"}\NormalTok{] }\OperatorTok{=} \BuiltInTok{float}\NormalTok{(approximate\_kl.detach())}
\NormalTok{            metrics[}\StringTok{"epochs\_completed"}\NormalTok{] }\OperatorTok{=}\NormalTok{ epoch}
\NormalTok{            metrics[}\StringTok{"kl\_early\_stop"}\NormalTok{] }\OperatorTok{=} \FloatTok{1.0}
            \ControlFlowTok{break}
\NormalTok{        unclipped }\OperatorTok{=}\NormalTok{ ratio }\OperatorTok{*}\NormalTok{ advantages}
\NormalTok{        clipped }\OperatorTok{=}\NormalTok{ torch.clamp(}
\NormalTok{            ratio,}
            \FloatTok{1.0} \OperatorTok{{-}}\NormalTok{ clip\_coefficient,}
            \FloatTok{1.0} \OperatorTok{+}\NormalTok{ clip\_coefficient,}
\NormalTok{        ) }\OperatorTok{*}\NormalTok{ advantages}
\NormalTok{        policy\_loss }\OperatorTok{=} \OperatorTok{{-}}\NormalTok{torch.minimum(unclipped, clipped).mean()}
\NormalTok{        values }\OperatorTok{=}\NormalTok{ critic(rollout.observations.reshape(}\OperatorTok{{-}}\DecValTok{1}\NormalTok{, OBSERVATION\_DIM)).reshape(}
\NormalTok{            rollout.rewards.shape}
\NormalTok{        )}
\NormalTok{        value\_loss }\OperatorTok{=} \FloatTok{0.5} \OperatorTok{*}\NormalTok{ (values }\OperatorTok{{-}}\NormalTok{ returns).square().mean()}
\NormalTok{        entropy }\OperatorTok{=}\NormalTok{ gaussian\_entropy(std)}
\NormalTok{        loss }\OperatorTok{=}\NormalTok{ (}
\NormalTok{            policy\_loss}
            \OperatorTok{+}\NormalTok{ value\_coefficient }\OperatorTok{*}\NormalTok{ value\_loss}
            \OperatorTok{{-}}\NormalTok{ entropy\_coefficient }\OperatorTok{*}\NormalTok{ entropy}
\NormalTok{        )}
        \ControlFlowTok{if} \KeywordTok{not}\NormalTok{ torch.isfinite(loss):}
            \ControlFlowTok{raise} \PreprocessorTok{RuntimeError}\NormalTok{(}\StringTok{"PPO loss became non{-}finite"}\NormalTok{)}
\NormalTok{        optimizer.zero\_grad(set\_to\_none}\OperatorTok{=}\VariableTok{True}\NormalTok{)}
\NormalTok{        loss.backward()}
\NormalTok{        gradient\_norm }\OperatorTok{=}\NormalTok{ nn.utils.clip\_grad\_norm\_(}
            \BuiltInTok{list}\NormalTok{(actor.parameters()) }\OperatorTok{+} \BuiltInTok{list}\NormalTok{(critic.parameters()),}
\NormalTok{            max\_grad\_norm,}
\NormalTok{        )}
        \ControlFlowTok{if} \KeywordTok{not}\NormalTok{ torch.isfinite(torch.as\_tensor(gradient\_norm)):}
            \ControlFlowTok{raise} \PreprocessorTok{RuntimeError}\NormalTok{(}\StringTok{"PPO gradient norm became non{-}finite"}\NormalTok{)}
\NormalTok{        optimizer.step()}
        \ControlFlowTok{if} \KeywordTok{not} \BuiltInTok{all}\NormalTok{(}
\NormalTok{            torch.isfinite(parameter).}\BuiltInTok{all}\NormalTok{()}
            \ControlFlowTok{for}\NormalTok{ parameter }\KeywordTok{in} \BuiltInTok{list}\NormalTok{(actor.parameters()) }\OperatorTok{+} \BuiltInTok{list}\NormalTok{(critic.parameters())}
\NormalTok{        ):}
            \ControlFlowTok{raise} \PreprocessorTok{RuntimeError}\NormalTok{(}\StringTok{"PPO parameters became non{-}finite"}\NormalTok{)}
        \ControlFlowTok{with}\NormalTok{ torch.no\_grad():}
\NormalTok{            updated\_means, \_, \_ }\OperatorTok{=}\NormalTok{ actor.forward\_sequence(}
\NormalTok{                rollout.observations, rollout.initial\_hidden, rollout.initial\_cell,}
\NormalTok{                rollout.episode\_starts,}
\NormalTok{            )}
\NormalTok{            updated\_log\_probs }\OperatorTok{=}\NormalTok{ gaussian\_log\_prob(rollout.actions, updated\_means, actor.action\_std())}
\NormalTok{            updated\_log\_ratio }\OperatorTok{=}\NormalTok{ updated\_log\_probs }\OperatorTok{{-}}\NormalTok{ rollout.old\_log\_probs}
\NormalTok{            final\_kl }\OperatorTok{=}\NormalTok{ (torch.exp(updated\_log\_ratio) }\OperatorTok{{-}} \DecValTok{1} \OperatorTok{{-}}\NormalTok{ updated\_log\_ratio).mean()}
        \ControlFlowTok{if} \KeywordTok{not}\NormalTok{ torch.isfinite(final\_kl):}
            \ControlFlowTok{raise} \PreprocessorTok{RuntimeError}\NormalTok{(}\StringTok{"post{-}update PPO KL became non{-}finite"}\NormalTok{)}
\NormalTok{        metrics }\OperatorTok{=}\NormalTok{ \{}
            \StringTok{"loss"}\NormalTok{: }\BuiltInTok{float}\NormalTok{(loss.detach()),}
            \StringTok{"policy\_loss"}\NormalTok{: }\BuiltInTok{float}\NormalTok{(policy\_loss.detach()),}
            \StringTok{"value\_loss"}\NormalTok{: }\BuiltInTok{float}\NormalTok{(value\_loss.detach()),}
            \StringTok{"entropy"}\NormalTok{: }\BuiltInTok{float}\NormalTok{(entropy.detach()),}
            \StringTok{"gradient\_norm"}\NormalTok{: }\BuiltInTok{float}\NormalTok{(gradient\_norm),}
            \StringTok{"approx\_kl"}\NormalTok{: }\BuiltInTok{float}\NormalTok{(final\_kl),}
            \StringTok{"epochs\_completed"}\NormalTok{: epoch }\OperatorTok{+} \DecValTok{1}\NormalTok{,}
            \StringTok{"kl\_early\_stop"}\NormalTok{: }\FloatTok{0.0}\NormalTok{,}
\NormalTok{        \}}
        \ControlFlowTok{if} \BuiltInTok{float}\NormalTok{(final\_kl) }\OperatorTok{\textgreater{}}\NormalTok{ target\_kl:}
\NormalTok{            metrics[}\StringTok{"kl\_early\_stop"}\NormalTok{] }\OperatorTok{=} \FloatTok{1.0}
            \ControlFlowTok{break}
    \ControlFlowTok{return}\NormalTok{ metrics}


\KeywordTok{def}\NormalTok{ actor\_weight\_change(}
\NormalTok{    actor: BasketballActor,}
\NormalTok{    initial\_state: }\BuiltInTok{dict}\NormalTok{[}\BuiltInTok{str}\NormalTok{, torch.Tensor],}
\NormalTok{) }\OperatorTok{{-}\textgreater{}} \BuiltInTok{float}\NormalTok{:}
\NormalTok{    total }\OperatorTok{=}\NormalTok{ torch.tensor(}\FloatTok{0.0}\NormalTok{)}
\NormalTok{    current }\OperatorTok{=}\NormalTok{ actor.state\_dict()}
    \ControlFlowTok{for}\NormalTok{ name, source }\KeywordTok{in}\NormalTok{ initial\_state.items():}
        \ControlFlowTok{if}\NormalTok{ name.startswith(}\StringTok{"obs\_normalizer."}\NormalTok{) }\KeywordTok{or}\NormalTok{ name }\OperatorTok{==} \StringTok{"distribution.std\_param"}\NormalTok{:}
            \ControlFlowTok{continue}
\NormalTok{        total }\OperatorTok{+=}\NormalTok{ (current[name].cpu() }\OperatorTok{{-}}\NormalTok{ source.cpu()).square().}\BuiltInTok{sum}\NormalTok{()}
    \ControlFlowTok{return} \BuiltInTok{float}\NormalTok{(torch.sqrt(total))}


\KeywordTok{def}\NormalTok{ save\_results(}
\NormalTok{    output: Path,}
\NormalTok{    actor: BasketballActor,}
\NormalTok{    critic: BasketballCritic,}
\NormalTok{    optimizer: torch.optim.Optimizer,}
    \OperatorTok{*}\NormalTok{,}
\NormalTok{    args: argparse.Namespace,}
\NormalTok{    source\_metadata: }\BuiltInTok{dict}\NormalTok{[}\BuiltInTok{str}\NormalTok{, Any],}
\NormalTok{    initial\_actor\_state: }\BuiltInTok{dict}\NormalTok{[}\BuiltInTok{str}\NormalTok{, torch.Tensor],}
\NormalTok{    local\_steps: }\BuiltInTok{int}\NormalTok{,}
\NormalTok{    update\_metrics: }\BuiltInTok{list}\NormalTok{[}\BuiltInTok{dict}\NormalTok{[}\BuiltInTok{str}\NormalTok{, }\BuiltInTok{float}\NormalTok{]],}
\NormalTok{) }\OperatorTok{{-}\textgreater{}} \BuiltInTok{dict}\NormalTok{[}\BuiltInTok{str}\NormalTok{, Any]:}
    \ImportTok{from}\NormalTok{ rlx.environments }\ImportTok{import}\NormalTok{ basketball }\ImportTok{as}\NormalTok{ basketball\_environment}

\NormalTok{    output.mkdir(parents}\OperatorTok{=}\VariableTok{True}\NormalTok{, exist\_ok}\OperatorTok{=}\VariableTok{True}\NormalTok{)}
\NormalTok{    checkpoint\_path }\OperatorTok{=}\NormalTok{ output }\OperatorTok{/} \StringTok{"checkpoint.pt"}
\NormalTok{    onnx\_path }\OperatorTok{=}\NormalTok{ output }\OperatorTok{/} \StringTok{"policy.onnx"}
\NormalTok{    checkpoint }\OperatorTok{=}\NormalTok{ \{}
        \StringTok{"format"}\NormalTok{: LOCAL\_CHECKPOINT\_FORMAT,}
        \StringTok{"actor\_state\_dict"}\NormalTok{: actor.state\_dict(),}
        \StringTok{"critic\_state\_dict"}\NormalTok{: critic.state\_dict(),}
        \StringTok{"optimizer\_state\_dict"}\NormalTok{: optimizer.state\_dict(),}
        \StringTok{"local\_steps"}\NormalTok{: local\_steps,}
        \StringTok{"total\_steps"}\NormalTok{: source\_metadata.get(}\StringTok{"source\_training\_steps"}\NormalTok{, }\DecValTok{0}\NormalTok{) }\OperatorTok{+}\NormalTok{ local\_steps,}
        \StringTok{"local\_updates"}\NormalTok{: }\BuiltInTok{int}\NormalTok{(update\_metrics[}\OperatorTok{{-}}\DecValTok{1}\NormalTok{][}\StringTok{"update"}\NormalTok{]) }\ControlFlowTok{if}\NormalTok{ update\_metrics }\ControlFlowTok{else} \DecValTok{0}\NormalTok{,}
        \StringTok{"source"}\NormalTok{: source\_metadata,}
        \StringTok{"normalization\_frozen"}\NormalTok{: }\VariableTok{True}\NormalTok{,}
        \StringTok{"actor\_observation\_dim"}\NormalTok{: OBSERVATION\_DIM,}
        \StringTok{"critic\_observation\_dim"}\NormalTok{: OBSERVATION\_DIM,}
\NormalTok{    \}}
\NormalTok{    torch.save(checkpoint, checkpoint\_path)}
\NormalTok{    export\_recurrent\_onnx(actor, onnx\_path)}
\NormalTok{    parity }\OperatorTok{=}\NormalTok{ compare\_actor\_to\_onnx(actor, onnx\_path)}
\NormalTok{    summary }\OperatorTok{=}\NormalTok{ \{}
        \StringTok{"created\_at"}\NormalTok{: datetime.now(timezone.utc).isoformat(),}
        \StringTok{"continuation\_kind"}\NormalTok{: }\StringTok{"actor warm{-}start with fresh local critic and optimizer"}\NormalTok{,}
        \StringTok{"full\_upstream\_resume"}\NormalTok{: }\VariableTok{False}\NormalTok{,}
        \StringTok{"source\_kind"}\NormalTok{: source\_metadata.get(}\StringTok{"source\_kind"}\NormalTok{, }\StringTok{"released\_actor"}\NormalTok{),}
        \StringTok{"source\_training\_steps"}\NormalTok{: source\_metadata.get(}\StringTok{"source\_training\_steps"}\NormalTok{, }\DecValTok{0}\NormalTok{),}
        \StringTok{"source\_checkpoint"}\NormalTok{: }\BuiltInTok{str}\NormalTok{(args.checkpoint.resolve()),}
        \StringTok{"source\_checkpoint\_sha256"}\NormalTok{: source\_metadata[}\StringTok{"source\_hashes"}\NormalTok{][}
            \StringTok{"checkpoint\_sha256"}
\NormalTok{        ],}
        \StringTok{"source\_onnx\_sha256"}\NormalTok{: source\_metadata[}\StringTok{"source\_hashes"}\NormalTok{].get(}\StringTok{"onnx\_sha256"}\NormalTok{),}
        \StringTok{"source\_iteration"}\NormalTok{: source\_metadata[}\StringTok{"source\_iteration"}\NormalTok{],}
        \StringTok{"checkpoint"}\NormalTok{: }\BuiltInTok{str}\NormalTok{(checkpoint\_path.resolve()),}
        \StringTok{"checkpoint\_sha256"}\NormalTok{: sha256\_file(checkpoint\_path),}
        \StringTok{"onnx"}\NormalTok{: }\BuiltInTok{str}\NormalTok{(onnx\_path.resolve()),}
        \StringTok{"onnx\_sha256"}\NormalTok{: sha256\_file(onnx\_path),}
        \StringTok{"onnx\_parity"}\NormalTok{: parity,}
        \StringTok{"local\_steps"}\NormalTok{: local\_steps,}
        \StringTok{"total\_steps"}\NormalTok{: source\_metadata.get(}\StringTok{"source\_training\_steps"}\NormalTok{, }\DecValTok{0}\NormalTok{) }\OperatorTok{+}\NormalTok{ local\_steps,}
        \StringTok{"local\_updates"}\NormalTok{: }\BuiltInTok{int}\NormalTok{(update\_metrics[}\OperatorTok{{-}}\DecValTok{1}\NormalTok{][}\StringTok{"update"}\NormalTok{]) }\ControlFlowTok{if}\NormalTok{ update\_metrics }\ControlFlowTok{else} \DecValTok{0}\NormalTok{,}
        \StringTok{"actor\_weight\_change\_l2"}\NormalTok{: actor\_weight\_change(actor, initial\_actor\_state),}
        \StringTok{"normalization\_frozen"}\NormalTok{: }\VariableTok{True}\NormalTok{,}
        \StringTok{"actor\_observation\_dim"}\NormalTok{: OBSERVATION\_DIM,}
        \StringTok{"ball\_state\_actor"}\NormalTok{: }\VariableTok{False}\NormalTok{,}
        \StringTok{"environment"}\NormalTok{: \{}
            \StringTok{"exact\_upstream\_recipe"}\NormalTok{: }\VariableTok{False}\NormalTok{,}
            \StringTok{"reward\_weights"}\NormalTok{: }\BuiltInTok{dict}\NormalTok{(basketball\_environment.REWARD\_WEIGHTS),}
            \StringTok{"reward\_scale"}\NormalTok{: CTRL\_DT,}
            \StringTok{"omitted\_upstream\_terms"}\NormalTok{: [}\StringTok{"pose"}\NormalTok{, }\StringTok{"angular\_momentum"}\NormalTok{, }\StringTok{"joint\_limits"}\NormalTok{, }\StringTok{"self\_collision"}\NormalTok{],}
            \StringTok{"hold\_levels"}\NormalTok{: }\BuiltInTok{list}\NormalTok{(basketball\_environment.HOLD\_LEVELS),}
            \StringTok{"source\_sha256"}\NormalTok{: sha256\_file(basketball\_environment.}\VariableTok{\_\_file\_\_}\NormalTok{),}
\NormalTok{        \},}
        \StringTok{"settings"}\NormalTok{: \{}
            \StringTok{"num\_envs"}\NormalTok{: args.num\_envs,}
            \StringTok{"steps"}\NormalTok{: args.steps,}
            \StringTok{"learning\_rate"}\NormalTok{: args.learning\_rate,}
            \StringTok{"initial\_std"}\NormalTok{: args.initial\_std,}
            \StringTok{"hold"}\NormalTok{: args.hold,}
            \StringTok{"curriculum"}\NormalTok{: args.curriculum,}
            \StringTok{"actuator"}\NormalTok{: args.actuator,}
            \StringTok{"seed"}\NormalTok{: args.seed,}
            \StringTok{"command"}\NormalTok{: args.command,}
            \StringTok{"pushes"}\NormalTok{: args.pushes,}
            \StringTok{"control\_dt"}\NormalTok{: CTRL\_DT,}
            \StringTok{"target\_kl"}\NormalTok{: args.target\_kl,}
            \StringTok{"obs\_noise"}\NormalTok{: }\VariableTok{False}\NormalTok{,}
            \StringTok{"domain\_rand"}\NormalTok{: }\VariableTok{False}\NormalTok{,}
            \StringTok{"randomized\_commands"}\NormalTok{: args.command }\KeywordTok{is} \VariableTok{None}\NormalTok{,}
\NormalTok{        \},}
        \StringTok{"updates"}\NormalTok{: update\_metrics,}
\NormalTok{    \}}
\NormalTok{    (output }\OperatorTok{/} \StringTok{"summary.json"}\NormalTok{).write\_text(}
\NormalTok{        json.dumps(summary, indent}\OperatorTok{=}\DecValTok{2}\NormalTok{, sort\_keys}\OperatorTok{=}\VariableTok{True}\NormalTok{, allow\_nan}\OperatorTok{=}\VariableTok{False}\NormalTok{) }\OperatorTok{+} \StringTok{"}\CharTok{\textbackslash{}n}\StringTok{"}
\NormalTok{    )}
    \ControlFlowTok{return}\NormalTok{ summary}


\KeywordTok{def}\NormalTok{ train(args: argparse.Namespace) }\OperatorTok{{-}\textgreater{}} \BuiltInTok{dict}\NormalTok{[}\BuiltInTok{str}\NormalTok{, Any]:}
    \ControlFlowTok{if}\NormalTok{ args.output.exists() }\KeywordTok{and}\NormalTok{ (}\KeywordTok{not}\NormalTok{ args.output.is\_dir() }\KeywordTok{or} \BuiltInTok{any}\NormalTok{(args.output.iterdir())):}
        \ControlFlowTok{raise} \PreprocessorTok{FileExistsError}\NormalTok{(}\SpecialStringTok{f"refusing to overwrite an existing run: }\SpecialCharTok{\{}\NormalTok{args}\SpecialCharTok{.}\NormalTok{output}\SpecialCharTok{\}}\SpecialStringTok{"}\NormalTok{)}
\NormalTok{    torch.set\_num\_threads(}\DecValTok{1}\NormalTok{)}
\NormalTok{    torch.manual\_seed(args.seed)}
\NormalTok{    np.random.seed(args.seed)}
\NormalTok{    actor, source\_metadata }\OperatorTok{=}\NormalTok{ load\_source\_actor(}
\NormalTok{        args.checkpoint,}
\NormalTok{        initial\_std}\OperatorTok{=}\NormalTok{args.initial\_std,}
\NormalTok{        verify\_hash}\OperatorTok{=}\VariableTok{True}\NormalTok{,}
\NormalTok{        allow\_local}\OperatorTok{=}\VariableTok{True}\NormalTok{,}
\NormalTok{    )}
\NormalTok{    initial\_actor\_state }\OperatorTok{=}\NormalTok{ \{}
\NormalTok{        name: value.detach().clone() }\ControlFlowTok{for}\NormalTok{ name, value }\KeywordTok{in}\NormalTok{ actor.state\_dict().items()}
\NormalTok{    \}}
\NormalTok{    critic }\OperatorTok{=}\NormalTok{ BasketballCritic(actor.obs\_normalizer)}
\NormalTok{    parameters }\OperatorTok{=} \BuiltInTok{list}\NormalTok{(actor.parameters()) }\OperatorTok{+}\NormalTok{ [}
\NormalTok{        parameter}
        \ControlFlowTok{for}\NormalTok{ name, parameter }\KeywordTok{in}\NormalTok{ critic.named\_parameters()}
        \ControlFlowTok{if} \KeywordTok{not}\NormalTok{ name.startswith(}\StringTok{"normalizer."}\NormalTok{)}
\NormalTok{    ]}
\NormalTok{    optimizer }\OperatorTok{=}\NormalTok{ torch.optim.Adam(parameters, lr}\OperatorTok{=}\NormalTok{args.learning\_rate)}
\NormalTok{    environments }\OperatorTok{=}\NormalTok{ make\_environments(args)}
\NormalTok{    generator }\OperatorTok{=}\NormalTok{ torch.Generator().manual\_seed(args.seed)}
\NormalTok{    reset\_counts }\OperatorTok{=}\NormalTok{ np.zeros(args.num\_envs, dtype}\OperatorTok{=}\NormalTok{np.int64)}
\NormalTok{    update\_metrics: }\BuiltInTok{list}\NormalTok{[}\BuiltInTok{dict}\NormalTok{[}\BuiltInTok{str}\NormalTok{, }\BuiltInTok{float}\NormalTok{]] }\OperatorTok{=}\NormalTok{ []}
\NormalTok{    local\_steps }\OperatorTok{=} \DecValTok{0}
    \ControlFlowTok{try}\NormalTok{:}
\NormalTok{        observations, \_ }\OperatorTok{=}\NormalTok{ reset\_environments(environments, seed}\OperatorTok{=}\NormalTok{args.seed)}
\NormalTok{        hidden, cell }\OperatorTok{=}\NormalTok{ actor.initial\_state(args.num\_envs)}
\NormalTok{        episode\_starts }\OperatorTok{=}\NormalTok{ torch.ones(args.num\_envs, dtype}\OperatorTok{=}\NormalTok{torch.}\BuiltInTok{bool}\NormalTok{)}
        \ControlFlowTok{for}\NormalTok{ update }\KeywordTok{in} \BuiltInTok{range}\NormalTok{(args.updates):}
\NormalTok{            rollout, observations, hidden, cell, episode\_starts }\OperatorTok{=}\NormalTok{ collect\_rollout(}
\NormalTok{                actor,}
\NormalTok{                critic,}
\NormalTok{                environments,}
\NormalTok{                observations,}
\NormalTok{                hidden,}
\NormalTok{                cell,}
\NormalTok{                episode\_starts,}
\NormalTok{                steps}\OperatorTok{=}\NormalTok{args.steps,}
\NormalTok{                generator}\OperatorTok{=}\NormalTok{generator,}
\NormalTok{                seed}\OperatorTok{=}\NormalTok{args.seed,}
\NormalTok{                reset\_counts}\OperatorTok{=}\NormalTok{reset\_counts,}
\NormalTok{            )}
\NormalTok{            metrics }\OperatorTok{=}\NormalTok{ ppo\_update(actor, critic, optimizer, rollout, target\_kl}\OperatorTok{=}\NormalTok{args.target\_kl)}
\NormalTok{            metrics[}\StringTok{"hold\_min"}\NormalTok{] }\OperatorTok{=} \BuiltInTok{min}\NormalTok{(environment.hold }\ControlFlowTok{for}\NormalTok{ environment }\KeywordTok{in}\NormalTok{ environments)}
\NormalTok{            metrics[}\StringTok{"hold\_max"}\NormalTok{] }\OperatorTok{=} \BuiltInTok{max}\NormalTok{(environment.hold }\ControlFlowTok{for}\NormalTok{ environment }\KeywordTok{in}\NormalTok{ environments)}
\NormalTok{            metrics[}\StringTok{"completed\_episodes"}\NormalTok{] }\OperatorTok{=} \BuiltInTok{int}\NormalTok{(reset\_counts.}\BuiltInTok{sum}\NormalTok{())}
\NormalTok{            local\_steps }\OperatorTok{+=}\NormalTok{ args.num\_envs }\OperatorTok{*}\NormalTok{ args.steps}
\NormalTok{            update\_metrics.append(\{}\StringTok{"update"}\NormalTok{: update }\OperatorTok{+} \DecValTok{1}\NormalTok{, }\OperatorTok{**}\NormalTok{metrics\})}
            \BuiltInTok{print}\NormalTok{(}
\NormalTok{                json.dumps(}
\NormalTok{                    \{}
                        \StringTok{"event"}\NormalTok{: }\StringTok{"basketball\_ppo\_update"}\NormalTok{,}
                        \StringTok{"update"}\NormalTok{: update }\OperatorTok{+} \DecValTok{1}\NormalTok{,}
                        \StringTok{"updates"}\NormalTok{: args.updates,}
                        \StringTok{"local\_steps"}\NormalTok{: local\_steps,}
                        \OperatorTok{**}\NormalTok{metrics,}
\NormalTok{                    \},}
\NormalTok{                    allow\_nan}\OperatorTok{=}\VariableTok{False}\NormalTok{,}
\NormalTok{                ),}
\NormalTok{                flush}\OperatorTok{=}\VariableTok{True}\NormalTok{,}
\NormalTok{            )}
            \ControlFlowTok{if}\NormalTok{ (update }\OperatorTok{+} \DecValTok{1}\NormalTok{) }\OperatorTok{\%}\NormalTok{ args.save\_interval }\OperatorTok{==} \DecValTok{0} \KeywordTok{and}\NormalTok{ update }\OperatorTok{+} \DecValTok{1} \OperatorTok{\textless{}}\NormalTok{ args.updates:}
\NormalTok{                save\_results(}
\NormalTok{                    args.output }\OperatorTok{/} \SpecialStringTok{f"update{-}}\SpecialCharTok{\{}\NormalTok{update }\OperatorTok{+} \DecValTok{1}\SpecialCharTok{:04d\}}\SpecialStringTok{"}\NormalTok{, actor, critic, optimizer,}
\NormalTok{                    args}\OperatorTok{=}\NormalTok{args, source\_metadata}\OperatorTok{=}\NormalTok{source\_metadata,}
\NormalTok{                    initial\_actor\_state}\OperatorTok{=}\NormalTok{initial\_actor\_state,}
\NormalTok{                    local\_steps}\OperatorTok{=}\NormalTok{local\_steps, update\_metrics}\OperatorTok{=}\NormalTok{update\_metrics,}
\NormalTok{                )}
    \ControlFlowTok{finally}\NormalTok{:}
        \ControlFlowTok{for}\NormalTok{ environment }\KeywordTok{in}\NormalTok{ environments:}
\NormalTok{            environment.close()}
    \ControlFlowTok{return}\NormalTok{ save\_results(}
\NormalTok{        args.output,}
\NormalTok{        actor,}
\NormalTok{        critic,}
\NormalTok{        optimizer,}
\NormalTok{        args}\OperatorTok{=}\NormalTok{args,}
\NormalTok{        source\_metadata}\OperatorTok{=}\NormalTok{source\_metadata,}
\NormalTok{        initial\_actor\_state}\OperatorTok{=}\NormalTok{initial\_actor\_state,}
\NormalTok{        local\_steps}\OperatorTok{=}\NormalTok{local\_steps,}
\NormalTok{        update\_metrics}\OperatorTok{=}\NormalTok{update\_metrics,}
\NormalTok{    )}


\KeywordTok{def}\NormalTok{ \_positive\_int(value: }\BuiltInTok{str}\NormalTok{) }\OperatorTok{{-}\textgreater{}} \BuiltInTok{int}\NormalTok{:}
\NormalTok{    parsed }\OperatorTok{=} \BuiltInTok{int}\NormalTok{(value)}
    \ControlFlowTok{if}\NormalTok{ parsed }\OperatorTok{\textless{}} \DecValTok{1}\NormalTok{:}
        \ControlFlowTok{raise}\NormalTok{ argparse.ArgumentTypeError(}\StringTok{"must be positive"}\NormalTok{)}
    \ControlFlowTok{return}\NormalTok{ parsed}


\KeywordTok{def}\NormalTok{ \_positive\_float(value: }\BuiltInTok{str}\NormalTok{) }\OperatorTok{{-}\textgreater{}} \BuiltInTok{float}\NormalTok{:}
\NormalTok{    parsed }\OperatorTok{=} \BuiltInTok{float}\NormalTok{(value)}
    \ControlFlowTok{if} \KeywordTok{not}\NormalTok{ math.isfinite(parsed) }\KeywordTok{or}\NormalTok{ parsed }\OperatorTok{\textless{}=} \FloatTok{0.0}\NormalTok{:}
        \ControlFlowTok{raise}\NormalTok{ argparse.ArgumentTypeError(}\StringTok{"must be finite and positive"}\NormalTok{)}
    \ControlFlowTok{return}\NormalTok{ parsed}


\KeywordTok{def}\NormalTok{ \_nonnegative\_float(value: }\BuiltInTok{str}\NormalTok{) }\OperatorTok{{-}\textgreater{}} \BuiltInTok{float}\NormalTok{:}
\NormalTok{    parsed }\OperatorTok{=} \BuiltInTok{float}\NormalTok{(value)}
    \ControlFlowTok{if} \KeywordTok{not}\NormalTok{ math.isfinite(parsed) }\KeywordTok{or}\NormalTok{ parsed }\OperatorTok{\textless{}} \FloatTok{0.0}\NormalTok{:}
        \ControlFlowTok{raise}\NormalTok{ argparse.ArgumentTypeError(}\StringTok{"must be finite and nonnegative"}\NormalTok{)}
    \ControlFlowTok{return}\NormalTok{ parsed}


\KeywordTok{def}\NormalTok{ parse\_args(argv: Sequence[}\BuiltInTok{str}\NormalTok{] }\OperatorTok{|} \VariableTok{None} \OperatorTok{=} \VariableTok{None}\NormalTok{) }\OperatorTok{{-}\textgreater{}}\NormalTok{ argparse.Namespace:}
\NormalTok{    parser }\OperatorTok{=}\NormalTok{ argparse.ArgumentParser(description}\OperatorTok{=}\NormalTok{\_\_doc\_\_)}
\NormalTok{    parser.add\_argument(}\StringTok{"{-}{-}checkpoint"}\NormalTok{, }\BuiltInTok{type}\OperatorTok{=}\NormalTok{Path, default}\OperatorTok{=}\NormalTok{DEFAULT\_CHECKPOINT)}
\NormalTok{    parser.add\_argument(}
        \StringTok{"{-}{-}output"}\NormalTok{,}
        \BuiltInTok{type}\OperatorTok{=}\NormalTok{Path,}
\NormalTok{        default}\OperatorTok{=}\NormalTok{Path(}\StringTok{"runs/basketball/local{-}recurrent{-}ppo"}\NormalTok{),}
\NormalTok{    )}
\NormalTok{    parser.add\_argument(}\StringTok{"{-}{-}num{-}envs"}\NormalTok{, }\BuiltInTok{type}\OperatorTok{=}\NormalTok{\_positive\_int, default}\OperatorTok{=}\DecValTok{4}\NormalTok{)}
\NormalTok{    parser.add\_argument(}\StringTok{"{-}{-}updates"}\NormalTok{, }\BuiltInTok{type}\OperatorTok{=}\NormalTok{\_positive\_int, default}\OperatorTok{=}\DecValTok{5}\NormalTok{)}
\NormalTok{    parser.add\_argument(}\StringTok{"{-}{-}steps"}\NormalTok{, }\BuiltInTok{type}\OperatorTok{=}\NormalTok{\_positive\_int, default}\OperatorTok{=}\DecValTok{32}\NormalTok{)}
\NormalTok{    parser.add\_argument(}\StringTok{"{-}{-}learning{-}rate"}\NormalTok{, }\BuiltInTok{type}\OperatorTok{=}\NormalTok{\_positive\_float, default}\OperatorTok{=}\FloatTok{2e{-}5}\NormalTok{)}
\NormalTok{    parser.add\_argument(}\StringTok{"{-}{-}target{-}kl"}\NormalTok{, }\BuiltInTok{type}\OperatorTok{=}\NormalTok{\_positive\_float, default}\OperatorTok{=}\FloatTok{0.02}\NormalTok{)}
\NormalTok{    parser.add\_argument(}\StringTok{"{-}{-}save{-}interval"}\NormalTok{, }\BuiltInTok{type}\OperatorTok{=}\NormalTok{\_positive\_int, default}\OperatorTok{=}\DecValTok{25}\NormalTok{)}
\NormalTok{    parser.add\_argument(}\StringTok{"{-}{-}initial{-}std"}\NormalTok{, }\BuiltInTok{type}\OperatorTok{=}\NormalTok{\_positive\_float, default}\OperatorTok{=}\FloatTok{0.03}\NormalTok{)}
\NormalTok{    parser.add\_argument(}\StringTok{"{-}{-}hold"}\NormalTok{, }\BuiltInTok{type}\OperatorTok{=}\NormalTok{\_nonnegative\_float, default}\OperatorTok{=}\FloatTok{0.0}\NormalTok{)}
\NormalTok{    parser.add\_argument(}
        \StringTok{"{-}{-}curriculum"}\NormalTok{,}
\NormalTok{        action}\OperatorTok{=}\NormalTok{argparse.BooleanOptionalAction,}
\NormalTok{        default}\OperatorTok{=}\VariableTok{False}\NormalTok{,}
\NormalTok{    )}
\NormalTok{    parser.add\_argument(}\StringTok{"{-}{-}actuator"}\NormalTok{, default}\OperatorTok{=}\StringTok{"bam"}\NormalTok{)}
\NormalTok{    parser.add\_argument(}\StringTok{"{-}{-}seed"}\NormalTok{, }\BuiltInTok{type}\OperatorTok{=}\BuiltInTok{int}\NormalTok{, default}\OperatorTok{=}\DecValTok{42}\NormalTok{)}
\NormalTok{    commands }\OperatorTok{=}\NormalTok{ parser.add\_mutually\_exclusive\_group()}
\NormalTok{    commands.add\_argument(}
        \StringTok{"{-}{-}command"}\NormalTok{,}
        \BuiltInTok{type}\OperatorTok{=}\BuiltInTok{float}\NormalTok{,}
\NormalTok{        nargs}\OperatorTok{=}\DecValTok{3}\NormalTok{,}
\NormalTok{        metavar}\OperatorTok{=}\NormalTok{(}\StringTok{"VX"}\NormalTok{, }\StringTok{"VY"}\NormalTok{, }\StringTok{"WZ"}\NormalTok{),}
\NormalTok{        default}\OperatorTok{=}\NormalTok{(}\FloatTok{0.08}\NormalTok{, }\FloatTok{0.0}\NormalTok{, }\FloatTok{0.0}\NormalTok{),}
\NormalTok{    )}
\NormalTok{    commands.add\_argument(}\StringTok{"{-}{-}randomized{-}commands"}\NormalTok{, action}\OperatorTok{=}\StringTok{"store\_true"}\NormalTok{)}
\NormalTok{    parser.add\_argument(}
        \StringTok{"{-}{-}pushes"}\NormalTok{,}
\NormalTok{        action}\OperatorTok{=}\NormalTok{argparse.BooleanOptionalAction,}
\NormalTok{        default}\OperatorTok{=}\VariableTok{False}\NormalTok{,}
\NormalTok{    )}
\NormalTok{    args }\OperatorTok{=}\NormalTok{ parser.parse\_args(argv)}
    \ControlFlowTok{if}\NormalTok{ args.randomized\_commands:}
\NormalTok{        args.command }\OperatorTok{=} \VariableTok{None}
    \ControlFlowTok{if} \KeywordTok{not}\NormalTok{ args.checkpoint.is\_file():}
\NormalTok{        parser.error(}\SpecialStringTok{f"checkpoint does not exist: }\SpecialCharTok{\{}\NormalTok{args}\SpecialCharTok{.}\NormalTok{checkpoint}\SpecialCharTok{\}}\SpecialStringTok{"}\NormalTok{)}
    \ControlFlowTok{if}\NormalTok{ args.command }\KeywordTok{is} \KeywordTok{not} \VariableTok{None} \KeywordTok{and} \KeywordTok{not} \BuiltInTok{all}\NormalTok{(math.isfinite(x) }\ControlFlowTok{for}\NormalTok{ x }\KeywordTok{in}\NormalTok{ args.command):}
\NormalTok{        parser.error(}\StringTok{"{-}{-}command values must be finite"}\NormalTok{)}
    \ControlFlowTok{return}\NormalTok{ args}


\KeywordTok{def}\NormalTok{ main(argv: Sequence[}\BuiltInTok{str}\NormalTok{] }\OperatorTok{|} \VariableTok{None} \OperatorTok{=} \VariableTok{None}\NormalTok{) }\OperatorTok{{-}\textgreater{}} \BuiltInTok{int}\NormalTok{:}
\NormalTok{    args }\OperatorTok{=}\NormalTok{ parse\_args(argv)}
\NormalTok{    summary }\OperatorTok{=}\NormalTok{ train(args)}
    \BuiltInTok{print}\NormalTok{(json.dumps(summary, allow\_nan}\OperatorTok{=}\VariableTok{False}\NormalTok{), flush}\OperatorTok{=}\VariableTok{True}\NormalTok{)}
    \ControlFlowTok{return} \DecValTok{0}


\ControlFlowTok{if} \VariableTok{\_\_name\_\_} \OperatorTok{==} \StringTok{"\_\_main\_\_"}\NormalTok{:}
    \ControlFlowTok{raise} \PreprocessorTok{SystemExit}\NormalTok{(main())}
\end{Highlighting}
\end{Shaded}

\subsection{\texorpdfstring{A4.
\texttt{rlx/\allowbreak{}rlx/\allowbreak{}environm\allowbreak{}ents/\allowbreak{}bridge.p\allowbreak{}y}}{A4. }}\label{a4.-rlxrlxenvironmentsbridge.py}

SHA-256:
\texttt{49bd21b6\allowbreak{}260841bb\allowbreak{}06ab1d04\allowbreak{}fcca18cd\allowbreak{}79c478e0\allowbreak{}484d8d72\allowbreak{}2775353a\allowbreak{}04a41f37\allowbreak{}}

\begin{Shaded}
\begin{Highlighting}[]
\CommentTok{"""Suspended narrow{-}bridge environment for Microduck locomotion prototypes."""}

\ImportTok{from}\NormalTok{ \_\_future\_\_ }\ImportTok{import}\NormalTok{ annotations}

\ImportTok{import}\NormalTok{ math}
\ImportTok{from}\NormalTok{ pathlib }\ImportTok{import}\NormalTok{ Path}
\ImportTok{from}\NormalTok{ types }\ImportTok{import}\NormalTok{ MappingProxyType}

\ImportTok{import}\NormalTok{ mujoco}
\ImportTok{import}\NormalTok{ numpy }\ImportTok{as}\NormalTok{ np}

\ImportTok{from}\NormalTok{ microduck\_local }\ImportTok{import}\NormalTok{ contract }\ImportTok{as}\NormalTok{ C}
\ImportTok{from}\NormalTok{ microduck\_local.walk\_env }\ImportTok{import}\NormalTok{ MicroduckWalkEnv}

\NormalTok{ROBOT\_DIR }\OperatorTok{=}\NormalTok{ Path(}\VariableTok{\_\_file\_\_}\NormalTok{).resolve().parents[}\DecValTok{1}\NormalTok{] }\OperatorTok{/} \StringTok{"mjlab\_microduck/robot/microduck"}

\NormalTok{CTRL\_DT }\OperatorTok{=}\NormalTok{ C.CTRL\_DT}
\NormalTok{GROUND\_Z }\OperatorTok{=} \FloatTok{0.0}
\NormalTok{PLATFORM\_TOP\_Z }\OperatorTok{=} \FloatTok{0.24}
\NormalTok{PLATFORM\_HALF\_LENGTH }\OperatorTok{=} \FloatTok{0.30}
\NormalTok{PLATFORM\_HALF\_WIDTH }\OperatorTok{=} \FloatTok{0.24}
\NormalTok{START\_PLATFORM\_X }\OperatorTok{=} \OperatorTok{{-}}\FloatTok{0.85}
\NormalTok{END\_PLATFORM\_X }\OperatorTok{=} \FloatTok{0.85}
\NormalTok{BRIDGE\_HALF\_LENGTH }\OperatorTok{=} \FloatTok{0.55}
\NormalTok{BRIDGE\_HALF\_WIDTH }\OperatorTok{=} \FloatTok{0.065}
\NormalTok{BRIDGE\_HALF\_THICKNESS }\OperatorTok{=} \FloatTok{0.015}
\NormalTok{BRIDGE\_MASS }\OperatorTok{=} \FloatTok{0.45}
\NormalTok{BRIDGE\_CENTER\_Z }\OperatorTok{=}\NormalTok{ PLATFORM\_TOP\_Z }\OperatorTok{{-}}\NormalTok{ BRIDGE\_HALF\_THICKNESS}
\NormalTok{FOOT\_CONTACT\_SETTLE\_M }\OperatorTok{=} \FloatTok{0.003}
\NormalTok{ASSISTANCE\_LEVELS }\OperatorTok{=}\NormalTok{ (}\FloatTok{1.0}\NormalTok{, }\FloatTok{0.6}\NormalTok{, }\FloatTok{0.3}\NormalTok{, }\FloatTok{0.1}\NormalTok{, }\FloatTok{0.0}\NormalTok{)}
\NormalTok{SPAWN\_X\_LEVELS }\OperatorTok{=}\NormalTok{ (}\FloatTok{0.0}\NormalTok{, }\OperatorTok{{-}}\FloatTok{0.30}\NormalTok{, }\OperatorTok{{-}}\FloatTok{0.55}\NormalTok{, }\OperatorTok{{-}}\FloatTok{0.75}\NormalTok{, START\_PLATFORM\_X }\OperatorTok{{-}} \FloatTok{0.05}\NormalTok{)}

\NormalTok{CURRICULUM\_KNOBS }\OperatorTok{=}\NormalTok{ MappingProxyType(\{}
    \StringTok{"assistance"}\NormalTok{: }\StringTok{"plank restoring force and level torque only"}\NormalTok{,}
    \StringTok{"spawn\_x"}\NormalTok{: }\StringTok{"progressive physical start position from plank to launch platform"}\NormalTok{,}
\NormalTok{\})}


\KeywordTok{def}\NormalTok{ \_suspension\_length(x: }\BuiltInTok{float}\NormalTok{, y: }\BuiltInTok{float}\NormalTok{, z: }\BuiltInTok{float}\NormalTok{) }\OperatorTok{{-}\textgreater{}} \BuiltInTok{float}\NormalTok{:}
    \ControlFlowTok{return}\NormalTok{ math.sqrt(x }\OperatorTok{*}\NormalTok{ x }\OperatorTok{+}\NormalTok{ y }\OperatorTok{*}\NormalTok{ y }\OperatorTok{+}\NormalTok{ z }\OperatorTok{*}\NormalTok{ z)}


\KeywordTok{def}\NormalTok{ bridge\_model(actuator: }\BuiltInTok{str} \OperatorTok{=} \StringTok{"xml"}\NormalTok{) }\OperatorTok{{-}\textgreater{}}\NormalTok{ mujoco.MjModel:}
    \CommentTok{"""Build a free suspended plank between two fixed safe platforms."""}
    \ControlFlowTok{if}\NormalTok{ actuator }\KeywordTok{not} \KeywordTok{in}\NormalTok{ (}\StringTok{"xml"}\NormalTok{, }\StringTok{"bam"}\NormalTok{):}
        \ControlFlowTok{raise} \PreprocessorTok{ValueError}\NormalTok{(}\StringTok{"actuator must be xml or bam"}\NormalTok{)}

\NormalTok{    spec }\OperatorTok{=}\NormalTok{ mujoco.MjSpec.from\_file(}\BuiltInTok{str}\NormalTok{(ROBOT\_DIR }\OperatorTok{/} \StringTok{"robot\_allcollisions.xml"}\NormalTok{))}
\NormalTok{    spec.option.timestep }\OperatorTok{=}\NormalTok{ C.PHYSICS\_DT}
\NormalTok{    spec.option.integrator }\OperatorTok{=}\NormalTok{ mujoco.mjtIntegrator.mjINT\_IMPLICITFAST}
\NormalTok{    spec.option.cone }\OperatorTok{=}\NormalTok{ mujoco.mjtCone.mjCONE\_PYRAMIDAL}
\NormalTok{    spec.option.iterations }\OperatorTok{=} \DecValTok{12}
\NormalTok{    spec.option.ls\_iterations }\OperatorTok{=} \DecValTok{24}

\NormalTok{    ground\_mat }\OperatorTok{=}\NormalTok{ spec.add\_material(name}\OperatorTok{=}\StringTok{"bridge\_ground\_mat"}\NormalTok{, rgba}\OperatorTok{=}\NormalTok{(}\FloatTok{0.10}\NormalTok{, }\FloatTok{0.12}\NormalTok{, }\FloatTok{0.14}\NormalTok{, }\DecValTok{1}\NormalTok{))}
\NormalTok{    platform\_mat }\OperatorTok{=}\NormalTok{ spec.add\_material(name}\OperatorTok{=}\StringTok{"bridge\_platform\_mat"}\NormalTok{, rgba}\OperatorTok{=}\NormalTok{(}\FloatTok{0.28}\NormalTok{, }\FloatTok{0.31}\NormalTok{, }\FloatTok{0.34}\NormalTok{, }\DecValTok{1}\NormalTok{))}
\NormalTok{    plank\_mat }\OperatorTok{=}\NormalTok{ spec.add\_material(name}\OperatorTok{=}\StringTok{"bridge\_plank\_mat"}\NormalTok{, rgba}\OperatorTok{=}\NormalTok{(}\FloatTok{0.52}\NormalTok{, }\FloatTok{0.31}\NormalTok{, }\FloatTok{0.13}\NormalTok{, }\DecValTok{1}\NormalTok{))}
\NormalTok{    cable\_mat }\OperatorTok{=}\NormalTok{ spec.add\_material(name}\OperatorTok{=}\StringTok{"bridge\_cable\_mat"}\NormalTok{, rgba}\OperatorTok{=}\NormalTok{(}\FloatTok{0.75}\NormalTok{, }\FloatTok{0.77}\NormalTok{, }\FloatTok{0.78}\NormalTok{, }\DecValTok{1}\NormalTok{))}
    \KeywordTok{del}\NormalTok{ ground\_mat, platform\_mat, plank\_mat, cable\_mat}

\NormalTok{    spec.worldbody.add\_geom(}
\NormalTok{        name}\OperatorTok{=}\StringTok{"floor"}\NormalTok{,}
        \BuiltInTok{type}\OperatorTok{=}\NormalTok{mujoco.mjtGeom.mjGEOM\_PLANE,}
\NormalTok{        pos}\OperatorTok{=}\NormalTok{(}\DecValTok{0}\NormalTok{, }\DecValTok{0}\NormalTok{, GROUND\_Z),}
\NormalTok{        size}\OperatorTok{=}\NormalTok{(}\DecValTok{0}\NormalTok{, }\DecValTok{0}\NormalTok{, }\FloatTok{0.05}\NormalTok{),}
\NormalTok{        material}\OperatorTok{=}\StringTok{"bridge\_ground\_mat"}\NormalTok{,}
\NormalTok{        friction}\OperatorTok{=}\NormalTok{(}\FloatTok{1.0}\NormalTok{, }\FloatTok{0.005}\NormalTok{, }\FloatTok{0.0001}\NormalTok{),}
\NormalTok{    )}
    \ControlFlowTok{for}\NormalTok{ name, x }\KeywordTok{in}\NormalTok{ (}
\NormalTok{        (}\StringTok{"start\_platform"}\NormalTok{, START\_PLATFORM\_X),}
\NormalTok{        (}\StringTok{"end\_platform"}\NormalTok{, END\_PLATFORM\_X),}
\NormalTok{    ):}
\NormalTok{        spec.worldbody.add\_geom(}
\NormalTok{            name}\OperatorTok{=}\NormalTok{name,}
            \BuiltInTok{type}\OperatorTok{=}\NormalTok{mujoco.mjtGeom.mjGEOM\_BOX,}
\NormalTok{            group}\OperatorTok{=}\DecValTok{2}\NormalTok{,}
\NormalTok{            pos}\OperatorTok{=}\NormalTok{(x, }\DecValTok{0}\NormalTok{, PLATFORM\_TOP\_Z }\OperatorTok{/} \DecValTok{2}\NormalTok{),}
\NormalTok{            size}\OperatorTok{=}\NormalTok{(PLATFORM\_HALF\_LENGTH, PLATFORM\_HALF\_WIDTH, PLATFORM\_TOP\_Z }\OperatorTok{/} \DecValTok{2}\NormalTok{),}
\NormalTok{            material}\OperatorTok{=}\StringTok{"bridge\_platform\_mat"}\NormalTok{,}
\NormalTok{            friction}\OperatorTok{=}\NormalTok{(}\FloatTok{1.1}\NormalTok{, }\FloatTok{0.005}\NormalTok{, }\FloatTok{0.0001}\NormalTok{),}
\NormalTok{            priority}\OperatorTok{=}\DecValTok{1}\NormalTok{,}
\NormalTok{        )}

\NormalTok{    plank }\OperatorTok{=}\NormalTok{ spec.worldbody.add\_body(}
\NormalTok{        name}\OperatorTok{=}\StringTok{"bridge\_plank"}\NormalTok{,}
\NormalTok{        pos}\OperatorTok{=}\NormalTok{(}\DecValTok{0}\NormalTok{, }\DecValTok{0}\NormalTok{, BRIDGE\_CENTER\_Z),}
\NormalTok{    )}
\NormalTok{    plank.add\_freejoint(name}\OperatorTok{=}\StringTok{"bridge\_freejoint"}\NormalTok{)}
\NormalTok{    plank.add\_geom(}
\NormalTok{        name}\OperatorTok{=}\StringTok{"bridge\_surface"}\NormalTok{,}
        \BuiltInTok{type}\OperatorTok{=}\NormalTok{mujoco.mjtGeom.mjGEOM\_BOX,}
\NormalTok{        group}\OperatorTok{=}\DecValTok{2}\NormalTok{,}
\NormalTok{        size}\OperatorTok{=}\NormalTok{(BRIDGE\_HALF\_LENGTH, BRIDGE\_HALF\_WIDTH, BRIDGE\_HALF\_THICKNESS),}
\NormalTok{        mass}\OperatorTok{=}\NormalTok{BRIDGE\_MASS,}
\NormalTok{        material}\OperatorTok{=}\StringTok{"bridge\_plank\_mat"}\NormalTok{,}
\NormalTok{        friction}\OperatorTok{=}\NormalTok{(}\FloatTok{1.15}\NormalTok{, }\FloatTok{0.006}\NormalTok{, }\FloatTok{0.0002}\NormalTok{),}
\NormalTok{        priority}\OperatorTok{=}\DecValTok{1}\NormalTok{,}
\NormalTok{        condim}\OperatorTok{=}\DecValTok{4}\NormalTok{,}
\NormalTok{    )}

\NormalTok{    anchor\_z }\OperatorTok{=} \FloatTok{0.77}
\NormalTok{    anchor\_y }\OperatorTok{=} \FloatTok{0.17}
\NormalTok{    attach\_y }\OperatorTok{=}\NormalTok{ BRIDGE\_HALF\_WIDTH }\OperatorTok{{-}} \FloatTok{0.008}
\NormalTok{    attach\_x }\OperatorTok{=}\NormalTok{ BRIDGE\_HALF\_LENGTH }\OperatorTok{{-}} \FloatTok{0.10}
\NormalTok{    vertical }\OperatorTok{=}\NormalTok{ anchor\_z }\OperatorTok{{-}}\NormalTok{ BRIDGE\_CENTER\_Z}
\NormalTok{    cable\_length }\OperatorTok{=}\NormalTok{ \_suspension\_length(}\FloatTok{0.0}\NormalTok{, anchor\_y }\OperatorTok{{-}}\NormalTok{ attach\_y, vertical)}
    \ControlFlowTok{for}\NormalTok{ longitudinal, x }\KeywordTok{in}\NormalTok{ ((}\StringTok{"start"}\NormalTok{, }\OperatorTok{{-}}\NormalTok{attach\_x), (}\StringTok{"end"}\NormalTok{, attach\_x)):}
\NormalTok{        spec.worldbody.add\_geom(}
\NormalTok{            name}\OperatorTok{=}\SpecialStringTok{f"bridge\_}\SpecialCharTok{\{}\NormalTok{longitudinal}\SpecialCharTok{\}}\SpecialStringTok{\_crossbar"}\NormalTok{, }\BuiltInTok{type}\OperatorTok{=}\NormalTok{mujoco.mjtGeom.mjGEOM\_BOX,}
\NormalTok{            pos}\OperatorTok{=}\NormalTok{(x, }\DecValTok{0}\NormalTok{, anchor\_z }\OperatorTok{+} \FloatTok{0.014}\NormalTok{), size}\OperatorTok{=}\NormalTok{(}\FloatTok{0.014}\NormalTok{, }\FloatTok{0.24}\NormalTok{, }\FloatTok{0.014}\NormalTok{),}
\NormalTok{            material}\OperatorTok{=}\StringTok{"bridge\_cable\_mat"}\NormalTok{, contype}\OperatorTok{=}\DecValTok{0}\NormalTok{, conaffinity}\OperatorTok{=}\DecValTok{0}\NormalTok{, group}\OperatorTok{=}\DecValTok{2}\NormalTok{,}
\NormalTok{        )}
        \ControlFlowTok{for}\NormalTok{ post\_side }\KeywordTok{in}\NormalTok{ (}\OperatorTok{{-}}\DecValTok{1}\NormalTok{, }\DecValTok{1}\NormalTok{):}
\NormalTok{            spec.worldbody.add\_geom(}
\NormalTok{                name}\OperatorTok{=}\SpecialStringTok{f"bridge\_}\SpecialCharTok{\{}\NormalTok{longitudinal}\SpecialCharTok{\}}\SpecialStringTok{\_post\_}\SpecialCharTok{\{}\NormalTok{post\_side}\SpecialCharTok{\}}\SpecialStringTok{"}\NormalTok{, }\BuiltInTok{type}\OperatorTok{=}\NormalTok{mujoco.mjtGeom.mjGEOM\_BOX,}
\NormalTok{                pos}\OperatorTok{=}\NormalTok{(x, post\_side }\OperatorTok{*} \FloatTok{0.225}\NormalTok{, anchor\_z }\OperatorTok{/} \DecValTok{2}\NormalTok{),}
\NormalTok{                size}\OperatorTok{=}\NormalTok{(}\FloatTok{0.014}\NormalTok{, }\FloatTok{0.014}\NormalTok{, anchor\_z }\OperatorTok{/} \DecValTok{2}\NormalTok{), material}\OperatorTok{=}\StringTok{"bridge\_cable\_mat"}\NormalTok{,}
\NormalTok{                contype}\OperatorTok{=}\DecValTok{0}\NormalTok{, conaffinity}\OperatorTok{=}\DecValTok{0}\NormalTok{, group}\OperatorTok{=}\DecValTok{2}\NormalTok{,}
\NormalTok{            )}
        \ControlFlowTok{for}\NormalTok{ lateral, anchor\_side, attach\_side }\KeywordTok{in}\NormalTok{ (}
\NormalTok{            (}\StringTok{"left"}\NormalTok{, anchor\_y, attach\_y),}
\NormalTok{            (}\StringTok{"right"}\NormalTok{, }\OperatorTok{{-}}\NormalTok{anchor\_y, }\OperatorTok{{-}}\NormalTok{attach\_y),}
\NormalTok{        ):}
\NormalTok{            name }\OperatorTok{=} \SpecialStringTok{f"bridge\_}\SpecialCharTok{\{}\NormalTok{longitudinal}\SpecialCharTok{\}}\SpecialStringTok{\_}\SpecialCharTok{\{}\NormalTok{lateral}\SpecialCharTok{\}}\SpecialStringTok{"}
\NormalTok{            spec.worldbody.add\_site(}
\NormalTok{                name}\OperatorTok{=}\SpecialStringTok{f"}\SpecialCharTok{\{}\NormalTok{name}\SpecialCharTok{\}}\SpecialStringTok{\_anchor"}\NormalTok{,}
\NormalTok{                pos}\OperatorTok{=}\NormalTok{(x, anchor\_side, anchor\_z),}
\NormalTok{                size}\OperatorTok{=}\NormalTok{(}\FloatTok{0.004}\NormalTok{,),}
\NormalTok{                rgba}\OperatorTok{=}\NormalTok{(}\FloatTok{0.75}\NormalTok{, }\FloatTok{0.77}\NormalTok{, }\FloatTok{0.78}\NormalTok{, }\DecValTok{1}\NormalTok{),}
\NormalTok{            )}
\NormalTok{            plank.add\_site(}
\NormalTok{                name}\OperatorTok{=}\SpecialStringTok{f"}\SpecialCharTok{\{}\NormalTok{name}\SpecialCharTok{\}}\SpecialStringTok{\_attach"}\NormalTok{,}
\NormalTok{                pos}\OperatorTok{=}\NormalTok{(x, attach\_side, }\DecValTok{0}\NormalTok{),}
\NormalTok{                size}\OperatorTok{=}\NormalTok{(}\FloatTok{0.004}\NormalTok{,),}
\NormalTok{                rgba}\OperatorTok{=}\NormalTok{(}\FloatTok{0.75}\NormalTok{, }\FloatTok{0.77}\NormalTok{, }\FloatTok{0.78}\NormalTok{, }\DecValTok{1}\NormalTok{),}
\NormalTok{            )}
\NormalTok{            tendon }\OperatorTok{=}\NormalTok{ spec.add\_tendon(}
\NormalTok{                name}\OperatorTok{=}\SpecialStringTok{f"}\SpecialCharTok{\{}\NormalTok{name}\SpecialCharTok{\}}\SpecialStringTok{\_cable"}\NormalTok{,}
\NormalTok{                stiffness}\OperatorTok{=}\FloatTok{320.0}\NormalTok{,}
\NormalTok{                damping}\OperatorTok{=}\FloatTok{2.0}\NormalTok{,}
\NormalTok{                springlength}\OperatorTok{=}\NormalTok{(}\FloatTok{0.0}\NormalTok{, cable\_length }\OperatorTok{{-}} \FloatTok{0.0015}\NormalTok{),}
\NormalTok{                limited}\OperatorTok{=}\VariableTok{True}\NormalTok{,}
                \BuiltInTok{range}\OperatorTok{=}\NormalTok{(}\FloatTok{0.0}\NormalTok{, cable\_length }\OperatorTok{+} \FloatTok{0.008}\NormalTok{),}
\NormalTok{                width}\OperatorTok{=}\FloatTok{0.0015}\NormalTok{,}
\NormalTok{                rgba}\OperatorTok{=}\NormalTok{(}\FloatTok{0.75}\NormalTok{, }\FloatTok{0.77}\NormalTok{, }\FloatTok{0.78}\NormalTok{, }\DecValTok{1}\NormalTok{),}
\NormalTok{                solref\_limit}\OperatorTok{=}\NormalTok{(}\FloatTok{0.015}\NormalTok{, }\FloatTok{1.0}\NormalTok{),}
\NormalTok{                solimp\_limit}\OperatorTok{=}\NormalTok{(}\FloatTok{0.90}\NormalTok{, }\FloatTok{0.95}\NormalTok{, }\FloatTok{0.001}\NormalTok{, }\FloatTok{0.5}\NormalTok{, }\FloatTok{2.0}\NormalTok{),}
\NormalTok{            )}
\NormalTok{            tendon.wrap\_site(}\SpecialStringTok{f"}\SpecialCharTok{\{}\NormalTok{name}\SpecialCharTok{\}}\SpecialStringTok{\_anchor"}\NormalTok{)}
\NormalTok{            tendon.wrap\_site(}\SpecialStringTok{f"}\SpecialCharTok{\{}\NormalTok{name}\SpecialCharTok{\}}\SpecialStringTok{\_attach"}\NormalTok{)}

\NormalTok{    spec.worldbody.add\_light(pos}\OperatorTok{=}\NormalTok{(}\DecValTok{0}\NormalTok{, }\OperatorTok{{-}}\FloatTok{1.5}\NormalTok{, }\DecValTok{3}\NormalTok{), }\BuiltInTok{dir}\OperatorTok{=}\NormalTok{(}\DecValTok{0}\NormalTok{, }\FloatTok{0.3}\NormalTok{, }\OperatorTok{{-}}\DecValTok{1}\NormalTok{),}
\NormalTok{                             diffuse}\OperatorTok{=}\NormalTok{(}\FloatTok{0.85}\NormalTok{, }\FloatTok{0.85}\NormalTok{, }\FloatTok{0.85}\NormalTok{), ambient}\OperatorTok{=}\NormalTok{(}\FloatTok{0.3}\NormalTok{, }\FloatTok{0.3}\NormalTok{, }\FloatTok{0.3}\NormalTok{))}
\NormalTok{    spec.add\_key(}
\NormalTok{        name}\OperatorTok{=}\StringTok{"STAND"}\NormalTok{,}
\NormalTok{        qpos}\OperatorTok{=}\NormalTok{np.concatenate((}
\NormalTok{            (}\DecValTok{0}\NormalTok{, }\DecValTok{0}\NormalTok{, }\FloatTok{0.12}\NormalTok{, }\DecValTok{1}\NormalTok{, }\DecValTok{0}\NormalTok{, }\DecValTok{0}\NormalTok{, }\DecValTok{0}\NormalTok{),}
\NormalTok{            C.DEFAULT\_POSE,}
\NormalTok{            (}\DecValTok{0}\NormalTok{, }\DecValTok{0}\NormalTok{, BRIDGE\_CENTER\_Z, }\DecValTok{1}\NormalTok{, }\DecValTok{0}\NormalTok{, }\DecValTok{0}\NormalTok{, }\DecValTok{0}\NormalTok{),}
\NormalTok{        )),}
\NormalTok{        ctrl}\OperatorTok{=}\NormalTok{C.DEFAULT\_POSE,}
\NormalTok{    )}
\NormalTok{    model }\OperatorTok{=}\NormalTok{ spec.}\BuiltInTok{compile}\NormalTok{()}
\NormalTok{    model.vis.global\_.offwidth }\OperatorTok{=} \DecValTok{1280}
\NormalTok{    model.vis.global\_.offheight }\OperatorTok{=} \DecValTok{960}
    \ControlFlowTok{return}\NormalTok{ model}


\KeywordTok{def}\NormalTok{ curriculum\_level(stage: }\BuiltInTok{int}\NormalTok{, }\OperatorTok{*}\NormalTok{, crossed: }\BuiltInTok{bool}\NormalTok{, elapsed\_s: }\BuiltInTok{float}\NormalTok{) }\OperatorTok{{-}\textgreater{}} \BuiltInTok{int}\NormalTok{:}
    \CommentTok{"""Advance after an honest crossing; retreat after an immediate failure."""}
    \ControlFlowTok{if}\NormalTok{ (}
        \KeywordTok{not} \DecValTok{0} \OperatorTok{\textless{}=}\NormalTok{ stage }\OperatorTok{\textless{}} \BuiltInTok{len}\NormalTok{(ASSISTANCE\_LEVELS)}
        \KeywordTok{or} \KeywordTok{not}\NormalTok{ math.isfinite(elapsed\_s)}
        \KeywordTok{or}\NormalTok{ elapsed\_s }\OperatorTok{\textless{}} \DecValTok{0}
\NormalTok{    ):}
        \ControlFlowTok{raise} \PreprocessorTok{ValueError}\NormalTok{(}\StringTok{"invalid curriculum stage or elapsed time"}\NormalTok{)}
    \ControlFlowTok{if}\NormalTok{ crossed:}
        \ControlFlowTok{return} \BuiltInTok{min}\NormalTok{(stage }\OperatorTok{+} \DecValTok{1}\NormalTok{, }\BuiltInTok{len}\NormalTok{(ASSISTANCE\_LEVELS) }\OperatorTok{{-}} \DecValTok{1}\NormalTok{)}
    \ControlFlowTok{if}\NormalTok{ elapsed\_s }\OperatorTok{\textless{}} \FloatTok{2.0}\NormalTok{:}
        \ControlFlowTok{return} \BuiltInTok{max}\NormalTok{(stage }\OperatorTok{{-}} \DecValTok{1}\NormalTok{, }\DecValTok{0}\NormalTok{)}
    \ControlFlowTok{return}\NormalTok{ stage}


\KeywordTok{class}\NormalTok{ BridgeEnv(MicroduckWalkEnv):}
    \CommentTok{"""Gymnasium environment for crossing a physically suspended narrow plank.}

\CommentTok{    The policy receives the standard 61{-}float Microduck observation only.}
\CommentTok{    Bridge pose, world progress, contacts, and crossing state are privileged}
\CommentTok{    metrics. Assistance acts on the plank, never on the robot, and the inherited}
\CommentTok{    locomotion reward is identical at every curriculum level.}
\CommentTok{    """}

    \KeywordTok{def} \FunctionTok{\_\_init\_\_}\NormalTok{(}
        \VariableTok{self}\NormalTok{,}
        \OperatorTok{*}\NormalTok{,}
\NormalTok{        max\_episode\_s: }\BuiltInTok{float} \OperatorTok{=} \FloatTok{20.0}\NormalTok{,}
\NormalTok{        seed: }\BuiltInTok{int} \OperatorTok{|} \VariableTok{None} \OperatorTok{=} \VariableTok{None}\NormalTok{,}
\NormalTok{        obs\_noise: }\BuiltInTok{bool} \OperatorTok{=} \VariableTok{False}\NormalTok{,}
\NormalTok{        domain\_rand: }\BuiltInTok{bool} \OperatorTok{=} \VariableTok{False}\NormalTok{,}
\NormalTok{        pushes: }\BuiltInTok{bool} \OperatorTok{=} \VariableTok{False}\NormalTok{,}
\NormalTok{        random\_yaw: }\BuiltInTok{bool} \OperatorTok{=} \VariableTok{False}\NormalTok{,}
\NormalTok{        actuator: }\BuiltInTok{str} \OperatorTok{=} \StringTok{"xml"}\NormalTok{,}
\NormalTok{        action\_delay: }\BuiltInTok{bool} \OperatorTok{=} \VariableTok{False}\NormalTok{,}
\NormalTok{        command: }\BuiltInTok{tuple}\NormalTok{[}\BuiltInTok{float}\NormalTok{, }\BuiltInTok{float}\NormalTok{, }\BuiltInTok{float}\NormalTok{] }\OperatorTok{|} \VariableTok{None} \OperatorTok{=}\NormalTok{ (}\FloatTok{0.25}\NormalTok{, }\FloatTok{0.0}\NormalTok{, }\FloatTok{0.0}\NormalTok{),}
\NormalTok{        navigation: }\BuiltInTok{bool} \OperatorTok{=} \VariableTok{True}\NormalTok{,}
\NormalTok{        assistance: }\BuiltInTok{float} \OperatorTok{|} \VariableTok{None} \OperatorTok{=} \VariableTok{None}\NormalTok{,}
\NormalTok{        curriculum: }\BuiltInTok{bool} \OperatorTok{=} \VariableTok{False}\NormalTok{,}
\NormalTok{        model: mujoco.MjModel }\OperatorTok{|} \VariableTok{None} \OperatorTok{=} \VariableTok{None}\NormalTok{,}
\NormalTok{    ) }\OperatorTok{{-}\textgreater{}} \VariableTok{None}\NormalTok{:}
        \ControlFlowTok{if} \KeywordTok{not}\NormalTok{ math.isfinite(max\_episode\_s) }\KeywordTok{or}\NormalTok{ max\_episode\_s }\OperatorTok{\textless{}=} \DecValTok{0}\NormalTok{:}
            \ControlFlowTok{raise} \PreprocessorTok{ValueError}\NormalTok{(}\StringTok{"episode length must be positive and finite"}\NormalTok{)}
        \ControlFlowTok{if}\NormalTok{ assistance }\KeywordTok{is} \VariableTok{None}\NormalTok{:}
\NormalTok{            assistance }\OperatorTok{=} \FloatTok{1.0} \ControlFlowTok{if}\NormalTok{ curriculum }\ControlFlowTok{else} \FloatTok{0.0}
        \ControlFlowTok{if} \KeywordTok{not}\NormalTok{ math.isfinite(assistance) }\KeywordTok{or} \KeywordTok{not} \DecValTok{0} \OperatorTok{\textless{}=}\NormalTok{ assistance }\OperatorTok{\textless{}=} \DecValTok{1}\NormalTok{:}
            \ControlFlowTok{raise} \PreprocessorTok{ValueError}\NormalTok{(}\StringTok{"assistance must be finite and between zero and one"}\NormalTok{)}
        \ControlFlowTok{if}\NormalTok{ curriculum }\KeywordTok{and}\NormalTok{ assistance }\KeywordTok{not} \KeywordTok{in}\NormalTok{ ASSISTANCE\_LEVELS:}
            \ControlFlowTok{raise} \PreprocessorTok{ValueError}\NormalTok{(}\StringTok{"curriculum assistance must be one of ASSISTANCE\_LEVELS"}\NormalTok{)}

        \VariableTok{self}\NormalTok{.fixed\_command }\OperatorTok{=} \VariableTok{self}\NormalTok{.\_validate\_command(command)}
        \VariableTok{self}\NormalTok{.navigation }\OperatorTok{=} \BuiltInTok{bool}\NormalTok{(navigation)}
        \VariableTok{self}\NormalTok{.assistance }\OperatorTok{=} \BuiltInTok{float}\NormalTok{(assistance)}
        \VariableTok{self}\NormalTok{.curriculum }\OperatorTok{=} \BuiltInTok{bool}\NormalTok{(curriculum)}
        \VariableTok{self}\NormalTok{.stage }\OperatorTok{=}\NormalTok{ (}
\NormalTok{            ASSISTANCE\_LEVELS.index(assistance)}
            \ControlFlowTok{if}\NormalTok{ assistance }\KeywordTok{in}\NormalTok{ ASSISTANCE\_LEVELS}
            \ControlFlowTok{else} \BuiltInTok{len}\NormalTok{(ASSISTANCE\_LEVELS) }\OperatorTok{{-}} \DecValTok{1}
\NormalTok{        )}
        \VariableTok{self}\NormalTok{.pushes }\OperatorTok{=} \BuiltInTok{bool}\NormalTok{(pushes)}
        \VariableTok{self}\NormalTok{.\_bridge\_random\_yaw }\OperatorTok{=} \BuiltInTok{bool}\NormalTok{(random\_yaw)}
        \VariableTok{self}\NormalTok{.\_episode\_finished }\OperatorTok{=} \VariableTok{False}
        \VariableTok{self}\NormalTok{.\_crossed }\OperatorTok{=} \VariableTok{False}
        \VariableTok{self}\NormalTok{.\_bridge\_contact\_observed }\OperatorTok{=} \VariableTok{False}
        \VariableTok{self}\NormalTok{.\_next\_push }\OperatorTok{=}\NormalTok{ math.inf}
        \VariableTok{self}\NormalTok{.\_last\_push\_force }\OperatorTok{=} \FloatTok{0.0}

        \BuiltInTok{super}\NormalTok{().}\FunctionTok{\_\_init\_\_}\NormalTok{(}
\NormalTok{            max\_episode\_s}\OperatorTok{=}\NormalTok{max\_episode\_s,}
\NormalTok{            command\_resample\_s}\OperatorTok{=}\FloatTok{5.0}\NormalTok{,}
\NormalTok{            obs\_noise}\OperatorTok{=}\NormalTok{obs\_noise,}
\NormalTok{            domain\_rand}\OperatorTok{=}\NormalTok{domain\_rand,}
\NormalTok{            action\_delay}\OperatorTok{=}\NormalTok{action\_delay,}
\NormalTok{            random\_yaw}\OperatorTok{=}\VariableTok{False}\NormalTok{,}
\NormalTok{            seed}\OperatorTok{=}\NormalTok{seed,}
\NormalTok{            actuator\_force}\OperatorTok{=}\NormalTok{actuator,}
\NormalTok{            model}\OperatorTok{=}\NormalTok{model }\ControlFlowTok{if}\NormalTok{ model }\KeywordTok{is} \KeywordTok{not} \VariableTok{None} \ControlFlowTok{else}\NormalTok{ bridge\_model(actuator),}
\NormalTok{            terminate\_on\_fall}\OperatorTok{=}\VariableTok{False}\NormalTok{,}
\NormalTok{            height\_termination}\OperatorTok{=}\VariableTok{False}\NormalTok{,}
\NormalTok{        )}

        \VariableTok{self}\NormalTok{.bridge\_body\_id }\OperatorTok{=} \VariableTok{self}\NormalTok{.model.body(}\StringTok{"bridge\_plank"}\NormalTok{).}\BuiltInTok{id}
        \VariableTok{self}\NormalTok{.bridge\_geom\_id }\OperatorTok{=} \VariableTok{self}\NormalTok{.model.geom(}\StringTok{"bridge\_surface"}\NormalTok{).}\BuiltInTok{id}
        \VariableTok{self}\NormalTok{.bridge\_qpos\_adr }\OperatorTok{=} \BuiltInTok{int}\NormalTok{(}\VariableTok{self}\NormalTok{.model.joint(}\StringTok{"bridge\_freejoint"}\NormalTok{).qposadr[}\DecValTok{0}\NormalTok{])}
        \VariableTok{self}\NormalTok{.bridge\_qvel\_adr }\OperatorTok{=} \BuiltInTok{int}\NormalTok{(}\VariableTok{self}\NormalTok{.model.joint(}\StringTok{"bridge\_freejoint"}\NormalTok{).dofadr[}\DecValTok{0}\NormalTok{])}
        \VariableTok{self}\NormalTok{.start\_platform\_geom }\OperatorTok{=} \VariableTok{self}\NormalTok{.model.geom(}\StringTok{"start\_platform"}\NormalTok{).}\BuiltInTok{id}
        \VariableTok{self}\NormalTok{.end\_platform\_geom }\OperatorTok{=} \VariableTok{self}\NormalTok{.model.geom(}\StringTok{"end\_platform"}\NormalTok{).}\BuiltInTok{id}
        \VariableTok{self}\NormalTok{.support\_geoms }\OperatorTok{=}\NormalTok{ \{}
            \VariableTok{self}\NormalTok{.start\_platform\_geom,}
            \VariableTok{self}\NormalTok{.bridge\_geom\_id,}
            \VariableTok{self}\NormalTok{.end\_platform\_geom,}
\NormalTok{        \}}
        \VariableTok{self}\NormalTok{.\_bridge\_nominal\_position }\OperatorTok{=}\NormalTok{ np.array((}\FloatTok{0.0}\NormalTok{, }\FloatTok{0.0}\NormalTok{, BRIDGE\_CENTER\_Z))}

    \AttributeTok{@staticmethod}
    \KeywordTok{def}\NormalTok{ \_validate\_command(command):}
        \ControlFlowTok{if}\NormalTok{ command }\KeywordTok{is} \VariableTok{None}\NormalTok{:}
            \ControlFlowTok{return} \VariableTok{None}
\NormalTok{        array }\OperatorTok{=}\NormalTok{ np.asarray(command, dtype}\OperatorTok{=}\NormalTok{np.float32)}
        \ControlFlowTok{if}\NormalTok{ array.shape }\OperatorTok{!=}\NormalTok{ (}\DecValTok{3}\NormalTok{,) }\KeywordTok{or} \KeywordTok{not}\NormalTok{ np.isfinite(array).}\BuiltInTok{all}\NormalTok{():}
            \ControlFlowTok{raise} \PreprocessorTok{ValueError}\NormalTok{(}\StringTok{"command must contain three finite numbers"}\NormalTok{)}
        \ControlFlowTok{return}\NormalTok{ array.copy()}

    \KeywordTok{def}\NormalTok{ set\_command(}\VariableTok{self}\NormalTok{, command: }\BuiltInTok{tuple}\NormalTok{[}\BuiltInTok{float}\NormalTok{, }\BuiltInTok{float}\NormalTok{, }\BuiltInTok{float}\NormalTok{]) }\OperatorTok{{-}\textgreater{}} \VariableTok{None}\NormalTok{:}
\NormalTok{        validated }\OperatorTok{=} \VariableTok{self}\NormalTok{.\_validate\_command(command)}
        \ControlFlowTok{if}\NormalTok{ validated }\KeywordTok{is} \VariableTok{None}\NormalTok{:}
            \ControlFlowTok{raise} \PreprocessorTok{ValueError}\NormalTok{(}\StringTok{"set\_command requires three numbers"}\NormalTok{)}
        \VariableTok{self}\NormalTok{.fixed\_command }\OperatorTok{=}\NormalTok{ validated}
        \VariableTok{self}\NormalTok{.twist\_cmd[:] }\OperatorTok{=}\NormalTok{ validated}

    \KeywordTok{def}\NormalTok{ \_sample\_commands(}\VariableTok{self}\NormalTok{) }\OperatorTok{{-}\textgreater{}} \VariableTok{None}\NormalTok{:}
        \ControlFlowTok{if} \VariableTok{self}\NormalTok{.fixed\_command }\KeywordTok{is} \KeywordTok{not} \VariableTok{None}\NormalTok{:}
            \VariableTok{self}\NormalTok{.twist\_cmd[:] }\OperatorTok{=} \VariableTok{self}\NormalTok{.fixed\_command}
        \ControlFlowTok{else}\NormalTok{:}
            \VariableTok{self}\NormalTok{.twist\_cmd[:] }\OperatorTok{=}\NormalTok{ (}
                \VariableTok{self}\NormalTok{.\_rng.uniform(}\FloatTok{0.12}\NormalTok{, }\FloatTok{0.24}\NormalTok{),}
                \VariableTok{self}\NormalTok{.\_rng.uniform(}\OperatorTok{{-}}\FloatTok{0.025}\NormalTok{, }\FloatTok{0.025}\NormalTok{),}
                \VariableTok{self}\NormalTok{.\_rng.uniform(}\OperatorTok{{-}}\FloatTok{0.12}\NormalTok{, }\FloatTok{0.12}\NormalTok{),}
\NormalTok{            )}
        \VariableTok{self}\NormalTok{.head\_cmd[:] }\OperatorTok{=} \DecValTok{0}
        \VariableTok{self}\NormalTok{.body\_cmd[:] }\OperatorTok{=} \DecValTok{0}

    \KeywordTok{def}\NormalTok{ \_update\_navigation\_command(}\VariableTok{self}\NormalTok{) }\OperatorTok{{-}\textgreater{}} \VariableTok{None}\NormalTok{:}
        \ControlFlowTok{if} \KeywordTok{not} \VariableTok{self}\NormalTok{.navigation }\KeywordTok{or} \VariableTok{self}\NormalTok{.fixed\_command }\KeywordTok{is} \VariableTok{None}\NormalTok{:}
            \ControlFlowTok{return}
\NormalTok{        rotation }\OperatorTok{=} \VariableTok{self}\NormalTok{.data.xmat[}\VariableTok{self}\NormalTok{.trunk\_body\_id].reshape(}\DecValTok{3}\NormalTok{, }\DecValTok{3}\NormalTok{)}
\NormalTok{        heading }\OperatorTok{=}\NormalTok{ rotation[:, }\DecValTok{0}\NormalTok{].astype(np.float64)}
\NormalTok{        heading[}\DecValTok{2}\NormalTok{] }\OperatorTok{=} \DecValTok{0}
\NormalTok{        norm }\OperatorTok{=} \BuiltInTok{float}\NormalTok{(np.linalg.norm(heading))}
\NormalTok{        heading }\OperatorTok{=}\NormalTok{ heading }\OperatorTok{/}\NormalTok{ norm }\ControlFlowTok{if}\NormalTok{ norm }\OperatorTok{\textgreater{}} \FloatTok{1e{-}9} \ControlFlowTok{else}\NormalTok{ np.array((}\FloatTok{1.0}\NormalTok{, }\FloatTok{0.0}\NormalTok{, }\FloatTok{0.0}\NormalTok{))}
\NormalTok{        side }\OperatorTok{=}\NormalTok{ np.array((}\OperatorTok{{-}}\NormalTok{heading[}\DecValTok{1}\NormalTok{], heading[}\DecValTok{0}\NormalTok{], }\FloatTok{0.0}\NormalTok{))}
\NormalTok{        reached\_destination }\OperatorTok{=}\NormalTok{ (}
            \VariableTok{self}\NormalTok{.data.xpos[}\VariableTok{self}\NormalTok{.trunk\_body\_id, }\DecValTok{0}\NormalTok{] }\OperatorTok{\textgreater{}=}\NormalTok{ END\_PLATFORM\_X }\OperatorTok{{-}} \FloatTok{0.05}
\NormalTok{        )}
\NormalTok{        desired\_world }\OperatorTok{=}\NormalTok{ np.array((}
            \FloatTok{0.0} \ControlFlowTok{if}\NormalTok{ reached\_destination }\ControlFlowTok{else} \BuiltInTok{float}\NormalTok{(}\VariableTok{self}\NormalTok{.fixed\_command[}\DecValTok{0}\NormalTok{]),}
            \BuiltInTok{float}\NormalTok{(np.clip(}\OperatorTok{{-}}\FloatTok{1.2} \OperatorTok{*} \VariableTok{self}\NormalTok{.data.xpos[}\VariableTok{self}\NormalTok{.trunk\_body\_id, }\DecValTok{1}\NormalTok{], }\OperatorTok{{-}}\FloatTok{0.12}\NormalTok{, }\FloatTok{0.12}\NormalTok{)),}
            \FloatTok{0.0}\NormalTok{,}
\NormalTok{        ))}
\NormalTok{        yaw\_error }\OperatorTok{=}\NormalTok{ math.atan2(}\BuiltInTok{float}\NormalTok{(heading[}\DecValTok{1}\NormalTok{]), }\BuiltInTok{float}\NormalTok{(heading[}\DecValTok{0}\NormalTok{]))}
        \VariableTok{self}\NormalTok{.twist\_cmd[:] }\OperatorTok{=}\NormalTok{ (}
            \BuiltInTok{float}\NormalTok{(desired\_world }\OperatorTok{@}\NormalTok{ heading),}
            \BuiltInTok{float}\NormalTok{(desired\_world }\OperatorTok{@}\NormalTok{ side),}
            \BuiltInTok{float}\NormalTok{(np.clip(}\OperatorTok{{-}}\FloatTok{1.5} \OperatorTok{*}\NormalTok{ yaw\_error, }\OperatorTok{{-}}\FloatTok{0.6}\NormalTok{, }\FloatTok{0.6}\NormalTok{)),}
\NormalTok{        )}

    \KeywordTok{def}\NormalTok{ \_get\_obs(}\VariableTok{self}\NormalTok{) }\OperatorTok{{-}\textgreater{}}\NormalTok{ np.ndarray:}
        \VariableTok{self}\NormalTok{.\_update\_navigation\_command()}
        \ControlFlowTok{return} \BuiltInTok{super}\NormalTok{().\_get\_obs()}

    \KeywordTok{def}\NormalTok{ reset(}\VariableTok{self}\NormalTok{, }\OperatorTok{*}\NormalTok{, seed: }\BuiltInTok{int} \OperatorTok{|} \VariableTok{None} \OperatorTok{=} \VariableTok{None}\NormalTok{, options: }\BuiltInTok{dict} \OperatorTok{|} \VariableTok{None} \OperatorTok{=} \VariableTok{None}\NormalTok{):}
        \ControlFlowTok{if} \VariableTok{self}\NormalTok{.curriculum }\KeywordTok{and} \VariableTok{self}\NormalTok{.\_episode\_finished:}
            \VariableTok{self}\NormalTok{.stage }\OperatorTok{=}\NormalTok{ curriculum\_level(}
                \VariableTok{self}\NormalTok{.stage,}
\NormalTok{                crossed}\OperatorTok{=}\VariableTok{self}\NormalTok{.\_crossed,}
\NormalTok{                elapsed\_s}\OperatorTok{=}\VariableTok{self}\NormalTok{.step\_count }\OperatorTok{*}\NormalTok{ CTRL\_DT,}
\NormalTok{            )}
            \VariableTok{self}\NormalTok{.assistance }\OperatorTok{=}\NormalTok{ ASSISTANCE\_LEVELS[}\VariableTok{self}\NormalTok{.stage]}

        \VariableTok{self}\NormalTok{.\_episode\_finished }\OperatorTok{=} \VariableTok{False}
        \VariableTok{self}\NormalTok{.\_crossed }\OperatorTok{=} \VariableTok{False}
        \VariableTok{self}\NormalTok{.\_bridge\_contact\_observed }\OperatorTok{=} \VariableTok{False}
        \BuiltInTok{super}\NormalTok{().reset(seed}\OperatorTok{=}\NormalTok{seed, options}\OperatorTok{=}\NormalTok{options)}

\NormalTok{        spawn\_x }\OperatorTok{=}\NormalTok{ (}
\NormalTok{            SPAWN\_X\_LEVELS[}\VariableTok{self}\NormalTok{.stage]}
            \ControlFlowTok{if} \VariableTok{self}\NormalTok{.curriculum}
            \ControlFlowTok{else}\NormalTok{ START\_PLATFORM\_X }\OperatorTok{{-}} \FloatTok{0.05}
\NormalTok{        )}
        \VariableTok{self}\NormalTok{.data.qpos[:}\DecValTok{3}\NormalTok{] }\OperatorTok{=}\NormalTok{ (}
\NormalTok{            spawn\_x,}
            \VariableTok{self}\NormalTok{.\_rng.uniform(}\OperatorTok{{-}}\FloatTok{0.015}\NormalTok{, }\FloatTok{0.015}\NormalTok{),}
            \VariableTok{self}\NormalTok{.stand\_z }\OperatorTok{+}\NormalTok{ PLATFORM\_TOP\_Z }\OperatorTok{{-}}\NormalTok{ FOOT\_CONTACT\_SETTLE\_M,}
\NormalTok{        )}
\NormalTok{        yaw }\OperatorTok{=}\NormalTok{ (}
            \VariableTok{self}\NormalTok{.\_rng.uniform(}\OperatorTok{{-}}\NormalTok{math.radians(}\DecValTok{5}\NormalTok{), math.radians(}\DecValTok{5}\NormalTok{))}
            \ControlFlowTok{if} \VariableTok{self}\NormalTok{.\_bridge\_random\_yaw}
            \ControlFlowTok{else} \FloatTok{0.0}
\NormalTok{        )}
        \VariableTok{self}\NormalTok{.data.qpos[}\DecValTok{3}\NormalTok{:}\DecValTok{7}\NormalTok{] }\OperatorTok{=}\NormalTok{ (math.cos(yaw }\OperatorTok{/} \DecValTok{2}\NormalTok{), }\DecValTok{0}\NormalTok{, }\DecValTok{0}\NormalTok{, math.sin(yaw }\OperatorTok{/} \DecValTok{2}\NormalTok{))}
        \VariableTok{self}\NormalTok{.data.qpos[}\VariableTok{self}\NormalTok{.joint\_qpos\_adr] }\OperatorTok{=}\NormalTok{ (}
\NormalTok{            C.DEFAULT\_POSE }\OperatorTok{+} \VariableTok{self}\NormalTok{.\_rng.uniform(}\OperatorTok{{-}}\FloatTok{0.025}\NormalTok{, }\FloatTok{0.025}\NormalTok{, C.NUM\_JOINTS)}
\NormalTok{        )}

\NormalTok{        roll }\OperatorTok{=} \VariableTok{self}\NormalTok{.\_rng.uniform(}\OperatorTok{{-}}\NormalTok{math.radians(}\FloatTok{1.5}\NormalTok{), math.radians(}\FloatTok{1.5}\NormalTok{))}
\NormalTok{        pitch }\OperatorTok{=} \VariableTok{self}\NormalTok{.\_rng.uniform(}\OperatorTok{{-}}\NormalTok{math.radians(}\FloatTok{0.8}\NormalTok{), math.radians(}\FloatTok{0.8}\NormalTok{))}
\NormalTok{        rotation }\OperatorTok{=}\NormalTok{ np.empty(}\DecValTok{4}\NormalTok{)}
\NormalTok{        roll\_quat }\OperatorTok{=}\NormalTok{ np.array((math.cos(roll }\OperatorTok{/} \DecValTok{2}\NormalTok{), math.sin(roll }\OperatorTok{/} \DecValTok{2}\NormalTok{), }\DecValTok{0}\NormalTok{, }\DecValTok{0}\NormalTok{))}
\NormalTok{        pitch\_quat }\OperatorTok{=}\NormalTok{ np.array((math.cos(pitch }\OperatorTok{/} \DecValTok{2}\NormalTok{), }\DecValTok{0}\NormalTok{, math.sin(pitch }\OperatorTok{/} \DecValTok{2}\NormalTok{), }\DecValTok{0}\NormalTok{))}
\NormalTok{        mujoco.mju\_mulQuat(rotation, pitch\_quat, roll\_quat)}
        \VariableTok{self}\NormalTok{.data.qpos[}
            \VariableTok{self}\NormalTok{.bridge\_qpos\_adr:}\VariableTok{self}\NormalTok{.bridge\_qpos\_adr }\OperatorTok{+} \DecValTok{7}
\NormalTok{        ] }\OperatorTok{=}\NormalTok{ (}\OperatorTok{*}\VariableTok{self}\NormalTok{.\_bridge\_nominal\_position, }\OperatorTok{*}\NormalTok{rotation)}
        \VariableTok{self}\NormalTok{.data.qvel[:] }\OperatorTok{=} \DecValTok{0}
        \VariableTok{self}\NormalTok{.data.qvel[}\VariableTok{self}\NormalTok{.bridge\_qvel\_adr }\OperatorTok{+} \DecValTok{3}\NormalTok{:}\VariableTok{self}\NormalTok{.bridge\_qvel\_adr }\OperatorTok{+} \DecValTok{6}\NormalTok{] }\OperatorTok{=}\NormalTok{ (}
            \VariableTok{self}\NormalTok{.\_rng.uniform(}\OperatorTok{{-}}\FloatTok{0.025}\NormalTok{, }\FloatTok{0.025}\NormalTok{, }\DecValTok{3}\NormalTok{)}
\NormalTok{        )}
        \VariableTok{self}\NormalTok{.data.xfrc\_applied[:] }\OperatorTok{=} \DecValTok{0}
        \VariableTok{self}\NormalTok{.data.ctrl[:] }\OperatorTok{=} \VariableTok{self}\NormalTok{.data.qpos[}\VariableTok{self}\NormalTok{.joint\_qpos\_adr]}
\NormalTok{        mujoco.mj\_forward(}\VariableTok{self}\NormalTok{.model, }\VariableTok{self}\NormalTok{.data)}
        \ControlFlowTok{if} \VariableTok{self}\NormalTok{.bam }\KeywordTok{is} \KeywordTok{not} \VariableTok{None}\NormalTok{:}
            \VariableTok{self}\NormalTok{.bam.reset(}\VariableTok{self}\NormalTok{.data.qpos[}\VariableTok{self}\NormalTok{.joint\_qpos\_adr])}
        \VariableTok{self}\NormalTok{.prev\_joint\_vel[:] }\OperatorTok{=} \VariableTok{self}\NormalTok{.\_joint\_vel()}
        \VariableTok{self}\NormalTok{.\_next\_push }\OperatorTok{=} \VariableTok{self}\NormalTok{.\_rng.uniform(}\FloatTok{1.5}\NormalTok{, }\FloatTok{3.0}\NormalTok{) }\ControlFlowTok{if} \VariableTok{self}\NormalTok{.pushes }\ControlFlowTok{else}\NormalTok{ math.inf}
        \VariableTok{self}\NormalTok{.\_last\_push\_force }\OperatorTok{=} \FloatTok{0.0}
        \ControlFlowTok{return} \VariableTok{self}\NormalTok{.\_get\_obs(), }\VariableTok{self}\NormalTok{.metrics()}

    \KeywordTok{def}\NormalTok{ \_apply\_physics\_curriculum(}\VariableTok{self}\NormalTok{) }\OperatorTok{{-}\textgreater{}} \VariableTok{None}\NormalTok{:}
        \VariableTok{self}\NormalTok{.data.xfrc\_applied[:] }\OperatorTok{=} \DecValTok{0}
        \VariableTok{self}\NormalTok{.\_last\_push\_force }\OperatorTok{=} \FloatTok{0.0}
        \ControlFlowTok{if} \VariableTok{self}\NormalTok{.assistance:}
\NormalTok{            pos }\OperatorTok{=} \VariableTok{self}\NormalTok{.data.xpos[}\VariableTok{self}\NormalTok{.bridge\_body\_id]}
\NormalTok{            vel }\OperatorTok{=} \VariableTok{self}\NormalTok{.data.qvel[}\VariableTok{self}\NormalTok{.bridge\_qvel\_adr:}\VariableTok{self}\NormalTok{.bridge\_qvel\_adr }\OperatorTok{+} \DecValTok{3}\NormalTok{]}
\NormalTok{            rotation }\OperatorTok{=} \VariableTok{self}\NormalTok{.data.xmat[}\VariableTok{self}\NormalTok{.bridge\_body\_id].reshape(}\DecValTok{3}\NormalTok{, }\DecValTok{3}\NormalTok{)}
\NormalTok{            angular\_world }\OperatorTok{=}\NormalTok{ (}
\NormalTok{                rotation}
                \OperatorTok{@} \VariableTok{self}\NormalTok{.data.qvel[}\VariableTok{self}\NormalTok{.bridge\_qvel\_adr }\OperatorTok{+} \DecValTok{3}\NormalTok{:}\VariableTok{self}\NormalTok{.bridge\_qvel\_adr }\OperatorTok{+} \DecValTok{6}\NormalTok{]}
\NormalTok{            )}
\NormalTok{            displacement }\OperatorTok{=}\NormalTok{ pos }\OperatorTok{{-}} \VariableTok{self}\NormalTok{.\_bridge\_nominal\_position}
\NormalTok{            restoring\_force }\OperatorTok{=} \OperatorTok{{-}}\FloatTok{45.0} \OperatorTok{*}\NormalTok{ displacement }\OperatorTok{{-}} \FloatTok{4.0} \OperatorTok{*}\NormalTok{ vel}
\NormalTok{            restoring\_force[}\DecValTok{2}\NormalTok{] }\OperatorTok{=} \OperatorTok{{-}}\FloatTok{25.0} \OperatorTok{*}\NormalTok{ displacement[}\DecValTok{2}\NormalTok{] }\OperatorTok{{-}} \FloatTok{2.5} \OperatorTok{*}\NormalTok{ vel[}\DecValTok{2}\NormalTok{]}
\NormalTok{            restoring\_torque }\OperatorTok{=}\NormalTok{ (}
                \FloatTok{3.5} \OperatorTok{*}\NormalTok{ np.cross(rotation[:, }\DecValTok{2}\NormalTok{], np.array((}\FloatTok{0.0}\NormalTok{, }\FloatTok{0.0}\NormalTok{, }\FloatTok{1.0}\NormalTok{)))}
                \OperatorTok{{-}} \FloatTok{0.18} \OperatorTok{*}\NormalTok{ angular\_world}
\NormalTok{            )}
            \VariableTok{self}\NormalTok{.data.xfrc\_applied[}\VariableTok{self}\NormalTok{.bridge\_body\_id, :}\DecValTok{3}\NormalTok{] }\OperatorTok{=}\NormalTok{ (}
                \VariableTok{self}\NormalTok{.assistance }\OperatorTok{*}\NormalTok{ restoring\_force}
\NormalTok{            )}
            \VariableTok{self}\NormalTok{.data.xfrc\_applied[}\VariableTok{self}\NormalTok{.bridge\_body\_id, }\DecValTok{3}\NormalTok{:] }\OperatorTok{=}\NormalTok{ (}
                \VariableTok{self}\NormalTok{.assistance }\OperatorTok{*}\NormalTok{ restoring\_torque}
\NormalTok{            )}

        \ControlFlowTok{if} \VariableTok{self}\NormalTok{.step\_count }\OperatorTok{*}\NormalTok{ CTRL\_DT }\OperatorTok{\textgreater{}=} \VariableTok{self}\NormalTok{.\_next\_push:}
            \VariableTok{self}\NormalTok{.\_last\_push\_force }\OperatorTok{=} \BuiltInTok{float}\NormalTok{(}\VariableTok{self}\NormalTok{.\_rng.choice((}\OperatorTok{{-}}\FloatTok{3.0}\NormalTok{, }\FloatTok{3.0}\NormalTok{)))}
            \VariableTok{self}\NormalTok{.data.xfrc\_applied[}\VariableTok{self}\NormalTok{.trunk\_body\_id, }\DecValTok{1}\NormalTok{] }\OperatorTok{=} \VariableTok{self}\NormalTok{.\_last\_push\_force}
            \VariableTok{self}\NormalTok{.\_next\_push }\OperatorTok{+=} \VariableTok{self}\NormalTok{.\_rng.uniform(}\FloatTok{1.5}\NormalTok{, }\FloatTok{3.0}\NormalTok{)}

    \KeywordTok{def}\NormalTok{ step(}\VariableTok{self}\NormalTok{, action):}
\NormalTok{        action }\OperatorTok{=}\NormalTok{ np.asarray(action, dtype}\OperatorTok{=}\NormalTok{np.float32)}
        \ControlFlowTok{if}\NormalTok{ action.shape }\OperatorTok{!=}\NormalTok{ (C.NUM\_JOINTS,) }\KeywordTok{or} \KeywordTok{not}\NormalTok{ np.isfinite(action).}\BuiltInTok{all}\NormalTok{():}
            \ControlFlowTok{raise} \PreprocessorTok{ValueError}\NormalTok{(}\StringTok{"actions must be finite with shape (14,)"}\NormalTok{)}
        \ControlFlowTok{if} \VariableTok{self}\NormalTok{.\_episode\_finished:}
            \ControlFlowTok{raise} \PreprocessorTok{RuntimeError}\NormalTok{(}\StringTok{"reset is required after a terminal episode"}\NormalTok{)}

        \VariableTok{self}\NormalTok{.\_apply\_physics\_curriculum()}
\NormalTok{        observation, reward, terminated, truncated, info }\OperatorTok{=} \BuiltInTok{super}\NormalTok{().step(action)}
\NormalTok{        metrics }\OperatorTok{=} \VariableTok{self}\NormalTok{.metrics()}
\NormalTok{        reasons }\OperatorTok{=}\NormalTok{ []}
        \ControlFlowTok{if}\NormalTok{ metrics[}\StringTok{"robot\_ground\_contact"}\NormalTok{]:}
\NormalTok{            reasons.append(}\StringTok{"robot\_ground\_contact"}\NormalTok{)}
        \ControlFlowTok{if}\NormalTok{ metrics[}\StringTok{"tilt\_deg"}\NormalTok{] }\OperatorTok{\textgreater{}} \DecValTok{70}\NormalTok{:}
\NormalTok{            reasons.append(}\StringTok{"tilt"}\NormalTok{)}
        \ControlFlowTok{if}\NormalTok{ metrics[}\StringTok{"trunk\_height\_m"}\NormalTok{] }\OperatorTok{\textless{}} \FloatTok{0.10}\NormalTok{:}
\NormalTok{            reasons.append(}\StringTok{"height"}\NormalTok{)}
        \ControlFlowTok{if}\NormalTok{ (}
            \BuiltInTok{abs}\NormalTok{(metrics[}\StringTok{"trunk\_lateral\_m"}\NormalTok{]) }\OperatorTok{\textgreater{}}\NormalTok{ PLATFORM\_HALF\_WIDTH }\OperatorTok{+} \FloatTok{0.12}
            \KeywordTok{and}\NormalTok{ metrics[}\StringTok{"supported\_feet"}\NormalTok{] }\OperatorTok{==} \DecValTok{0}
\NormalTok{        ):}
\NormalTok{            reasons.append(}\StringTok{"fell\_off\_side"}\NormalTok{)}
        \ControlFlowTok{if} \KeywordTok{not}\NormalTok{ np.isfinite(}\VariableTok{self}\NormalTok{.data.qpos).}\BuiltInTok{all}\NormalTok{() }\KeywordTok{or} \KeywordTok{not}\NormalTok{ np.isfinite(}\VariableTok{self}\NormalTok{.data.qvel).}\BuiltInTok{all}\NormalTok{():}
\NormalTok{            reasons.append(}\StringTok{"nonfinite\_physics"}\NormalTok{)}

\NormalTok{        terminated }\OperatorTok{=}\NormalTok{ terminated }\KeywordTok{or} \BuiltInTok{bool}\NormalTok{(reasons)}
        \VariableTok{self}\NormalTok{.\_episode\_finished }\OperatorTok{=}\NormalTok{ terminated }\KeywordTok{or}\NormalTok{ truncated}
\NormalTok{        info.update(metrics)}
\NormalTok{        info[}\StringTok{"termination\_reasons"}\NormalTok{] }\OperatorTok{=}\NormalTok{ reasons}
        \ControlFlowTok{if} \VariableTok{self}\NormalTok{.\_episode\_finished:}
\NormalTok{            info[}\StringTok{"episode\_rewards"}\NormalTok{] }\OperatorTok{=} \BuiltInTok{dict}\NormalTok{(}\VariableTok{self}\NormalTok{.reward\_sums)}
        \ControlFlowTok{return}\NormalTok{ observation, reward, terminated, truncated, info}

    \KeywordTok{def}\NormalTok{ \_contact\_state(}\VariableTok{self}\NormalTok{) }\OperatorTok{{-}\textgreater{}} \BuiltInTok{dict}\NormalTok{:}
\NormalTok{        feet }\OperatorTok{=} \VariableTok{self}\NormalTok{.foot\_geoms}
\NormalTok{        foot\_supports: }\BuiltInTok{dict}\NormalTok{[}\BuiltInTok{str}\NormalTok{, }\BuiltInTok{set}\NormalTok{[}\BuiltInTok{int}\NormalTok{]] }\OperatorTok{=}\NormalTok{ \{}\StringTok{"left"}\NormalTok{: }\BuiltInTok{set}\NormalTok{(), }\StringTok{"right"}\NormalTok{: }\BuiltInTok{set}\NormalTok{()\}}
\NormalTok{        foot\_ground }\OperatorTok{=}\NormalTok{ \{}\StringTok{"left"}\NormalTok{: }\VariableTok{False}\NormalTok{, }\StringTok{"right"}\NormalTok{: }\VariableTok{False}\NormalTok{\}}
\NormalTok{        robot\_ground }\OperatorTok{=} \VariableTok{False}
\NormalTok{        robot\_bridge }\OperatorTok{=} \VariableTok{False}
\NormalTok{        nonfoot\_support }\OperatorTok{=} \VariableTok{False}
\NormalTok{        foot\_geom\_ids }\OperatorTok{=} \BuiltInTok{set}\NormalTok{(feet.values())}
        \ControlFlowTok{for}\NormalTok{ contact }\KeywordTok{in} \VariableTok{self}\NormalTok{.data.contact:}
\NormalTok{            geom1 }\OperatorTok{=} \BuiltInTok{int}\NormalTok{(contact.geom1)}
\NormalTok{            geom2 }\OperatorTok{=} \BuiltInTok{int}\NormalTok{(contact.geom2)}
\NormalTok{            pair }\OperatorTok{=}\NormalTok{ \{geom1, geom2\}}
            \ControlFlowTok{for}\NormalTok{ side, foot }\KeywordTok{in}\NormalTok{ feet.items():}
                \ControlFlowTok{if}\NormalTok{ foot }\KeywordTok{not} \KeywordTok{in}\NormalTok{ pair:}
                    \ControlFlowTok{continue}
\NormalTok{                other }\OperatorTok{=}\NormalTok{ geom2 }\ControlFlowTok{if}\NormalTok{ geom1 }\OperatorTok{==}\NormalTok{ foot }\ControlFlowTok{else}\NormalTok{ geom1}
                \ControlFlowTok{if}\NormalTok{ other }\KeywordTok{in} \VariableTok{self}\NormalTok{.support\_geoms:}
\NormalTok{                    foot\_supports[side].add(other)}
                \ControlFlowTok{if}\NormalTok{ other }\OperatorTok{==} \VariableTok{self}\NormalTok{.floor\_geom:}
\NormalTok{                    foot\_ground[side] }\OperatorTok{=} \VariableTok{True}
            \ControlFlowTok{if} \VariableTok{self}\NormalTok{.floor\_geom }\KeywordTok{in}\NormalTok{ pair:}
\NormalTok{                other }\OperatorTok{=}\NormalTok{ geom2 }\ControlFlowTok{if}\NormalTok{ geom1 }\OperatorTok{==} \VariableTok{self}\NormalTok{.floor\_geom }\ControlFlowTok{else}\NormalTok{ geom1}
\NormalTok{                body\_id }\OperatorTok{=} \BuiltInTok{int}\NormalTok{(}\VariableTok{self}\NormalTok{.model.geom\_bodyid[other])}
                \ControlFlowTok{if}\NormalTok{ body\_id }\KeywordTok{not} \KeywordTok{in}\NormalTok{ (}\DecValTok{0}\NormalTok{, }\VariableTok{self}\NormalTok{.bridge\_body\_id):}
\NormalTok{                    robot\_ground }\OperatorTok{=} \VariableTok{True}
            \ControlFlowTok{if} \VariableTok{self}\NormalTok{.bridge\_geom\_id }\KeywordTok{in}\NormalTok{ pair:}
\NormalTok{                other }\OperatorTok{=}\NormalTok{ geom2 }\ControlFlowTok{if}\NormalTok{ geom1 }\OperatorTok{==} \VariableTok{self}\NormalTok{.bridge\_geom\_id }\ControlFlowTok{else}\NormalTok{ geom1}
\NormalTok{                body\_id }\OperatorTok{=} \BuiltInTok{int}\NormalTok{(}\VariableTok{self}\NormalTok{.model.geom\_bodyid[other])}
                \ControlFlowTok{if}\NormalTok{ body\_id }\KeywordTok{not} \KeywordTok{in}\NormalTok{ (}\DecValTok{0}\NormalTok{, }\VariableTok{self}\NormalTok{.bridge\_body\_id):}
\NormalTok{                    robot\_bridge }\OperatorTok{=} \VariableTok{True}
            \ControlFlowTok{if}\NormalTok{ pair }\OperatorTok{\&} \VariableTok{self}\NormalTok{.support\_geoms:}
\NormalTok{                support }\OperatorTok{=} \BuiltInTok{next}\NormalTok{(}\BuiltInTok{iter}\NormalTok{(pair }\OperatorTok{\&} \VariableTok{self}\NormalTok{.support\_geoms))}
\NormalTok{                other }\OperatorTok{=}\NormalTok{ geom2 }\ControlFlowTok{if}\NormalTok{ geom1 }\OperatorTok{==}\NormalTok{ support }\ControlFlowTok{else}\NormalTok{ geom1}
\NormalTok{                body\_id }\OperatorTok{=} \BuiltInTok{int}\NormalTok{(}\VariableTok{self}\NormalTok{.model.geom\_bodyid[other])}
                \ControlFlowTok{if}\NormalTok{ body\_id }\KeywordTok{not} \KeywordTok{in}\NormalTok{ (}\DecValTok{0}\NormalTok{, }\VariableTok{self}\NormalTok{.bridge\_body\_id) }\KeywordTok{and}\NormalTok{ other }\KeywordTok{not} \KeywordTok{in}\NormalTok{ foot\_geom\_ids:}
\NormalTok{                    nonfoot\_support }\OperatorTok{=} \VariableTok{True}
        \ControlFlowTok{return}\NormalTok{ \{}
            \StringTok{"foot\_supports"}\NormalTok{: foot\_supports,}
            \StringTok{"foot\_ground"}\NormalTok{: foot\_ground,}
            \StringTok{"robot\_ground"}\NormalTok{: robot\_ground,}
            \StringTok{"robot\_bridge"}\NormalTok{: robot\_bridge,}
            \StringTok{"nonfoot\_support"}\NormalTok{: nonfoot\_support,}
\NormalTok{        \}}

    \KeywordTok{def}\NormalTok{ \_foot\_contacts(}\VariableTok{self}\NormalTok{) }\OperatorTok{{-}\textgreater{}} \BuiltInTok{dict}\NormalTok{[}\BuiltInTok{str}\NormalTok{, }\BuiltInTok{bool}\NormalTok{]:}
\NormalTok{        contacts }\OperatorTok{=} \VariableTok{self}\NormalTok{.\_contact\_state()[}\StringTok{"foot\_supports"}\NormalTok{]}
        \ControlFlowTok{return}\NormalTok{ \{side: }\BuiltInTok{bool}\NormalTok{(surfaces) }\ControlFlowTok{for}\NormalTok{ side, surfaces }\KeywordTok{in}\NormalTok{ contacts.items()\}}

    \KeywordTok{def}\NormalTok{ metrics(}\VariableTok{self}\NormalTok{) }\OperatorTok{{-}\textgreater{}} \BuiltInTok{dict}\NormalTok{:}
\NormalTok{        contacts }\OperatorTok{=} \VariableTok{self}\NormalTok{.\_contact\_state()}
\NormalTok{        foot\_supports }\OperatorTok{=}\NormalTok{ contacts[}\StringTok{"foot\_supports"}\NormalTok{]}
\NormalTok{        feet\_on\_bridge }\OperatorTok{=}\NormalTok{ (}
            \BuiltInTok{int}\NormalTok{(}\VariableTok{self}\NormalTok{.bridge\_geom\_id }\KeywordTok{in}\NormalTok{ foot\_supports[}\StringTok{"left"}\NormalTok{])}
            \OperatorTok{+} \BuiltInTok{int}\NormalTok{(}\VariableTok{self}\NormalTok{.bridge\_geom\_id }\KeywordTok{in}\NormalTok{ foot\_supports[}\StringTok{"right"}\NormalTok{])}
\NormalTok{        )}
        \VariableTok{self}\NormalTok{.\_bridge\_contact\_observed }\OperatorTok{=}\NormalTok{ (}
            \VariableTok{self}\NormalTok{.\_bridge\_contact\_observed }\KeywordTok{or}\NormalTok{ feet\_on\_bridge }\OperatorTok{\textgreater{}} \DecValTok{0}
\NormalTok{        )}
\NormalTok{        left\_end }\OperatorTok{=} \VariableTok{self}\NormalTok{.end\_platform\_geom }\KeywordTok{in}\NormalTok{ foot\_supports[}\StringTok{"left"}\NormalTok{]}
\NormalTok{        right\_end }\OperatorTok{=} \VariableTok{self}\NormalTok{.end\_platform\_geom }\KeywordTok{in}\NormalTok{ foot\_supports[}\StringTok{"right"}\NormalTok{]}
\NormalTok{        full\_traversal }\OperatorTok{=} \BuiltInTok{float}\NormalTok{(}\VariableTok{self}\NormalTok{.data.xpos[}\VariableTok{self}\NormalTok{.trunk\_body\_id, }\DecValTok{0}\NormalTok{]) }\OperatorTok{\textgreater{}=}\NormalTok{ (}
\NormalTok{            END\_PLATFORM\_X }\OperatorTok{{-}} \FloatTok{0.08}
\NormalTok{        )}
\NormalTok{        accepted }\OperatorTok{=}\NormalTok{ (}
\NormalTok{            full\_traversal}
            \KeywordTok{and}\NormalTok{ left\_end}
            \KeywordTok{and}\NormalTok{ right\_end}
            \KeywordTok{and} \KeywordTok{not} \BuiltInTok{any}\NormalTok{(contacts[}\StringTok{"foot\_ground"}\NormalTok{].values())}
            \KeywordTok{and} \KeywordTok{not}\NormalTok{ contacts[}\StringTok{"robot\_ground"}\NormalTok{]}
            \KeywordTok{and} \KeywordTok{not}\NormalTok{ contacts[}\StringTok{"nonfoot\_support"}\NormalTok{]}
            \KeywordTok{and} \VariableTok{self}\NormalTok{.\_bridge\_contact\_observed}
\NormalTok{        )}
        \VariableTok{self}\NormalTok{.\_crossed }\OperatorTok{=} \VariableTok{self}\NormalTok{.\_crossed }\KeywordTok{or}\NormalTok{ accepted}

\NormalTok{        rotation }\OperatorTok{=} \VariableTok{self}\NormalTok{.data.xmat[}\VariableTok{self}\NormalTok{.bridge\_body\_id].reshape(}\DecValTok{3}\NormalTok{, }\DecValTok{3}\NormalTok{)}
\NormalTok{        plank\_up }\OperatorTok{=}\NormalTok{ rotation[:, }\DecValTok{2}\NormalTok{]}
\NormalTok{        roll }\OperatorTok{=}\NormalTok{ math.atan2(}\BuiltInTok{float}\NormalTok{(rotation[}\DecValTok{2}\NormalTok{, }\DecValTok{1}\NormalTok{]), }\BuiltInTok{float}\NormalTok{(rotation[}\DecValTok{2}\NormalTok{, }\DecValTok{2}\NormalTok{]))}
\NormalTok{        pitch }\OperatorTok{=}\NormalTok{ math.atan2(}
            \BuiltInTok{float}\NormalTok{(}\OperatorTok{{-}}\NormalTok{rotation[}\DecValTok{2}\NormalTok{, }\DecValTok{0}\NormalTok{]),}
\NormalTok{            math.sqrt(}\BuiltInTok{float}\NormalTok{(rotation[}\DecValTok{2}\NormalTok{, }\DecValTok{1}\NormalTok{] }\OperatorTok{**} \DecValTok{2} \OperatorTok{+}\NormalTok{ rotation[}\DecValTok{2}\NormalTok{, }\DecValTok{2}\NormalTok{] }\OperatorTok{**} \DecValTok{2}\NormalTok{)),}
\NormalTok{        )}
\NormalTok{        gravity }\OperatorTok{=} \VariableTok{self}\NormalTok{.\_projected\_gravity()}
        \ControlFlowTok{return}\NormalTok{ \{}
            \StringTok{"elapsed\_s"}\NormalTok{: }\VariableTok{self}\NormalTok{.step\_count }\OperatorTok{*}\NormalTok{ CTRL\_DT,}
            \StringTok{"assistance"}\NormalTok{: }\VariableTok{self}\NormalTok{.assistance,}
            \StringTok{"curriculum\_stage"}\NormalTok{: }\VariableTok{self}\NormalTok{.stage,}
            \StringTok{"command"}\NormalTok{: }\VariableTok{self}\NormalTok{.twist\_cmd.tolist(),}
            \StringTok{"trunk\_x\_m"}\NormalTok{: }\BuiltInTok{float}\NormalTok{(}\VariableTok{self}\NormalTok{.data.xpos[}\VariableTok{self}\NormalTok{.trunk\_body\_id, }\DecValTok{0}\NormalTok{]),}
            \StringTok{"trunk\_lateral\_m"}\NormalTok{: }\BuiltInTok{float}\NormalTok{(}\VariableTok{self}\NormalTok{.data.xpos[}\VariableTok{self}\NormalTok{.trunk\_body\_id, }\DecValTok{1}\NormalTok{]),}
            \StringTok{"trunk\_height\_m"}\NormalTok{: }\BuiltInTok{float}\NormalTok{(}\VariableTok{self}\NormalTok{.data.xpos[}\VariableTok{self}\NormalTok{.trunk\_body\_id, }\DecValTok{2}\NormalTok{]),}
            \StringTok{"tilt\_deg"}\NormalTok{: math.degrees(}
\NormalTok{                math.acos(}\BuiltInTok{float}\NormalTok{(np.clip(}\OperatorTok{{-}}\NormalTok{gravity[}\DecValTok{2}\NormalTok{], }\OperatorTok{{-}}\FloatTok{1.0}\NormalTok{, }\FloatTok{1.0}\NormalTok{)))}
\NormalTok{            ),}
            \StringTok{"bridge\_position"}\NormalTok{: }\VariableTok{self}\NormalTok{.data.xpos[}\VariableTok{self}\NormalTok{.bridge\_body\_id].tolist(),}
            \StringTok{"bridge\_roll\_deg"}\NormalTok{: math.degrees(roll),}
            \StringTok{"bridge\_pitch\_deg"}\NormalTok{: math.degrees(pitch),}
            \StringTok{"bridge\_up\_z"}\NormalTok{: }\BuiltInTok{float}\NormalTok{(plank\_up[}\DecValTok{2}\NormalTok{]),}
            \StringTok{"bridge\_speed\_mps"}\NormalTok{: }\BuiltInTok{float}\NormalTok{(}
\NormalTok{                np.linalg.norm(}
                    \VariableTok{self}\NormalTok{.data.qvel[}
                        \VariableTok{self}\NormalTok{.bridge\_qvel\_adr:}\VariableTok{self}\NormalTok{.bridge\_qvel\_adr }\OperatorTok{+} \DecValTok{3}
\NormalTok{                    ]}
\NormalTok{                )}
\NormalTok{            ),}
            \StringTok{"left\_support"}\NormalTok{: }\BuiltInTok{sorted}\NormalTok{(foot\_supports[}\StringTok{"left"}\NormalTok{]),}
            \StringTok{"right\_support"}\NormalTok{: }\BuiltInTok{sorted}\NormalTok{(foot\_supports[}\StringTok{"right"}\NormalTok{]),}
            \StringTok{"supported\_feet"}\NormalTok{: }\BuiltInTok{int}\NormalTok{(}\BuiltInTok{bool}\NormalTok{(foot\_supports[}\StringTok{"left"}\NormalTok{]))}
            \OperatorTok{+} \BuiltInTok{int}\NormalTok{(}\BuiltInTok{bool}\NormalTok{(foot\_supports[}\StringTok{"right"}\NormalTok{])),}
            \StringTok{"feet\_on\_bridge"}\NormalTok{: feet\_on\_bridge,}
            \StringTok{"feet\_on\_end\_platform"}\NormalTok{: }\BuiltInTok{int}\NormalTok{(left\_end) }\OperatorTok{+} \BuiltInTok{int}\NormalTok{(right\_end),}
            \StringTok{"foot\_ground\_contacts"}\NormalTok{: }\BuiltInTok{int}\NormalTok{(contacts[}\StringTok{"foot\_ground"}\NormalTok{][}\StringTok{"left"}\NormalTok{])}
            \OperatorTok{+} \BuiltInTok{int}\NormalTok{(contacts[}\StringTok{"foot\_ground"}\NormalTok{][}\StringTok{"right"}\NormalTok{]),}
            \StringTok{"robot\_ground\_contact"}\NormalTok{: contacts[}\StringTok{"robot\_ground"}\NormalTok{],}
            \StringTok{"robot\_bridge\_contact"}\NormalTok{: contacts[}\StringTok{"robot\_bridge"}\NormalTok{],}
            \StringTok{"nonfoot\_support"}\NormalTok{: contacts[}\StringTok{"nonfoot\_support"}\NormalTok{],}
            \StringTok{"full\_traversal"}\NormalTok{: full\_traversal,}
            \StringTok{"bridge\_contact\_observed"}\NormalTok{: }\VariableTok{self}\NormalTok{.\_bridge\_contact\_observed,}
            \StringTok{"crossing\_accepted"}\NormalTok{: }\VariableTok{self}\NormalTok{.\_crossed,}
            \StringTok{"push\_force\_n"}\NormalTok{: }\VariableTok{self}\NormalTok{.\_last\_push\_force,}
\NormalTok{        \}}

    \KeywordTok{def}\NormalTok{ recipe\_metrics(}\VariableTok{self}\NormalTok{) }\OperatorTok{{-}\textgreater{}} \BuiltInTok{dict}\NormalTok{[}\BuiltInTok{str}\NormalTok{, }\BuiltInTok{float}\NormalTok{]:}
        \CommentTok{"""Numeric{-}only privileged metrics for training logs and evaluation."""}
\NormalTok{        metrics }\OperatorTok{=} \VariableTok{self}\NormalTok{.metrics()}
        \ControlFlowTok{return}\NormalTok{ \{}
            \StringTok{"bridge\_crossed"}\NormalTok{: }\BuiltInTok{float}\NormalTok{(metrics[}\StringTok{"crossing\_accepted"}\NormalTok{]),}
            \StringTok{"bridge\_progress\_m"}\NormalTok{: }\BuiltInTok{float}\NormalTok{(}
\NormalTok{                metrics[}\StringTok{"trunk\_x\_m"}\NormalTok{] }\OperatorTok{{-}}\NormalTok{ (START\_PLATFORM\_X }\OperatorTok{{-}} \FloatTok{0.05}\NormalTok{)}
\NormalTok{            ),}
            \StringTok{"bridge\_assistance"}\NormalTok{: }\BuiltInTok{float}\NormalTok{(}\VariableTok{self}\NormalTok{.assistance),}
            \StringTok{"robot\_floor\_contact"}\NormalTok{: }\BuiltInTok{float}\NormalTok{(metrics[}\StringTok{"robot\_ground\_contact"}\NormalTok{]),}
            \StringTok{"nonfoot\_support"}\NormalTok{: }\BuiltInTok{float}\NormalTok{(metrics[}\StringTok{"nonfoot\_support"}\NormalTok{]),}
            \StringTok{"feet\_support"}\NormalTok{: }\BuiltInTok{float}\NormalTok{(metrics[}\StringTok{"supported\_feet"}\NormalTok{]),}
            \StringTok{"bridge\_contact\_observed"}\NormalTok{: }\BuiltInTok{float}\NormalTok{(}
\NormalTok{                metrics[}\StringTok{"bridge\_contact\_observed"}\NormalTok{]}
\NormalTok{            ),}
            \StringTok{"feet\_on\_bridge"}\NormalTok{: }\BuiltInTok{float}\NormalTok{(metrics[}\StringTok{"feet\_on\_bridge"}\NormalTok{]),}
            \StringTok{"feet\_on\_end\_platform"}\NormalTok{: }\BuiltInTok{float}\NormalTok{(metrics[}\StringTok{"feet\_on\_end\_platform"}\NormalTok{]),}
            \StringTok{"bridge\_roll\_deg"}\NormalTok{: }\BuiltInTok{float}\NormalTok{(metrics[}\StringTok{"bridge\_roll\_deg"}\NormalTok{]),}
            \StringTok{"bridge\_pitch\_deg"}\NormalTok{: }\BuiltInTok{float}\NormalTok{(metrics[}\StringTok{"bridge\_pitch\_deg"}\NormalTok{]),}
            \StringTok{"bridge\_speed\_mps"}\NormalTok{: }\BuiltInTok{float}\NormalTok{(metrics[}\StringTok{"bridge\_speed\_mps"}\NormalTok{]),}
            \StringTok{"robot\_tilt\_deg"}\NormalTok{: }\BuiltInTok{float}\NormalTok{(metrics[}\StringTok{"tilt\_deg"}\NormalTok{]),}
            \StringTok{"robot\_lateral\_m"}\NormalTok{: }\BuiltInTok{float}\NormalTok{(metrics[}\StringTok{"trunk\_lateral\_m"}\NormalTok{]),}
\NormalTok{        \}}
\end{Highlighting}
\end{Shaded}

\subsection{\texorpdfstring{A5.
\texttt{rlx/\allowbreak{}rlx/\allowbreak{}environm\allowbreak{}ents/\allowbreak{}bridge\_\allowbreak{}evaluati\allowbreak{}on.py}}{A5. }}\label{a5.-rlxrlxenvironmentsbridge_evaluation.py}

SHA-256:
\texttt{dd3b1169\allowbreak{}33777d78\allowbreak{}9e860e6a\allowbreak{}1398204d\allowbreak{}dfb60ae2\allowbreak{}8e77cf3f\allowbreak{}35127504\allowbreak{}9936078a\allowbreak{}}

\begin{Shaded}
\begin{Highlighting}[]
\CommentTok{"""Strict per{-}episode acceptance for an unassisted suspended{-}bridge crossing."""}

\ImportTok{from}\NormalTok{ \_\_future\_\_ }\ImportTok{import}\NormalTok{ annotations}

\ImportTok{from}\NormalTok{ dataclasses }\ImportTok{import}\NormalTok{ asdict, dataclass, field}
\ImportTok{import}\NormalTok{ math}
\ImportTok{from}\NormalTok{ numbers }\ImportTok{import}\NormalTok{ Integral, Real}
\ImportTok{from}\NormalTok{ typing }\ImportTok{import}\NormalTok{ Any}


\NormalTok{CONTROL\_HZ }\OperatorTok{=} \DecValTok{50}
\NormalTok{REQUIRED\_METRICS }\OperatorTok{=}\NormalTok{ (}
    \StringTok{"bridge\_crossed"}\NormalTok{,}
    \StringTok{"bridge\_progress\_m"}\NormalTok{,}
    \StringTok{"bridge\_assistance"}\NormalTok{,}
    \StringTok{"robot\_floor\_contact"}\NormalTok{,}
    \StringTok{"nonfoot\_support"}\NormalTok{,}
    \StringTok{"feet\_support"}\NormalTok{,}
    \StringTok{"bridge\_contact\_observed"}\NormalTok{,}
\NormalTok{)}
\NormalTok{\_BINARY\_METRICS }\OperatorTok{=}\NormalTok{ (}
    \StringTok{"bridge\_crossed"}\NormalTok{,}
    \StringTok{"robot\_floor\_contact"}\NormalTok{,}
    \StringTok{"nonfoot\_support"}\NormalTok{,}
    \StringTok{"bridge\_contact\_observed"}\NormalTok{,}
\NormalTok{)}


\AttributeTok{@dataclass}\NormalTok{(frozen}\OperatorTok{=}\VariableTok{True}\NormalTok{)}
\KeywordTok{class}\NormalTok{ BridgeEvaluationCriteria:}
\NormalTok{    version: }\BuiltInTok{int} \OperatorTok{=} \DecValTok{1}
\NormalTok{    required\_steps: }\BuiltInTok{int} \OperatorTok{=} \DecValTok{1000}
\NormalTok{    control\_hz: }\BuiltInTok{int} \OperatorTok{=}\NormalTok{ CONTROL\_HZ}
\NormalTok{    min\_progress\_m: }\BuiltInTok{float} \OperatorTok{=} \FloatTok{1.60}
\NormalTok{    required\_end\_feet: }\BuiltInTok{int} \OperatorTok{=} \DecValTok{2}
\NormalTok{    requires\_unassisted: }\BuiltInTok{bool} \OperatorTok{=} \VariableTok{True}
\NormalTok{    requires\_no\_falls: }\BuiltInTok{bool} \OperatorTok{=} \VariableTok{True}
\NormalTok{    all\_episodes\_required: }\BuiltInTok{bool} \OperatorTok{=} \VariableTok{True}

    \KeywordTok{def}\NormalTok{ \_\_post\_init\_\_(}\VariableTok{self}\NormalTok{) }\OperatorTok{{-}\textgreater{}} \VariableTok{None}\NormalTok{:}
        \ControlFlowTok{if}\NormalTok{ (}
            \BuiltInTok{isinstance}\NormalTok{(}\VariableTok{self}\NormalTok{.required\_steps, }\BuiltInTok{bool}\NormalTok{)}
            \KeywordTok{or} \KeywordTok{not} \BuiltInTok{isinstance}\NormalTok{(}\VariableTok{self}\NormalTok{.required\_steps, Integral)}
            \KeywordTok{or} \VariableTok{self}\NormalTok{.required\_steps }\OperatorTok{\textless{}} \DecValTok{1}
\NormalTok{        ):}
            \ControlFlowTok{raise} \PreprocessorTok{ValueError}\NormalTok{(}\StringTok{"bridge required\_steps must be a positive integer"}\NormalTok{)}

    \KeywordTok{def}\NormalTok{ to\_dict(}\VariableTok{self}\NormalTok{) }\OperatorTok{{-}\textgreater{}} \BuiltInTok{dict}\NormalTok{[}\BuiltInTok{str}\NormalTok{, Any]:}
        \ControlFlowTok{return}\NormalTok{ asdict(}\VariableTok{self}\NormalTok{)}


\KeywordTok{def}\NormalTok{ \_numeric\_metrics(metrics: Any) }\OperatorTok{{-}\textgreater{}} \BuiltInTok{dict}\NormalTok{[}\BuiltInTok{str}\NormalTok{, }\BuiltInTok{float}\NormalTok{] }\OperatorTok{|} \VariableTok{None}\NormalTok{:}
    \ControlFlowTok{if} \KeywordTok{not} \BuiltInTok{isinstance}\NormalTok{(metrics, }\BuiltInTok{dict}\NormalTok{):}
        \ControlFlowTok{return} \VariableTok{None}
\NormalTok{    values: }\BuiltInTok{dict}\NormalTok{[}\BuiltInTok{str}\NormalTok{, }\BuiltInTok{float}\NormalTok{] }\OperatorTok{=}\NormalTok{ \{\}}
    \ControlFlowTok{for}\NormalTok{ key }\KeywordTok{in}\NormalTok{ REQUIRED\_METRICS:}
\NormalTok{        value }\OperatorTok{=}\NormalTok{ metrics.get(key)}
        \ControlFlowTok{if} \BuiltInTok{isinstance}\NormalTok{(value, }\BuiltInTok{bool}\NormalTok{) }\KeywordTok{or} \KeywordTok{not} \BuiltInTok{isinstance}\NormalTok{(value, Real):}
            \ControlFlowTok{return} \VariableTok{None}
\NormalTok{        converted }\OperatorTok{=} \BuiltInTok{float}\NormalTok{(value)}
        \ControlFlowTok{if} \KeywordTok{not}\NormalTok{ math.isfinite(converted):}
            \ControlFlowTok{return} \VariableTok{None}
\NormalTok{        values[key] }\OperatorTok{=}\NormalTok{ converted}
    \ControlFlowTok{return}\NormalTok{ values}


\AttributeTok{@dataclass}
\KeywordTok{class}\NormalTok{ \_Episode:}
\NormalTok{    env\_index: }\BuiltInTok{int}
\NormalTok{    index: }\BuiltInTok{int}
\NormalTok{    steps: }\BuiltInTok{int} \OperatorTok{=} \DecValTok{0}
\NormalTok{    measured\_steps: }\BuiltInTok{int} \OperatorTok{=} \DecValTok{0}
\NormalTok{    max\_progress\_m: }\BuiltInTok{float} \OperatorTok{=} \FloatTok{0.0}
\NormalTok{    crossed: }\BuiltInTok{bool} \OperatorTok{=} \VariableTok{False}
\NormalTok{    crossing\_step: }\BuiltInTok{int} \OperatorTok{|} \VariableTok{None} \OperatorTok{=} \VariableTok{None}
\NormalTok{    start\_support\_observed: }\BuiltInTok{bool} \OperatorTok{=} \VariableTok{False}
\NormalTok{    end\_feet\_at\_crossing: }\BuiltInTok{int} \OperatorTok{=} \DecValTok{0}
\NormalTok{    assisted\_steps: }\BuiltInTok{int} \OperatorTok{=} \DecValTok{0}
\NormalTok{    floor\_contact\_steps: }\BuiltInTok{int} \OperatorTok{=} \DecValTok{0}
\NormalTok{    nonfoot\_support\_steps: }\BuiltInTok{int} \OperatorTok{=} \DecValTok{0}
\NormalTok{    previous\_crossed: }\BuiltInTok{float} \OperatorTok{|} \VariableTok{None} \OperatorTok{=} \VariableTok{None}
\NormalTok{    bridge\_contact\_observed: }\BuiltInTok{bool} \OperatorTok{=} \VariableTok{False}
\NormalTok{    failures: }\BuiltInTok{set}\NormalTok{[}\BuiltInTok{str}\NormalTok{] }\OperatorTok{=}\NormalTok{ field(default\_factory}\OperatorTok{=}\BuiltInTok{set}\NormalTok{)}

    \KeywordTok{def}\NormalTok{ observe(}
        \VariableTok{self}\NormalTok{,}
\NormalTok{        metrics: }\BuiltInTok{dict}\NormalTok{[}\BuiltInTok{str}\NormalTok{, Any],}
\NormalTok{        criteria: BridgeEvaluationCriteria,}
\NormalTok{    ) }\OperatorTok{{-}\textgreater{}} \VariableTok{None}\NormalTok{:}
        \VariableTok{self}\NormalTok{.steps }\OperatorTok{+=} \DecValTok{1}
\NormalTok{        values }\OperatorTok{=}\NormalTok{ \_numeric\_metrics(metrics)}
        \ControlFlowTok{if}\NormalTok{ values }\KeywordTok{is} \VariableTok{None}\NormalTok{:}
            \VariableTok{self}\NormalTok{.failures.add(}\StringTok{"missing, malformed, or non{-}finite bridge metrics"}\NormalTok{)}
            \ControlFlowTok{return}
        \ControlFlowTok{if} \BuiltInTok{any}\NormalTok{(values[key] }\KeywordTok{not} \KeywordTok{in}\NormalTok{ (}\FloatTok{0.0}\NormalTok{, }\FloatTok{1.0}\NormalTok{) }\ControlFlowTok{for}\NormalTok{ key }\KeywordTok{in}\NormalTok{ \_BINARY\_METRICS):}
            \VariableTok{self}\NormalTok{.failures.add(}\StringTok{"invalid binary bridge metrics"}\NormalTok{)}
            \ControlFlowTok{return}
        \ControlFlowTok{if}\NormalTok{ (}
\NormalTok{            values[}\StringTok{"bridge\_progress\_m"}\NormalTok{] }\OperatorTok{\textless{}} \OperatorTok{{-}}\FloatTok{0.5}
            \KeywordTok{or}\NormalTok{ values[}\StringTok{"bridge\_assistance"}\NormalTok{] }\OperatorTok{\textless{}} \DecValTok{0}
            \KeywordTok{or}\NormalTok{ values[}\StringTok{"bridge\_assistance"}\NormalTok{] }\OperatorTok{\textgreater{}} \DecValTok{1}
            \KeywordTok{or}\NormalTok{ values[}\StringTok{"feet\_support"}\NormalTok{] }\KeywordTok{not} \KeywordTok{in}\NormalTok{ (}\FloatTok{0.0}\NormalTok{, }\FloatTok{1.0}\NormalTok{, }\FloatTok{2.0}\NormalTok{)}
\NormalTok{        ):}
            \VariableTok{self}\NormalTok{.failures.add(}\StringTok{"invalid physical bridge metrics"}\NormalTok{)}
            \ControlFlowTok{return}
        \ControlFlowTok{if} \VariableTok{self}\NormalTok{.previous\_crossed }\OperatorTok{==} \FloatTok{1.0} \KeywordTok{and}\NormalTok{ values[}\StringTok{"bridge\_crossed"}\NormalTok{] }\OperatorTok{==} \FloatTok{0.0}\NormalTok{:}
            \VariableTok{self}\NormalTok{.failures.add(}\StringTok{"bridge\_crossed regressed within an episode"}\NormalTok{)}

        \VariableTok{self}\NormalTok{.measured\_steps }\OperatorTok{+=} \DecValTok{1}
        \VariableTok{self}\NormalTok{.max\_progress\_m }\OperatorTok{=} \BuiltInTok{max}\NormalTok{(}
            \VariableTok{self}\NormalTok{.max\_progress\_m,}
\NormalTok{            values[}\StringTok{"bridge\_progress\_m"}\NormalTok{],}
\NormalTok{        )}
        \VariableTok{self}\NormalTok{.start\_support\_observed }\OperatorTok{=}\NormalTok{ (}
            \VariableTok{self}\NormalTok{.start\_support\_observed}
            \KeywordTok{or}\NormalTok{ (}
\NormalTok{                values[}\StringTok{"bridge\_progress\_m"}\NormalTok{] }\OperatorTok{\textless{}=} \FloatTok{0.10}
                \KeywordTok{and}\NormalTok{ values[}\StringTok{"feet\_support"}\NormalTok{] }\OperatorTok{\textgreater{}=} \FloatTok{1.0}
\NormalTok{            )}
\NormalTok{        )}
        \VariableTok{self}\NormalTok{.assisted\_steps }\OperatorTok{+=} \BuiltInTok{int}\NormalTok{(values[}\StringTok{"bridge\_assistance"}\NormalTok{] }\OperatorTok{\textgreater{}} \FloatTok{0.0}\NormalTok{)}
        \VariableTok{self}\NormalTok{.floor\_contact\_steps }\OperatorTok{+=} \BuiltInTok{int}\NormalTok{(values[}\StringTok{"robot\_floor\_contact"}\NormalTok{] }\OperatorTok{==} \FloatTok{1.0}\NormalTok{)}
        \VariableTok{self}\NormalTok{.nonfoot\_support\_steps }\OperatorTok{+=} \BuiltInTok{int}\NormalTok{(values[}\StringTok{"nonfoot\_support"}\NormalTok{] }\OperatorTok{==} \FloatTok{1.0}\NormalTok{)}
        \VariableTok{self}\NormalTok{.bridge\_contact\_observed }\OperatorTok{=}\NormalTok{ (}
            \VariableTok{self}\NormalTok{.bridge\_contact\_observed}
            \KeywordTok{or}\NormalTok{ values[}\StringTok{"bridge\_contact\_observed"}\NormalTok{] }\OperatorTok{==} \FloatTok{1.0}
\NormalTok{        )}

\NormalTok{        crossing\_now }\OperatorTok{=}\NormalTok{ values[}\StringTok{"bridge\_crossed"}\NormalTok{] }\OperatorTok{==} \FloatTok{1.0} \KeywordTok{and} \KeywordTok{not} \VariableTok{self}\NormalTok{.crossed}
        \ControlFlowTok{if}\NormalTok{ crossing\_now:}
            \VariableTok{self}\NormalTok{.crossed }\OperatorTok{=} \VariableTok{True}
            \VariableTok{self}\NormalTok{.crossing\_step }\OperatorTok{=} \VariableTok{self}\NormalTok{.steps}
            \VariableTok{self}\NormalTok{.end\_feet\_at\_crossing }\OperatorTok{=} \BuiltInTok{int}\NormalTok{(values[}\StringTok{"feet\_support"}\NormalTok{])}
            \ControlFlowTok{if} \VariableTok{self}\NormalTok{.end\_feet\_at\_crossing }\OperatorTok{!=}\NormalTok{ criteria.required\_end\_feet:}
                \VariableTok{self}\NormalTok{.failures.add(}\StringTok{"crossing lacked two{-}foot landing support"}\NormalTok{)}
            \ControlFlowTok{if} \KeywordTok{not} \VariableTok{self}\NormalTok{.bridge\_contact\_observed:}
                \VariableTok{self}\NormalTok{.failures.add(}
                    \StringTok{"crossing lacked suspended{-}plank foot{-}contact history"}
\NormalTok{                )}
        \VariableTok{self}\NormalTok{.previous\_crossed }\OperatorTok{=}\NormalTok{ values[}\StringTok{"bridge\_crossed"}\NormalTok{]}

    \KeywordTok{def}\NormalTok{ finish(}
        \VariableTok{self}\NormalTok{,}
\NormalTok{        criteria: BridgeEvaluationCriteria,}
        \OperatorTok{*}\NormalTok{,}
\NormalTok{        terminated: }\BuiltInTok{bool}\NormalTok{,}
\NormalTok{        truncated: }\BuiltInTok{bool}\NormalTok{,}
\NormalTok{    ) }\OperatorTok{{-}\textgreater{}} \BuiltInTok{dict}\NormalTok{[}\BuiltInTok{str}\NormalTok{, Any]:}
\NormalTok{        failures }\OperatorTok{=} \BuiltInTok{set}\NormalTok{(}\VariableTok{self}\NormalTok{.failures)}
\NormalTok{        complete }\OperatorTok{=}\NormalTok{ (}
\NormalTok{            truncated}
            \KeywordTok{and} \KeywordTok{not}\NormalTok{ terminated}
            \KeywordTok{and} \VariableTok{self}\NormalTok{.steps }\OperatorTok{\textgreater{}=}\NormalTok{ criteria.required\_steps}
\NormalTok{        )}
        \ControlFlowTok{if}\NormalTok{ terminated:}
\NormalTok{            failures.add(}\StringTok{"episode terminated or fell before evaluation completion"}\NormalTok{)}
        \ControlFlowTok{if} \KeywordTok{not}\NormalTok{ complete:}
\NormalTok{            failures.add(}
                \SpecialStringTok{f"episode did not complete the required "}
                \SpecialStringTok{f"}\SpecialCharTok{\{}\NormalTok{criteria}\SpecialCharTok{.}\NormalTok{required\_steps }\OperatorTok{/}\NormalTok{ criteria}\SpecialCharTok{.}\NormalTok{control\_hz}\SpecialCharTok{:g\}}\SpecialStringTok{{-}second horizon"}
\NormalTok{            )}
        \ControlFlowTok{if} \VariableTok{self}\NormalTok{.steps }\OperatorTok{==} \DecValTok{0} \KeywordTok{or} \VariableTok{self}\NormalTok{.measured\_steps }\OperatorTok{!=} \VariableTok{self}\NormalTok{.steps:}
\NormalTok{            failures.add(}\StringTok{"incomplete bridge metric coverage"}\NormalTok{)}
        \ControlFlowTok{if} \KeywordTok{not} \VariableTok{self}\NormalTok{.start\_support\_observed:}
\NormalTok{            failures.add(}\StringTok{"launch{-}platform foot support was not observed"}\NormalTok{)}
        \ControlFlowTok{if} \KeywordTok{not} \VariableTok{self}\NormalTok{.crossed:}
\NormalTok{            failures.add(}\StringTok{"contact{-}validated end{-}to{-}end crossing was not observed"}\NormalTok{)}
        \ControlFlowTok{if} \KeywordTok{not} \VariableTok{self}\NormalTok{.bridge\_contact\_observed:}
\NormalTok{            failures.add(}\StringTok{"suspended{-}plank foot contact was not observed"}\NormalTok{)}
        \ControlFlowTok{if} \VariableTok{self}\NormalTok{.max\_progress\_m }\OperatorTok{\textless{}}\NormalTok{ criteria.min\_progress\_m:}
\NormalTok{            failures.add(}\StringTok{"bridge progress remained below end{-}to{-}end target"}\NormalTok{)}
        \ControlFlowTok{if} \VariableTok{self}\NormalTok{.assisted\_steps:}
\NormalTok{            failures.add(}\StringTok{"bridge assistance was active during evaluation"}\NormalTok{)}
        \ControlFlowTok{if} \VariableTok{self}\NormalTok{.floor\_contact\_steps:}
\NormalTok{            failures.add(}\StringTok{"robot contacted the floor"}\NormalTok{)}
        \ControlFlowTok{if} \VariableTok{self}\NormalTok{.nonfoot\_support\_steps:}
\NormalTok{            failures.add(}\StringTok{"non{-}foot robot support contact occurred"}\NormalTok{)}

        \ControlFlowTok{return}\NormalTok{ \{}
            \StringTok{"env\_index"}\NormalTok{: }\VariableTok{self}\NormalTok{.env\_index,}
            \StringTok{"episode\_index"}\NormalTok{: }\VariableTok{self}\NormalTok{.index,}
            \StringTok{"steps"}\NormalTok{: }\VariableTok{self}\NormalTok{.steps,}
            \StringTok{"measured\_steps"}\NormalTok{: }\VariableTok{self}\NormalTok{.measured\_steps,}
            \StringTok{"terminated"}\NormalTok{: }\BuiltInTok{bool}\NormalTok{(terminated),}
            \StringTok{"truncated"}\NormalTok{: }\BuiltInTok{bool}\NormalTok{(truncated),}
            \StringTok{"complete"}\NormalTok{: complete,}
            \StringTok{"max\_progress\_m"}\NormalTok{: }\VariableTok{self}\NormalTok{.max\_progress\_m,}
            \StringTok{"crossed"}\NormalTok{: }\VariableTok{self}\NormalTok{.crossed,}
            \StringTok{"crossing\_step"}\NormalTok{: }\VariableTok{self}\NormalTok{.crossing\_step,}
            \StringTok{"start\_support\_observed"}\NormalTok{: }\VariableTok{self}\NormalTok{.start\_support\_observed,}
            \StringTok{"end\_feet\_at\_crossing"}\NormalTok{: }\VariableTok{self}\NormalTok{.end\_feet\_at\_crossing,}
            \StringTok{"bridge\_contact\_observed"}\NormalTok{: }\VariableTok{self}\NormalTok{.bridge\_contact\_observed,}
            \StringTok{"assisted\_steps"}\NormalTok{: }\VariableTok{self}\NormalTok{.assisted\_steps,}
            \StringTok{"floor\_contact\_steps"}\NormalTok{: }\VariableTok{self}\NormalTok{.floor\_contact\_steps,}
            \StringTok{"nonfoot\_support\_steps"}\NormalTok{: }\VariableTok{self}\NormalTok{.nonfoot\_support\_steps,}
            \StringTok{"passed"}\NormalTok{: }\KeywordTok{not}\NormalTok{ failures,}
            \StringTok{"failures"}\NormalTok{: }\BuiltInTok{sorted}\NormalTok{(failures),}
\NormalTok{        \}}


\KeywordTok{class}\NormalTok{ BridgeEvaluation:}
    \CommentTok{"""Accumulate strict unassisted crossing evidence per vector lane."""}

\NormalTok{    criteria: BridgeEvaluationCriteria}
\NormalTok{    \_current: }\BuiltInTok{list}\NormalTok{[\_Episode]}

    \KeywordTok{def} \FunctionTok{\_\_init\_\_}\NormalTok{(}\VariableTok{self}\NormalTok{, num\_envs: }\BuiltInTok{int}\NormalTok{, required\_steps: }\BuiltInTok{int}\NormalTok{) }\OperatorTok{{-}\textgreater{}} \VariableTok{None}\NormalTok{:}
        \ControlFlowTok{if}\NormalTok{ (}
            \BuiltInTok{isinstance}\NormalTok{(num\_envs, }\BuiltInTok{bool}\NormalTok{)}
            \KeywordTok{or} \KeywordTok{not} \BuiltInTok{isinstance}\NormalTok{(num\_envs, Integral)}
            \KeywordTok{or}\NormalTok{ num\_envs }\OperatorTok{\textless{}} \DecValTok{1}
\NormalTok{        ):}
            \ControlFlowTok{raise} \PreprocessorTok{ValueError}\NormalTok{(}\StringTok{"num\_envs must be positive"}\NormalTok{)}
        \ControlFlowTok{if}\NormalTok{ (}
            \BuiltInTok{isinstance}\NormalTok{(required\_steps, }\BuiltInTok{bool}\NormalTok{)}
            \KeywordTok{or} \KeywordTok{not} \BuiltInTok{isinstance}\NormalTok{(required\_steps, Integral)}
            \KeywordTok{or}\NormalTok{ required\_steps }\OperatorTok{\textless{}} \DecValTok{1}
\NormalTok{        ):}
            \ControlFlowTok{raise} \PreprocessorTok{ValueError}\NormalTok{(}\StringTok{"required\_steps must be a positive integer"}\NormalTok{)}
        \VariableTok{self}\NormalTok{.criteria }\OperatorTok{=}\NormalTok{ BridgeEvaluationCriteria(required\_steps}\OperatorTok{=}\BuiltInTok{int}\NormalTok{(required\_steps))}
        \VariableTok{self}\NormalTok{.\_current }\OperatorTok{=}\NormalTok{ [\_Episode(index, }\DecValTok{0}\NormalTok{) }\ControlFlowTok{for}\NormalTok{ index }\KeywordTok{in} \BuiltInTok{range}\NormalTok{(}\BuiltInTok{int}\NormalTok{(num\_envs))]}
        \VariableTok{self}\NormalTok{.\_episodes: }\BuiltInTok{list}\NormalTok{[}\BuiltInTok{dict}\NormalTok{[}\BuiltInTok{str}\NormalTok{, Any]] }\OperatorTok{=}\NormalTok{ []}

    \KeywordTok{def}\NormalTok{ observe(}
        \VariableTok{self}\NormalTok{,}
\NormalTok{        env\_index: }\BuiltInTok{int}\NormalTok{,}
\NormalTok{        metrics: }\BuiltInTok{dict}\NormalTok{[}\BuiltInTok{str}\NormalTok{, Any],}
\NormalTok{        terminated: }\BuiltInTok{bool} \OperatorTok{=} \VariableTok{False}\NormalTok{,}
\NormalTok{        truncated: }\BuiltInTok{bool} \OperatorTok{=} \VariableTok{False}\NormalTok{,}
\NormalTok{    ) }\OperatorTok{{-}\textgreater{}} \VariableTok{None}\NormalTok{:}
        \ControlFlowTok{if} \KeywordTok{not} \DecValTok{0} \OperatorTok{\textless{}=}\NormalTok{ env\_index }\OperatorTok{\textless{}} \BuiltInTok{len}\NormalTok{(}\VariableTok{self}\NormalTok{.\_current):}
            \ControlFlowTok{raise} \PreprocessorTok{IndexError}\NormalTok{(}\SpecialStringTok{f"env\_index out of range: }\SpecialCharTok{\{}\NormalTok{env\_index}\SpecialCharTok{\}}\SpecialStringTok{"}\NormalTok{)}
\NormalTok{        episode }\OperatorTok{=} \VariableTok{self}\NormalTok{.\_current[env\_index]}
\NormalTok{        episode.observe(metrics, }\VariableTok{self}\NormalTok{.criteria)}
        \ControlFlowTok{if}\NormalTok{ terminated }\KeywordTok{or}\NormalTok{ truncated:}
            \VariableTok{self}\NormalTok{.\_episodes.append(}
\NormalTok{                episode.finish(}
                    \VariableTok{self}\NormalTok{.criteria,}
\NormalTok{                    terminated}\OperatorTok{=}\BuiltInTok{bool}\NormalTok{(terminated),}
\NormalTok{                    truncated}\OperatorTok{=}\BuiltInTok{bool}\NormalTok{(truncated),}
\NormalTok{                )}
\NormalTok{            )}
            \VariableTok{self}\NormalTok{.\_current[env\_index] }\OperatorTok{=}\NormalTok{ \_Episode(env\_index, episode.index }\OperatorTok{+} \DecValTok{1}\NormalTok{)}

    \KeywordTok{def}\NormalTok{ report(}\VariableTok{self}\NormalTok{) }\OperatorTok{{-}\textgreater{}} \BuiltInTok{dict}\NormalTok{[}\BuiltInTok{str}\NormalTok{, Any]:}
\NormalTok{        episodes }\OperatorTok{=}\NormalTok{ [}
            \OperatorTok{*}\VariableTok{self}\NormalTok{.\_episodes,}
            \OperatorTok{*}\NormalTok{(}
\NormalTok{                episode.finish(}
                    \VariableTok{self}\NormalTok{.criteria,}
\NormalTok{                    terminated}\OperatorTok{=}\VariableTok{False}\NormalTok{,}
\NormalTok{                    truncated}\OperatorTok{=}\VariableTok{False}\NormalTok{,}
\NormalTok{                )}
                \ControlFlowTok{for}\NormalTok{ episode }\KeywordTok{in} \VariableTok{self}\NormalTok{.\_current}
                \ControlFlowTok{if}\NormalTok{ episode.steps}
\NormalTok{            ),}
\NormalTok{        ]}
\NormalTok{        covered\_lanes }\OperatorTok{=}\NormalTok{ \{episode[}\StringTok{"env\_index"}\NormalTok{] }\ControlFlowTok{for}\NormalTok{ episode }\KeywordTok{in}\NormalTok{ episodes\}}
\NormalTok{        failures }\OperatorTok{=}\NormalTok{ \{}
\NormalTok{            failure}
            \ControlFlowTok{for}\NormalTok{ episode }\KeywordTok{in}\NormalTok{ episodes}
            \ControlFlowTok{for}\NormalTok{ failure }\KeywordTok{in}\NormalTok{ episode[}\StringTok{"failures"}\NormalTok{]}
\NormalTok{        \}}
        \ControlFlowTok{if} \BuiltInTok{set}\NormalTok{(}\BuiltInTok{range}\NormalTok{(}\BuiltInTok{len}\NormalTok{(}\VariableTok{self}\NormalTok{.\_current))) }\OperatorTok{{-}}\NormalTok{ covered\_lanes }\KeywordTok{or} \KeywordTok{not}\NormalTok{ episodes:}
\NormalTok{            failures.add(}\StringTok{"missing evaluation lane coverage"}\NormalTok{)}
        \ControlFlowTok{return}\NormalTok{ \{}
            \StringTok{"passed"}\NormalTok{: }\BuiltInTok{bool}\NormalTok{(episodes) }\KeywordTok{and} \KeywordTok{not}\NormalTok{ failures,}
            \StringTok{"failures"}\NormalTok{: }\BuiltInTok{sorted}\NormalTok{(failures),}
            \StringTok{"episodes"}\NormalTok{: episodes,}
            \StringTok{"criteria"}\NormalTok{: }\VariableTok{self}\NormalTok{.criteria.to\_dict(),}
            \StringTok{"completed\_episodes"}\NormalTok{: }\BuiltInTok{sum}\NormalTok{(}
                \BuiltInTok{bool}\NormalTok{(episode[}\StringTok{"complete"}\NormalTok{]) }\ControlFlowTok{for}\NormalTok{ episode }\KeywordTok{in}\NormalTok{ episodes}
\NormalTok{            ),}
            \StringTok{"passed\_episodes"}\NormalTok{: }\BuiltInTok{sum}\NormalTok{(}
                \BuiltInTok{bool}\NormalTok{(episode[}\StringTok{"passed"}\NormalTok{]) }\ControlFlowTok{for}\NormalTok{ episode }\KeywordTok{in}\NormalTok{ episodes}
\NormalTok{            ),}
            \StringTok{"measurement\_scope"}\NormalTok{: (}
                \StringTok{"50 Hz per{-}episode numeric recipe metrics; crossing credit "}
                \StringTok{"requires full progress and contact{-}validated two{-}foot landing"}
\NormalTok{            ),}
            \StringTok{"hardware\_safe"}\NormalTok{: }\VariableTok{False}\NormalTok{,}
\NormalTok{        \}}


\VariableTok{\_\_all\_\_} \OperatorTok{=}\NormalTok{ [}
    \StringTok{"BridgeEvaluation"}\NormalTok{,}
    \StringTok{"BridgeEvaluationCriteria"}\NormalTok{,}
    \StringTok{"REQUIRED\_METRICS"}\NormalTok{,}
\NormalTok{]}
\end{Highlighting}
\end{Shaded}

\subsection{\texorpdfstring{A6.
\texttt{rlx/\allowbreak{}rlx/\allowbreak{}models/\allowbreak{}bridge\_\allowbreak{}bootstra\allowbreak{}p.py}}{A6. }}\label{a6.-rlxrlxmodelsbridge_bootstrap.py}

SHA-256:
\texttt{f1dd30fb\allowbreak{}c474dcd0\allowbreak{}0080e9df\allowbreak{}a1026dcf\allowbreak{}1528cb94\allowbreak{}6e640a94\allowbreak{}86588256\allowbreak{}6584ca7a\allowbreak{}}

\begin{Shaded}
\begin{Highlighting}[]
\CommentTok{"""Create a bridge{-}training warm start from the shipped walking actor."""}

\ImportTok{from}\NormalTok{ \_\_future\_\_ }\ImportTok{import}\NormalTok{ annotations}

\ImportTok{import}\NormalTok{ hashlib}
\ImportTok{import}\NormalTok{ math}
\ImportTok{from}\NormalTok{ pathlib }\ImportTok{import}\NormalTok{ Path}
\ImportTok{from}\NormalTok{ typing }\ImportTok{import}\NormalTok{ Any}

\ImportTok{import}\NormalTok{ numpy }\ImportTok{as}\NormalTok{ np}


\NormalTok{WALK\_POLICY }\OperatorTok{=}\NormalTok{ (}
\NormalTok{    Path(}\VariableTok{\_\_file\_\_}\NormalTok{).resolve().parents[}\DecValTok{3}\NormalTok{] }\OperatorTok{/} \StringTok{"microduck"} \OperatorTok{/} \StringTok{"policies"} \OperatorTok{/} \StringTok{"alpha\_walking.onnx"}
\NormalTok{)}
\NormalTok{EXPECTED\_NODE\_TYPES }\OperatorTok{=}\NormalTok{ (}
    \StringTok{"Sub"}\NormalTok{,}
    \StringTok{"Div"}\NormalTok{,}
    \StringTok{"Gemm"}\NormalTok{,}
    \StringTok{"Elu"}\NormalTok{,}
    \StringTok{"Gemm"}\NormalTok{,}
    \StringTok{"Elu"}\NormalTok{,}
    \StringTok{"Gemm"}\NormalTok{,}
    \StringTok{"Elu"}\NormalTok{,}
    \StringTok{"Gemm"}\NormalTok{,}
\NormalTok{)}
\NormalTok{EXPECTED\_INITIALIZERS }\OperatorTok{=}\NormalTok{ \{}
    \StringTok{"obs\_normalizer.\_mean"}\NormalTok{: (}\DecValTok{1}\NormalTok{, }\DecValTok{61}\NormalTok{),}
    \StringTok{"onnx::Div\_24"}\NormalTok{: (}\DecValTok{1}\NormalTok{, }\DecValTok{61}\NormalTok{),}
    \StringTok{"mlp.0.weight"}\NormalTok{: (}\DecValTok{512}\NormalTok{, }\DecValTok{61}\NormalTok{),}
    \StringTok{"mlp.0.bias"}\NormalTok{: (}\DecValTok{512}\NormalTok{,),}
    \StringTok{"mlp.2.weight"}\NormalTok{: (}\DecValTok{256}\NormalTok{, }\DecValTok{512}\NormalTok{),}
    \StringTok{"mlp.2.bias"}\NormalTok{: (}\DecValTok{256}\NormalTok{,),}
    \StringTok{"mlp.4.weight"}\NormalTok{: (}\DecValTok{128}\NormalTok{, }\DecValTok{256}\NormalTok{),}
    \StringTok{"mlp.4.bias"}\NormalTok{: (}\DecValTok{128}\NormalTok{,),}
    \StringTok{"mlp.6.weight"}\NormalTok{: (}\DecValTok{14}\NormalTok{, }\DecValTok{128}\NormalTok{),}
    \StringTok{"mlp.6.bias"}\NormalTok{: (}\DecValTok{14}\NormalTok{,),}
\NormalTok{\}}
\NormalTok{PARITY\_SAMPLES }\OperatorTok{=} \DecValTok{32}
\NormalTok{PARITY\_THRESHOLD }\OperatorTok{=} \FloatTok{1e{-}4}


\KeywordTok{def}\NormalTok{ \_fixed\_shape(value\_info: Any) }\OperatorTok{{-}\textgreater{}} \BuiltInTok{tuple}\NormalTok{[}\BuiltInTok{int}\NormalTok{, ...]:}
    \ControlFlowTok{return} \BuiltInTok{tuple}\NormalTok{(dimension.dim\_value }\ControlFlowTok{for}\NormalTok{ dimension }\KeywordTok{in}\NormalTok{ value\_info.}\BuiltInTok{type}\NormalTok{.tensor\_type.shape.dim)}


\KeywordTok{def}\NormalTok{ \_validated\_source() }\OperatorTok{{-}\textgreater{}} \BuiltInTok{tuple}\NormalTok{[Any, }\BuiltInTok{dict}\NormalTok{[}\BuiltInTok{str}\NormalTok{, np.ndarray], }\BuiltInTok{str}\NormalTok{]:}
    \ImportTok{import}\NormalTok{ onnx}
    \ImportTok{from}\NormalTok{ onnx }\ImportTok{import}\NormalTok{ numpy\_helper}

    \ControlFlowTok{if} \KeywordTok{not}\NormalTok{ WALK\_POLICY.is\_file():}
        \ControlFlowTok{raise} \PreprocessorTok{FileNotFoundError}\NormalTok{(}\SpecialStringTok{f"walking source policy not found: }\SpecialCharTok{\{}\NormalTok{WALK\_POLICY}\SpecialCharTok{\}}\SpecialStringTok{"}\NormalTok{)}
\NormalTok{    source\_hash }\OperatorTok{=}\NormalTok{ hashlib.sha256(WALK\_POLICY.read\_bytes()).hexdigest()}
\NormalTok{    source }\OperatorTok{=}\NormalTok{ onnx.load(WALK\_POLICY)}
\NormalTok{    onnx.checker.check\_model(source)}

    \ControlFlowTok{if} \BuiltInTok{len}\NormalTok{(source.graph.}\BuiltInTok{input}\NormalTok{) }\OperatorTok{!=} \DecValTok{1} \KeywordTok{or}\NormalTok{ source.graph.}\BuiltInTok{input}\NormalTok{[}\DecValTok{0}\NormalTok{].name }\OperatorTok{!=} \StringTok{"obs"}\NormalTok{:}
        \ControlFlowTok{raise} \PreprocessorTok{ValueError}\NormalTok{(}\StringTok{"Walking teacher input contract changed"}\NormalTok{)}
    \ControlFlowTok{if}\NormalTok{ \_fixed\_shape(source.graph.}\BuiltInTok{input}\NormalTok{[}\DecValTok{0}\NormalTok{]) }\OperatorTok{!=}\NormalTok{ (}\DecValTok{1}\NormalTok{, }\DecValTok{61}\NormalTok{):}
        \ControlFlowTok{raise} \PreprocessorTok{ValueError}\NormalTok{(}\StringTok{"Walking teacher observation shape changed"}\NormalTok{)}
    \ControlFlowTok{if} \BuiltInTok{len}\NormalTok{(source.graph.output) }\OperatorTok{!=} \DecValTok{1} \KeywordTok{or}\NormalTok{ source.graph.output[}\DecValTok{0}\NormalTok{].name }\OperatorTok{!=} \StringTok{"actions"}\NormalTok{:}
        \ControlFlowTok{raise} \PreprocessorTok{ValueError}\NormalTok{(}\StringTok{"Walking teacher output contract changed"}\NormalTok{)}
    \ControlFlowTok{if}\NormalTok{ \_fixed\_shape(source.graph.output[}\DecValTok{0}\NormalTok{]) }\OperatorTok{!=}\NormalTok{ (}\DecValTok{1}\NormalTok{, }\DecValTok{14}\NormalTok{):}
        \ControlFlowTok{raise} \PreprocessorTok{ValueError}\NormalTok{(}\StringTok{"Walking teacher action shape changed"}\NormalTok{)}
    \ControlFlowTok{if} \BuiltInTok{tuple}\NormalTok{(node.op\_type }\ControlFlowTok{for}\NormalTok{ node }\KeywordTok{in}\NormalTok{ source.graph.node) }\OperatorTok{!=}\NormalTok{ EXPECTED\_NODE\_TYPES:}
        \ControlFlowTok{raise} \PreprocessorTok{ValueError}\NormalTok{(}\StringTok{"Walking teacher graph changed; exact transfer needs a parity review"}\NormalTok{)}

\NormalTok{    tensors }\OperatorTok{=}\NormalTok{ \{}
\NormalTok{        value.name: np.asarray(numpy\_helper.to\_array(value))}
        \ControlFlowTok{for}\NormalTok{ value }\KeywordTok{in}\NormalTok{ source.graph.initializer}
\NormalTok{    \}}
\NormalTok{    shapes }\OperatorTok{=}\NormalTok{ \{name: value.shape }\ControlFlowTok{for}\NormalTok{ name, value }\KeywordTok{in}\NormalTok{ tensors.items()\}}
    \ControlFlowTok{if}\NormalTok{ shapes }\OperatorTok{!=}\NormalTok{ EXPECTED\_INITIALIZERS:}
        \ControlFlowTok{raise} \PreprocessorTok{ValueError}\NormalTok{(}\StringTok{"Walking teacher initializer contract changed"}\NormalTok{)}
    \ControlFlowTok{if} \BuiltInTok{any}\NormalTok{(value.dtype }\OperatorTok{!=}\NormalTok{ np.float32 }\ControlFlowTok{for}\NormalTok{ value }\KeywordTok{in}\NormalTok{ tensors.values()):}
        \ControlFlowTok{raise} \PreprocessorTok{ValueError}\NormalTok{(}\StringTok{"Walking teacher initializers must be float32"}\NormalTok{)}

\NormalTok{    mean }\OperatorTok{=}\NormalTok{ tensors[}\StringTok{"obs\_normalizer.\_mean"}\NormalTok{]}
\NormalTok{    standard\_deviation }\OperatorTok{=}\NormalTok{ tensors[}\StringTok{"onnx::Div\_24"}\NormalTok{]}
    \ControlFlowTok{if} \KeywordTok{not}\NormalTok{ np.}\BuiltInTok{all}\NormalTok{(np.isfinite(mean)):}
        \ControlFlowTok{raise} \PreprocessorTok{ValueError}\NormalTok{(}\StringTok{"Walking teacher normalizer mean is invalid"}\NormalTok{)}
    \ControlFlowTok{if} \KeywordTok{not}\NormalTok{ np.}\BuiltInTok{all}\NormalTok{(np.isfinite(standard\_deviation)) }\KeywordTok{or} \KeywordTok{not}\NormalTok{ np.}\BuiltInTok{all}\NormalTok{(}
\NormalTok{        standard\_deviation }\OperatorTok{\textgreater{}} \DecValTok{0}
\NormalTok{    ):}
        \ControlFlowTok{raise} \PreprocessorTok{ValueError}\NormalTok{(}\StringTok{"Walking teacher normalizer standard deviation is invalid"}\NormalTok{)}
    \ControlFlowTok{return}\NormalTok{ source, tensors, source\_hash}


\KeywordTok{def}\NormalTok{ \_representative\_observations(}
\NormalTok{    mean: np.ndarray, standard\_deviation: np.ndarray, seed: }\BuiltInTok{int}
\NormalTok{) }\OperatorTok{{-}\textgreater{}}\NormalTok{ np.ndarray:}
\NormalTok{    rng }\OperatorTok{=}\NormalTok{ np.random.default\_rng(seed)}
\NormalTok{    normalized }\OperatorTok{=}\NormalTok{ rng.uniform(}\OperatorTok{{-}}\FloatTok{3.0}\NormalTok{, }\FloatTok{3.0}\NormalTok{, size}\OperatorTok{=}\NormalTok{(PARITY\_SAMPLES }\OperatorTok{{-}} \DecValTok{1}\NormalTok{, }\DecValTok{61}\NormalTok{)).astype(}
\NormalTok{        np.float32}
\NormalTok{    )}
\NormalTok{    normalized }\OperatorTok{=}\NormalTok{ np.concatenate(}
\NormalTok{        (np.zeros((}\DecValTok{1}\NormalTok{, }\DecValTok{61}\NormalTok{), dtype}\OperatorTok{=}\NormalTok{np.float32), normalized), axis}\OperatorTok{=}\DecValTok{0}
\NormalTok{    )}
\NormalTok{    observations }\OperatorTok{=}\NormalTok{ mean.reshape(}\DecValTok{1}\NormalTok{, }\DecValTok{61}\NormalTok{) }\OperatorTok{+}\NormalTok{ normalized }\OperatorTok{*}\NormalTok{ standard\_deviation.reshape(}
        \DecValTok{1}\NormalTok{, }\DecValTok{61}
\NormalTok{    )}
    \ControlFlowTok{if} \KeywordTok{not}\NormalTok{ np.}\BuiltInTok{all}\NormalTok{(np.isfinite(observations)):}
        \ControlFlowTok{raise} \PreprocessorTok{RuntimeError}\NormalTok{(}\StringTok{"Seeded parity observations are not finite"}\NormalTok{)}
    \ControlFlowTok{return}\NormalTok{ observations.astype(np.float32)}


\KeywordTok{def}\NormalTok{ initialize\_bridge(}
\NormalTok{    output: Path, }\OperatorTok{*}\NormalTok{, seed: }\BuiltInTok{int} \OperatorTok{=} \DecValTok{7}\NormalTok{, initial\_std: }\BuiltInTok{float} \OperatorTok{=} \FloatTok{0.03}
\NormalTok{) }\OperatorTok{{-}\textgreater{}}\NormalTok{ Path:}
    \CommentTok{"""Transfer alpha\_walking exactly into an MLX actor with fresh training state."""}
    \ImportTok{import}\NormalTok{ mlx.core }\ImportTok{as}\NormalTok{ mx}
    \ImportTok{import}\NormalTok{ onnxruntime }\ImportTok{as}\NormalTok{ ort}

    \ImportTok{from}\NormalTok{ rlx.models.microduck }\ImportTok{import}\NormalTok{ (}
\NormalTok{        create\_actor\_critic,}
\NormalTok{        normalize\_observations,}
\NormalTok{        save\_checkpoint,}
\NormalTok{    )}

    \ControlFlowTok{if} \KeywordTok{not}\NormalTok{ math.isfinite(initial\_std) }\KeywordTok{or}\NormalTok{ initial\_std }\OperatorTok{\textless{}=} \DecValTok{0}\NormalTok{:}
        \ControlFlowTok{raise} \PreprocessorTok{ValueError}\NormalTok{(}\StringTok{"initial\_std must be finite and positive"}\NormalTok{)}

\NormalTok{    \_, tensors, source\_hash }\OperatorTok{=}\NormalTok{ \_validated\_source()}
\NormalTok{    mean }\OperatorTok{=}\NormalTok{ tensors[}\StringTok{"obs\_normalizer.\_mean"}\NormalTok{].reshape(}\DecValTok{61}\NormalTok{).copy()}
\NormalTok{    standard\_deviation }\OperatorTok{=}\NormalTok{ tensors[}\StringTok{"onnx::Div\_24"}\NormalTok{].reshape(}\DecValTok{61}\NormalTok{).copy()}
\NormalTok{    variance }\OperatorTok{=}\NormalTok{ np.square(standard\_deviation, dtype}\OperatorTok{=}\NormalTok{np.float32)}

\NormalTok{    mx.random.seed(seed)}
\NormalTok{    network }\OperatorTok{=}\NormalTok{ create\_actor\_critic(initial\_std}\OperatorTok{=}\NormalTok{initial\_std)}
\NormalTok{    actor\_weights }\OperatorTok{=}\NormalTok{ [}
\NormalTok{        (}
            \SpecialStringTok{f"actor\_mean.layers.}\SpecialCharTok{\{}\NormalTok{index}\SpecialCharTok{\}}\SpecialStringTok{.}\SpecialCharTok{\{}\NormalTok{kind}\SpecialCharTok{\}}\SpecialStringTok{"}\NormalTok{,}
\NormalTok{            mx.array(tensors[}\SpecialStringTok{f"mlp.}\SpecialCharTok{\{}\NormalTok{index}\SpecialCharTok{\}}\SpecialStringTok{.}\SpecialCharTok{\{}\NormalTok{kind}\SpecialCharTok{\}}\SpecialStringTok{"}\NormalTok{].copy()),}
\NormalTok{        )}
        \ControlFlowTok{for}\NormalTok{ index }\KeywordTok{in}\NormalTok{ (}\DecValTok{0}\NormalTok{, }\DecValTok{2}\NormalTok{, }\DecValTok{4}\NormalTok{, }\DecValTok{6}\NormalTok{)}
        \ControlFlowTok{for}\NormalTok{ kind }\KeywordTok{in}\NormalTok{ (}\StringTok{"weight"}\NormalTok{, }\StringTok{"bias"}\NormalTok{)}
\NormalTok{    ]}
\NormalTok{    network.load\_weights(actor\_weights, strict}\OperatorTok{=}\VariableTok{False}\NormalTok{)}
\NormalTok{    mx.}\BuiltInTok{eval}\NormalTok{(network.parameters())}

\NormalTok{    observations }\OperatorTok{=}\NormalTok{ \_representative\_observations(mean, standard\_deviation, seed)}
\NormalTok{    session }\OperatorTok{=}\NormalTok{ ort.InferenceSession(}
        \BuiltInTok{str}\NormalTok{(WALK\_POLICY), providers}\OperatorTok{=}\NormalTok{[}\StringTok{"CPUExecutionProvider"}\NormalTok{]}
\NormalTok{    )}
\NormalTok{    expected }\OperatorTok{=}\NormalTok{ np.concatenate(}
\NormalTok{        [session.run([}\StringTok{"actions"}\NormalTok{], \{}\StringTok{"obs"}\NormalTok{: item[}\VariableTok{None}\NormalTok{]\})[}\DecValTok{0}\NormalTok{] }\ControlFlowTok{for}\NormalTok{ item }\KeywordTok{in}\NormalTok{ observations],}
\NormalTok{        axis}\OperatorTok{=}\DecValTok{0}\NormalTok{,}
\NormalTok{    )}
\NormalTok{    normalized }\OperatorTok{=}\NormalTok{ normalize\_observations(}
\NormalTok{        mx.array(observations), mean, variance, epsilon}\OperatorTok{=}\FloatTok{0.0}\NormalTok{, clip}\OperatorTok{=}\FloatTok{1e6}
\NormalTok{    )}
\NormalTok{    actual }\OperatorTok{=}\NormalTok{ np.asarray(network.deterministic(normalized), dtype}\OperatorTok{=}\NormalTok{np.float32)}
\NormalTok{    parity\_error }\OperatorTok{=} \BuiltInTok{float}\NormalTok{(np.}\BuiltInTok{max}\NormalTok{(np.}\BuiltInTok{abs}\NormalTok{(actual }\OperatorTok{{-}}\NormalTok{ expected)))}
    \ControlFlowTok{if} \KeywordTok{not}\NormalTok{ math.isfinite(parity\_error) }\KeywordTok{or}\NormalTok{ parity\_error }\OperatorTok{\textgreater{}}\NormalTok{ PARITY\_THRESHOLD:}
        \ControlFlowTok{raise} \PreprocessorTok{RuntimeError}\NormalTok{(}
            \SpecialStringTok{f"Exact walking actor transfer failed parity: }\SpecialCharTok{\{}\NormalTok{parity\_error}\SpecialCharTok{\}}\SpecialStringTok{"}
\NormalTok{        )}
    \ControlFlowTok{if}\NormalTok{ hashlib.sha256(WALK\_POLICY.read\_bytes()).hexdigest() }\OperatorTok{!=}\NormalTok{ source\_hash:}
        \ControlFlowTok{raise} \PreprocessorTok{RuntimeError}\NormalTok{(}\StringTok{"Walking source policy changed during initialization"}\NormalTok{)}

    \ControlFlowTok{return}\NormalTok{ save\_checkpoint(}
\NormalTok{        output,}
\NormalTok{        network,}
\NormalTok{        mean,}
\NormalTok{        variance,}
        \FloatTok{1.0}\NormalTok{,}
\NormalTok{        epsilon}\OperatorTok{=}\FloatTok{0.0}\NormalTok{,}
\NormalTok{        clip}\OperatorTok{=}\FloatTok{1e6}\NormalTok{,}
\NormalTok{        metadata}\OperatorTok{=}\NormalTok{\{}
            \StringTok{"recipe"}\NormalTok{: }\StringTok{"bridge"}\NormalTok{,}
            \StringTok{"steps"}\NormalTok{: }\DecValTok{0}\NormalTok{,}
            \StringTok{"bootstrap"}\NormalTok{: (}
                \StringTok{"exact alpha\_walking actor and normalizer transfer; "}
                \StringTok{"fresh critic and optimizer"}
\NormalTok{            ),}
            \StringTok{"initialization\_method"}\NormalTok{: }\StringTok{"exact\_weight\_transfer"}\NormalTok{,}
            \StringTok{"continuation\_kind"}\NormalTok{: (}
                \StringTok{"actor warm{-}start with fresh local critic and optimizer"}
\NormalTok{            ),}
            \StringTok{"source\_policy"}\NormalTok{: }\BuiltInTok{str}\NormalTok{(WALK\_POLICY.resolve()),}
            \StringTok{"source\_policy\_sha256"}\NormalTok{: source\_hash,}
            \StringTok{"walking\_policy\_sha256"}\NormalTok{: source\_hash,}
            \StringTok{"parity\_samples"}\NormalTok{: }\BuiltInTok{len}\NormalTok{(observations),}
            \StringTok{"parity\_max\_absolute\_action\_error"}\NormalTok{: parity\_error,}
            \StringTok{"parity\_threshold"}\NormalTok{: PARITY\_THRESHOLD,}
            \StringTok{"critic\_initialization"}\NormalTok{: }\StringTok{"fresh\_seeded\_mlx"}\NormalTok{,}
            \StringTok{"optimizer\_initialization"}\NormalTok{: }\StringTok{"fresh\_by\_training\_caller"}\NormalTok{,}
            \StringTok{"freeze\_observation\_normalization\_by\_caller"}\NormalTok{: }\VariableTok{True}\NormalTok{,}
            \StringTok{"seed"}\NormalTok{: seed,}
            \StringTok{"initial\_std"}\NormalTok{: initial\_std,}
            \StringTok{"trained"}\NormalTok{: }\VariableTok{False}\NormalTok{,}
\NormalTok{        \},}
\NormalTok{    )}
\end{Highlighting}
\end{Shaded}

\subsection{\texorpdfstring{A7.
\texttt{rlx/\allowbreak{}rlx/\allowbreak{}algorith\allowbreak{}ms/\allowbreak{}ppo.py}}{A7. }}\label{a7.-rlxrlxalgorithmsppo.py}

SHA-256:
\texttt{126783ef\allowbreak{}9600ad7b\allowbreak{}27030d45\allowbreak{}d87bdf9a\allowbreak{}4f890cb0\allowbreak{}98393167\allowbreak{}17ad46f5\allowbreak{}778b2f83\allowbreak{}}

\begin{Shaded}
\begin{Highlighting}[]
\ImportTok{from}\NormalTok{ dataclasses }\ImportTok{import}\NormalTok{ dataclass}
\ImportTok{from}\NormalTok{ time }\ImportTok{import}\NormalTok{ perf\_counter}
\ImportTok{from}\NormalTok{ typing }\ImportTok{import}\NormalTok{ Callable, Optional}

\ImportTok{import}\NormalTok{ numpy }\ImportTok{as}\NormalTok{ np}

\ImportTok{import}\NormalTok{ mlx.core }\ImportTok{as}\NormalTok{ mx}
\ImportTok{import}\NormalTok{ mlx.nn }\ImportTok{as}\NormalTok{ nn}
\ImportTok{import}\NormalTok{ mlx.optimizers }\ImportTok{as}\NormalTok{ optim}
\ImportTok{from}\NormalTok{ mlx.utils }\ImportTok{import}\NormalTok{ tree\_flatten, tree\_map}

\ImportTok{from}\NormalTok{ rlx.environments.environment }\ImportTok{import}\NormalTok{ Environment}
\ImportTok{from}\NormalTok{ rlx.buffers.rollout\_buffer }\ImportTok{import}\NormalTok{ RolloutBuffer}
\ImportTok{from}\NormalTok{ rlx.utils }\ImportTok{import}\NormalTok{ compute\_generalized\_advantage\_estimate, flatten}


\NormalTok{LOG\_RATIO\_MIN }\OperatorTok{=} \OperatorTok{{-}}\FloatTok{20.0}
\NormalTok{LOG\_RATIO\_MAX }\OperatorTok{=} \FloatTok{20.0}
\NormalTok{VALUE\_LOSS\_DELTA }\OperatorTok{=} \FloatTok{10.0}


\AttributeTok{@dataclass}
\KeywordTok{class}\NormalTok{ PPOConfig:}
\NormalTok{    num\_envs: }\BuiltInTok{int} \OperatorTok{=} \DecValTok{4096}
\NormalTok{    num\_steps: }\BuiltInTok{int} \OperatorTok{=} \DecValTok{16}
\NormalTok{    gamma: }\BuiltInTok{float} \OperatorTok{=} \FloatTok{0.99}
\NormalTok{    gae\_lambda: }\BuiltInTok{float} \OperatorTok{=} \FloatTok{0.95}
\NormalTok{    num\_minibatches: }\BuiltInTok{int} \OperatorTok{=} \DecValTok{16}
\NormalTok{    update\_epochs: }\BuiltInTok{int} \OperatorTok{=} \DecValTok{2}
\NormalTok{    normalize\_advantages: }\BuiltInTok{bool} \OperatorTok{=} \VariableTok{True}
\NormalTok{    clip\_coefficient: }\BuiltInTok{float} \OperatorTok{=} \FloatTok{0.2}
\NormalTok{    clip\_value\_loss: }\BuiltInTok{bool} \OperatorTok{=} \VariableTok{True}
\NormalTok{    entropy\_coefficient: }\BuiltInTok{float} \OperatorTok{=} \FloatTok{0.01}
\NormalTok{    value\_coefficient: }\BuiltInTok{float} \OperatorTok{=} \FloatTok{0.5}
\NormalTok{    max\_grad\_norm: }\BuiltInTok{float} \OperatorTok{=} \FloatTok{0.5}


\AttributeTok{@dataclass}
\KeywordTok{class}\NormalTok{ PPO:}
\NormalTok{    config: PPOConfig}
\NormalTok{    env: Environment}
\NormalTok{    network: nn.Module}
\NormalTok{    optimizer: optim.Optimizer}
    \BuiltInTok{buffer}\NormalTok{: RolloutBuffer}
\NormalTok{    key: mx.array}
\NormalTok{    step: }\BuiltInTok{int} \OperatorTok{=} \DecValTok{0}

    \KeywordTok{def}\NormalTok{ \_\_post\_init\_\_(}\VariableTok{self}\NormalTok{):}
        \VariableTok{self}\NormalTok{.\_constrain\_actor\_log\_std()}
\NormalTok{        mx.}\BuiltInTok{eval}\NormalTok{(}\VariableTok{self}\NormalTok{.network.state)}
        \VariableTok{self}\NormalTok{.\_assert\_finite(}\StringTok{"model parameters"}\NormalTok{, }\VariableTok{self}\NormalTok{.network.parameters())}
\NormalTok{        state }\OperatorTok{=}\NormalTok{ [}\VariableTok{self}\NormalTok{.network.state, }\VariableTok{self}\NormalTok{.optimizer.state]}
        \VariableTok{self}\NormalTok{.update\_step }\OperatorTok{=}\NormalTok{ mx.}\BuiltInTok{compile}\NormalTok{(}\VariableTok{self}\NormalTok{.update\_step, inputs}\OperatorTok{=}\NormalTok{state, outputs}\OperatorTok{=}\NormalTok{state)}
        \VariableTok{self}\NormalTok{.\_update\_step\_with\_metrics }\OperatorTok{=}\NormalTok{ mx.}\BuiltInTok{compile}\NormalTok{(}
            \VariableTok{self}\NormalTok{.\_update\_step\_with\_metrics,}
\NormalTok{            inputs}\OperatorTok{=}\NormalTok{state,}
\NormalTok{            outputs}\OperatorTok{=}\NormalTok{state,}
\NormalTok{        )}

    \KeywordTok{def}\NormalTok{ \_constrain\_actor\_log\_std(}\VariableTok{self}\NormalTok{):}
\NormalTok{        constrain }\OperatorTok{=} \BuiltInTok{getattr}\NormalTok{(}\VariableTok{self}\NormalTok{.network, }\StringTok{"constrain\_actor\_log\_std"}\NormalTok{, }\VariableTok{None}\NormalTok{)}
        \ControlFlowTok{if}\NormalTok{ constrain }\KeywordTok{is} \KeywordTok{not} \VariableTok{None}\NormalTok{:}
\NormalTok{            constrain()}

    \AttributeTok{@staticmethod}
    \KeywordTok{def}\NormalTok{ \_all\_finite(tree):}
\NormalTok{        checks }\OperatorTok{=}\NormalTok{ [mx.}\BuiltInTok{all}\NormalTok{(mx.isfinite(value)) }\ControlFlowTok{for}\NormalTok{ \_, value }\KeywordTok{in}\NormalTok{ tree\_flatten(tree)]}
        \ControlFlowTok{return}\NormalTok{ mx.}\BuiltInTok{all}\NormalTok{(mx.stack(checks))}

    \AttributeTok{@classmethod}
    \KeywordTok{def}\NormalTok{ \_assert\_finite(cls, label, tree):}
\NormalTok{        check }\OperatorTok{=}\NormalTok{ cls.\_all\_finite(tree)}
\NormalTok{        mx.}\BuiltInTok{eval}\NormalTok{(check)}
        \ControlFlowTok{if} \KeywordTok{not} \BuiltInTok{bool}\NormalTok{(np.asarray(check)):}
            \ControlFlowTok{raise} \PreprocessorTok{RuntimeError}\NormalTok{(}\SpecialStringTok{f"PPO }\SpecialCharTok{\{}\NormalTok{label}\SpecialCharTok{\}}\SpecialStringTok{ contains non{-}finite values"}\NormalTok{)}

    \AttributeTok{@staticmethod}
    \KeywordTok{def}\NormalTok{ \_raise\_if\_non\_finite(label, check):}
        \ControlFlowTok{if} \KeywordTok{not} \BuiltInTok{bool}\NormalTok{(np.asarray(check)):}
            \ControlFlowTok{raise} \PreprocessorTok{RuntimeError}\NormalTok{(}\SpecialStringTok{f"PPO }\SpecialCharTok{\{}\NormalTok{label}\SpecialCharTok{\}}\SpecialStringTok{ contains non{-}finite values"}\NormalTok{)}

    \KeywordTok{def}\NormalTok{ \_gradient\_failure\_details(}
        \VariableTok{self}\NormalTok{,}
\NormalTok{        observations,}
\NormalTok{        actions,}
\NormalTok{        old\_log\_probabilities,}
\NormalTok{        old\_values,}
\NormalTok{        advantages,}
\NormalTok{        returns,}
\NormalTok{    ) }\OperatorTok{{-}\textgreater{}} \BuiltInTok{str}\NormalTok{:}
        \CommentTok{"""Recompute one failed minibatch eagerly and name unstable tensors."""}
\NormalTok{        loss, grads }\OperatorTok{=}\NormalTok{ nn.value\_and\_grad(}\VariableTok{self}\NormalTok{.network, }\VariableTok{self}\NormalTok{.loss\_fn)(}
            \VariableTok{self}\NormalTok{.network,}
\NormalTok{            observations,}
\NormalTok{            actions,}
\NormalTok{            old\_log\_probabilities,}
\NormalTok{            old\_values,}
\NormalTok{            advantages,}
\NormalTok{            returns,}
\NormalTok{        )}
\NormalTok{        flat\_grads }\OperatorTok{=}\NormalTok{ tree\_flatten(grads)}
\NormalTok{        mx.}\BuiltInTok{eval}\NormalTok{(loss, }\OperatorTok{*}\NormalTok{(value }\ControlFlowTok{for}\NormalTok{ \_, value }\KeywordTok{in}\NormalTok{ flat\_grads))}
\NormalTok{        bad }\OperatorTok{=}\NormalTok{ [}
\NormalTok{            name }\ControlFlowTok{for}\NormalTok{ name, value }\KeywordTok{in}\NormalTok{ flat\_grads}
            \ControlFlowTok{if} \KeywordTok{not}\NormalTok{ np.isfinite(np.asarray(value)).}\BuiltInTok{all}\NormalTok{()}
\NormalTok{        ]}

        \KeywordTok{def}\NormalTok{ summary(name, value):}
\NormalTok{            array }\OperatorTok{=}\NormalTok{ np.asarray(value, dtype}\OperatorTok{=}\NormalTok{np.float64)}
\NormalTok{            finite }\OperatorTok{=}\NormalTok{ array[np.isfinite(array)]}
            \ControlFlowTok{if}\NormalTok{ finite.size:}
\NormalTok{                bounds }\OperatorTok{=} \SpecialStringTok{f"}\SpecialCharTok{\{}\NormalTok{finite}\SpecialCharTok{.}\BuiltInTok{min}\NormalTok{()}\SpecialCharTok{:.6g\}}\SpecialStringTok{..}\SpecialCharTok{\{}\NormalTok{finite}\SpecialCharTok{.}\BuiltInTok{max}\NormalTok{()}\SpecialCharTok{:.6g\}}\SpecialStringTok{"}
            \ControlFlowTok{else}\NormalTok{:}
\NormalTok{                bounds }\OperatorTok{=} \StringTok{"no finite values"}
            \ControlFlowTok{return} \SpecialStringTok{f"}\SpecialCharTok{\{}\NormalTok{name}\SpecialCharTok{\}}\SpecialStringTok{[}\SpecialCharTok{\{}\NormalTok{bounds}\SpecialCharTok{\}}\SpecialStringTok{; nonfinite=}\SpecialCharTok{\{}\NormalTok{array}\SpecialCharTok{.}\NormalTok{size }\OperatorTok{{-}}\NormalTok{ finite}\SpecialCharTok{.}\NormalTok{size}\SpecialCharTok{\}}\SpecialStringTok{]"}

\NormalTok{        inputs }\OperatorTok{=} \StringTok{", "}\NormalTok{.join(}
\NormalTok{            summary(name, value)}
            \ControlFlowTok{for}\NormalTok{ name, value }\KeywordTok{in}\NormalTok{ (}
\NormalTok{                (}\StringTok{"observations"}\NormalTok{, observations),}
\NormalTok{                (}\StringTok{"actions"}\NormalTok{, actions),}
\NormalTok{                (}\StringTok{"old\_log\_probabilities"}\NormalTok{, old\_log\_probabilities),}
\NormalTok{                (}\StringTok{"old\_values"}\NormalTok{, old\_values),}
\NormalTok{                (}\StringTok{"advantages"}\NormalTok{, advantages),}
\NormalTok{                (}\StringTok{"returns"}\NormalTok{, returns),}
\NormalTok{            )}
\NormalTok{        )}
\NormalTok{        tensors }\OperatorTok{=} \StringTok{", "}\NormalTok{.join(bad) }\ControlFlowTok{if}\NormalTok{ bad }\ControlFlowTok{else} \StringTok{"not reproduced during eager diagnosis"}
        \ControlFlowTok{return} \SpecialStringTok{f"gradient tensors=}\SpecialCharTok{\{}\NormalTok{tensors}\SpecialCharTok{\}}\SpecialStringTok{; }\SpecialCharTok{\{}\NormalTok{inputs}\SpecialCharTok{\}}\SpecialStringTok{"}

    \KeywordTok{def}\NormalTok{ warmup(}\VariableTok{self}\NormalTok{, num\_steps: }\BuiltInTok{int}\NormalTok{):}
        \ControlFlowTok{pass}

    \KeywordTok{def}\NormalTok{ \_truncation\_values(}\VariableTok{self}\NormalTok{, terminated, truncated, info):}
        \CommentTok{"""Evaluate only final timeout states, never autoreset observations."""}
\NormalTok{        selected }\OperatorTok{=}\NormalTok{ np.flatnonzero(np.asarray(truncated) }\OperatorTok{\&} \OperatorTok{\textasciitilde{}}\NormalTok{np.asarray(terminated))}
\NormalTok{        result }\OperatorTok{=}\NormalTok{ mx.zeros(}\VariableTok{self}\NormalTok{.config.num\_envs)}
        \ControlFlowTok{if} \KeywordTok{not}\NormalTok{ selected.size:}
            \ControlFlowTok{return}\NormalTok{ result}
        \ControlFlowTok{if} \StringTok{"terminal\_observation"} \KeywordTok{not} \KeywordTok{in}\NormalTok{ info }\KeywordTok{or} \StringTok{"\_terminal\_observation"} \KeywordTok{not} \KeywordTok{in}\NormalTok{ info:}
            \ControlFlowTok{raise} \PreprocessorTok{ValueError}\NormalTok{(}\StringTok{"PPO timeout requires terminal\_observation and its validity mask"}\NormalTok{)}
\NormalTok{        terminals }\OperatorTok{=}\NormalTok{ mx.array(info[}\StringTok{"terminal\_observation"}\NormalTok{])}
\NormalTok{        mask }\OperatorTok{=}\NormalTok{ np.asarray(info[}\StringTok{"\_terminal\_observation"}\NormalTok{])}
\NormalTok{        expected }\OperatorTok{=}\NormalTok{ (}\VariableTok{self}\NormalTok{.config.num\_envs, }\OperatorTok{*}\VariableTok{self}\NormalTok{.env.observation\_space.shape)}
        \ControlFlowTok{if}\NormalTok{ terminals.shape }\OperatorTok{!=}\NormalTok{ expected }\KeywordTok{or}\NormalTok{ mask.shape }\OperatorTok{!=}\NormalTok{ (}\VariableTok{self}\NormalTok{.config.num\_envs,) }\KeywordTok{or}\NormalTok{ mask.dtype }\OperatorTok{!=}\NormalTok{ np.bool\_:}
            \ControlFlowTok{raise} \PreprocessorTok{ValueError}\NormalTok{(}\StringTok{"PPO terminal observations or validity mask have invalid shape or dtype"}\NormalTok{)}
        \ControlFlowTok{if} \KeywordTok{not}\NormalTok{ np.}\BuiltInTok{all}\NormalTok{(mask[selected]):}
            \ControlFlowTok{raise} \PreprocessorTok{ValueError}\NormalTok{(}\StringTok{"PPO timeout is missing a valid terminal observation"}\NormalTok{)}
\NormalTok{        indices }\OperatorTok{=}\NormalTok{ mx.array(selected.astype(np.int32))}
\NormalTok{        \_, values }\OperatorTok{=} \VariableTok{self}\NormalTok{.network(terminals[indices])}
        \VariableTok{self}\NormalTok{.\_assert\_finite(}\StringTok{"timeout critic values"}\NormalTok{, \{}\StringTok{"values"}\NormalTok{: values\})}
        \ControlFlowTok{return}\NormalTok{ result.at[indices].add(values.squeeze(}\OperatorTok{{-}}\DecValTok{1}\NormalTok{))}

    \KeywordTok{def}\NormalTok{ train(}
        \VariableTok{self}\NormalTok{,}
\NormalTok{        num\_steps: }\BuiltInTok{int}\NormalTok{,}
\NormalTok{        callback: Optional[Callable] }\OperatorTok{=} \VariableTok{None}\NormalTok{,}
        \OperatorTok{*}\NormalTok{,}
\NormalTok{        observer: Optional[Callable[[}\BuiltInTok{dict}\NormalTok{], }\VariableTok{None}\NormalTok{]] }\OperatorTok{=} \VariableTok{None}\NormalTok{,}
\NormalTok{    ):}
        \CommentTok{"""Train to the step target, optionally observing completed phases.}

\CommentTok{        The episode callback keeps its \textasciigrave{}\textasciigrave{}callback(info, step)\textasciigrave{}\textasciigrave{} convention.}
\CommentTok{        Observer dictionaries contain per{-}phase \textasciigrave{}\textasciigrave{}phase\textasciigrave{}\textasciigrave{}, \textasciigrave{}\textasciigrave{}steps\textasciigrave{}\textasciigrave{} and wall}
\CommentTok{        \textasciigrave{}\textasciigrave{}seconds\textasciigrave{}\textasciigrave{}; updates also contain \textasciigrave{}\textasciigrave{}optimizer\_steps\textasciigrave{}\textasciigrave{} and \textasciigrave{}\textasciigrave{}mean\_loss\textasciigrave{}\textasciigrave{}}
\CommentTok{        plus the mean actual{-}minibatch \textasciigrave{}\textasciigrave{}policy\_loss\textasciigrave{}\textasciigrave{}, \textasciigrave{}\textasciigrave{}value\_loss\textasciigrave{}\textasciigrave{},}
\CommentTok{        \textasciigrave{}\textasciigrave{}entropy\textasciigrave{}\textasciigrave{}, \textasciigrave{}\textasciigrave{}approximate\_kl\textasciigrave{}\textasciigrave{}, \textasciigrave{}\textasciigrave{}clip\_fraction\textasciigrave{}\textasciigrave{} and}
\CommentTok{        \textasciigrave{}\textasciigrave{}explained\_variance\textasciigrave{}\textasciigrave{}.}
\CommentTok{        Collection includes buffer reset and rollout; update includes last{-}value}
\CommentTok{        evaluation, GAE and deferred critic work. Observer time is excluded;}
\CommentTok{        these intervals do not isolate device execution or transfer costs.}
\CommentTok{        Observer exceptions propagate. A failed phase emits no completion event.}
\CommentTok{        Observers must not mutate training state or consume training RNGs.}
\CommentTok{        """}
        \VariableTok{self}\NormalTok{.key, reset\_key }\OperatorTok{=}\NormalTok{ mx.random.split(}\VariableTok{self}\NormalTok{.key)}
\NormalTok{        keys }\OperatorTok{=}\NormalTok{ mx.random.split(reset\_key, }\VariableTok{self}\NormalTok{.config.num\_envs)}
\NormalTok{        observation, state, \_ }\OperatorTok{=} \VariableTok{self}\NormalTok{.env.reset(keys)}
\NormalTok{        mx.}\BuiltInTok{eval}\NormalTok{(observation, }\OperatorTok{*}\NormalTok{state.values())}

\NormalTok{        next\_done }\OperatorTok{=}\NormalTok{ mx.zeros(}\VariableTok{self}\NormalTok{.config.num\_envs)}
        \ControlFlowTok{while} \VariableTok{self}\NormalTok{.step }\OperatorTok{\textless{}}\NormalTok{ num\_steps:}
            \ControlFlowTok{if}\NormalTok{ observer }\KeywordTok{is} \KeywordTok{not} \VariableTok{None}\NormalTok{:}
\NormalTok{                collection\_step }\OperatorTok{=} \VariableTok{self}\NormalTok{.step}
\NormalTok{                collection\_start }\OperatorTok{=}\NormalTok{ perf\_counter()}
            \VariableTok{self}\NormalTok{.}\BuiltInTok{buffer}\NormalTok{.reset()}
\NormalTok{            timeout\_values }\OperatorTok{=}\NormalTok{ []}
            \ControlFlowTok{for}\NormalTok{ \_ }\KeywordTok{in} \BuiltInTok{range}\NormalTok{(}\VariableTok{self}\NormalTok{.config.num\_steps):}
                \VariableTok{self}\NormalTok{.\_constrain\_actor\_log\_std()}
\NormalTok{                distribution, value }\OperatorTok{=} \VariableTok{self}\NormalTok{.network(observation)}
\NormalTok{                action }\OperatorTok{=}\NormalTok{ distribution.sample()}
\NormalTok{                log\_probability }\OperatorTok{=}\NormalTok{ distribution.log\_prob(action)}

                \VariableTok{self}\NormalTok{.key, step\_key }\OperatorTok{=}\NormalTok{ mx.random.split(}\VariableTok{self}\NormalTok{.key)}
\NormalTok{                keys }\OperatorTok{=}\NormalTok{ mx.random.split(step\_key, }\VariableTok{self}\NormalTok{.config.num\_envs)}
\NormalTok{                next\_observation, state, reward, terminated, truncated, info }\OperatorTok{=}\NormalTok{ (}
                    \VariableTok{self}\NormalTok{.env.step(keys, state, action)}
\NormalTok{                )}
\NormalTok{                mx.}\BuiltInTok{eval}\NormalTok{(}
\NormalTok{                    next\_observation, reward, terminated, truncated, }\OperatorTok{*}\NormalTok{state.values()}
\NormalTok{                )}

\NormalTok{                timeout\_values.append(}\VariableTok{self}\NormalTok{.\_truncation\_values(terminated, truncated, info))}
                \ControlFlowTok{if}\NormalTok{ callback }\KeywordTok{is} \KeywordTok{not} \VariableTok{None} \KeywordTok{and} \StringTok{"episode"} \KeywordTok{in}\NormalTok{ info:}
\NormalTok{                    callback(info, }\VariableTok{self}\NormalTok{.step)}

                \VariableTok{self}\NormalTok{.}\BuiltInTok{buffer}\NormalTok{.add(}
\NormalTok{                    observation,}
\NormalTok{                    next\_observation,}
\NormalTok{                    action,}
\NormalTok{                    reward,}
\NormalTok{                    terminated,}
\NormalTok{                    truncated,}
\NormalTok{                    value}\OperatorTok{=}\NormalTok{value,}
\NormalTok{                    log\_prob}\OperatorTok{=}\NormalTok{log\_probability,}
\NormalTok{                )}

\NormalTok{                observation }\OperatorTok{=}\NormalTok{ next\_observation}
\NormalTok{                next\_done }\OperatorTok{=}\NormalTok{ mx.logical\_or(terminated, truncated).astype(}
\NormalTok{                    mx.float32}
\NormalTok{                )}
                \VariableTok{self}\NormalTok{.step }\OperatorTok{+=} \VariableTok{self}\NormalTok{.config.num\_envs}

            \ControlFlowTok{if}\NormalTok{ observer }\KeywordTok{is} \KeywordTok{not} \VariableTok{None}\NormalTok{:}
\NormalTok{                collection\_seconds }\OperatorTok{=}\NormalTok{ perf\_counter() }\OperatorTok{{-}}\NormalTok{ collection\_start}
\NormalTok{                phase\_steps }\OperatorTok{=} \VariableTok{self}\NormalTok{.step }\OperatorTok{{-}}\NormalTok{ collection\_step}
\NormalTok{                observer(\{}
                    \StringTok{"phase"}\NormalTok{: }\StringTok{"collection"}\NormalTok{,}
                    \StringTok{"steps"}\NormalTok{: phase\_steps,}
                    \StringTok{"seconds"}\NormalTok{: collection\_seconds,}
\NormalTok{                \})}
\NormalTok{                update\_start }\OperatorTok{=}\NormalTok{ perf\_counter()}

\NormalTok{            \_, last\_value }\OperatorTok{=} \VariableTok{self}\NormalTok{.network(observation)}
\NormalTok{            advantages }\OperatorTok{=}\NormalTok{ compute\_generalized\_advantage\_estimate(}
                \VariableTok{self}\NormalTok{.}\BuiltInTok{buffer}\NormalTok{.rewards,}
                \VariableTok{self}\NormalTok{.}\BuiltInTok{buffer}\NormalTok{.values,}
                \VariableTok{self}\NormalTok{.}\BuiltInTok{buffer}\NormalTok{.terminations,}
\NormalTok{                last\_value.squeeze(}\OperatorTok{{-}}\DecValTok{1}\NormalTok{),}
\NormalTok{                next\_done,}
                \VariableTok{self}\NormalTok{.config.gamma,}
                \VariableTok{self}\NormalTok{.config.gae\_lambda,}
\NormalTok{                truncations}\OperatorTok{=}\VariableTok{self}\NormalTok{.}\BuiltInTok{buffer}\NormalTok{.truncations,}
\NormalTok{                truncation\_values}\OperatorTok{=}\NormalTok{mx.stack(timeout\_values),}
\NormalTok{            )}
\NormalTok{            returns }\OperatorTok{=}\NormalTok{ advantages }\OperatorTok{+} \VariableTok{self}\NormalTok{.}\BuiltInTok{buffer}\NormalTok{.values}

            \ControlFlowTok{if}\NormalTok{ observer }\KeywordTok{is} \VariableTok{None}\NormalTok{:}
                \VariableTok{self}\NormalTok{.update(advantages, returns)}
            \ControlFlowTok{else}\NormalTok{:}
\NormalTok{                metrics }\OperatorTok{=} \VariableTok{self}\NormalTok{.update(advantages, returns, collect\_metrics}\OperatorTok{=}\VariableTok{True}\NormalTok{)}
\NormalTok{                update\_seconds }\OperatorTok{=}\NormalTok{ perf\_counter() }\OperatorTok{{-}}\NormalTok{ update\_start}
\NormalTok{                observer(\{}
                    \StringTok{"phase"}\NormalTok{: }\StringTok{"update"}\NormalTok{,}
                    \StringTok{"steps"}\NormalTok{: phase\_steps,}
                    \StringTok{"seconds"}\NormalTok{: update\_seconds,}
                    \OperatorTok{**}\NormalTok{metrics,}
\NormalTok{                \})}

    \KeywordTok{def}\NormalTok{ loss\_fn(}
        \VariableTok{self}\NormalTok{,}
\NormalTok{        network,}
\NormalTok{        observations,}
\NormalTok{        actions,}
\NormalTok{        old\_log\_probabilities,}
\NormalTok{        old\_values,}
\NormalTok{        advantages,}
\NormalTok{        returns,}
\NormalTok{    ):}
        \ControlFlowTok{return} \VariableTok{self}\NormalTok{.\_loss\_and\_metrics(}
\NormalTok{            network,}
\NormalTok{            observations,}
\NormalTok{            actions,}
\NormalTok{            old\_log\_probabilities,}
\NormalTok{            old\_values,}
\NormalTok{            advantages,}
\NormalTok{            returns,}
\NormalTok{        )[}\DecValTok{0}\NormalTok{]}

    \KeywordTok{def}\NormalTok{ \_loss\_and\_metrics(}
        \VariableTok{self}\NormalTok{,}
\NormalTok{        network,}
\NormalTok{        observations,}
\NormalTok{        actions,}
\NormalTok{        old\_log\_probabilities,}
\NormalTok{        old\_values,}
\NormalTok{        advantages,}
\NormalTok{        returns,}
\NormalTok{    ):}
\NormalTok{        distribution, new\_values }\OperatorTok{=}\NormalTok{ network(observations)}
\NormalTok{        new\_log\_probabilities }\OperatorTok{=}\NormalTok{ distribution.log\_prob(actions)}
\NormalTok{        entropy }\OperatorTok{=}\NormalTok{ distribution.entropy()}

        \CommentTok{\# PPO\textquotesingle{}s clipped objective does not need an unbounded probability ratio.}
        \CommentTok{\# Large but finite log{-}probability differences can overflow exp() during}
        \CommentTok{\# backpropagation before gradient clipping has a chance to act.}
\NormalTok{        log\_ratio }\OperatorTok{=}\NormalTok{ mx.clip(}
\NormalTok{            new\_log\_probabilities }\OperatorTok{{-}}\NormalTok{ old\_log\_probabilities,}
\NormalTok{            LOG\_RATIO\_MIN,}
\NormalTok{            LOG\_RATIO\_MAX,}
\NormalTok{        )}
\NormalTok{        ratio }\OperatorTok{=}\NormalTok{ mx.exp(log\_ratio)}

        \ControlFlowTok{if} \VariableTok{self}\NormalTok{.config.normalize\_advantages:}
\NormalTok{            advantages }\OperatorTok{=}\NormalTok{ (advantages }\OperatorTok{{-}}\NormalTok{ mx.mean(advantages)) }\OperatorTok{/}\NormalTok{ (}
\NormalTok{                mx.std(advantages) }\OperatorTok{+} \FloatTok{1e{-}8}
\NormalTok{            )}

\NormalTok{        policy\_loss }\OperatorTok{=}\NormalTok{ mx.maximum(}
            \OperatorTok{{-}}\NormalTok{advantages }\OperatorTok{*}\NormalTok{ ratio,}
            \OperatorTok{{-}}\NormalTok{advantages}
            \OperatorTok{*}\NormalTok{ mx.clip(}
\NormalTok{                ratio,}
                \DecValTok{1} \OperatorTok{{-}} \VariableTok{self}\NormalTok{.config.clip\_coefficient,}
                \DecValTok{1} \OperatorTok{+} \VariableTok{self}\NormalTok{.config.clip\_coefficient,}
\NormalTok{            ),}
\NormalTok{        )}
\NormalTok{        policy\_loss }\OperatorTok{=}\NormalTok{ mx.mean(policy\_loss)}

\NormalTok{        new\_values }\OperatorTok{=}\NormalTok{ new\_values.squeeze(}\OperatorTok{{-}}\DecValTok{1}\NormalTok{)}

        \KeywordTok{def}\NormalTok{ value\_error\_loss(predictions):}
\NormalTok{            errors }\OperatorTok{=}\NormalTok{ predictions }\OperatorTok{{-}}\NormalTok{ returns}
\NormalTok{            absolute\_errors }\OperatorTok{=}\NormalTok{ mx.}\BuiltInTok{abs}\NormalTok{(errors)}
\NormalTok{            quadratic }\OperatorTok{=}\NormalTok{ mx.minimum(absolute\_errors, VALUE\_LOSS\_DELTA)}
\NormalTok{            linear }\OperatorTok{=}\NormalTok{ absolute\_errors }\OperatorTok{{-}}\NormalTok{ quadratic}
            \ControlFlowTok{return} \FloatTok{0.5} \OperatorTok{*}\NormalTok{ mx.square(quadratic) }\OperatorTok{+}\NormalTok{ VALUE\_LOSS\_DELTA }\OperatorTok{*}\NormalTok{ linear}

        \ControlFlowTok{if} \VariableTok{self}\NormalTok{.config.clip\_value\_loss:}
\NormalTok{            value\_loss\_unclipped }\OperatorTok{=}\NormalTok{ value\_error\_loss(new\_values)}
\NormalTok{            clipped\_values }\OperatorTok{=}\NormalTok{ old\_values }\OperatorTok{+}\NormalTok{ mx.clip(}
\NormalTok{                new\_values }\OperatorTok{{-}}\NormalTok{ old\_values,}
                \OperatorTok{{-}}\VariableTok{self}\NormalTok{.config.clip\_coefficient,}
                \VariableTok{self}\NormalTok{.config.clip\_coefficient,}
\NormalTok{            )}
\NormalTok{            value\_loss\_clipped }\OperatorTok{=}\NormalTok{ value\_error\_loss(clipped\_values)}
\NormalTok{            value\_loss }\OperatorTok{=}\NormalTok{ mx.mean(}
\NormalTok{                mx.maximum(value\_loss\_unclipped, value\_loss\_clipped)}
\NormalTok{            )}
        \ControlFlowTok{else}\NormalTok{:}
\NormalTok{            value\_loss }\OperatorTok{=}\NormalTok{ mx.mean(value\_error\_loss(new\_values))}

\NormalTok{        entropy\_loss }\OperatorTok{=}\NormalTok{ mx.mean(entropy)}

\NormalTok{        loss }\OperatorTok{=}\NormalTok{ (}
\NormalTok{            policy\_loss}
            \OperatorTok{{-}} \VariableTok{self}\NormalTok{.config.entropy\_coefficient }\OperatorTok{*}\NormalTok{ entropy\_loss}
            \OperatorTok{+} \VariableTok{self}\NormalTok{.config.value\_coefficient }\OperatorTok{*}\NormalTok{ value\_loss}
\NormalTok{        )}
\NormalTok{        approximate\_kl }\OperatorTok{=}\NormalTok{ mx.mean((ratio }\OperatorTok{{-}} \DecValTok{1}\NormalTok{) }\OperatorTok{{-}}\NormalTok{ log\_ratio)}
\NormalTok{        clip\_fraction }\OperatorTok{=}\NormalTok{ mx.mean(}
\NormalTok{            mx.}\BuiltInTok{abs}\NormalTok{(ratio }\OperatorTok{{-}} \DecValTok{1}\NormalTok{) }\OperatorTok{\textgreater{}} \VariableTok{self}\NormalTok{.config.clip\_coefficient}
\NormalTok{        )}
\NormalTok{        value\_errors }\OperatorTok{=}\NormalTok{ returns }\OperatorTok{{-}}\NormalTok{ new\_values}
        \ControlFlowTok{return}\NormalTok{ (}
\NormalTok{            loss,}
\NormalTok{            policy\_loss,}
\NormalTok{            value\_loss,}
\NormalTok{            entropy\_loss,}
\NormalTok{            approximate\_kl,}
\NormalTok{            clip\_fraction,}
\NormalTok{            mx.}\BuiltInTok{sum}\NormalTok{(returns),}
\NormalTok{            mx.}\BuiltInTok{sum}\NormalTok{(mx.square(returns)),}
\NormalTok{            mx.}\BuiltInTok{sum}\NormalTok{(value\_errors),}
\NormalTok{            mx.}\BuiltInTok{sum}\NormalTok{(mx.square(value\_errors)),}
\NormalTok{        )}

    \KeywordTok{def}\NormalTok{ \_apply\_update(}\VariableTok{self}\NormalTok{, loss, grads):}
\NormalTok{        loss\_finite }\OperatorTok{=}\NormalTok{ mx.}\BuiltInTok{all}\NormalTok{(mx.isfinite(loss))}
\NormalTok{        gradients\_finite }\OperatorTok{=} \VariableTok{self}\NormalTok{.\_all\_finite(grads)}
\NormalTok{        update\_finite }\OperatorTok{=}\NormalTok{ mx.logical\_and(loss\_finite, gradients\_finite)}
\NormalTok{        grads }\OperatorTok{=}\NormalTok{ tree\_map(}
            \KeywordTok{lambda}\NormalTok{ grad: mx.where(update\_finite, grad, mx.zeros\_like(grad)),}
\NormalTok{            grads,}
\NormalTok{        )}
\NormalTok{        grads, \_ }\OperatorTok{=}\NormalTok{ optim.clip\_grad\_norm(grads, }\VariableTok{self}\NormalTok{.config.max\_grad\_norm)}
        \VariableTok{self}\NormalTok{.optimizer.update(}\VariableTok{self}\NormalTok{.network, grads)}
        \VariableTok{self}\NormalTok{.\_constrain\_actor\_log\_std()}
\NormalTok{        parameters\_finite }\OperatorTok{=} \VariableTok{self}\NormalTok{.\_all\_finite(}\VariableTok{self}\NormalTok{.network.parameters())}
        \ControlFlowTok{return}\NormalTok{ loss\_finite, gradients\_finite, parameters\_finite}

    \KeywordTok{def}\NormalTok{ update\_step(}
        \VariableTok{self}\NormalTok{,}
\NormalTok{        observations,}
\NormalTok{        actions,}
\NormalTok{        old\_log\_probabilities,}
\NormalTok{        old\_values,}
\NormalTok{        advantages,}
\NormalTok{        returns,}
\NormalTok{    ):}
\NormalTok{        loss, grads }\OperatorTok{=}\NormalTok{ nn.value\_and\_grad(}\VariableTok{self}\NormalTok{.network, }\VariableTok{self}\NormalTok{.loss\_fn)(}
            \VariableTok{self}\NormalTok{.network,}
\NormalTok{            observations,}
\NormalTok{            actions,}
\NormalTok{            old\_log\_probabilities,}
\NormalTok{            old\_values,}
\NormalTok{            advantages,}
\NormalTok{            returns,}
\NormalTok{        )}
\NormalTok{        loss\_finite, gradients\_finite, parameters\_finite }\OperatorTok{=} \VariableTok{self}\NormalTok{.\_apply\_update(}
\NormalTok{            loss, grads}
\NormalTok{        )}
        \ControlFlowTok{return}\NormalTok{ loss, loss\_finite, gradients\_finite, parameters\_finite}

    \KeywordTok{def}\NormalTok{ \_update\_step\_with\_metrics(}
        \VariableTok{self}\NormalTok{,}
\NormalTok{        observations,}
\NormalTok{        actions,}
\NormalTok{        old\_log\_probabilities,}
\NormalTok{        old\_values,}
\NormalTok{        advantages,}
\NormalTok{        returns,}
\NormalTok{    ):}
\NormalTok{        metrics, grads }\OperatorTok{=}\NormalTok{ nn.value\_and\_grad(}
            \VariableTok{self}\NormalTok{.network, }\VariableTok{self}\NormalTok{.\_loss\_and\_metrics}
\NormalTok{        )(}
            \VariableTok{self}\NormalTok{.network,}
\NormalTok{            observations,}
\NormalTok{            actions,}
\NormalTok{            old\_log\_probabilities,}
\NormalTok{            old\_values,}
\NormalTok{            advantages,}
\NormalTok{            returns,}
\NormalTok{        )}
\NormalTok{        loss\_finite, gradients\_finite, parameters\_finite }\OperatorTok{=} \VariableTok{self}\NormalTok{.\_apply\_update(}
\NormalTok{            metrics[}\DecValTok{0}\NormalTok{], grads}
\NormalTok{        )}
        \ControlFlowTok{return}\NormalTok{ (}
            \OperatorTok{*}\NormalTok{metrics,}
\NormalTok{            loss\_finite,}
\NormalTok{            gradients\_finite,}
\NormalTok{            parameters\_finite,}
\NormalTok{        )}

    \KeywordTok{def}\NormalTok{ update(}
        \VariableTok{self}\NormalTok{,}
\NormalTok{        advantages: mx.array,}
\NormalTok{        returns: mx.array,}
        \OperatorTok{*}\NormalTok{,}
\NormalTok{        collect\_metrics: }\BuiltInTok{bool} \OperatorTok{=} \VariableTok{False}\NormalTok{,}
\NormalTok{    ):}
        \CommentTok{"""Update minibatches, returning metrics only when requested.}

\CommentTok{        Metrics count completed optimizer steps and average the values produced}
\CommentTok{        by their actual minibatch objective evaluations.}
\CommentTok{        Metrics require at least one optimizer step; exceptions propagate without}
\CommentTok{        returning partial metrics. The default return value remains \textasciigrave{}\textasciigrave{}None\textasciigrave{}\textasciigrave{}.}
\CommentTok{        """}
        \ControlFlowTok{if}\NormalTok{ collect\_metrics:}
\NormalTok{            optimizer\_steps }\OperatorTok{=} \DecValTok{0}
\NormalTok{            totals }\OperatorTok{=}\NormalTok{ \{}
                \StringTok{"mean\_loss"}\NormalTok{: }\FloatTok{0.0}\NormalTok{,}
                \StringTok{"policy\_loss"}\NormalTok{: }\FloatTok{0.0}\NormalTok{,}
                \StringTok{"value\_loss"}\NormalTok{: }\FloatTok{0.0}\NormalTok{,}
                \StringTok{"entropy"}\NormalTok{: }\FloatTok{0.0}\NormalTok{,}
                \StringTok{"approximate\_kl"}\NormalTok{: }\FloatTok{0.0}\NormalTok{,}
                \StringTok{"clip\_fraction"}\NormalTok{: }\FloatTok{0.0}\NormalTok{,}
\NormalTok{            \}}
\NormalTok{            return\_sum }\OperatorTok{=} \FloatTok{0.0}
\NormalTok{            return\_square\_sum }\OperatorTok{=} \FloatTok{0.0}
\NormalTok{            value\_error\_sum }\OperatorTok{=} \FloatTok{0.0}
\NormalTok{            value\_error\_square\_sum }\OperatorTok{=} \FloatTok{0.0}
\NormalTok{            metric\_samples }\OperatorTok{=} \DecValTok{0}
\NormalTok{        observations }\OperatorTok{=}\NormalTok{ flatten(}\VariableTok{self}\NormalTok{.}\BuiltInTok{buffer}\NormalTok{.observations)}
\NormalTok{        actions }\OperatorTok{=}\NormalTok{ flatten(}\VariableTok{self}\NormalTok{.}\BuiltInTok{buffer}\NormalTok{.actions)}
\NormalTok{        log\_probabilities }\OperatorTok{=}\NormalTok{ flatten(}\VariableTok{self}\NormalTok{.}\BuiltInTok{buffer}\NormalTok{.log\_probs)}
\NormalTok{        values }\OperatorTok{=}\NormalTok{ flatten(}\VariableTok{self}\NormalTok{.}\BuiltInTok{buffer}\NormalTok{.values)}
\NormalTok{        advantages }\OperatorTok{=}\NormalTok{ flatten(advantages)}
\NormalTok{        returns }\OperatorTok{=}\NormalTok{ flatten(returns)}
        \VariableTok{self}\NormalTok{.\_assert\_finite(}
            \StringTok{"update inputs"}\NormalTok{,}
\NormalTok{            \{}
                \StringTok{"observations"}\NormalTok{: observations,}
                \StringTok{"actions"}\NormalTok{: actions,}
                \StringTok{"log\_probabilities"}\NormalTok{: log\_probabilities,}
                \StringTok{"values"}\NormalTok{: values,}
                \StringTok{"advantages"}\NormalTok{: advantages,}
                \StringTok{"returns"}\NormalTok{: returns,}
\NormalTok{            \},}
\NormalTok{        )}

\NormalTok{        batch\_size }\OperatorTok{=} \VariableTok{self}\NormalTok{.config.num\_steps }\OperatorTok{*} \VariableTok{self}\NormalTok{.config.num\_envs}
\NormalTok{        minibatch\_size }\OperatorTok{=}\NormalTok{ batch\_size }\OperatorTok{//} \VariableTok{self}\NormalTok{.config.num\_minibatches}

        \ControlFlowTok{for}\NormalTok{ \_ }\KeywordTok{in} \BuiltInTok{range}\NormalTok{(}\VariableTok{self}\NormalTok{.config.update\_epochs):}
\NormalTok{            permutation }\OperatorTok{=}\NormalTok{ np.random.permutation(batch\_size)}
            \ControlFlowTok{for}\NormalTok{ start }\KeywordTok{in} \BuiltInTok{range}\NormalTok{(}\DecValTok{0}\NormalTok{, batch\_size, minibatch\_size):}
\NormalTok{                minibatch\_indices }\OperatorTok{=}\NormalTok{ mx.array(}
\NormalTok{                    permutation[start : start }\OperatorTok{+}\NormalTok{ minibatch\_size]}
\NormalTok{                )}
                \ControlFlowTok{if}\NormalTok{ collect\_metrics:}
\NormalTok{                    result }\OperatorTok{=} \VariableTok{self}\NormalTok{.\_update\_step\_with\_metrics(}
\NormalTok{                        observations[minibatch\_indices],}
\NormalTok{                        actions[minibatch\_indices],}
\NormalTok{                        log\_probabilities[minibatch\_indices],}
\NormalTok{                        values[minibatch\_indices],}
\NormalTok{                        advantages[minibatch\_indices],}
\NormalTok{                        returns[minibatch\_indices],}
\NormalTok{                    )}
\NormalTok{                    metric\_values }\OperatorTok{=}\NormalTok{ result[:}\DecValTok{6}\NormalTok{]}
\NormalTok{                    metric\_moments }\OperatorTok{=}\NormalTok{ result[}\DecValTok{6}\NormalTok{:}\DecValTok{10}\NormalTok{]}
\NormalTok{                    loss\_finite, gradients\_finite, parameters\_finite }\OperatorTok{=}\NormalTok{ result[}\DecValTok{10}\NormalTok{:]}
\NormalTok{                    loss }\OperatorTok{=}\NormalTok{ metric\_values[}\DecValTok{0}\NormalTok{]}
                \ControlFlowTok{else}\NormalTok{:}
\NormalTok{                    loss, loss\_finite, gradients\_finite, parameters\_finite }\OperatorTok{=}\NormalTok{ (}
                        \VariableTok{self}\NormalTok{.update\_step(}
\NormalTok{                            observations[minibatch\_indices],}
\NormalTok{                            actions[minibatch\_indices],}
\NormalTok{                            log\_probabilities[minibatch\_indices],}
\NormalTok{                            values[minibatch\_indices],}
\NormalTok{                            advantages[minibatch\_indices],}
\NormalTok{                            returns[minibatch\_indices],}
\NormalTok{                        )}
\NormalTok{                    )}
\NormalTok{                mx.}\BuiltInTok{eval}\NormalTok{(}
\NormalTok{                    loss,}
\NormalTok{                    loss\_finite,}
\NormalTok{                    gradients\_finite,}
\NormalTok{                    parameters\_finite,}
                    \VariableTok{self}\NormalTok{.network.state,}
                    \VariableTok{self}\NormalTok{.optimizer.state,}
\NormalTok{                )}
                \VariableTok{self}\NormalTok{.\_raise\_if\_non\_finite(}\StringTok{"loss"}\NormalTok{, loss\_finite)}
                \ControlFlowTok{if} \KeywordTok{not} \BuiltInTok{bool}\NormalTok{(np.asarray(gradients\_finite)):}
\NormalTok{                    details }\OperatorTok{=} \VariableTok{self}\NormalTok{.\_gradient\_failure\_details(}
\NormalTok{                        observations[minibatch\_indices],}
\NormalTok{                        actions[minibatch\_indices],}
\NormalTok{                        log\_probabilities[minibatch\_indices],}
\NormalTok{                        values[minibatch\_indices],}
\NormalTok{                        advantages[minibatch\_indices],}
\NormalTok{                        returns[minibatch\_indices],}
\NormalTok{                    )}
                    \ControlFlowTok{raise} \PreprocessorTok{RuntimeError}\NormalTok{(}
                        \SpecialStringTok{f"PPO gradients contains non{-}finite values: }\SpecialCharTok{\{}\NormalTok{details}\SpecialCharTok{\}}\SpecialStringTok{"}
\NormalTok{                    )}
                \VariableTok{self}\NormalTok{.\_raise\_if\_non\_finite(}
                    \StringTok{"model parameters"}\NormalTok{, parameters\_finite}
\NormalTok{                )}
                \ControlFlowTok{if}\NormalTok{ collect\_metrics:}
                    \ControlFlowTok{for}\NormalTok{ name, value }\KeywordTok{in} \BuiltInTok{zip}\NormalTok{(totals, metric\_values):}
\NormalTok{                        totals[name] }\OperatorTok{+=} \BuiltInTok{float}\NormalTok{(value.item())}
\NormalTok{                    return\_sum }\OperatorTok{+=} \BuiltInTok{float}\NormalTok{(metric\_moments[}\DecValTok{0}\NormalTok{].item())}
\NormalTok{                    return\_square\_sum }\OperatorTok{+=} \BuiltInTok{float}\NormalTok{(metric\_moments[}\DecValTok{1}\NormalTok{].item())}
\NormalTok{                    value\_error\_sum }\OperatorTok{+=} \BuiltInTok{float}\NormalTok{(metric\_moments[}\DecValTok{2}\NormalTok{].item())}
\NormalTok{                    value\_error\_square\_sum }\OperatorTok{+=} \BuiltInTok{float}\NormalTok{(metric\_moments[}\DecValTok{3}\NormalTok{].item())}
\NormalTok{                    metric\_samples }\OperatorTok{+=} \BuiltInTok{int}\NormalTok{(returns[minibatch\_indices].size)}
\NormalTok{                    optimizer\_steps }\OperatorTok{+=} \DecValTok{1}

        \ControlFlowTok{if}\NormalTok{ collect\_metrics:}
            \ControlFlowTok{if}\NormalTok{ optimizer\_steps }\OperatorTok{==} \DecValTok{0}\NormalTok{:}
                \ControlFlowTok{raise} \PreprocessorTok{ValueError}\NormalTok{(}\StringTok{"PPO update metrics require at least one optimizer step"}\NormalTok{)}
\NormalTok{            return\_variance }\OperatorTok{=}\NormalTok{ (}
\NormalTok{                return\_square\_sum }\OperatorTok{/}\NormalTok{ metric\_samples}
                \OperatorTok{{-}}\NormalTok{ (return\_sum }\OperatorTok{/}\NormalTok{ metric\_samples) }\OperatorTok{**} \DecValTok{2}
\NormalTok{            )}
\NormalTok{            value\_error\_variance }\OperatorTok{=}\NormalTok{ (}
\NormalTok{                value\_error\_square\_sum }\OperatorTok{/}\NormalTok{ metric\_samples}
                \OperatorTok{{-}}\NormalTok{ (value\_error\_sum }\OperatorTok{/}\NormalTok{ metric\_samples) }\OperatorTok{**} \DecValTok{2}
\NormalTok{            )}
\NormalTok{            explained\_variance }\OperatorTok{=}\NormalTok{ (}
                \DecValTok{1} \OperatorTok{{-}}\NormalTok{ value\_error\_variance }\OperatorTok{/}\NormalTok{ return\_variance}
                \ControlFlowTok{if}\NormalTok{ return\_variance }\OperatorTok{\textgreater{}} \FloatTok{1e{-}8}
                \ControlFlowTok{else} \FloatTok{0.0}
\NormalTok{            )}
            \ControlFlowTok{return}\NormalTok{ \{}
                \StringTok{"optimizer\_steps"}\NormalTok{: optimizer\_steps,}
                \OperatorTok{**}\NormalTok{\{}
\NormalTok{                    name: total }\OperatorTok{/}\NormalTok{ optimizer\_steps}
                    \ControlFlowTok{for}\NormalTok{ name, total }\KeywordTok{in}\NormalTok{ totals.items()}
\NormalTok{                \},}
                \StringTok{"explained\_variance"}\NormalTok{: explained\_variance,}
\NormalTok{            \}}

    \KeywordTok{def}\NormalTok{ evaluate(}\VariableTok{self}\NormalTok{, num\_steps: }\BuiltInTok{int}\NormalTok{, callback: Optional[Callable] }\OperatorTok{=} \VariableTok{None}\NormalTok{):}
        \VariableTok{self}\NormalTok{.key, reset\_key }\OperatorTok{=}\NormalTok{ mx.random.split(}\VariableTok{self}\NormalTok{.key)}
\NormalTok{        keys }\OperatorTok{=}\NormalTok{ mx.random.split(reset\_key, }\VariableTok{self}\NormalTok{.config.num\_envs)}
\NormalTok{        observation, state, \_ }\OperatorTok{=} \VariableTok{self}\NormalTok{.env.reset(keys)}
\NormalTok{        mx.}\BuiltInTok{eval}\NormalTok{(observation, }\OperatorTok{*}\NormalTok{state.values())}

        \ControlFlowTok{for}\NormalTok{ \_ }\KeywordTok{in} \BuiltInTok{range}\NormalTok{(}\DecValTok{0}\NormalTok{, num\_steps, }\VariableTok{self}\NormalTok{.config.num\_envs):}
            \VariableTok{self}\NormalTok{.\_constrain\_actor\_log\_std()}
\NormalTok{            distribution, \_ }\OperatorTok{=} \VariableTok{self}\NormalTok{.network(observation)}
\NormalTok{            action }\OperatorTok{=}\NormalTok{ distribution.sample()}

            \VariableTok{self}\NormalTok{.key, step\_key }\OperatorTok{=}\NormalTok{ mx.random.split(}\VariableTok{self}\NormalTok{.key)}
\NormalTok{            keys }\OperatorTok{=}\NormalTok{ mx.random.split(step\_key, }\VariableTok{self}\NormalTok{.config.num\_envs)}
\NormalTok{            observation, state, reward, terminated, truncated, info }\OperatorTok{=}\NormalTok{ (}
                \VariableTok{self}\NormalTok{.env.step(keys, state, action)}
\NormalTok{            )}
\NormalTok{            mx.}\BuiltInTok{eval}\NormalTok{(observation, reward, terminated, truncated, }\OperatorTok{*}\NormalTok{state.values())}

            \ControlFlowTok{if}\NormalTok{ callback }\KeywordTok{is} \KeywordTok{not} \VariableTok{None} \KeywordTok{and} \StringTok{"episode"} \KeywordTok{in}\NormalTok{ info:}
\NormalTok{                callback(info, }\VariableTok{self}\NormalTok{.step)}
\end{Highlighting}
\end{Shaded}

\subsection{Source inventory /
源码清单}\label{source-inventory-ux6e90ux7801ux6e05ux5355}

{\def\LTcaptype{none} % do not increment counter
\begin{longtable}[]{@{}
  >{\raggedright\arraybackslash}p{(\linewidth - 4\tabcolsep) * \real{0.3000}}
  >{\raggedleft\arraybackslash}p{(\linewidth - 4\tabcolsep) * \real{0.4000}}
  >{\raggedright\arraybackslash}p{(\linewidth - 4\tabcolsep) * \real{0.3000}}@{}}
\toprule\noalign{}
\begin{minipage}[b]{\linewidth}\raggedright
File / 文件
\end{minipage} & \begin{minipage}[b]{\linewidth}\raggedleft
Bytes / 字节
\end{minipage} & \begin{minipage}[b]{\linewidth}\raggedright
Printed / 全文
\end{minipage} \\
\midrule\noalign{}
\endhead
\bottomrule\noalign{}
\endlastfoot
\texttt{rlx/\allowbreak{}rlx/\allowbreak{}environm\allowbreak{}ents/\allowbreak{}basketba\allowbreak{}ll.py}
& 15053 & True \\
\texttt{rlx/\allowbreak{}rlx/\allowbreak{}models/\allowbreak{}basketba\allowbreak{}ll.py}
& 12553 & True \\
\texttt{rlx/\allowbreak{}examples\allowbreak{}/\allowbreak{}ppo\_\allowbreak{}microduc\allowbreak{}k\_\allowbreak{}basketba\allowbreak{}ll.py}
& 25019 & True \\
\texttt{rlx/\allowbreak{}rlx/\allowbreak{}environm\allowbreak{}ents/\allowbreak{}bridge.p\allowbreak{}y}
& 23794 & True \\
\texttt{rlx/\allowbreak{}rlx/\allowbreak{}environm\allowbreak{}ents/\allowbreak{}bridge\_\allowbreak{}evaluati\allowbreak{}on.py}
& 10122 & True \\
\texttt{rlx/\allowbreak{}rlx/\allowbreak{}models/\allowbreak{}bridge\_\allowbreak{}bootstra\allowbreak{}p.py}
& 6821 & True \\
\texttt{rlx/\allowbreak{}rlx/\allowbreak{}algorith\allowbreak{}ms/\allowbreak{}ppo.py}
& 21802 & True \\
\texttt{rlx/\allowbreak{}examples\allowbreak{}/\allowbreak{}ppo\_\allowbreak{}microduc\allowbreak{}k\_\allowbreak{}balance.\allowbreak{}py}
& 37084 & False \\
\texttt{rlx/\allowbreak{}examples\allowbreak{}/\allowbreak{}ppo\_\allowbreak{}microduc\allowbreak{}k\_\allowbreak{}studio.p\allowbreak{}y}
& 59467 & False \\
\texttt{rlx/\allowbreak{}scripts/\allowbreak{}eval\_\allowbreak{}basketba\allowbreak{}ll\_\allowbreak{}local.py\allowbreak{}}
& 29570 & False \\
\texttt{rlx/\allowbreak{}rlx/\allowbreak{}environm\allowbreak{}ents/\allowbreak{}microduc\allowbreak{}k\_\allowbreak{}recipes.\allowbreak{}py}
& 48106 & False \\
\texttt{rlx/\allowbreak{}rlx/\allowbreak{}environm\allowbreak{}ents/\allowbreak{}microduc\allowbreak{}k.py}
& 16607 & False \\
\texttt{rlx/\allowbreak{}rlx/\allowbreak{}models/\allowbreak{}microduc\allowbreak{}k.py}
& 7786 & False \\
\texttt{rlx/\allowbreak{}rlx/\allowbreak{}export/\allowbreak{}microduc\allowbreak{}k\_\allowbreak{}onnx.py}
& 5599 & False \\
\texttt{rlx/\allowbreak{}rlx/\allowbreak{}buffers/\allowbreak{}rollout\_\allowbreak{}buffer.p\allowbreak{}y}
& 3063 & False \\
\texttt{rlx/\allowbreak{}rlx/\allowbreak{}utils/\allowbreak{}utils.py\allowbreak{}}
& 3138 & False \\
\texttt{rlx/\allowbreak{}rlx/\allowbreak{}utils/\allowbreak{}distribu\allowbreak{}tions.py\allowbreak{}}
& 2399 & False \\
\texttt{rlx/\allowbreak{}rlx/\allowbreak{}utils/\allowbreak{}\_\allowbreak{}\_\allowbreak{}init\_\allowbreak{}\_\allowbreak{}.py}
& 522 & False \\
\texttt{microduc\allowbreak{}k\_\allowbreak{}local/\allowbreak{}src/\allowbreak{}microduc\allowbreak{}k\_\allowbreak{}local/\allowbreak{}contract\allowbreak{}.py}
& 4067 & False \\
\texttt{microduc\allowbreak{}k\_\allowbreak{}local/\allowbreak{}src/\allowbreak{}microduc\allowbreak{}k\_\allowbreak{}local/\allowbreak{}walk\_\allowbreak{}env.py}
& 38041 & False \\
\texttt{microduc\allowbreak{}k\_\allowbreak{}local/\allowbreak{}src/\allowbreak{}microduc\allowbreak{}k\_\allowbreak{}local/\allowbreak{}studio\_\allowbreak{}policies\allowbreak{}.py}
& 16512 & False \\
\texttt{microduc\allowbreak{}k\_\allowbreak{}local/\allowbreak{}src/\allowbreak{}microduc\allowbreak{}k\_\allowbreak{}local/\allowbreak{}viz\_\allowbreak{}server.p\allowbreak{}y}
& 185964 & False \\
\texttt{duck-\allowbreak{}viewer/\allowbreak{}lib/\allowbreak{}experime\allowbreak{}nts.ts}
& 37262 & False \\
\texttt{duck-\allowbreak{}viewer/\allowbreak{}lib/\allowbreak{}rlx-\allowbreak{}job.ts}
& 46969 & False \\
\texttt{duck-\allowbreak{}viewer/\allowbreak{}lib/\allowbreak{}scene-\allowbreak{}layout.t\allowbreak{}s}
& 645 & False \\
\texttt{scripts/\allowbreak{}train-\allowbreak{}basketba\allowbreak{}ll.sh}
& 256 & False \\
\texttt{scripts/\allowbreak{}train-\allowbreak{}bridge.s\allowbreak{}h}
& 368 & False \\
\texttt{scripts/\allowbreak{}verify-\allowbreak{}balance-\allowbreak{}studio-\allowbreak{}api.mjs}
& 3756 & False \\
\texttt{scripts/\allowbreak{}verify-\allowbreak{}seven-\allowbreak{}cases.mj\allowbreak{}s}
& 5701 & False \\
\texttt{rlx/\allowbreak{}pyprojec\allowbreak{}t.toml} & 3578 & False \\
\texttt{microduc\allowbreak{}k\_\allowbreak{}local/\allowbreak{}pyprojec\allowbreak{}t.toml}
& 2379 & False \\
\texttt{duck-\allowbreak{}viewer/\allowbreak{}package.\allowbreak{}json}
& 826 & False \\
\texttt{microduc\allowbreak{}k-\allowbreak{}playgrou\allowbreak{}nd/\allowbreak{}artifact\allowbreak{}s/\allowbreak{}basketba\allowbreak{}ll/\allowbreak{}NOTICE}
& 763 & False \\
\texttt{microduc\allowbreak{}k-\allowbreak{}playgrou\allowbreak{}nd/\allowbreak{}artifact\allowbreak{}s/\allowbreak{}basketba\allowbreak{}ll/\allowbreak{}LICENSE}
& 11351 & False \\
\texttt{rlx/\allowbreak{}tests/\allowbreak{}test\_\allowbreak{}basketba\allowbreak{}ll.py}
& 4998 & False \\
\texttt{rlx/\allowbreak{}tests/\allowbreak{}test\_\allowbreak{}basketba\allowbreak{}ll\_\allowbreak{}evaluati\allowbreak{}on.py}
& 7185 & False \\
\texttt{rlx/\allowbreak{}tests/\allowbreak{}test\_\allowbreak{}basketba\allowbreak{}ll\_\allowbreak{}policy.p\allowbreak{}y}
& 16678 & False \\
\texttt{rlx/\allowbreak{}tests/\allowbreak{}test\_\allowbreak{}bridge.p\allowbreak{}y}
& 10383 & False \\
\texttt{rlx/\allowbreak{}tests/\allowbreak{}test\_\allowbreak{}bridge\_\allowbreak{}bootstra\allowbreak{}p.py}
& 2546 & False \\
\texttt{rlx/\allowbreak{}tests/\allowbreak{}test\_\allowbreak{}bridge\_\allowbreak{}evaluati\allowbreak{}on.py}
& 5319 & False \\
\texttt{rlx/\allowbreak{}tests/\allowbreak{}test\_\allowbreak{}balance\_\allowbreak{}adapter.\allowbreak{}py}
& 18789 & False \\
\end{longtable}
}

\end{document}
